\documentclass[a4paper,12pt]{article}


\usepackage{mathtools, bbm, bm,array,amsmath,amsthm}
\usepackage{tikz}
\usepackage{caption}
\usepackage{environ}
\usepackage{graphicx}
\usetikzlibrary{shapes}
\usetikzlibrary{decorations.markings}
\usetikzlibrary{patterns}


\makeatletter
\newcommand\footnoteref[1]{\protected@xdef\@thefnmark{\ref{#1}}\@footnotemark}
\makeatother



\newtheorem{theorem}{Theorem}%  
\newtheorem{proposition}[theorem]{Proposition}% 
\newtheorem{lemma}[theorem]{Lemma}% 
\newtheorem{corollary}[theorem]{Corollary}% 

\newtheorem{definition}[theorem]{Definition}%

\raggedbottom

\let\originalleft\left
\let\originalright\right
\renewcommand{\left}{\mathopen{}\mathclose\bgroup\originalleft}
\renewcommand{\right}{\aftergroup\egroup\originalright}


\newcommand{\semcol}{\! ; \!}


% Commands
\newcommand{\C}{\mathbb{C}}
\newcommand{\K}{\mathbb{K}}
\newcommand{\N}{\mathbb{N}}
\newcommand{\Q}{\mathbb{Q}}
\newcommand{\R}{\mathbb{R}}
\newcommand{\Z}{\mathbb{Z}}
\newcommand{\X}{\mathbb{X}}
\newcommand{\FF}{\mathbb{F}}
\newcommand{\diff}{\mathop{}\!d}
\newcommand{\ecfup}{M^{\uparrow}}
\newcommand{\volup}{V^{\uparrow}}
\newcommand{\symbolup}{\chi}
\newcommand{\weylsymb}{\symbolnot_{\text{w}}}
\newcommand{\symbolnotsobw}{\symbolnot_{\mathcal{W}}}
\newcommand{\symbolnotsch}{\symbolnot_{\mathcal{S}}}
\newcommand{\ballsch}{\mathcal{B}^{(s)}_{\mathcal{S}}}
\newcommand{\ballweight}{\mathcal{B}^{(s,r)}_{\mathcal{W}}}
\newcommand{\tsch}{T_{\mathcal{S}}}
\newcommand{\tweight}{T_{\mathcal{W}}}
\newcommand{\sobsch}{W^{s,2}_{\mathcal{S}}}
\newcommand{\sobweight}{W^{s, r, 2}_{\mathcal{W}}}
\newcommand{\psidosymbclass}{\Gamma^m_\rho}
\newcommand{\hyposymbclass}{\text{H}\Gamma^{m_-,m_+}_\rho}
\newcommand{\hyposymbclassinv}{\mathrm{H}\Gamma^{-m_+,-m_-}_\rho}
\newcommand{\psidoopclass}{\Psi^m_\rho}
\newcommand{\hyponegclass}{\mathrm{H}\Gamma^-}
\newcommand{\hypoposclass}{\mathrm{H}\Gamma^+}
\newcommand{\symbolnot}{\sigma}


\DeclareMathAlphabet{\pazocal}{OMS}{zplm}{m}{n}


\usepackage{pdfsync}

\usepackage{amsmath}
\usepackage{amsfonts}
\usepackage{amssymb}
\usepackage{amsthm}
\usepackage{mathrsfs}
\newtheorem*{lemma*}{Lemma}
\newtheorem*{proposition*}{Proposition}
\newtheorem*{theorem*}{Theorem}
\usepackage[normalem]{ulem}
\usepackage{extarrows}
\usepackage{xcolor}
\usepackage[margin=1cm]{caption}
\usepackage{subcaption}
\usepackage[inline]{enumitem}


\newlength{\mylabelwidth}
\settowidth{\mylabelwidth}{\textbf{(H2)}} % <-- set widest label

\newenvironment{citemize}
  {\begin{list}{}%
     {\setlength{\labelwidth}{\mylabelwidth}%
      \setlength{\leftmargin}{\dimexpr\labelwidth+\labelsep\relax}%
      \renewcommand{\makelabel}[1]{\makebox[\labelwidth][c]{##1}}}}
  {\end{list}}


\usepackage[numbers, sort, compress]{natbib}
\bibliographystyle{IEEEtranS_edit}



\setcounter{tocdepth}{2}



\DeclareUnicodeCharacter{025B}{\ensuremath{\varepsilon}}



\usepackage{framed}
\newenvironment{customabstract}{
    \begin{center}
    \begin{minipage}{0.9\textwidth}
}{%
    \end{minipage}
    \end{center}
    \vspace{1cm}
}




\usepackage[linkcolor=blue,colorlinks=true,citecolor=blue,filecolor=black]{hyperref}
\usepackage{geometry}
 \geometry{
 a4paper,
 total={170mm,257mm},
 left=20mm,
 top=20mm,
 }
\usepackage{setspace}

\newcommand\blindfootnote[1]{%
  \begingroup
  \renewcommand\thefootnote{}\footnote{#1}%
  \addtocounter{footnote}{-1}%
  \endgroup
}


\title{Entropy and Minimax Risk of Hypoelliptic Pseudodifferential Operators}
\date{}
\author{Thomas Allard \\ tallard@ethz.ch  \and Helmut Bölcskei \\ hboelcskei@ethz.ch}




\begin{document}




\maketitle





\abstract{
\noindent
We characterize the %metric 
entropy and
minimax risk of a broad class of compact pseudodifferential operators.
Under suitable decay and regularity conditions on the symbol, 
we combine a Weyl-type asymptotic relation between the 
eigenvalue-counting function and the phase-space volume of the symbol 
with a general correspondence between spectral
quantities, entropy, and minimax risk for compact operators.
This approach yields explicit
asymptotic formulae for both entropy and minimax risk directly in terms
of the symbol.
% Under suitable decay and regularity conditions on the symbol, 
% we establish a Weyl-type asymptotic relation between the 
% eigenvalue-counting function and the phase-space volume of the symbol.
% This relation, combined with the general correspondence between spectral
% quantities, entropy, and minimax risk for compact operators, yields explicit
% asymptotic formulae for both metric entropy and minimax risk directly in terms
% of the symbol.
As an application, we derive sharp entropy and minimax risk asymptotics for
unit balls in Sobolev spaces on unbounded domains, thereby extending
Pinsker’s theorem for Sobolev classes beyond the bounded-domain setting,
and showing that the sharp asymptotic constants are determined by
phase-space geometry rather than domain geometry.
}


\vspace{1cm}

\noindent
{\em H. B\"{o}lcskei dedicates this paper to Prof. Thomas Kailath on the occasion of his 90th birthday.}










\section{Introduction}


Metric entropy provides a natural measure of the massiveness of classes of mathematical objects—such as functions, dynamical systems, or statistical estimators—by quantifying the number of bits required to approximate their elements uniformly to a given precision. 
%Metric entropy provides a natural measure of the number of device for counting the number of bits 
%required to describe mathematical objects (signals, function classes, statistical estimators) uniformly within a prescribed precision level.
This concept plays a central role in a wide range of mathematical fields, including approximation theory \cite{lorentzApproximationFunctions1966,lorentzMetricEntropyApproximation1966,lorentzConstructiveApproximationAdvanced1996,cohenTreeApproximationOptimal2001, cohen2022optimal}, 
harmonic analysis \cite{donohoUnconditionalBasesAre1993,donohoUnconditionalBasesBitLevel1996,donohoDataCompressionHarmonic1998,Grohs2015},
high-dimensional statistics and probability \cite{Vershynin_2018,wainwrightHighDimensionalStatistics2019}, 
and machine learning \cite{elbrachterDeepNeuralNetwork2021,hutterMetricEntropyLimits2022,20.500.11850/734099}. The theory of
metric entropy itself has also attracted renewed interest in recent years  \cite{expdecaypaper,polydecaypaper,ab2025compactoperators,allard2025metricentropyminimaxrisk,20.500.11850/734099}; this line of work goes beyond purely scaling-based considerations and develops a more quantitative framework for complexity estimation.

Minimax risk in non-parametric estimation provides a complementary notion of complexity, quantifying the intrinsic difficulty of recovering elements of a class from noisy observations; see, e.g., \cite{tsybakovIntroductionNonparametricEstimation2009,johnstone2019estimation,nussbaum1999minimax}. While metric entropy and minimax risk are closely related, existing characterizations typically treat them through separate analytical frameworks and under restrictive structural assumptions. The present paper develops a unified framework for characterizing the 
% metric 
entropy and minimax risk of a large class of compact pseudodifferential operators. Explicit asymptotic characterizations are derived in terms of the associated operator symbols. These quantities characterize operator compactness by describing, respectively, the effective size of the operator image of the unit ball and the intrinsic difficulty of recovering elements of this image from noisy measurements \cite{ab2025compactoperators,allard2025metricentropyminimaxrisk}. 


%The present paper develops a joint theory for the characterization of the metric entropy and the minimax risk of a large class of compact %pseudodifferential operators. Specifically, we derive explicit asymptotic characterizations in terms of the operator symbols.
%These quantities measure the degree of operator compactness by quantifying the effective size of their images of unit balls and, respectively, the %intrinsic difficulty of recovering these images from noisy measurements \cite{ab2025compactoperators,allard2025metricentropyminimaxrisk}. 
%We derive explicit asymptotic characterizations for a large family of compact pseudodifferential operators in terms of their symbols.


%In the present paper, we develop a joint theory for the entropy and minimax risk of a large family of compact pseudodifferential operators, yielding %explicit asymptotic characterizations in terms of their symbols. For a compact operator $T$, the associated quantities
%Minimax risk in non-parametric estimation is also a natural device to study such mathematical objects; see, e.g., %\cite{tsybakovIntroductionNonparametricEstimation2009,johnstone2019estimation,nussbaum1999minimax}.
%We provide, in the present paper, a joint theory for the entropy and minimax risk of a large family of compact pseudodifferential  operators.
%$H_{T}$ and $R_{T}$ provide quantitative measures of its compactness \cite{ab2025compactoperators,allard2025metricentropyminimaxrisk}.


%These operator-level quantities also allow one to quantify the metric entropy and minimax risk of sets obtained as the image of a compact operator.
%The point of view of compact operators allows us to characterize the metric entropy and minimax risk of mathematical objects that can be expressed as the %image of a (compact) operator. These operator-level quantities also allow one to quantify the metric entropy and minimax risk of sets obtained as the %image of a compact operator.

We build on the well-known observation that images of pseudodifferential operators are often localized in phase space, with localization properties governed by the associated operator symbols; see, for example, the discussion of uncertainty principles and phase-space concentration in \cite{fefferman_uncertainty_1983}. This observation motivates characterizing the metric entropy and minimax risk of operator images as a natural and effective way to study the complexity of sets of phase-space-localized signals. Such a viewpoint was adopted by the authors in \cite[Section~3.1]{ab2025compactoperators}, where the metric entropy of signal classes subject to time–frequency concentration constraints is analyzed via the ellipsoidal structure of the image of the unit ball under the Landau–Pollak–Slepian operator. This approach yields improvements over the previously best-known entropy characterizations for both the Landau–Pollak–Slepian operator and Sobolev balls. 
%The same viewpoint was adopted by the authors in \cite[Section~3.1]{ab2025compactoperators}, where the metric entropy of signal classes subject to time–%frequency concentration constraints is analyzed via the ellipsoidal structure of the image of the unit ball under the Landau-Slepian-Pollak operator, %yielding improvements over the previously best-known entropy characterizations for the Landau–Pollak–Slepian operator and for Sobolev balls. 

The central contribution of the present paper is to expand this localization-based methodology by dispensing with explicit ellipsoidal structure and developing a symbol-based framework that applies to broad classes of compact pseudodifferential operators, thereby enabling a unified analysis of % metric 
entropy and minimax risk under general phase-space localization conditions encoded by symbol decay and regularity.

\iffalse
Images of pseudodifferential operators are often well localized in phase space, with localization properties governed by their symbols; see, for example, the discussion of uncertainty principles and phase-space concentration in \cite{fefferman_uncertainty_1983}. This observation makes the characterization of the metric entropy of operator images a natural and effective tool for studying the complexity of phase-space localized signals. 
This vantage point has already been adopted by the authors of the present paper in \cite[Section~3.1]{ab2025compactoperators}, where the metric entropy of signal classes subject to time–frequency concentration constraints is analyzed via the ellipsoidal structure of the image of the unit ball under the associated operator, yielding refinements of the previously best-known entropy characterizations for the Landau–Pollak–Slepian operator and for Sobolev balls. The present work extends this analysis to more general forms of phase-space localization by treating broader families of pseudodifferential operators.
\fi


\iffalse 
The characterization of the metric entropy of operator images through pseudodifferential operators is particularly effective for studying signals that are well localized in phase space. Such signals can often be realized as elements of the image of a pseudodifferential operator whose symbol encodes the relevant localization properties. In the extreme case of signals that are strictly time- and band-limited, the metric entropy of the corresponding operator image has already been analyzed in \cite[Section 3.1]{ab2025compactoperators}. The present work extends this analysis to more general forms of phase-space localization by treating broader families of pseudodifferential operators.

This viewpoint becomes particularly relevant when studying signals that are well-localized in the phase space.
In the extreme case of signals that are strictly time- and band-limited, 
the analysis of its metric entropy has already been carried out in \cite[Section~3.1]{ab2025compactoperators}.
The present work can be seen as an extension of this result to more general types of localization in the phase space.
\fi
%For compact pseudodifferential operators whose symbols exhibit sufficient decay and regularity, 
For compact pseudodifferential operators with sufficiently regular symbols, phase-space localization at the symbol level is reflected in the spectral 
properties % distribution -> 'spectral distribution' has a different meaning
of the operator; see, e.g., \cite{fefferman_uncertainty_1983} for the underlying microlocal principle.
%Phase-space localization at the symbol level is reflected in the spectral distribution of the operator \cite{fefferman_uncertainty_1983}. 
A spectral formulation is therefore natural.  In the analysis of % metric
entropy and minimax risk for compact linear operators, it suffices to restrict attention to operators that are positive and self-adjoint. This follows from the fact that any compact operator admits a polar decomposition, and that metric entropy and minimax risk depend only on the geometry of the image of the unit ball, which is preserved under the isometry induced by the polar decomposition 
%restricting attention to positive self-adjoint operators entails no loss of generality 
(see, e.g., \cite{ab2025compactoperators}, \cite{prosserEEntropyECapacityCertain1966},  or \cite[Chapter~3.4]{carlEntropyCompactnessApproximation1990}).
%This restriction entails no loss of generality, as follows from standard polar decomposition arguments (see, e.g., \cite{ab2025compactoperators}, %\cite{prosserEEntropyECapacityCertain1966},  or \cite[Chapter~3.4]{carlEntropyCompactnessApproximation1990}).
Positive self-adjoint compact operators $T$ admit a spectral decomposition; the corresponding nonzero eigenvalues $\{\lambda_n\}_{n\in\N^*}$, listed in non-increasing order and counted with % according to
multiplicity, are nonnegative and tend to zero.
The spectral properties of $T$ are typically characterized through the eigenvalue-counting function
\begin{equation}\label{eq:definition-ecf}
    M_T(\lambda) \coloneqq  \# \left\{n\in\N^* \mid \lambda_n \geq \lambda \right\},
    \quad \text{for all } \lambda>0.
\end{equation}
It is well established that the entropy of a positive self-adjoint compact linear operator is governed by its spectral properties;
see, e.g., \cite{ab2025compactoperators,prosserEEntropyECapacityCertain1966,carlInequalitiesEigenvaluesEntropy1980,carl1981entropy,konigEigenvalueDistributionCompact1986,edmundsFunctionSpacesEntropy1996}.
Expressing the operator in its eigenbasis turns these spectral properties into a geometric description of the image of the unit ball under the operator. In particular, this image is an ellipsoid whose semi-axes are given by the eigenvalues $\{\lambda_n\}_{n\in\N^*}$, see \cite{ab2025compactoperators}.
As shown in \cite[Theorems~2, 4, and 5]{allard2025metricentropyminimaxrisk}, this ellipsoidal structure yields 
sharp asymptotic characterizations for both the % metric 
entropy $H_T$ and the minimax risk $R_T$ of $T$ via the type-$\tau$ (for $\tau \geq 1$) integrals \begin{equation}\label{eq:type-tau-integrals}
    I_\tau (\varepsilon) \coloneqq \int_{\varepsilon}^{\infty} \frac{M_T(\lambda)}{\lambda^\tau}\diff \lambda, \quad 
    \varepsilon > 0.
\end{equation}
%$I_\tau$ (for $\tau\geq 1$)
%to be characterized by the type-$\tau$ integrals $I_\tau$ (for $\tau\geq 1$) which are defined according to
%\begin{equation}\label{eq:type-tau-integrals}
%    I_\tau (\varepsilon) \coloneqq \int_{\varepsilon}^{\infty} \frac{M_T(\lambda)}{\lambda^\tau}\diff \lambda,
%    \quad \text{for all } \varepsilon>0.
%\end{equation}
Specifically, the minimax risk $R_T$ %for the ellipsoid with semi-axes $\{\lambda_n\}_{\N^*}$ 
is asymptotically determined by the type-$2$ and type-$3$ integrals according to
\begin{equation}\label{eq:integral-ecf-risk}
    R_T (\kappa)
    \sim \kappa^2 \varepsilon_\kappa I_2(\varepsilon_\kappa),
    \quad \text{as } \kappa \rightarrow 0,
\end{equation}
where the critical radius $\varepsilon_\kappa$ is defined as the unique solution of
\begin{equation}\label{eq:def-critical-radius}
    \kappa^2 \left(2I_3(\varepsilon_\kappa) - \frac{I_2(\varepsilon_\kappa)}{\varepsilon_\kappa}\right)=1,
    \quad \text{for all } \kappa > 0.
\end{equation}
Moreover, under a mild regularity condition on the eigenvalue-counting function $M_T$,
the % metric 
entropy $H_T$ is asymptotically equivalent to the type-$1$ integral, that is, 
\begin{equation}\label{eq:integral-ecf-entropy}
    H_T(\varepsilon) 
    \sim  I_1(\varepsilon),
    \quad \text{as } \varepsilon \rightarrow 0.
\end{equation}


The relations \eqref{eq:integral-ecf-risk}--\eqref{eq:integral-ecf-entropy} reduce the characterization of the asymptotic behavior of the % metric 
entropy and minimax risk of a compact operator to that of its asymptotic spectral distribution. 
This reduction has already proved effective in \cite[Section~3.2]{ab2025compactoperators} and \cite[Section~5]{allard2025metricentropyminimaxrisk}, where sharp asymptotic characterizations of metric entropy and minimax risk for unit balls in Sobolev spaces on bounded domains were derived.
In that setting, the operator $T$ is related to the inverse of the Laplacian, and the required asymptotics of the eigenvalue-counting function follow from the Weyl law (see, e.g., \cite[Chapter~9.5]{shubinInvitationPartialDifferential2020}) and its Riesz-means counterpart (see \cite{frank2024riesz}).
The goal of the present paper is to extend this spectral-reduction strategy to operator classes for which the asymptotic behavior of the eigenvalue-counting function can be derived from symbol-level information, notably compact pseudodifferential operators.

An important consequence of our results is a conceptual extension of Pinsker’s
theorem \cite{pinsker1980optimal} to unbounded domains. Pinsker’s original result
and its subsequent extensions are formulated in settings where compactness
arises from working on bounded domains, leading to compact embeddings of Sobolev
balls into $L^2$. We show that after restoring compactness on $\R^d$ through
spatial confinement, the same minimax principle continues to govern statistical
estimation. In this setting, the sharp asymptotic constants are determined by
phase-space geometry rather than by domain geometry, indicating that Pinsker’s
theorem reflects a structural feature of nonparametric estimation rather than an
artifact of bounded domains.

Pseudodifferential operators are linear operators acting on the Schwartz space
\[
\mathcal{S}(\R^d)
=
\left\{ f \in C^\infty(\R^d) :
\sup_{x\in\R^d} \left| x^\alpha \partial^\beta f(x) \right| < \infty
\text{ for all } \alpha, \beta \in \N^d \right\},
\]
and are of the form
\begin{equation}\label{eq:kohn-nirenberg-def}
    T_{\symbolnot} f (x)
    =
    \int_{\R^d}
    \symbolnot(x, \omega)\, \hat f (\omega)\, e^{2i\pi x \cdot \omega}\, d\omega,
\end{equation}
where $\symbolnot \in C^\infty(\R^d\times\R^d)$ is the associated symbol and
$x \cdot \omega$ denotes the Euclidean inner product in $\R^d$.
%Pseudodifferential operators %---sometimes also referred to as time-frequency localization operators---
%are linear operators, 
%\textcolor{blue}{acting on the Schwartz space $\mathcal{S}(\R^d)$,}
%of the form
%\begin{equation}\label{eq:kohn-nirenberg-def}
%    T_{\symbolnot} f (x) = \int_{\R^d} \symbolnot(x, \omega) \hat f (\omega) \, e^{2i\pi x \cdot \omega} \diff \omega,
%\end{equation}
%where $\symbolnot\in C^\infty(\R^d\times\R^d)$ is the symbol associated with $T_{\symbolnot}$ and $x \cdot \omega$ denotes the Euclidean inner product in %$\R^d$. 
Under appropriate regularity assumptions on $\symbolnot$, such operators 
can be extended compactly to % act boundedly on 
$L^2(\R^d)$ and admit a well-developed spectral theory. Pseudodifferential operators were originally introduced in mathematical physics in connection with quantization; see, e.g.,
%as a tool for addressing the quantization problem, see, e.g., 
\cite{weyl1950theory,wong1998weyl}.
They subsequently became central tools in harmonic analysis (see \cite{steinharmonic} and \cite[Chapter~2]{follandharmonic}), time-frequency analysis (see \cite[Chapter 14]{grochenig_foundations_2001}), 
electrical engineering in the context of linear-time-varying systems \cite{zadeh1950determination,zadeh1950frequency}, 
and wavelet theory (\cite[Chapter 2.10]{meyerWaveletsOperators1993}). 
In pure mathematics, pseudodifferential operators play a fundamental role in partial differential equations and microlocal analysis \cite{Taylor+1981,shubin_pseudodifferential_2001,hormander2007analysis,hormanderAnalysisLinearPartial2009,treves2013psido,hinzmicrolocal}, and they form a key ingredient in the proof of the Atiyah-Singer index theorem (see, e.g., \cite{landweber2005k}).
%However, as the name suggests, it is undoubtedly within partial differential equation (PDE) theory that pseudodifferential operators have had the most %profound impact. 
% For instance, every elliptic differential operator has a parametrix---an inverse modulo small terms---which is a pseudodifferential operator.
% Consequently, pseudodifferential operators permit to solve, at least formally, elliptic linear differential equations.
%We refer to \cite{Taylor+1981,shubin_pseudodifferential_2001,hormander2007analysis,hormanderAnalysisLinearPartial2009,treves2013psido,hinzmicrolocal} for %comprehensive treatments of pseudodifferential operators in PDE theory and microlocal analysis.

The mapping \eqref{eq:kohn-nirenberg-def} that associates a linear operator to a
symbol is known as the Kohn--Nirenberg quantization.
In addition to \eqref{eq:kohn-nirenberg-def}, two other classical quantization
schemes are commonly used, namely the Weyl quantization and the right
quantization; these are reviewed in Appendix~\ref{sec:review-lit}.
An important aspect of our analysis is that the results do not depend on the
particular choice of quantization; this is formalized in
Theorem~\ref{thm:main-result}.
Accordingly, we often write $H_\sigma$, $M_\sigma$, and $R_\sigma$ for the % metric
entropy, eigenvalue-counting function, and minimax risk of $T_\sigma$, omitting
explicit reference to the quantization when it is irrelevant.

\iffalse 
\textcolor{blue}{
    The process \eqref{eq:kohn-nirenberg-def} that associates a linear operator to a symbol is called the Kohn-Nirenberg quantization.
    Besides \eqref{eq:kohn-nirenberg-def}, there are two other classical ways to associate a symbol to an operator; namely, the Weyl quantization and the right quantization.
    We review them in Appendix~\ref{sec:review-lit}.
    Importantly, all the results presented in the present work are independent of the choice of quantization; this fact is formulated precisely in Theorem~\ref{thm:main-result}.
    For this reason, we write $H_\symbolnot$, $M_\symbolnot$, and $R_\symbolnot$ for the metric entropy $H_{T_{\symbolnot}}$, the eigenvalue-counting function $M_{T_{\symbolnot}}$, and the minimax risk $R_{T_{\symbolnot}}$ of the operator $T_{\symbolnot}$, respectively.
}
\fi 


%The function $\symbolnot$ appearing in \eqref{eq:kohn-nirenberg-def} is commonly referred to as the symbol of the pseudodifferential operator %$T_{\symbolnot}$,
%and the operator $T_{\symbolnot}$ in \eqref{eq:kohn-nirenberg-def} is traditionally called the Kohn-Nirenberg representation of $\symbolnot$.
Many properties of pseudodifferential operators relevant to spectral asymptotics are naturally expressed at the level of the symbol. In particular, for symbols that are positive and decay at infinity, 
a central quantity in our analysis is the phase-space volume above level 
%Many properties of pseudodifferential operators are naturally expressed in terms of their symbols.
%A central quantity in our analysis is the phase-space volume above level 
$\lambda>0$, defined by
%For instance, the volume above level $\lambda>0$ in the phase space occupied by the operator, a recurrent quantity in our analysis, is defined in terms %of $\symbolnot$ according to 
\begin{equation}\label{eq:definition-volume-function}
    V_{\symbolnot}(\lambda) 
    \coloneqq \int_{\mathbb{R}^{2d}} \mathbbm{1}_{\{\symbolnot(x, \omega)\, > \, \lambda\}} \diff  x\diff \omega.
\end{equation}
The function $V_{\symbolnot}$ records the phase-space distribution of the symbol.
%and, under the decay conditions specified in the main results below,
%provides a quantitative description of this decay.
% We write $H_\symbolnot$, $M_\symbolnot$, and $R_\symbolnot$ for the metric entropy $H_{T_{\symbolnot}}$, the eigenvalue-counting function $M_{T_{\symbolnot}}$, and the minimax risk $R_{T_{\symbolnot}}$ of the operator $T_{\symbolnot}$, respectively.
Under the regularity conditions specified in the main results below, the behavior of $V_{\symbolnot}(\lambda)$ as $\lambda \rightarrow 0$ quantifies the decay of $\sigma$. 
% As established in Theorem~\ref{thm:main-result}, this 
As recalled in Theorem~\ref{thm:Dauge-Robert}, this gives rise to a Weyl-type asymptotic for the spectral distribution of $T_{\symbolnot}$, namely % is governed by the phase-space volume $V_{\symbolnot}$ in the sense of %captures the asymptotic spectral properties of $T_{\symbolnot}$ according %to
\begin{equation}\label{eq:volume-ecf-relation}
    M_{\symbolnot}(\lambda)  \sim V_{\symbolnot}(\lambda), \quad 
    \lambda \rightarrow 0. 
\end{equation}
Relation \eqref{eq:volume-ecf-relation} serves as the link between symbol-level quantities and the general % metric 
entropy and minimax-risk characterizations in 
%, the decay properties of the symbol $\symbolnot$ determine the asymptotic spectral properties of  $T_{\symbolnot}$ and thus, by the previous discussion, %the asymptotic properties of its entropy and minimax risk.
%Concretely, when plugged in 
\eqref{eq:integral-ecf-risk}--\eqref{eq:integral-ecf-entropy}. Combining these results leads to the main conclusions of the paper,
%\eqref{eq:volume-ecf-relation} in \eqref{eq:integral-ecf-risk}--\eqref{eq:integral-ecf-entropy}, yields the main results of the paper, 
stated in Theorem~\ref{thm:main-result} and Corollaries~\ref{cor:ME-Main}--\ref{cor:MR-Main}. In particular, the % metric 
entropy satisfies
%the asymptotic relation \eqref{eq:volume-ecf-relation} yields the main contributions of the paper, that is, 
%the asymptotic formula for the metric entropy of $T_{\symbolnot}$
 \begin{equation}\label{eq:main-res-entropy}
    H_{\symbolnot}(\varepsilon) 
    \sim \int_{\R^{2d}} \ln_{+}\left(\frac{\symbolnot(x, \omega)}{\varepsilon}\right) \diff  x \diff \omega, 
    \quad \varepsilon \to 0,
\end{equation}
and the minimax risk has the asymptotic behavior 
\begin{equation}\label{eq:main-res-risk-1}
    R_\symbolnot (\kappa)
    \sim \kappa^2 \int_{\R^{2d}} \left(1-\frac{\varepsilon_\kappa}{\symbolnot(x, \omega)}\right)_+ \diff  x \diff \omega, 
    \quad  \kappa \to 0,
\end{equation}
where the critical radius $\varepsilon_\kappa$ is determined by 
\begin{equation}\label{eq:main-res-risk-2}
    \kappa^2 
    \int_{\R^{2d}} \frac{1}{\symbolnot(x, \omega) \, \varepsilon_\kappa }\left(1-\frac{\varepsilon_\kappa}{\symbolnot(x, \omega)}\right)_+ \diff  x \diff \omega =1, 
    \quad  \kappa >0.
\end{equation}
%The detailed derivations of \eqref{eq:main-res-entropy}--\eqref{eq:main-res-risk-2} are provided in Corollaries~\ref{cor:ME-Main} and \ref{cor:MR-Main} %respectively.















% Variations of \eqref{eq:volume-ecf-relation} for elliptic, non-compact operators are well known in the literature; see, e.g., \cite[Chapter~29]{hormanderAnalysisLinearPartial2009} or \cite[Corollary~16.1 and Theorem~30.1]{shubin_pseudodifferential_2001}. These results describe the asymptotic behavior of large eigenvalues of elliptic operators whose symbols grow at infinity. The Weyl-type relation \eqref{eq:volume-ecf-relation} established here differs from these classical results in several respects. First, it applies to compact pseudodifferential operators and, under the regularity conditions specified in the main results below, describes the asymptotic behavior of small eigenvalues. Second, the relation is formulated in terms of the superlevel-set volume $V_\sigma(\lambda)$, which is tailored to the compact setting and does not arise in the standard elliptic theory. Third---and most importantly for the present work---relation \eqref{eq:volume-ecf-relation} expresses the eigenvalue-counting function $M_\sigma(\lambda)$ explicitly as a phase-space volume. 
% This identification allows the general % metric 
% entropy and minimax risk formulae in \eqref{eq:integral-ecf-risk}--\eqref{eq:integral-ecf-entropy} to be reduced to phase-space integrals.




%Third---and most importantly for the present work---relation (8) is derived in a form that is directly compatible with entropy and minimax-risk analysis, %allowing it to be combined with the general characterizations \eqref{eq:integral-ecf-risk}--\eqref{eq:integral-ecf-entropy}.

% \color{red}
To put \eqref{eq:volume-ecf-relation} into context, we review the classical
Weyl-type theory for elliptic operators in Appendix~\ref{sec:review-lit}.
In Section~\ref{sec:main-results}, we adapt these ideas to the compact
symbol-decay setting considered here and use them to establish the
entropy–volume relation of Theorem~\ref{thm:main-result}. Section~\ref{sec:main-results}
% \color{black}
also contains the principal conceptual contributions of the paper, namely the asymptotic formulae \eqref{eq:main-res-entropy}–\eqref{eq:main-res-risk-2} together with their interpretation. Finally, Section~\ref{sec:appli-sobolev-spaces} applies these results to Sobolev spaces on unbounded domains, complementing the corresponding results for bounded domains obtained in \cite[Section~3.2]{ab2025compactoperators} and establishing Pinsker’s theorem in the unbounded setting.


\iffalse 
We review such standard results in Appendix~\ref{sec:review-lit}, 
and show how to transfer them to our setting in Section~\ref{sec:main-results}, 
thereby establishing an integrated version of \eqref{eq:volume-ecf-relation} in Theorem~\ref{thm:main-result}---our main technical contribution.
Section~\ref{sec:main-results} also contains the main conceptual contributions, namely \eqref{eq:main-res-entropy}--\eqref{eq:main-res-risk-2} and their interpretation.
Finally, Section~\ref{sec:appli-sobolev-spaces} applies the findings of Section~\ref{sec:main-results} to the study of Sobolev spaces on unbounded domain, 
thereby complementing the analogous result for Sobolev spaces on bounded domains obtained in \cite[Section~3.2]{ab2025compactoperators} and establishing Pinsker's theorem for unbounded domains.
\fi 



\paragraph{Notation.}
We write $\N$ for the set of natural numbers including zero,
$\N^*$ for the set of natural numbers excluding zero,
$\R$ for the real numbers, and $\R^*_+$ for the positive real numbers.
For $d \in \N^*$, we denote by $\omega_d$ the volume of the unit ball in $\R^d$,
and by $\|z\|_2$ the Euclidean norm of $z \in \R^d$.

$L^2(\R^d)$ and $C^\infty(\R^d)$ designate, respectively, the Lebesgue space of square-integrable functions and the space of infinitely differentiable functions on $\R^d$.
For $f \in C^\infty(\R^d)$, we write $\nabla f$ for its gradient and $\partial_i f$ for its partial derivative with respect to the $i$-th coordinate, where $i \in \{1,\dots,d\}$. A multi-index $\alpha = (\alpha_1,\dots,\alpha_d) \in \N^d$ defines the partial derivative
$\partial^\alpha = \partial_1^{\alpha_1} \cdots \partial_d^{\alpha_d}$,
with order $|\alpha| = \alpha_1 + \cdots + \alpha_d$.
The Fourier transform of a function $f \in L^2(\R^d)$ is denoted by $\hat f \in L^2(\R^d)$.

When comparing the asymptotic behavior of functions
$f,g \colon \R_+^* \to \R_+^*$ as $x \to 0$, we write
$f(x) = o_{x \to 0}(g(x))$ if $\lim_{x \to 0} f(x)/g(x) = 0$, and
$f(x) = O_{x \to 0}(g(x))$ if there exists a constant $C>0$ such that
$\limsup_{x \to 0} f(x)/g(x) \le C$.
Further, $f(x) \sim g(x)$ as $x \to 0$ if $\lim_{x \to 0} f(x)/g(x) = 1$.
% \textcolor{red}{
We write $f(x) \lesssim g(x)$ as $x \to 0$ if there exist constants
$C>0$ and $x_0>0$ such that $f(x) \le C\, g(x)$ for all $x \in (0,x_0)$,
and $f(x) \gtrsim g(x)$ as $x \to 0$ if $g(x) \lesssim f(x)$.
% }
Moreover, $f(x) \asymp g(x)$ as $x \to 0$
if there exist constants $0<c\le C<\infty$ and $x_0>0$ so that $
c\,g(x) \le f(x) \le C\,g(x), \, x \in (0,x_0)$.



Finally, $\ln(\cdot)$ denotes the natural logarithm,
$\ln_+(x) = \max\{0,\ln(x)\}$ for $x \ge 0$ (with the convention $\ln_+(0)=0$),
$(x)_+ = \max\{0,x\}$ for $x \in \R$,
and $\mathbbm{1}_X(\cdot)$ stands for the indicator function of the set $X$.



\iffalse 
\paragraph{Notation.}
We write $\N$ for the set of natural numbers including zero, 
$\N^*$ for the set of natural numbers excluding zero, 
$\R$ for the real numbers, and $\R^*_+$ for the positive real numbers.
For $d\in\N^*$, we denote by $\omega_d$ the volume of the unit ball in $\R^d$,
 and by $\|z\|_2$ the Euclidean norm of $z\in\R^d$.
$L^2(\R^d)$ and $C^\infty(\R^d)$ stand for both the Lebesgue space of square-integrable functions and, respectively, the space of infinitely differentiable functions, with domain $\R^d$.
We write $\nabla f $ for the gradient of the function $f\in C^\infty(\R^d)$ and $\partial_i f $ for its partial derivative with respect to the $i$-th variable,
with $i \in \{1, \dots, d\}$.
A multi-index $\alpha=(\alpha_1, \dots, \alpha_d)\in \N^d$ has the associated partial derivative composition $\partial^\alpha=\partial_1^{\alpha_1}\cdots\partial_d^{\alpha_d}$ and magnitude $|\alpha|=\alpha_1+\dots+\alpha_d$.
The Fourier transform of $f\in L^2(\R^d)$ is denoted by  $\hat f \in L^2(\R^d)$.
When comparing the asymptotic behavior of the functions $f, g\colon \R_+^*\to \R_+^*$ as $x \to 0$, we use $f(x) = o_{x \to 0}(g(x))$ to express that $\lim_{x \to 0} {f(x)}/{g(x)} =0$ and $f(x) = O_{x \to 0}(g(x))$ to indicate the existence of a constant $C>0$ such that $\lim_{x \to 0} {f(x)}/{g(x)} \leq C$.
We further write $f(x) \sim g(x)$, as $x\to 0$, for $\lim_{x \to 0} {f(x)}/{g(x)} =1$. 
Finally, $\ln(\cdot)$ is the natural logarithm, 
$\ln_+(x) = \max\{0,\ln(x)\}$ for all $x\geq 0$ (with the convention $\ln_+(0)=0$), $(x)_+ = \max\{0, x\}$ for all $x\in\R$,
and $\mathbbm{1}_{X}(\cdot)$ is the indicator function of the set $X$.
\fi 








\section{Entropy and Minimax Risk of Compact Hypoelliptic Pseudodifferential Operators}\label{sec:main-results}

We recall the definition of metric entropy for compact sets and the associated notion of entropy for compact linear operators.
Let $\mathcal{H}$ be a separable real Hilbert space and let $\mathcal{K} \subset \mathcal{H}$ be a compact set.
For $\varepsilon > 0$, an {$\varepsilon$-covering} of $\mathcal{K}$ is a finite set
$\{x_1,\dots,x_N\} \subset \mathcal{H}$ with the property that, for every
$x \in \mathcal{K}$, there exists an index $i \in \{1,\dots,N\}$ such that
\[
\|x - x_i\|_{\mathcal{H}} \le \varepsilon .
\]
The {$\varepsilon$-covering number} $N(\varepsilon;\mathcal{K},\|\cdot\|_{\mathcal{H}})$ is defined as the cardinality of a smallest such $\varepsilon$-covering; the {metric entropy} of $\mathcal{K}$ is the natural logarithm of the $\varepsilon$-covering number.

Given a compact linear operator $T\colon \mathcal{H} \to \mathcal{H}$, the {entropy} of $T$ is defined as the metric entropy of the closure of the image of the unit ball $\mathcal{B}_{\mathcal{H}} \subset \mathcal{H}$ under $T$, namely
\begin{equation}
\label{eq:def-entropy-operator}
    H_T(\varepsilon)
    \coloneqq
    \ln N \left(\varepsilon; \overline{T(\mathcal{B}_{\mathcal{H}})}, \|\cdot\|_{\mathcal{H}}\right).
\end{equation}
Since $T$ is compact, the set $\overline{T(\mathcal{B}_{\mathcal{H}})}$ is compact in $\mathcal{H}$, and the right-hand side of \eqref{eq:def-entropy-operator} is therefore finite for every $\varepsilon>0$.


% \iffalse
% Before proceeding, we recall the definition of metric entropy of a set and of the related notion of entropy of a compact linear operator.
% Let $\mathcal{H}$ be a separable real Hilbert space and let $\mathcal{K}\subset\mathcal{H}$ be a compact set.
% For $\varepsilon>0$, an {$\varepsilon$-covering} of $\mathcal{K}$ is a set $\{x_1,...\, ,x_N\} \subseteq \mathcal{H}$  such that,
% for each $x \in \mathcal{K}$, there exists an $i\in \{1, \dots, N\}$ so that $\|  x-x_i \|_{\mathcal{H}}\leq \varepsilon$. 
% The $\varepsilon$-covering number $N(\varepsilon ; \mathcal{K}, \| \cdot\|_{\mathcal{H}})$ is the cardinality of a smallest such $\varepsilon$-covering;
% the {metric entropy} of $\mathcal{K}$ is defined as the natural logarithm of the $\varepsilon$-covering number.
% The {entropy} of a compact linear operator $T\colon \mathcal{H}\to \mathcal{H}$ is defined as the metric entropy of the closure of the image of $\mathcal{B}_{\mathcal{H}}$, the unit ball in $\mathcal{H}$, under $T$, that is, 
% \begin{equation}\label{eq:def-entropy-operator}
%     H_T \left(\varepsilon \right)
%     \coloneqq \ln N \left(\varepsilon ; \overline{T(\mathcal{B}_{\mathcal{H}})}, \|\cdot\|_{\mathcal{H}} \right).
% \end{equation}
% The right-hand side of \eqref{eq:def-entropy-operator} is well-defined and finite as the set $\overline{T(\mathcal{B}_{\mathcal{H}})}$ is compact, 
% thanks to $T$ being compact.
% \fi 
The symbols treated in the classical literature on spectral asymptotics for
pseudo\-differential operators are predominantly drawn from the class
$\hypoposclass(\R^{2d})$ of positive-order hypoelliptic symbols;
see Theorem~\ref{thm-shubin} and the review of the standard spectral theory for
hypoelliptic pseudodifferential operators in
Appendix~\ref{sec:review-lit}. In this regime, the associated operators are necessarily
non-compact. Consequently, their entropy is not well defined and their minimax
risk does not converge to zero. 
%By contrast, operators whose symbols lie in 
% $\hyposymbclass(\R^{2d})$ and satisfy
% $m_- \le m_+ < 0$ 
%$\hyponegclass(\R^{2d})$ are compact; see, for example,
%\cite[Theorem~24.4]{shubin_pseudodifferential_2001}. 
By contrast, operators whose symbols lie in $\hyponegclass(\R^{2d})$
can be extended to compact operators on $L^2(\R^d)$; % are compact; 
see, for example, \cite[Theorem~24.4]{shubin_pseudodifferential_2001}.
% \color{red}
The results developed here build on spectral asymptotics for such compact
hypoelliptic operators, recalled in Theorem~\ref{thm:Dauge-Robert}.
%which extend classical Weyl-type results such as
%Theorem~\ref{thm-shubin}. 
%\textcolor{blue}{
%    The results developed here build on
%    spectral asymptotics for such compact hypoelliptic operators}---obtained in the spirit of Theorem~\ref{thm-shubin} and recalled in %Theorem~\ref{thm:Dauge-Robert}---to derive asymptotic characterizations of their entropy and minimax risk.
% The results developed here build on
% spectral asymptotics for such compact hypoelliptic operators—obtained in the spirit of
% Theorem~\ref{thm-shubin}—to derive asymptotic characterizations of their entropy and minimax risk.
\color{black}
% \iffalse 
% The symbols considered in the classical literature on spectral asymptotics for
% pseudo\-differential operators typically belong to the class
% $\mathrm{H}\Gamma^{m_-,m_+}_\rho(\R^{2d})$ of hypoelliptic symbols with parameters
% $m_+,m_->0$; see Theorem~\ref{thm-shubin} and the review of standard spectral theory for
% hypoelliptic pseudodifferential operators in Appendix~\ref{sec:review-lit}. In this regime, the
% operators associated with such symbols are necessarily non-compact. As a
% consequence, their entropy is not well defined and their minimax risk does not
% converge to zero. By contrast, operators whose symbols lie in the class $\hyposymbclass(\R^{2d})$ with
% $m_{-} \leq m_{+} < 0$ are compact; see, for example,
% \cite[Theorem~24.4]{shubin_pseudodifferential_2001}. The results developed here build on
% spectral asymptotics for such compact hypoelliptic operators—obtained in the spirit of
% Theorem~\ref{thm-shubin}—to derive asymptotic characterizations of their entropy.
% \fi 


% \iffalse 
% The symbols considered in the literature on spectral asymptotics for pseudodifferential operators typically belong to the class $\hyposymbclass(\R^{2d})$ of hypoelliptic operators with parameters $m_+,m_->0$; see Theorem~\ref{thm-shubin} and the review on  standard spectral theory for hypoelliptic operators in Appendix~\ref{sec:review-lit}.
% In particular, the corresponding operators can never be compact: their entropy is not well-defined and their minimax risk do not converge to zero.
% On the other hand, operators having their symbols in the class $\hyposymbclass(\R^{2d})$, with negative exponents $m_+,m_-<0$, are necessarily compact (see, e.g., \cite[Theorem~24.4]{shubin_pseudodifferential_2001}).
% Our main results then leverage spectral asymptotics of compact hypoelliptic operators, 
% obtained on the template of Theorem~\ref{thm-shubin}, to provide an asymptotic characterization of their entropy.
% \fi 

The asymptotic formulae derived below are obtained under the assumption that the
volume function $V_\sigma$ is regularly varying at zero with negative index.
This assumption excludes highly irregular decay behavior while encompassing
the polynomial decay rates that arise naturally in hypoelliptic
pseudodifferential calculus, as well as mild logarithmic perturbations thereof. Its role is to ensure sufficient
regularity of the associated eigenvalue-counting function to allow for a precise
evaluation of the leading-order asymptotics in the integral characterizations of entropy
and minimax risk.

More concretely, regular variation of the volume function $V_\sigma$ implies
condition~(RC) in \cite[Theorem~2]{allard2025metricentropyminimaxrisk}, %.
% Combined with the spectral asymptotics established above, this condition allows
allowing one to pass from spectral information to the entropy and minimax-risk asymptotics
in \eqref{eq:main-res-entropy}--\eqref{eq:main-res-risk-2}.
Background on regular variation and a detailed proof of the implication
\textcolor{blue}{`}$V_\sigma$ regularly varying $\Rightarrow$ condition~(RC)\textcolor{blue}{'} can be found in
\cite{binghamRegularVariation1987} and in
\cite[Appendix~C]{allard2025metricentropyminimaxrisk}.





% \iffalse
% To deduce asymptotics on the entropy,
% one actually needs to further impose regular variations on the volume $V_{\symbolnot}$.
% This ensures that the condition (RC) in \cite[Theorem~2]{allard2025metricentropyminimaxrisk} is satisfied.
% We refer to \cite{binghamRegularVariation1987} and \cite[Appendix~C]{allard2025metricentropyminimaxrisk} for more details on regular variations and how they ensure (RC) is satisfied.
% \fi 

\begin{theorem}\label{thm:main-result}
Let $\sigma\in\hyponegclass(\R^{2d})$ be a strictly positive hypoelliptic symbol, and
assume that the associated volume function $V_\sigma$ is regularly varying at
zero. % with negative index. 
Then,
\begin{equation}\label{eq:main-res-volume-integral}
    H_\sigma(\varepsilon)
    \sim
    \int_{\varepsilon}^{\infty}
        \frac{V_\sigma(\lambda)}{\lambda}\,\diff\lambda,
    \qquad \varepsilon\to0,
\end{equation}
where $H_\sigma$ denotes the entropy of the operator associated with $\sigma$,
independently of whether $\sigma$ is quantized using the left, right, or Weyl
quantization. Moreover, if the operator $T_{\sigma^{-1}}$ is invertible, then
\begin{equation}\label{eq:main-result-part2}
    H_{T^{-1}_{\sigma^{-1}}}(\varepsilon)
    \sim
    H_\sigma(\varepsilon),
    \qquad \varepsilon\to0.
\end{equation}
\end{theorem}

% \iffalse
% \begin{theorem}\label{thm:main-result}
% % Let $\rho\in(0,1]$ and let $m_+,m_->0$ satisfy $m_+\ge m_-$.
% % Let $\sigma\in\hyposymbclassinv(\R^{2d})$ be a strictly positive hypoelliptic symbol,
% Let $\sigma\in\hyponegclass(\R^{2d})$ be a strictly positive hypoelliptic symbol
% and assume that the associated volume function $V_\sigma$ is regularly varying at
% zero. Then
% \begin{equation}\label{eq:main-res-volume-integral}
%     H_\sigma(\varepsilon)
%     \sim
%     \int_{\varepsilon}^{\infty}
%         \frac{V_\sigma(\lambda)}{\lambda}\,\diff\lambda,
%     \qquad \varepsilon\to0,
% \end{equation}
% \textcolor{blue}{where $H_\sigma$ denotes the entropy of either the left, right, or Weyl quantization of $\sigma$.
% Additionally, if $T_{\sigma^{-1}}$ is invertible,
% then }
% \begin{equation}\label{eq:main-result-part2}
%     H_{T^{-1}_{\sigma^{-1}}}(\varepsilon)
%     \sim
%     H_\sigma(\varepsilon),
%     \qquad \varepsilon\to0.
% \end{equation}
% \end{theorem}
% \fi 



% \iffalse 
% \begin{theorem}\label{thm:main-result}
% Let $\rho \in (0,1]$, let $m_+,m_->0$ be such that $m_+ \ge m_-$, and 
% \textcolor{blue}{let
% $\sigma \in \hyposymbclassinv(\R^{2d})$ be a strictly positive hypoelliptic symbol such}
% % Assume that there exist constants $R>0$ and $C>0$ such that
% % \begin{equation}\label{eq:assumption-main-result}
% %     |z \cdot \nabla \sigma(z)| \ge C\,\sigma(z),
% %     \qquad \text{for all } z \in \R^{2d} \text{ with } \|z\|_2 \ge R.
% % \end{equation}
% % Assume in addition 
% that the associated volume function $V_\sigma$ is regularly
% varying at zero. Then
% \begin{equation}\label{eq:main-res-volume-integral}
%     H_\sigma(\varepsilon)
%     \sim \int_{\varepsilon}^\infty \frac{V_\sigma(\lambda)}{\lambda}\,\diff\lambda,
%     \qquad \varepsilon \to 0.
% \end{equation}
% \textcolor{blue}{Moreover, if the operator $T_{\symbolnot^{-1}}$ is invertible, 
% then we have $H_{T^{-1}_{\symbolnot^{-1}}}(\varepsilon) 
% \sim H_\sigma(\varepsilon)$, as $\varepsilon \to 0$.}
% \end{theorem}
% \fi 


% \begin{theorem}\label{thm:main-result}
% Let $\rho \in (0,1]$, let $m_+,m_->0$ be such that $m_+ \ge m_-$, and let
% $\sigma \in C^\infty(\R^{2d})$ be such that the associated volume function $V_\sigma$ is regularly
% varying at zero. 
% If at least one of the following assumptions is satisfied
% \begin{citemize}
%     \item [(A1)] $\sigma \in \hyposymbclassinv(\R^{2d})$  and $m_- = m_+$;
    
    

%     \item [(A2)] $\sigma \in \hyposymbclassinv(\R^{2d})$  is strictly positive and there exist constants $R>0$ and $C>0$ such that
%     \begin{equation}\label{eq:assumption-main-result}
%         |z \cdot \nabla \sigma(z)| \ge C\,\sigma(z),
%         \qquad \text{for all } z \in \R^{2d} \text{ with } \|z\|_2 \ge R;
%     \end{equation}
    
%     \item [(A3)] $\sigma$ is factorially hypoelliptic in the sense of \eqref{eq:fact-hypo1}--\eqref{eq:fact-hypo2}; 
% \end{citemize}
% then, 
% \begin{equation}\label{eq:main-res-volume-integral}
%     H_\sigma(\varepsilon)
%     \sim \int_{\varepsilon}^\infty \frac{V_\sigma(\lambda)}{\lambda}\,\diff\lambda,
%     \qquad \varepsilon \to 0.
% \end{equation}
% \end{theorem}

\begin{proof}
See Section~\ref{sec:proof-main-result}.
\end{proof}





\noindent
The proof of Theorem~\ref{thm:main-result} is carried out within the standard
hypoelliptic pseudodifferential operator framework summarized in
Appendix~\ref{sec:review-lit}, following \cite{shubin_pseudodifferential_2001}.
The strict positivity assumption on the symbol $\sigma$ in
Theorem~\ref{thm:main-result} is not essential, but it simplifies the proof by
ensuring that the symbol is invertible. A strategy for removing this assumption
proceeds as follows. By hypoellipticity, the symbol $\sigma$ has a fixed sign
outside a sufficiently large ball in phase space; without loss of generality,
this sign may be taken to be positive. One may therefore modify $\sigma$ inside
this ball so as to make it strictly positive everywhere, without affecting the
asymptotic behavior of the associated volume function $V_\sigma$, and hence
without altering the resulting entropy asymptotics. Consequently, the additional (strict) positivity assumption in
Theorem~\ref{thm:main-result} entails no loss of generality for the problems
considered here. Moreover, all symbols arising in the applications discussed in
Section~\ref{sec:appli-sobolev-spaces} are strictly positive. 

% \iffalse 
% The strict positivity assumption on the symbol $\symbolnot$ in Theorem~\ref{thm:main-result} is superfluous, but simplifies the proof as it ensures that the symbol is invertible.
% A proof strategy to remove this assumption would go as follows.
% First, note that, by hypoellipticity, we are guaranteed that $\symbolnot$ has fixed sign outside a ball of radius $R$ in the phase space, which, without loss of generality, can be assumed to be positive.
% Therefore, one could modify the symbol $\symbolnot$ within that ball (to make it positive) without modifying the asymptotics of the corresponding volume function and, incidentally, of the corresponding entropy.
% However, all the symbols considered in the applications presented in Section~\ref{sec:appli-sobolev-spaces} are positive, implying that the extra positivity assumption is not problematic for the present work.
% \fi 


% \iffalse
% It follows from its proof that Theorem~\ref{thm:main-result} remains valid if
% $\sigma$ is taken to be a left or right symbol, rather than a Weyl symbol.
% Indeed, by \cite[Theorem~23.3]{shubin_pseudodifferential_2001}, the Weyl symbol
% associated with an operator whose left symbol is
% $\sigma \in \hyposymbclassinv(\R^{2d})$ can be written in the form
% $\sigma(1+r)$, where $r \in \Gamma^{-2\rho}_\rho(\R^{2d})$.
% Equivalently, the corresponding operators satisfy
% $T_L = T_W(1+R)$, with $R \in \Psi^{-2\rho}_\rho$, where $T_L$ and $T_W$
% denote the operators with left and Weyl symbols given by $\sigma$, respectively.

% As is clear from Steps~3–5 in the proof of Theorem~\ref{thm:main-result}, the
% presence of the factor $1+R$ does not affect the leading-order behavior of the
% entropy. In particular, the entropy of $T_L$ has the same asymptotic scaling as
% that of $T_W$. Consequently, the asymptotic behavior of $H_{T_L}(\varepsilon)$
% can be deduced directly from that of $H_{T_W}(\varepsilon)$ given by
% Theorem~\ref{thm:main-result}. An analogous argument applies to operators whose
% symbol $\sigma$ is taken in the right quantization.
% \fi 



% % \begin{theorem}\label{thm:main-result-2}
% %     % Let $\rho\in(0,1]$ and let $m_+,m_->0$ satisfy $m_+\ge m_-$.
% %     Let $\sigma\in\hyponegclass(\R^{2d})$ be a strictly positive hypoelliptic symbol such that $T_{\sigma^{-1}}$ is invertible,
% %     and assume that the associated volume function $V_\sigma$ is regularly varying at
% %     zero. Then
% % \[
% %     H_{T^{-1}_{\sigma^{-1}}}(\varepsilon)
% %     \sim
% %     H_\sigma(\varepsilon),
% %     \qquad \varepsilon\to0.
% % \]
% % \end{theorem}







% % \begin{proof}
% % See Section~\ref{sec:proof-main-result-2}.
% % \end{proof}




% % \noindent
% % %%%T7: Please scrutinize this paragraph.
% % The entropy asymptotics described in
% % Theorems~\ref{thm:main-result} and~\ref{thm:main-result-2}
% % are robust with respect to the choice of quantization.
% % This robustness ultimately follows from an operator perturbation argument,
% % which is incarnated at the symbol level in the proof of
% % Theorem~\ref{thm:main-result} and at the operator level in the proof of
% % Theorem~\ref{thm:main-result-2}.
% % By \cite[Theorem~23.3]{shubin_pseudodifferential_2001}, if an operator has left
% % symbol $\sigma \in \hyposymbclassinv(\R^{2d})$ for some $\rho\in(0,1]$ and
% % $m_+\ge m_->0$, then the corresponding Weyl symbol $\sigma_W$ satisfies
% % \begin{equation}\label{eq:symbol-level}
% % \sigma_W - \sigma \in \Gamma^{-m_+ - 2\rho}_\rho(\R^{2d}).
% % \end{equation}
% % Equivalently, the operators associated with the left and Weyl quantizations are
% % related by a strictly lower-order perturbation,
% % \begin{equation}\label{eq:operator-level}
% % T_L = T_W + S,
% % \qquad S \in \Psi^{-m_+ - 2\rho}_\rho,
% % \end{equation}
% % where $T_L$ and $T_W$ denote the operators with left and Weyl symbols $\sigma$,
% % respectively.
% % In the proof of Theorem~\ref{thm:main-result}, such lower-order perturbations at the
% % symbol level \eqref{eq:symbol-level} do not affect the leading-order asymptotics of the associated
% % phase-space volume, and hence leave the entropy asymptotics unchanged.
% % In contrast, the proof of Theorem~\ref{thm:main-result-2} proceeds at the operator
% % level, using the decomposition~\eqref{eq:operator-level} together with
% % entropy-number estimates to show that additive compact operators of strictly lower order
% % do not contribute at leading order.


% \iffalse 
% The robustness of the entropy asymptotics with respect to the choice of
% quantization arises through different mechanisms in
% Theorems~\ref{thm:main-result} and~\ref{thm:main-result-2}.
% In the proof of Theorem~\ref{thm:main-result}, the entropy is characterized
% directly in terms of the phase-space volume function associated with the symbol.
% Lower-order modifications of the symbol, such as those induced by switching
% between Weyl, left, or right quantization, do not affect the leading-order
% asymptotics of this volume and therefore leave the entropy asymptotics unchanged.

% In contrast, the robustness established in
% Theorem~\ref{thm:main-result-2} is made explicit at the operator level.
% By \cite[Theorem~23.3]{shubin_pseudodifferential_2001}, if an operator has left
% symbol $\sigma \in \hyposymbclassinv(\R^{2d})$ for some $\rho\in(0,1]$ and
% $m_+\ge m_->0$, then the corresponding Weyl symbol $\sigma_W$ satisfies
% \[
% \sigma_W - \sigma \in \Gamma^{-m_+ - 2\rho}_\rho(\R^{2d}).
% \]
% As a consequence, the associated operators differ by a strictly lower-order
% perturbation,
% \[
% T_L = T_W + S,
% \qquad S \in \Psi^{-m_+ - 2\rho}_\rho,
% \]
% where $T_L$ and $T_W$ denote the operators with left and Weyl symbols $\sigma$,
% respectively. Since $S$ is compact and of strictly lower order, the
% entropy-number comparison argument used in the proof of
% Theorem~\ref{thm:main-result-2} shows that such perturbations contribute only
% lower-order corrections and do not affect the leading-order entropy asymptotics.
% An entirely analogous argument applies to the right quantization.


% The entropy asymptotics described in
% Theorems~\ref{thm:main-result} and~\ref{thm:main-result-2}
% are robust with respect to the choice of quantization.
% In particular, the leading-order entropy scaling remains unchanged if the symbol
% \(\sigma\) is taken in the left or right quantization rather than in the Weyl
% quantization.
% By \cite[Theorem~23.3]{shubin_pseudodifferential_2001}, if an operator has left
% symbol \(\sigma \in \hyposymbclassinv(\R^{2d})\) for some \(\rho\in(0,1]\) and
% \(m_+\ge m_->0\), then the corresponding Weyl symbol \(\sigma_W\) satisfies
% \[
% \sigma_W - \sigma \in \Gamma^{-m_+ - 2\rho}_\rho(\R^{2d}).
% \]
% As a consequence, the associated operators differ by a strictly lower-order
% perturbation,
% \[
% T_L = T_W + S,
% \qquad S \in \Psi^{-m_+ - 2\rho}_\rho,
% \]
% where \(T_L\) and \(T_W\) denote the operators with left and Weyl symbols
% \(\sigma\), respectively. In particular, \(S\) is compact and of strictly lower
% order than \(T_W\).

% The robustness of the entropy asymptotics established in
% Theorems~\ref{thm:main-result} and~\ref{thm:main-result-2} now stems from
% the fact that lower-order perturbations do not contribute at leading order.
% This mechanism appears in different forms in the two proofs.
% In Theorem~\ref{thm:main-result}, the entropy is characterized directly in terms
% of the volume function associated with the symbol, and lower-order modifications
% of the symbol do not affect the leading-order behavior of this volume.
% In Theorem~\ref{thm:main-result-2}, robustness is made explicit through the
% entropy-number comparison arising from the operator decompositions
% \(T_* = T_\sigma(I+R_1)+R'_{-\infty}\) and \(T_\sigma = T_*(I+R_0)\), where the
% operators \(R_0\), \(R_1\), and \(R'_{-\infty}\) are compact and of strictly lower
% order.
% \fi 




% \iffalse 
% The entropy asymptotics described in 
% Theorems~\ref{thm:main-result} and \ref{thm:main-result-2}
% are robust with
% respect to the choice of quantization. In particular, the leading-order entropy
% scaling remains unchanged if $\sigma$ is taken in the left or right quantization
% rather than in the Weyl quantization. By \cite[Theorem~23.3]{shubin_pseudodifferential_2001},
% if an operator has left symbol $\sigma \in \hyposymbclassinv(\R^{2d})$, 
% for some $\rho\in(0,1]$ and $m_+\geq m_->0$,
% then the
% corresponding Weyl symbol $\sigma_W$ satisfies
% \[
% \sigma_W - \sigma \in \Gamma^{-m_+ - 2\rho}_\rho(\R^{2d}).
% \]
% Accordingly, the corresponding operators differ by a lower-order perturbation,
% in the sense that
% \[
% T_L = T_W + S,
% \qquad S \in \Psi^{-m_+ - 2\rho}_\rho,
% \]
% where $T_L$ and $T_W$ denote the operators with left and Weyl symbols given by
% $\sigma$, respectively. In particular, $S$ is a compact operator of strictly lower order.


% As can be deduced from the arguments in 
% % Steps~3--5 
% \textcolor{blue}{the proof of
% Theorem~\ref{thm:main-result-2}}, perturbations of strictly lower order do not affect
% the leading-order asymptotic scaling of the entropy. It follows that \(H_{T_L}(\varepsilon)\) exhibits the same asymptotic scaling as
% \(H_{T_W}(\varepsilon)\) in the limit \(\varepsilon \to 0\), and its behavior can
% therefore be inferred from the Weyl-quantized setting treated in
% Theorem~\ref{thm:main-result}. An analogous argument applies when the symbol
% \(\sigma\) is quantized using the right quantization.
% \fi 


The following heuristic motivates the use of \eqref{eq:volume-ecf-relation} in the
derivation of the entropy and minimax-risk asymptotics
\eqref{eq:main-res-entropy} and \eqref{eq:main-res-risk-1}. Guided by the
uncertainty principle, one expects the eigenfunctions of $T_\sigma$ to be
essentially localized in phase space on regions of unit volume. Within such a
localized region centered at a point $(x_0,\omega_0)\in\R^{2d}$, the symbol
$\sigma$ may be regarded as approximately constant, with value
$\sigma(x_0,\omega_0)$.

Under this approximation, the number of eigenvalues exceeding a level
$\lambda>0$ is expected to be proportional to the number of disjoint unit-volume
regions of phase space centered at points $(x,\omega)$ for which
$\sigma(x,\omega)>\lambda$. This leads naturally to the volume-based relation
\eqref{eq:volume-ecf-relation}. A related heuristic for elliptic operators is
discussed by Fefferman in \cite{fefferman_uncertainty_1983}. Of course, the
argument above is purely heuristic; a central contribution of the present work is
to identify precise conditions under which the insertion of
\eqref{eq:volume-ecf-relation} into the entropy and minimax-risk integrals can be
rigorously justified in the compact, symbol-decay setting considered here.

% \iffalse 
% The following heuristic motivates the use of \eqref{eq:volume-ecf-relation} in the
% derivation of the entropy and minimax-risk asymptotics
% \eqref{eq:main-res-entropy} and \eqref{eq:main-res-risk-1}. Guided by the
% uncertainty principle, one expects the eigenfunctions of $T_\sigma$ to be
% essentially localized in phase space on regions of unit volume. Within such a
% localized region centered at a point $(x_0,\omega_0)\in\R^{2d}$, the symbol
% $\sigma$ may be regarded as approximately constant, with value
% $\sigma(x_0,\omega_0)$.

% Under this approximation, the number of eigenvalues exceeding a level
% $\lambda>0$ is expected to be proportional to the number of disjoint unit-volume
% regions of phase space centered at points $(x,\omega)$ for which
% $\sigma(x,\omega)>\lambda$. This leads naturally to the volume-based relation
% \eqref{eq:volume-ecf-relation}. A related heuristic for elliptic operators is
% discussed by Fefferman in \cite{fefferman_uncertainty_1983}. Of course, the
% argument above is purely heuristic; a central contribution of the present work is
% to identify precise conditions under which the insertion of
% \eqref{eq:volume-ecf-relation} into the entropy and minimax-risk integrals can be
% rigorously justified in the compact, symbol-decay setting considered here.


% The following heuristic motivates the use of \eqref{eq:volume-ecf-relation} in the
% derivation of the entropy and minimax-risk asymptotics \eqref{eq:main-res-entropy}
% and \eqref{eq:main-res-risk-1}. Guided by the uncertainty principle, one expects
% the eigenfunctions of $T_\sigma$ to be essentially localized in phase space on
% regions of unit volume. For concreteness, consider such a region centered at a
% point $(x_0,\omega_0)\in\R^{2d}$, and let $\psi_0$ denote a corresponding
% eigenfunction. Within this localized region, the symbol $\sigma$ may be regarded
% as approximately constant, with value $\sigma(x_0,\omega_0)$.
% Under this approximation, the number of eigenvalues exceeding a level
% $\lambda>0$ is expected to be proportional to the number of disjoint
% unit-volume regions of phase space centered at points $(x,\omega)$ for
% which $\sigma(x,\omega)>\lambda$. This leads naturally to the volume-based relation
% \eqref{eq:volume-ecf-relation}. A related heuristic for elliptic operators is discussed by Fefferman in \cite{fefferman_uncertainty_1983}.
% Of course, the argument above is purely heuristic; a central contribution of the
% present work is to identify precise conditions under which the insertion of
% \eqref{eq:volume-ecf-relation} into the entropy and minimax-risk integrals can be
% rigorously justified in the compact, symbol-decay setting considered here.
% \fi 

% \iffalse
% The heuristic behind the integrated version of \eqref{eq:volume-ecf-relation} we use in the proof of Theorem~\ref{thm:main-result} is the following.
% According to the uncertainty principle, 
% the eigenfunctions of $T_{\symbolnot}$ are localized within a portion of the phase space with unit volume. 
% For concreteness, assume that this portion of the phase space is centered around a point $(x_0, \omega_0)$, and let $\psi_0$ denote the corresponding eigenfunction.
% In this portion of the phase space, the symbol $\symbolnot$ is approximately constant equal to $\symbolnot(x_0, \omega_0)$.
% Therefore, the number of eigenvalues above $\lambda$ approximates the number of unit-volume portions of the phase space centered around points $(x_0, \omega_0)$ for which $\symbolnot(x_0, \omega_0)>\lambda$; 
% this is expressed by \eqref{eq:volume-ecf-relation}.
% A similar heuristic for elliptic operators is discussed by Fefferman in \cite{fefferman_uncertainty_1983}.
% Obviously, this heuristic argument does not serve as proof, and providing explicit conditions under which the integrated form of \eqref{eq:volume-ecf-relation} holds is a central contribution of the present work.
% \fi 

The volume principle underlying \eqref{eq:volume-ecf-relation} is reminiscent of the
approximate diagonalization results in underspread operator theory
\cite{kailath1963timevariant, bello1969measurement, pfander2006measurement,
durisi2008noncoherent, heckel2013identification, matz1998time, matz2013time}, in that
both frameworks interpret the Weyl symbol as an effective time--frequency transfer
function and relate spectral quantities to its pointwise values through a
phase-space localization principle.
%The volume principle underlying \eqref{eq:volume-ecf-relation} is reminiscent of the approximate diagonalization results in underspread operator theory %\cite{kailath1963timevariant, bello1969measurement, pfander2006measurement, durisi2008noncoherent, heckel2013identification, matz1998time, matz2013time}, %where the Weyl symbol acts as an effective time–frequency transfer function and spectral quantities can be related to pointwise symbol values. 
In that setting, one expects entropy formulae of the same log-integral form as \eqref{eq:main-res-entropy}, paralleling Shannon-type expressions for frequency-selective channels. Developing a rigorous analogue of Theorem~\ref{thm:main-result} for underspread operators is an interesting direction for future work.

% \iffalse 
% An alternative approach could be based on
% \emph{underspread operators}, in analogy with the theory developed in
% \cite{shubin_pseudodifferential_2001} for hypoelliptic operators.

% Underspread operators are characterized by the property that their spreading
% function—the Fourier transform of the Weyl symbol—has finite
% $\|\cdot\|_1$-norm; see \cite{matz1998time,kozek1996matched,Kozek1998,matz2013time}
% for a general treatment of this class. In particular, \cite[Theorem~III.6]{matz1998time}
% provides asymptotic characterizations of the eigenvalue-counting function in terms
% of the Weyl symbol for underspread operators whose eigenvalues tend to infinity.

% Using transfer arguments similar to those employed in the proof of
% Theorem~\ref{thm:main-result}, one may then expect to obtain an analogue of
% Theorem~\ref{thm:main-result} for compact operators arising as inverses of
% underspread operators. A rigorous development of this extension is beyond the
% scope of the present work and is left for future investigation.

% The technical requirements on hypoelliptic pseudodifferential operators for the proof of Theorem~\ref{thm:main-result} are provided in Appendix~\ref{sec:review-lit} and are based on the standard reference \cite{shubin_pseudodifferential_2001}.
% Note that one could develop a similar theory for underspread operators---parallel to the one presented in \cite{shubin_pseudodifferential_2001} for hypoelliptic operators.
% Underpsread operators are operators whose spreading function (the Fourier transform of the Weyl symbol) has finite $\|\cdot\|_1$-norm; 
% see \cite{matz1998time,kozek1996matched,Kozek1998,matz2013time} for a general theory of underspread operators.
% Specifically, based on \cite[Theorem~III.6]{matz1998time}, one can characterize the asymptotics of the eingenvalue-counting function in terms of the Weyl symbol for underpsread operators whose eigenvalues tend to infinity.
% Then,
% one gets a version of Theorem~\ref{thm:main-result} for compact operators that are inverses of underspread operators
%  using transfer techniques similar to the one developed in the proof of Theorem~\ref{thm:main-result}.
% Formalizing this discussion is left for future work.
% \fi 



We now present a reformulation of Theorem~\ref{thm:main-result} that makes the
connection between entropy and the symbol explicit. Again, the
symbol $\sigma$ in Corollary~\ref{cor:ME-Main} may be taken to be the left, right,
or Weyl symbol of the associated operator.

\begin{corollary}\label{cor:ME-Main}
% Under the assumptions of Theorem~\ref{thm:main-result}
Let $\sigma\in\hyponegclass(\R^{2d})$ be a strictly positive hypoelliptic symbol
and assume that the associated volume function $V_\sigma$ is regularly varying at
zero with negative index. Then, the entropy admits the
asymptotic representation
\begin{equation}\label{eq:main-res-entropy-corr}
    H_{\sigma}(\varepsilon)
    \sim \int_{\R^{2d}} \ln_{+}\!\left(\frac{\sigma(x,\omega)}{\varepsilon}\right)
    \diff x\,\diff\omega,
    \qquad \varepsilon \to 0.
\end{equation}
\end{corollary}

% \iffalse 
% We now turn to a reformulation of Theorem~\ref{thm:main-result} that links entropy and symbol.
% Once again, $\symbolnot$ in Corollary~\ref{cor:ME-Main} below can represent either the left, right, or Weyl symbol of the associated operator.


% \begin{corollary}\label{cor:ME-Main}
%     Under the assumptions of Theorem~\ref{thm:main-result}, one gets
%     \begin{equation}\label{eq:main-res-entropy-corr}
%         H_{\symbolnot}(\varepsilon) 
%         \sim \int_{\R^{2d}} \ln_{+}\left(\frac{\symbolnot(x, \omega)}{\varepsilon}\right) \diff  x \diff \omega, 
%         \quad \text{as } \varepsilon \to 0.
%     \end{equation}
% \end{corollary}
% \fi 



\begin{proof}
    See Section~\ref{sec:proof-ME-Main}.
\end{proof}


\noindent
Recall that, after normalization by a factor $\ln(2)$, the metric entropy of a
compact set admits an information-theoretic interpretation as the number of bits
required to uniformly encode the set with accuracy $\varepsilon$.
From this perspective, relation \eqref{eq:main-res-entropy-corr} suggests the
following interpretation.
The phase space $\R^{2d}$ may be partitioned into regions of unit volume, centered
at points $(x,\omega)$, each contributing to the total entropy the number of bits
required to quantize the local symbol value $\sigma(x,\omega)$.
At the heuristic level, this contribution is given by
$\lceil \log_{+}(\sigma(x,\omega)/\varepsilon)\rceil$, which is equivalent in the
limit $\varepsilon\to0$ to $\log_{+}(\sigma(x,\omega)/\varepsilon)$.
Corollary~\ref{cor:ME-Main} shows that this phase-space heuristic is in fact
correct and yields a sharp characterization of the leading-order term of the
entropy.

% \iffalse 
% Recall that, upon dividing it by a factor $\ln(2)$, one can interpret the metric entropy of a compact set as the number of bits required to uniformly encode the set to accuracy $\varepsilon$.
% The relation \eqref{eq:main-res-entropy-corr} thus allows the following interpretation.
% The phase space $\R^{2d}$ can be divided into regions of unit volume, centered around points $(x_0, \omega_0)$, 
% each contributing to the entropy the number of bits required to quantize $\symbolnot(x_0, \omega_0)$, 
% that is, $\lceil\log_{+}({\symbolnot(x_0, \omega_0)}/{\varepsilon})\rceil$.
% Indeed, recall that the metric entropy of the interval of length $A>0$ is $\log(\lceil A/(2\varepsilon))\rceil$ (see, e.g., \cite[Section~2]{expdecaypaper}), which is asymptotically equivalent to $\lceil\log_{+}({\symbolnot(x_0, \omega_0)}/{\varepsilon})\rceil$ when $A=\symbolnot(x_0, \omega_0)$.
% The point made by Corollary~\ref{cor:ME-Main} is that thus this heuristic is correct and provides a sharp characterization of the entropy's leading term.
% \fi 


Relation \eqref{eq:main-res-entropy-corr} admits a direct phase-space
interpretation in terms of activation and deactivation of degrees of freedom.
Specifically, only those phase-space cells for which the local symbol value
$\sigma(x,\omega)$ exceeds the threshold $\varepsilon$ contribute to the entropy,
while cells with $\sigma(x,\omega)\le\varepsilon$ are effectively inactive.
For each active cell, the contribution to the entropy is logarithmic in the
ratio $\sigma(x,\omega)/\varepsilon$. From this perspective, Corollary~\ref{cor:ME-Main} shows that the overall entropy is
obtained by aggregating the contributions of all active phase-space cells.

% \iffalse 
% We can also interpret \eqref{eq:main-res-entropy-corr} as an instance of a water-filling formula: 
% the way bits are allocated in the phase space mimics the way water distributes itself in a vessel.
% This gives Corollary~\ref{cor:ME-Main} a channel-capacity interpretation, where $\varepsilon$ plays the role of the signal-to-noise ratio;
% compare with the capacity formula for an additive Gaussian noise channel \cite[Chapter~9.5]{cover1999elements}.
% We refer to \cite[Chapters~9 and 10]{cover1999elements} for a general background on water-filling arguments and to \cite{donohoCountingBitsKolmogorov2000} for a discussion in the context of metric entropy.
% Note that, for a general compact operator $T$ with eigenvalues $\{\lambda_n\}_{\N^*}$, one has the following discrete water-filling representation
% \begin{equation}\label{eq:water-filling}
%     H_{T}(\varepsilon) 
%     \sim \sum_{n\in\N^*} \ln_+\left(\frac{\lambda_n}{\varepsilon}\right),
%     \quad \text{as } \varepsilon\to 0;
% \end{equation}
% see \cite[Theorem~2]{allard2025metricentropyminimaxrisk} and the subsequent discussion.
% Corollary~\ref{cor:ME-Main} can thus be interpreted as a continuous version of \eqref{eq:water-filling}.
% \fi 

Minimax risk has been studied extensively in nonparametric estimation for function
classes and compact subsets of Hilbert spaces, in particular through the analysis
of the spectral decay of associated compact operators, such as covariance
operators, embedding operators, or diagonal operators arising from basis
expansions; see, for example,
\cite{ibragimov1981statistical,pinsker1980optimal,tsybakovIntroductionNonparametricEstimation2009}.
In these settings, the function class is typically represented implicitly as the
image of a unit ball under a compact linear operator, and minimax risk is
characterized in terms of the eigenvalues or singular values of that operator.

In contrast, the present framework treats minimax risk as a functional of the
operator itself and aims to characterize its asymptotic behavior directly at the
symbol level. This operator-centric perspective makes it possible to link minimax
risk explicitly to phase-space geometry and spectral asymptotics, in parallel with
the corresponding entropy analysis developed above. 
%We therefore introduce the
%operator-level definition of minimax risk.

To make this operator-level perspective precise, let
$T\colon \mathcal{H}\to\mathcal{H}$ be a compact linear operator on a separable
Hilbert space $\mathcal{H}$. The minimax risk of $T$ is defined over the compact
set
\(
\mathcal{K}=\overline{T(\mathcal{B}_{\mathcal{H}})}\subset\mathcal{H}.
\)
Specifically, for a noise level $\kappa>0$, it is given by
\begin{equation*}
    % R_\kappa(T)
    R_T(\kappa)
    \coloneqq
    \inf_{\hat x_\kappa}
    \sup_{x\in \overline{T(\mathcal{B}_{\mathcal{H}})}}
    \mathbb{E}_{y\sim x}\!\left[
        \bigl\|\hat x_\kappa(y)-x\bigr\|_{\mathcal{H}}^{2}
    \right],
\end{equation*}
where the observation model is
\(
y = x + \kappa \xi,
\)
with $\xi$ a Gaussian random element in $\mathcal{H}$ whose coordinates with
respect to some (and hence any) orthonormal basis of $\mathcal{H}$ are i.i.d.\
standard normal random variables. Here, $\mathbb{E}_{y\sim x}$ is expectation with respect to the law of $y$ for fixed
$x$, and $\hat x_\kappa$ is a measurable estimator based on the observation $y$.

The asymptotic behavior of the minimax risk then follows from
Theorem~\ref{thm:main-result} via the general correspondence between metric entropy
and minimax risk developed in \cite{allard2025metricentropyminimaxrisk}.

% \iffalse 
% We now turn to the corresponding result for the minimax risk.
% In direct analogy with the entropy, the minimax risk of an operator $T$ is defined
% as the minimax risk associated with the compact set
% $\mathcal{K}=\overline{T(\mathcal{B}_{\mathcal{H}})}\subset\mathcal{H}$.
% Specifically, it is given by
% \begin{equation*}
%     R_\kappa(T)
%     \coloneqq
%     \inf_{\hat x_\kappa}
%     \sup_{x\in \overline{T(\mathcal{B}_{\mathcal{H}})}}
%     \mathbb{E}_{y\sim x}\!\left[
%         \bigl\|\hat x_\kappa(y)-x\bigr\|_{\mathcal{H}}^{2}
%     \right],
% \end{equation*}
% where the observation model is $y=x+\kappa\xi$, with $\xi$ a Gaussian random
% element in $\mathcal{H}$ with i.i.d.\ standard normal coordinates, and
% $\hat x_\kappa$ denotes an estimator mapping the noisy observation $y$ to an
% estimate of $x$.

% The asymptotic behavior of the minimax risk can be obtained from
% Theorem~\ref{thm:main-result} by invoking the general correspondence between
% metric entropy and minimax risk developed in
% \cite{allard2025metricentropyminimaxrisk}.



% We now state the corresponding result for minimax risk.
% Similarly to entropy, we define the minimax risk of an operator to be the corresponding quantity for the compact set $\mathcal{K} = \overline{T(\mathcal{B}_{\mathcal{H}})}$.
% Specifically, it is defined as 
% \begin{equation*}
%     R_\kappa(T) \coloneqq \inf_{\hat x_\kappa} \sup_{x\in \overline{T(\mathcal{B}_{\mathcal{H}})}} \mathbb{E}_{y\sim x}\left[\left\|\hat x_\kappa(y) - x\right\|^2_{\mathcal{H}}\right],
% \end{equation*}
% where the notation $y\sim x$ indicates that $y = x + \kappa \xi$, where the components of $\xi$ are i.i.d. standard Gaussian, 
% and  $\hat x_\kappa$ is an estimator mapping the noisy observation $y$ to the true parameter $x$.
% The main result about the minimax risk can be derived from Theorem~\ref{thm:main-result} using the correspondence between entropy and minimax risk developed in \cite{allard2025metricentropyminimaxrisk}.
% \fi 



\begin{corollary}\label{cor:MR-Main}
% Under the assumptions of Theorem~\ref{thm:main-result}
Let $\sigma\in\hyponegclass(\R^{2d})$ be a strictly positive hypoelliptic symbol
and assume that the associated volume function $V_\sigma$ is regularly varying at
zero with negative index. Then, the minimax risk satisfies
\begin{equation}\label{eq:main-res-risk-1-corr}
    R_{\sigma}(\kappa)
    \sim
    \kappa^{2}
    \int_{\R^{2d}}
        \left(1-\frac{\varepsilon_\kappa}{\sigma(x,\omega)}\right)_{+}
        \,\diff x\,\diff\omega,
    \qquad \kappa \to 0,
\end{equation}
where the critical radius $\varepsilon_\kappa$ is determined (implicitly) by
\begin{equation}\label{eq:main-res-risk-2-corr}
    \kappa^{2}
    \int_{\R^{2d}}
        \frac{1}{\sigma(x,\omega)\,\varepsilon_\kappa}
        \left(1-\frac{\varepsilon_\kappa}{\sigma(x,\omega)}\right)_{+}
        \,\diff x\,\diff\omega
    = 1,
    \qquad \kappa>0.
\end{equation}
\end{corollary}

\begin{proof}
See Section~\ref{sec:proof-MR-Main}.
\end{proof}






\section{Applications to Sobolev Spaces}\label{sec:appli-sobolev-spaces}

We now apply the results developed in the preceding sections to characterize the
metric entropy of unit balls in Sobolev spaces on $\R^d$. This extends the
corresponding analysis for Sobolev spaces on bounded domains presented in
\cite[Section~3.2]{ab2025compactoperators}. In contrast to the bounded-domain
setting, the unit ball of a classical Sobolev space on $\R^d$—as defined, for
example, in \cite[Definition~3.1]{Taylor+1981} or
\cite[Section~9.1]{brezisFunctionalAnalysisSobolev2011}—is not compact with
respect to the $L^2(\R^d)$ metric due to translation invariance
(cf.\ \cite[Example~6.11]{adamsSobolevSpaces2003}). As a result, its metric
entropy is infinite for every $\varepsilon>0$.

Compactness-restoring variants of Sobolev spaces on $\R^d$, such as those
involving confining potentials or weights, are classical and have been studied
extensively from the perspectives of functional analysis and spectral theory.
% What appears to be largely absent, however, are explicit characterizations of
% metric entropy for these spaces with respect to the $L^2$ metric.

% The operator-based framework developed above provides a natural way to address
% this gap. 
By modifying the Sobolev structure so as to restore compactness while
preserving the underlying notion of regularity, one obtains function classes
that fall within the scope of the present theory. Two such modifications are
particularly natural and widely used in practice: Sobolev spaces with a
confining potential (Theorem~\ref{thm:sob-sch}) and weighted Sobolev spaces
(Theorem~\ref{thm:sob-weighted}). In both cases, the resulting spaces admit
compact embeddings into $L^2(\R^d)$ and can be treated directly within our framework.


% \iffalse 
% We now apply the results developed in the previous sections to characterize the metric entropy of unit balls in Sobolev spaces on $\R^d$.
% This extends the corresponding result for Sobolev spaces on bounded domains provided in \cite[Section~3.2]{ab2025compactoperators}.
% The unit ball in  classical Sobolev spaces on $\R^d$ (as defined, e.g., in \cite[Definition~3.1]{Taylor+1981} or \cite[Section~9.1]{brezisFunctionalAnalysisSobolev2011}) is never compact under the $L^2(\R^d)$-metric (cf. \cite[Example~6.11]{adamsSobolevSpaces2003}), 
% so that its metric entropy is infinite for nonzero values of $\varepsilon$.
% Therefore, we consider two variants that are widely used in practice,  
% namely Sobolev spaces with a potential (Theorem~\ref{thm:sob-sch})
% and weighted Sobolev spaces (Theorem~\ref{thm:sob-weighted}).
% \fi 

We begin with Sobolev spaces equipped with a confining potential.
Fix $s>0$, let $u\colon\R^d\to\R_+^*$ be a potential function, and define
\begin{equation}\label{eq:introduce-operator-sch}
    T_{\mathcal{S}}
    =
    (I+U-\Delta)^{s/2},
\end{equation}
where $I$ is the identity operator, $U$ acts by pointwise multiplication with $u(\cdot)$,
and $-\Delta$ is the Laplacian operator.
Throughout, $T_{\mathcal{S}}$ is understood as the self-adjoint operator on
$L^2(\R^d)$ obtained as the closure of its action on $C_0^\infty(\R^d)$.
We note that $I+U-\Delta$ is a Schrödinger operator; see, for example,
\cite[Chapter~11.2]{lieb2001analysis}.

When $s$ is an even integer, say $s/2=k\in\N^*$, the operator $T_{\mathcal{S}}$
coincides with the $k$-fold composition of $I+U-\Delta$ with itself and is a
differential operator of order $2k$. Its principal Kohn--Nirenberg symbol is
\begin{equation}\label{eq:KN-symbol-sch}
    \sigma_{\mathcal{S}}(x,\omega)
    =
    \bigl(1+u(x)+(2\pi\|\omega\|_2)^2\bigr)^{s/2}.
\end{equation}
For general $s>0$, the operator $T_{\mathcal{S}}$ is defined via functional
calculus and is a pseudodifferential operator whose principal Kohn--Nirenberg
symbol is given by \eqref{eq:KN-symbol-sch}; see, for instance,
\cite[Definition~3.13]{hinzmicrolocal}. As a concrete and widely studied example, we work with the quadratic
potential (cf.\ \cite[Section~25.3]{shubin_pseudodifferential_2001}),
\[
    u(x)=c\,\|x\|_2^2,
    \qquad x\in\R^d,
\]
with a fixed constant $c>0$.

The Sobolev space of order $s$ induced by $T_{\mathcal{S}}$ is defined as
\[
\mathcal{H}_{\mathcal{S}}^{(s)}
\;\coloneqq\;
\overline{C_0^\infty(\R^d)}^{\,\|\cdot\|_{\mathcal{S}}^{(s)}},
\qquad
\|\varphi\|_{\mathcal{S}}^{(s)}
\coloneqq
\|T_{\mathcal{S}}\varphi\|_{L^2(\R^d)},
\quad \varphi\in C_0^\infty(\R^d),
\]
where the closure is taken with respect to the topology induced by
$\|\cdot\|_{\mathcal{S}}^{(s)}$, so that
$\mathcal{H}_{\mathcal{S}}^{(s)}$ embeds continuously into $L^2(\R^d)$.
Equivalently, $\mathcal{H}_{\mathcal{S}}^{(s)}$ coincides with the domain of the
% closed 
operator $T_{\mathcal{S}}$ on $L^2(\R^d)$, endowed with the norm
\[
\|f\|_{\mathcal{S}}^{(s)}=\|T_{\mathcal{S}}f\|_{L^2(\R^d)}.
\]
In this sense, $\mathcal{H}_{\mathcal{S}}^{(s)}$ may be viewed as a Sobolev space
of order $s$ with a confining potential, encoding $s$ degrees of smoothness
together with spatial localization.
The corresponding unit ball is
\begin{equation}\label{eq:definition-ballsch}
    \mathcal{B}_{\mathcal{S}}^{(s)}
    \coloneqq
    \bigl\{f\in\mathcal{H}_{\mathcal{S}}^{(s)}:
    \|T_{\mathcal{S}}f\|_{L^2(\R^d)}\le 1\bigr\}.
\end{equation}

We are now in a position to characterize the metric entropy of
$\mathcal{B}_{\mathcal{S}}^{(s)}$.

\begin{theorem}\label{thm:sob-sch}
Let $d\in\N^*$ and $s,c>0$.
The metric entropy of $\mathcal{B}_{\mathcal{S}}^{(s)}$ obeys
\begin{equation}\label{eq:pinsker-entropy}
    H\!\left(\varepsilon;\mathcal{B}_{\mathcal{S}}^{(s)}\right)
    \sim
    \frac{s\,\omega_{2d}}{2d(2\pi\sqrt{c})^{d}}\,
    \varepsilon^{-2d/s},
    \qquad \varepsilon\to0,
\end{equation}
and the corresponding minimax risk admits the asymptotic expansion
\begin{equation}\label{eq:pinsker-unbounded}
    R_\kappa \left(\mathcal{B}_{\mathcal{S}}^{(s)}\right)
    \sim
    \frac{d+s}{d}
    \left(
        \frac{d\,s\,\omega_{2d}\,\kappa^2}
        {(2\pi\sqrt{c})^{d}(d+s)(2d+s)}
    \right)^{\frac{s}{d+s}},
    \qquad \kappa\to0.
\end{equation}
\end{theorem}

\begin{proof}
    See Section~\ref{sec:proof-thm-sob-sch}.
\end{proof}

\noindent
Theorem~\ref{thm:sob-sch}, in particular the minimax-risk asymptotic
\eqref{eq:pinsker-unbounded}, extends Pinsker’s theorem to Sobolev spaces on
unbounded domains. Pinsker’s original result was established in
\cite{pinsker1980optimal}; general reviews can be found in
\cite{nussbaum1999minimax,johnstone2019estimation,tsybakovIntroductionNonparametricEstimation2009},
and an extension to arbitrary bounded domains is provided in
\cite[Theorem~9]{allard2025metricentropyminimaxrisk}. The result obtained here
shows that Pinsker’s principle is not confined to bounded domains: after
restoring compactness through spatial confinement, the minimax risk remains
governed by the same structural mechanism identified by Pinsker, with the sharp
asymptotic constant now determined by phase-space volume rather than boundary
effects. This indicates that Pinsker’s theorem reflects an underlying geometric
principle of statistical estimation, rather than an artifact of bounded domains
or specific boundary conditions.

% \iffalse 
% Theorem~\ref{thm:sob-sch},
% and more specifically \eqref{eq:pinsker-unbounded},
% extend Pinsker's theorem to Sobolev spaces on unbounded domains.
% We refer to \cite{pinsker1980optimal} for the original statement of Pinsker's theorem, to \cite{nussbaum1999minimax,johnstone2019estimation,tsybakovIntroductionNonparametricEstimation2009} for a review, and to \cite[Theorem~9]{allard2025metricentropyminimaxrisk} for a version on arbitrary bounded domains.
% \fi 

As a second application, we consider weighted Sobolev spaces, defined following
\cite[Section~3.3]{hinzmicrolocal}.
Fix $s,r,c>0$, set $u(x)=c\|x\|_2^2$ for $x\in\R^d$, and let
\begin{equation*}
    T_{\mathcal{W}}
    \coloneqq
    (I-\Delta)^{s/2}(I+U)^{r/2},
\end{equation*}
where $I$ is the identity operator, $U$ acts by pointwise multiplication with $u(\cdot)$,
and $-\Delta$ is the Laplacian operator.
As before, $T_{\mathcal{W}}$ is interpreted as a pseudodifferential operator; in
this case, its principal Kohn--Nirenberg symbol is given by
\begin{equation}\label{eq:symb-sob-weight-inverse}
    \sigma_{\mathcal{W}}(x,\omega)
    =
    \bigl(1+(2\pi\|\omega\|_2)^2\bigr)^{s/2}
    \bigl(1+c\|x\|_2^2\bigr)^{r/2}.
\end{equation}





% \iffalse 
% For the second application, we consider weighted Sobolev spaces,
% which we define following \cite[Section~3.3]{hinzmicrolocal}.
% Fix $s, r, c>0$, set $u(x)= c\|x\|_2^2$, for all $x\in\R^d$, and consider the operator
% \begin{equation*}
%     \tweight \coloneqq (I-\Delta )^{s/2} \left(I+U \right)^{r/2},
% \end{equation*}
% where $I$ denotes the identity operator,
% $U$ denotes the pointwise multiplication by $u(\cdot)$,
% and $- \Delta$ denotes the Laplacian. 
% Similarly to what has been done for $\tsch$, we interpret $\tweight$ as the pseudodifferential operator with right symbol
% \begin{equation}\label{eq:symb-sob-weight-inverse}
%     \symbolnotsobw (x, \omega)
%     = {\left(1+(2\pi \|\omega\|_2)^2\right)^{s/2}}{\left(1+c\|x\|_2^2\right)^{r/2}}.
% \end{equation}
% \fi 




The weighted Sobolev space of regularity $s$ and weight exponent $r$ induced by
$T_{\mathcal{W}}$ is defined as
\[
\mathcal{H}_{\mathcal{W}}^{(s,r)}
\;\coloneqq\;
\overline{C_0^\infty(\R^d)}^{\,\|\cdot\|_{\mathcal{W}}^{(s,r)}},
\qquad
\|\varphi\|_{\mathcal{W}}^{(s,r)}
\coloneqq
\|T_{\mathcal{W}}\varphi\|_{L^2(\R^d)},
\quad \varphi\in C_0^\infty(\R^d),
\]
where the closure is taken with respect to the topology induced by
$\|\cdot\|_{\mathcal{W}}^{(s,r)}$, so that
$\mathcal{H}_{\mathcal{W}}^{(s,r)}$ embeds continuously into $L^2(\R^d)$.
Equivalently, $\mathcal{H}_{\mathcal{W}}^{(s,r)}$ coincides with the domain of the
% closed 
operator $T_{\mathcal{W}}$ on $L^2(\R^d)$, endowed with the norm
\[
\|f\|_{\mathcal{W}}^{(s,r)}=\|T_{\mathcal{W}}f\|_{L^2(\R^d)}.
\]
In this sense, $\mathcal{H}_{\mathcal{W}}^{(s,r)}$ may be viewed as a weighted
Sobolev space of regularity $s$, where the weight exponent $r$ enforces spatial
localization.
The corresponding unit ball is
\begin{equation}\label{eq:definition-ballweight}
    \mathcal{B}_{\mathcal{W}}^{(s,r)}
    \coloneqq
    \bigl\{f\in\mathcal{H}_{\mathcal{W}}^{(s,r)}:
    \|T_{\mathcal{W}}f\|_{L^2(\R^d)}\le 1\bigr\}.
\end{equation}

We are now ready to characterize the metric entropy of
$\mathcal{B}_{\mathcal{W}}^{(s,r)}$.

\begin{theorem}\label{thm:sob-weighted}
Let $d \in \N^*$ and $s,r,c>0$. 
The metric entropy of $\mathcal{B}_{\mathcal{W}}^{(s,r)}$ obeys
\begin{equation*}
    H \left(\varepsilon;\mathcal{B}_{\mathcal{W}}^{(s,r)}, \|\cdot\|_{L^2(\R^d)}\right)
    \sim \frac{\omega_d^2}{(2\pi\sqrt{c})^{d}}
    \begin{dcases}
        \Xi_{r,s,d}\, \varepsilon^{-\frac{d}{\min\{r,s\}}}, 
        \quad &\text{if } r\neq s,\\[.25cm]           
        \varepsilon^{-\frac{d}{s}}\ln\left(\varepsilon^{-1}\right), \quad &\text{if } r=s,
    \end{dcases}
    \quad  \varepsilon\to 0,
\end{equation*}
where, for $r \neq s$,
\begin{equation}\label{eq:expression-constant-sobw}
     \Xi_{r,s,d}
     \coloneqq \frac{\min\{r,s\} \, \Gamma\left(\frac{d \, |s-r| }{2 \min\{r,s\}}\right) \Gamma \left( \frac{d}{2} \right)}{2\Gamma \left(\frac{d\max\{r,s\} }{2\min\{r,s\} }\right)}, 
\end{equation}
and $\Gamma$ denotes Euler's Gamma function.
% and define $f_{s,r,d}\colon \R_+^* \to \R_+^*$ by
% \[
% f_{s,r,d}(\varepsilon)=
% \frac{\omega_{d}^2}{d(2 \pi\sqrt{c})^{d}}
% \begin{cases}
% \dfrac{s^2}{r-s}\,\varepsilon^{-d/s}, & r>s, \\[.25cm]
% d\,\varepsilon^{-d/s}\ln(\varepsilon^{-1}), & r=s, \\[.25cm]
% \dfrac{r^2}{s-r}\,\varepsilon^{-d/r}, & r<s,
% \end{cases}
% \]
% As $\varepsilon \to 0$,
% \[
% f_{s,r,1}(\varepsilon)
% \;\lesssim\;
% H \left(\varepsilon;\mathcal{B}_{\mathcal{W}}^{(s,r)},
% \|\cdot\|_{L^2(\mathbb R)}\right)
% \;\lesssim\;
% 2^{\frac{s+r}{2\min\{r,s\}}}\,
% f_{s,r,1}(\varepsilon),
% \quad \text{if } d=1,
% \]
% and
% \[
% 2^{-\frac{(d-2)(s+r)}{2\min\{r,s\}}}\,
% f_{s,r,d}(\varepsilon)
% \;\lesssim\;
% H \left(\varepsilon;\mathcal{B}_{\mathcal{W}}^{(s,r)},
% \|\cdot\|_{L^2(\mathbb R^d)}\right)
% \;\lesssim\;
% f_{s,r,d}(\varepsilon),
% \quad \text{if } d\ge2.
% \]
\end{theorem}
\color{black}

% \iffalse 
% \begin{theorem}\label{thm:sob-weighted}
% Let $d \in \N^*$ and $s,r>0$.
% The metric entropy of $\mathcal{B}_{\mathcal{W}}^{(s,r)}$ satisfies
% \begin{equation*}
%     f_{s,r,1}(\varepsilon)
%     \leq
%     H \left(\varepsilon;\mathcal{B}_{\mathcal{W}}^{(s,r)},\|\cdot\|_{L^2(\R)}\right)
%     \leq
%     2^{\frac{s+r}{\min\{r,s\}}}\, f_{s,r,1}(\varepsilon),
% \end{equation*}
% if $d=1$, and
% \begin{equation*}
%     2^{-\frac{(d-2)(s+r)}{\min\{r,s\}}}\, f_{s,r,d}(\varepsilon)
%     \leq
%     H \left(\varepsilon;\mathcal{B}_{\mathcal{W}}^{(s,r)},\|\cdot\|_{L^2(\R^d)}\right)
%     \leq
%     f_{s,r,d}(\varepsilon),
% \end{equation*}
% if $d\ge2$, where $f_{s,r,d}\colon\R_+^*\to\R_+^*$, $d\in\N^*$, satisfies
% \begin{equation*}
%     f_{s,r,d}(\varepsilon)
%     \sim
%     \frac{\omega_{d}^2}{d(2\pi)^d}
%     \begin{cases}
%         \dfrac{s^2}{r-s}\,\varepsilon^{-d/s},
%         & \text{if } r>s, \\[.25cm]
%         d\,\varepsilon^{-d/s}\ln(\varepsilon^{-1}),
%         & \text{if } r=s, \\[.25cm]
%         \dfrac{r^2}{s-r}\,\varepsilon^{-d/r},
%         & \text{if } r<s,
%     \end{cases}
%     \qquad \text{as } \varepsilon \to 0.
% \end{equation*}
% \end{theorem}
% %%%T5[OK]: I extended the definition of f to d=1 as we also need it for d=1.
% %%%T5[OK]: I removed the ln(2) factor, please doublecheck.
% \fi 

\begin{proof}
    See Section~\ref{sec:proof-sob-weighted}.
\end{proof}

% \noindent
% Theorem~\ref{thm:sob-weighted} characterizes the metric entropy of the unit ball
% $\mathcal{B}_{\mathcal{W}}^{(s,r)}$ in the weighted Sobolev space
% $\mathcal{H}_{\mathcal{W}}^{(s,r)}$ up to a dimension-dependent multiplicative constant,
% namely $2^{\frac{s+r}{2\min\{r,s\}}}$ if $d=1$ and
% $2^{-\frac{(d-2)(s+r)}{2\min\{r,s\}}}$ if $d \ge 2$.
% These factors arise from the comparison step in the proof, where the volume
% function associated with the symbol is bounded by that of a simpler auxiliary
% quantity.
% Importantly, the result identifies the correct asymptotic regimes
% $r>s$, $r=s$, and $r<s$, together with the corresponding leading-order entropy
% behavior as $\varepsilon\to0$, thereby capturing the essential interplay between
% regularity and spatial localization.
% In particular, for $d\ge2$ the upper bound and for $d=1$ the lower bound
% capture the correct leading-order asymptotic term.
% For $d=2$, both bounds coincide and the entropy is determined exactly
% to first order.

% \iffalse 
% Theorem~\ref{thm:sob-weighted} characterizes the metric entropy of the unit ball in the weighted Sobolev space up to a multiplicative constant $2^{-\frac{(d-2)(s+r)}{\min\{r,s\}}}$.
% This is likely an artifact of our proof technique, and we expect the upper bound in the case $d\geq 2$, respectively the lower bound if $d=1$, to provide the correct asymptotics.
% Note that the result is sharp for $d=2$.
% \fi 






% \color{red}


\noindent
The expression \eqref{eq:expression-constant-sobw} for the constant appearing in
Theorem~\ref{thm:sob-weighted} can often be simplified by exploiting standard
identities for the Gamma function \cite{artin}. 
In particular, when the dimension is even, say $d=2k$ with $k\in\N^*$,
the Gamma factors reduce to factorials and one obtains
\[
\Xi_{r,s,2k}
=
\frac{(k-1)!\,\min\{r,s\}^{\,k+1}}
{2\prod_{j=1}^{k}\bigl(k\max\{r,s\}-j\min\{r,s\}\bigr)},
\qquad r\neq s .
\]
For example, in dimension $d=2$ this yields
\[
\Xi_{r,s,2}
=
\frac{\min\{r,s\}^2}{2\, (\max\{r,s\}-\min\{r,s\})}.
\]

\iffalse
The expression \eqref{eq:expression-constant-sobw} for the constant appearing in Theorem~\ref{thm:sob-weighted} can often be simplified
using the properties of the Gamma function \cite{artin}.
For instance, in even dimensions, that is, when $d=2k$ for some $k\in\N^*$, one obtains 
\begin{equation*}
    \Xi_{r,s,2k}
    = \dfrac{(k-1)!\, \min\{r,s\}^{k+1}}{2\prod_{j=1}^{k}\left(k\max\{r,s\}-j\min\{r,s\}\right)},
    \quad \text{e.g.,} \quad 
    \Xi_{r,s,2} = \frac{\min\{r,s\}^2}{2(\max\{r,s\}-\min\{r,s\})},
\end{equation*}
for all $s \neq r$.
\fi 

\color{black}

\iffalse
Theorem~\ref{thm:sob-weighted} also yields sharp entropy bounds for compact
embeddings between weighted Sobolev spaces defined via the operators
$T_{\mathcal{W}}$.
Specifically, for $s_1>s_2$ and $r_1>r_2$, consider the inclusion map
\[
I\colon \mathcal{H}_{\mathcal{W}}^{(s_1,r_1)} \longrightarrow
\mathcal{H}_{\mathcal{W}}^{(s_2,r_2)},
\]
%%%T10[OK]: Again, this may not be the right reference, see the corr. comment in the proof of Theorem 5.
which is compact by Rellich’s theorem
(see \cite[Theorem~2.17 \textcolor{blue}{and Exercise~2.12}]{hinzmicrolocal}).
The case treated in Theorem~\ref{thm:sob-weighted} corresponds to the special
choice $s_2=r_2=0$.
\fi 

% \color{red}
Theorem~\ref{thm:sob-weighted} also yields sharp entropy asymptotics for
compact embeddings between weighted Sobolev spaces defined through the
operators $T_{\mathcal W}$. In particular, for parameters $s_1>s_2$ and
$r_1>r_2$, consider the inclusion map
\[
I \colon \mathcal{H}_{\mathcal W}^{(s_1,r_1)}
\longrightarrow
\mathcal{H}_{\mathcal W}^{(s_2,r_2)} ,
\]
which is compact by Rellich’s theorem
(see \cite[Theorem~2.17 and Exercise~2.12]{hinzmicrolocal}).
The special case treated in Theorem~\ref{thm:sob-weighted}
corresponds to the choice $s_2=r_2=0$. The same analysis as in the proof of Theorem~\ref{thm:sob-weighted} applies to the
embedding operator $I$ with the parameters $(r,s)$ replaced by
$(r_1-r_2,s_1-s_2)$. Consequently, 
\begin{equation*}
    H_I(\varepsilon)
    \sim \frac{\omega_d^2}{(2\pi\sqrt{c})^{d}}
    \begin{dcases}
        \Xi_{r_1-r_2,s_1-s_2,d}\,
        \varepsilon^{-\frac{d}{\min\{r_1-r_2,s_1-s_2\}}},
        & \text{if } r_1+s_2 \neq r_2+s_1,\\[.25cm]
        \varepsilon^{-\frac{d}{r_1-r_2}}\ln \left(\varepsilon^{-1}\right),
        & \text{if } r_1+s_2 = r_2+s_1,
    \end{dcases}
    \qquad \varepsilon\to 0,
\end{equation*}
where the constant $\Xi_{r,s,d}$ is defined as in
Theorem~\ref{thm:sob-weighted}. The order of the entropy numbers for such embeddings has been studied
in a more general setting—for instance for weighted Besov and
Triebel–Lizorkin spaces \cite[Theorem~4.3.2]{edmundsFunctionSpacesEntropy1996}.
However, these results do not identify the leading constants.
The present approach provides these constants explicitly in the weighted
Sobolev setting.

\iffalse
\color{blue}
Applying the same argument as above, one obtains explicit entropy bounds for the
embedding $I$.
Specifically, 
%%%T10[OK]:There was a factor 2 missing in the numerator of the exponent, I believe, I changed this in the proof of Theorem 5 throughout and think it should %changed here as well?
\begin{equation*}
    H_I(\varepsilon)
    \sim \frac{\omega_d^2}{(2\pi\sqrt{c})^{d}}
    \begin{dcases}
        \Xi_{r_1-r_2,s_1-s_2,d}\, \varepsilon^{-\frac{d}{\min\{r_1-r_2,s_1-s_2\}}}, 
        \quad &\text{if } r_1+s_2 \neq r_2+s_1,\\[.25cm]           
        \varepsilon^{-\frac{d}{s}}\ln\left(\varepsilon^{-1}\right), \quad &\text{if } r_1+s_2 = r_2+s_1,
    \end{dcases}
    \quad  \varepsilon\to 0,
\end{equation*}
where $\Xi_{r,s,d}$ is defined as in Theorem~\ref{thm:sob-weighted}.
The order of the entropy numbers for such embeddings has been investigated in
greater generality—for instance for weighted Besov and Triebel–Lizorkin spaces
\cite[Theorem~4.3.2]{edmundsFunctionSpacesEntropy1996}—but without identification
of the leading constants.
The present approach provides these constants explicitly in the weighted Sobolev
setting.
\fi 

\color{black}


\section{Proofs}

\subsection{Proof of Theorem~\ref{thm:main-result}}\label{sec:proof-main-result}

We begin by recalling that, by definition of $\hyponegclass(\R^{2d})$
(see \eqref{eq:negative-order-def-class}), there exist $\rho\in(0,1]$ and real numbers $m_+,m_-$ with $m_+\ge m_->0$, all of which we fix throughout the proof,
such that $\symbolnot \in \hyposymbclassinv(\R^{2d})$.
Moreover, by \cite[Lemma~25.1,~1)a)]{shubin_pseudodifferential_2001},
$\symbolnot \in \hyposymbclassinv(\R^{2d})$ implies that
$\symbolnot^{-1}\in \hyposymbclass(\R^{2d})$.

% \iffalse
% \textcolor{blue}{
%     We first note that, by definition of $\hyponegclass(\R^{2d})$ (recall \eqref{eq:negative-order-def-class}),
%     there exists $\rho\in(0,1]$ and  $m_+\ge m_->0$ (which will be fixed throughout the proof)
%     such that  $\symbolnot \in \hyposymbclassinv(\R^{2d})$.
%     \textcolor{blue}{
%     Moreover, by \cite[Lemma~25.1 1) a)]{shubin_pseudodifferential_2001},  $\symbolnot \in \hyposymbclassinv(\R^{2d})$ implies $\symbolnot^{-1}\in \hyposymbclass(\R^{2d})$.}
% }
% \fi 


To derive the entropy asymptotics, we first reduce the Weyl quantization
$T_\sigma$ to a positive self-adjoint operator.
By the polar decomposition \cite[VIII.3.11]{conway},
there exists a partial isometry $U$ such that
$T_\sigma = U |T_\sigma|$.
Since $U$ acts isometrically on the range of $|T_\sigma|$,
it maps the image of the unit ball under $|T_\sigma|$
onto the image of the unit ball under $T_\sigma$.
Consequently,
\begin{equation}\label{eq:equality-entropies-absolute-value}
    H_{T_\sigma}(\varepsilon)
    =
    H_{|T_\sigma|}(\varepsilon),
    \qquad \varepsilon>0.
\end{equation}





Our strategy is to invoke the spectral asymptotics of
Dauge and Robert stated in Theorem~\ref{thm:Dauge-Robert}.
To this end, we choose the weight functions
\begin{equation}\label{eq:weight-choice-hypo}
    \phi(z)
    = \varphi(z)
    = \bigl(1+\|z\|_2^2\bigr)^{\rho/2},
    \qquad
    w(z)=\sigma(z),
    \qquad z\in\R^{2d}.
\end{equation}
Once it is verified that this choice satisfies conditions
\emph{(H1)}, \emph{(H2)}, \emph{(W)}, and \emph{(N)}, and that
$\sigma\in\Gamma(\R^{2d};w,\phi,\varphi)$,
Theorem~\ref{thm:Dauge-Robert} yields
\begin{equation}\label{eq:spec-asymptotics}
M_{|T_\sigma|}(\lambda)
\sim
V_\sigma(\lambda),
\qquad \lambda\to0.
\end{equation}
Since $w=\sigma$, the assumption
$V_w(\lambda)=O(V_\sigma(\lambda))$ required in
Theorem~\ref{thm:Dauge-Robert} holds trivially.
Because $V_\sigma$ is assumed to be regularly varying at zero,
it follows that $M_{|T_\sigma|}$ is regularly varying at zero as well.


% \color{red}
Since $|T_\sigma|$ is a positive compact operator,
it admits an orthonormal eigenbasis
$\{e_n\}_{n\in\N^*}$ with corresponding eigenvalues
$\{\lambda_n\}_{n\in\N^*}$ listed in nonincreasing order
and converging to zero.
In this basis, the image of the unit ball under $|T_\sigma|$
is the ellipsoid
\[
\mathcal E_\lambda
=
\Bigl\{
x=\sum_{n=1}^\infty x_n e_n :
\sum_{n=1}^\infty \frac{x_n^2}{\lambda_n^2}\le 1
\Bigr\}.
\]
The semi-axis-counting function of the ellipsoid $\mathcal E_\lambda$
coincides with the eigenvalue-counting function $M_{|T_\sigma|}(\lambda)$ of $|T_\sigma|$.
Hence, the entropy $H_{|T_\sigma|}(\varepsilon)$ is precisely the
metric entropy of the ellipsoid $\mathcal E_\lambda$.
%%%T12[OK]: the implication in the next sentence, i.e., reg. var at 0 with neg. index implies RC is not completely straightforward, perhaps a brief explanation.
Since $M_{|T_\sigma|}(\lambda)\sim V_\sigma(\lambda)$ as
$\lambda\to0$ and $V_\sigma$ is regularly varying at zero,
the regularity condition \emph{(RC)} required in
\cite[Theorem~2]{allard2025metricentropyminimaxrisk}
is satisfied (cf. the discussion following \cite[Lemma~10]{allard2025metricentropyminimaxrisk}).
We may therefore apply \cite[Theorem~2]{allard2025metricentropyminimaxrisk}
to conclude that
\[
H_{|T_\sigma|}(\varepsilon)
\sim
\int_{\varepsilon}^{\infty}
\frac{M_{|T_\sigma|}(\lambda)}{\lambda}\,d\lambda,
\qquad \varepsilon\to0.
\]
Combining this with \eqref{eq:equality-entropies-absolute-value} and \eqref{eq:spec-asymptotics} yields
\begin{equation}\label{eq:mainr-result-for-Weyl-1}
H_{T_\sigma}(\varepsilon)
=
H_{|T_\sigma|}(\varepsilon)
\sim
\int_{\varepsilon}^{\infty}
\frac{V_\sigma(\lambda)}{\lambda}\,d\lambda,
\qquad \varepsilon\to0,
\end{equation}
which is precisely \eqref{eq:main-res-volume-integral}.
It remains to verify conditions \emph{(H1)}–\emph{(N)}
for the choice \eqref{eq:weight-choice-hypo}.

\iffalse 
\color{blue}
% To pass from spectral asymptotics to entropy, 
To obtain asymptotics on the entropy,
we first reduce the Weyl quantization $T_\sigma$ of $\sigma$ to a positive
self-adjoint operator. By the polar decomposition
\cite[VIII.3.11]{conway}, there exists a partial isometry mapping the image of the
unit ball under \(|T_\sigma|\) onto the image of the unit ball under \(T_\sigma\).
Consequently,
\[
    H_{T_\sigma}(\varepsilon)\sim H_{|T_\sigma|}(\varepsilon),
    \qquad \varepsilon\to0.
\]
% Moreover, the eigenvalues of \(|T_\sigma|\) are given by
% \(\{|\lambda_n|\}_{n\in\N^*}\), where \(\{\lambda_n\}_{n\in\N^*}\) denotes the
% eigenvalues of \(T_\sigma\).
%%%T9[OK]: Please check the following argument and the reference to Shubin, was not sure but thought that this would be somewhere in Ch. 25 or 26.
% Since $\sigma$ is real-valued and bounded away from zero outside a compact subset
% of phase space, the operator $T_\sigma$ is elliptic and uniformly positive at
% infinity. In particular, $0$ lies outside the essential spectrum of $T_\sigma$,
% and only finitely many eigenvalues can be non-positive; see
% \cite[Chapters~25–26]{shubin_pseudodifferential_2001}. As a result, taking absolute
% values affects only finitely many eigenvalues, and hence does not change the
% eigenvalue-counting function at leading order. Therefore,
% \[
%     M_{|T_\sigma|}(\lambda)
%     \sim
%     M_\sigma(\lambda)
%     \sim
%     V_\sigma(\lambda),
%     \qquad \lambda \to 0.
% \]
%
Our strategy is to invoke the spectral asymptotics of Dauge and Robert stated in
Theorem~\ref{thm:Dauge-Robert}.
To this end, we select the weight functions
\begin{equation}\label{eq:weight-choice-hypo}
    \phi(z)
    = \varphi(z)
    = \bigl(1+\|z\|_2^2\bigr)^{\rho/2},
    \qquad
    w(z)=\sigma(z),
    \qquad z\in\R^{2d}.
\end{equation}
Once it is verified that this choice satisfies conditions
\emph{(H1)}, \emph{(H2)}, \emph{(W)}, and \emph{(N)}, and that
\(\sigma\in\Gamma(\R^{2d};w,\phi,\varphi)\),
Theorem~\ref{thm:Dauge-Robert} applies and yields
\[
    M_{|T_\sigma|} (\lambda)\sim V_\sigma(\lambda),
    \qquad \lambda\to0.
\]
For the present choice \(w=\sigma\), the assumption \(V_w(\lambda)=O(V_\sigma(\lambda))\) required in
Theorem~\ref{thm:Dauge-Robert} is automatically satisfied.
In particular, \(M_{|T_\sigma|}\) is regularly varying at zero.
We may therefore apply \cite[Theorem~2]{allard2025metricentropyminimaxrisk}
(see also \eqref{eq:type-tau-integrals} and \eqref{eq:integral-ecf-entropy}) to the
positive compact operator \(|T_\sigma|\) to conclude that
\begin{equation}\label{eq:mainr-result-for-Weyl-1}
    H_{T_\sigma}(\varepsilon)
    \sim H_{|T_\sigma|}(\varepsilon)  
    \sim \int_{\varepsilon}^{\infty}
        \frac{M_{|T_\sigma|}(\lambda)}{\lambda}\,\diff\lambda \nonumber 
    \sim \int_{\varepsilon}^{\infty}
        \frac{V_\sigma(\lambda)}{\lambda}\,\diff\lambda,
    \qquad \varepsilon\to0,%\label{eq:mainr-result-for-Weyl-2} 
\end{equation}
% \begin{align}
%     H_\sigma(\varepsilon)
%     &\sim H_{|T_\sigma|}(\varepsilon) \label{eq:mainr-result-for-Weyl-1} \\
%     &\sim \int_{\varepsilon}^{\infty}
%         \frac{M_{|T_\sigma|}(\lambda)}{\lambda}\,\diff\lambda \nonumber \\
%     &\sim \int_{\varepsilon}^{\infty}
%         \frac{M_\sigma(\lambda)}{\lambda}\,\diff\lambda
%      \sim \int_{\varepsilon}^{\infty}
%         \frac{V_\sigma(\lambda)}{\lambda}\,\diff\lambda,
%     \qquad \varepsilon\to0,\label{eq:mainr-result-for-Weyl-2} 
% \end{align}
which is precisely \eqref{eq:main-res-volume-integral}.
It remains to verify conditions \emph{(H1)}–\emph{(N)} for the choice
\eqref{eq:weight-choice-hypo}.
\color{black}
\fi 

% \iffalse 
% \color{blue}
% % \textbf{Proof of \eqref{eq:main-res-volume-integral}.}
% Our strategy for the first part of the statement will be to apply the result by Dauge and Robert presented in Theorem~\ref{thm:Dauge-Robert}.
% To this end, we choose the weights according to 
% \begin{equation}\label{eq:weight-choice-hypo}
%     \phi(z)
%     = \varphi(z) 
%     = \left(1+\|z\|_2^2\right)^{\frac{\rho}{2}},
%     \quad \text{and} \quad     w(z)
%     = \sigma(z),
%     \quad \text{for all } z \in \R^{2d}.
% \end{equation}
% Upon showing that, for the choice \eqref{eq:weight-choice-hypo}, 
% the conditions (H1), (H2), (W), and (N) are satisfied and that $\sigma\in\Gamma(\R^{2d}; w, \phi, \varphi)$ with $V_w (\lambda) = O_{\lambda\to\infty} \left(V_\sigma (\lambda)\right)$ (which are both trivial for our choice $w=\sigma$),
% we can apply Theorem~\ref{thm:Dauge-Robert} to obtain $M_{\symbolnot}(\lambda)  \sim V_{\symbolnot}(\lambda)$, as $\lambda\to 0$.
% Now, by polar decomposition \cite[VIII.3.11]{conway}, there exists a partial isometry mapping the image of the unit ball through the operator $|T_\sigma|$ to the image of the unit ball through $T_\sigma$, 
% so that $H_\sigma(\varepsilon)
% \sim
% H_{|T_\sigma|}(\varepsilon)$, as $\varepsilon\to 0$.
% Moreover, the eigenvalues of $|T_\sigma|$ are given by $\{|\lambda_n|\}_{n\in\N^*}$, where $\{\lambda_n\}_{n\in\N^*}$ denotes the eigenvalues of $T_\sigma$,
% which means that $M_{|T_\sigma|}(\lambda)  \sim M_{\symbolnot}(\lambda)  \sim V_{\symbolnot}(\lambda)$, as $\lambda\to 0$.
% This implies that $M_{|T_\sigma|}$ is regularly varying, and one can apply \cite[Theorem~2]{allard2025metricentropyminimaxrisk}
% (also see \eqref{eq:type-tau-integrals} and \eqref{eq:integral-ecf-entropy})
% to the positive self-adjoint compact operator $|T_\sigma|$
% to get 
% \begin{equation*}
%     H_\sigma(\varepsilon)
%     \sim
%     H_{|T_\sigma|}(\varepsilon)
%     \sim \int_{\varepsilon}^\infty \frac{M_{|T_\sigma|}(\lambda)}{\lambda}\,\diff\lambda
%     \sim \int_{\varepsilon}^\infty \frac{M_\sigma(\lambda)}{\lambda}\,\diff\lambda
%     \sim \int_{\varepsilon}^\infty \frac{V_\sigma(\lambda)}{\lambda}\,\diff\lambda,
%     \qquad \varepsilon \to 0,
% \end{equation*}
% which is the desired result \eqref{eq:main-res-volume-integral}.
% \fi

\color{black}
\textbf{Verifying (H1).}
To establish that $\phi^{-1}$ and $\varphi^{-1}$ are $(\phi,\varphi)$-continuous,
we start from the elementary inequality
\[
1+(t_1+t_2)^2 \le \bigl((1+t_1^2)^{1/2}+t_2\bigr)^2,
\qquad t_1,t_2\ge 0,
\]
which follows by direct expansion. Applying this with
$t_1=\|z\|_2$ and $t_2=\|z'\|_2$, and using the triangle inequality
$\|z+z'\|_2 \le \|z\|_2+\|z'\|_2$, we obtain
\begin{equation}\label{eq:verifyH1-DR-cont-1}
\bigl(1+\|z+z'\|_2^2\bigr)^{1/2}
\le
\bigl(1+\|z\|_2^2\bigr)^{1/2}+\|z'\|_2,
\qquad z,z'\in\R^{2d}.
\end{equation}
Replacing \(z\) by \(z+z'\) and \(z'\) by \(-z'\) in
\eqref{eq:verifyH1-DR-cont-1}, yields the
corresponding lower bound
\begin{equation}\label{eq:verifyH1-DR-cont-2}
\bigl(1+\|z+z'\|_2^2\bigr)^{1/2}
\ge
\bigl(1+\|z\|_2^2\bigr)^{1/2}-\|z'\|_2,
\qquad z,z'\in\R^{2d}.
\end{equation}
Fix $c>0$ and take $x,x',\omega,\omega'\in\R^d$ such that
\[
\frac{\|x'\|_2}{\varphi(x,\omega)}+\frac{\|\omega'\|_2}{\phi(x,\omega)}\le c.
\]
Setting $z=(x,\omega)$ and $z'=(x',\omega')$, and using $\phi(z)=\varphi(z)=(1+\|z\|_2^2)^{\rho/2}$, we obtain
\[
\|z'\|_2
\le \|x'\|_2+\|\omega'\|_2
\le c\,(1+\|z\|_2^2)^{\rho/2}
\le c\,(1+\|z\|_2^2)^{1/2},
\]
since $\rho\le 1$. Combining this with
\eqref{eq:verifyH1-DR-cont-1}–\eqref{eq:verifyH1-DR-cont-2} yields
\[
1-c
\le
\frac{(1+\|z+z'\|_2^2)^{1/2}}{(1+\|z\|_2^2)^{1/2}}
\le
1+c.
\]
Raising both sides to the power $\rho$ and choosing $c=1/2$ shows that
\[
C^{-1}
\le
\frac{\phi^{-1}(z+z')}{\phi^{-1}(z)}
\le
C,
\qquad C=2^\rho,
\]
which proves that $\phi^{-1}$ is $(\phi,\varphi)$-continuous. Since $\phi=\varphi$ by \eqref{eq:weight-choice-hypo}, we have $\phi^{-1}=\varphi^{-1}$.
Therefore, the same estimate holds for $\varphi^{-1}$, and both functions are
$(\phi,\varphi)$-continuous.

We next verify that $\phi^{-1}$ and $\varphi^{-1}$ are $(1,1)$-temperate.
By Peetre’s inequality (see \cite[(2.21)]{treves_introduction_1980}),
\begin{equation}\label{eq:Peetre}
(1+\|z\|_2^2)^{1/2}
\le
\sqrt{2}\,(1+\|z'\|_2^2)^{1/2}(1+\|z+z'\|_2^2)^{1/2},
\qquad z,z'\in\R^{2d}.
\end{equation}
Rewriting \eqref{eq:Peetre} gives
\begin{equation}\label{eq:DR-HA-hypo-tempstep}
(1+\|z+z'\|_2^2)^{-\rho/2}
\le
2^{\rho/2}\,
(1+\|z\|_2^2)^{-\rho/2}
(1+\|z'\|_2^2)^{\rho/2}.
\end{equation}
For $z=(x,\omega)$ and $z'=(x',\omega')$, using
$1+\|z'\|_2^2\le (1+\|x'\|_2+\|\omega'\|_2)^2$, we obtain
\[
\phi^{-1}(x+x',\omega+\omega')
\le
2^{\rho/2}\,\phi^{-1}(x,\omega)\,
\bigl(1+\|x'\|_2+\|\omega'\|_2\bigr)^\rho.
\]
Since $\phi=\varphi$ by \eqref{eq:weight-choice-hypo}, we have $\phi^{-1}=\varphi^{-1}$.
Hence the same bound holds for $\varphi^{-1}$, and both functions are
$(1,1)$-temperate.

Finally, since $\rho>0$, we have $\phi^{-1}(z)=\varphi^{-1}(z)\le 1$ for all
$z\in\R^{2d}$, and hence both functions are bounded on $\R^{2d}$.
Together with the $(\phi,\varphi)$-continuity and $(1,1)$-temperateness
established above, this verifies condition \emph{(H1)}.

% \iffalse 
% \textbf{Verifying (H1).}
% To show that $\phi^{-1}, \varphi^{-1}$ are $(\phi,\varphi)$-continuous, we first start from the inequality $1+(t_1+t_2)^2 \leq ((1+t_1^2)^{1/2}+t_2)^2$, valid for all $t_1,t_2\geq 0$, which can be verified  by developing both sides.
% Applying this inequality to $t_1=\|z\|_2$ and $t_2=\|z'\|_2$, we obtain
% \begin{equation}\label{eq:verifyH1-DR-cont-1}
%     \left(1+\|z + z'\|_2^2\right)^{\frac{1}{2}}
%     \leq \left(1+\|z \|_2^2\right)^{\frac{1}{2}} + \|z'\|_2,
%     \quad \text{for all } z,z'\in\R^{2d}.
% \end{equation}
% Substituting $z$ by $z+z'$ and $z'$ by $-z'$ in \eqref{eq:verifyH1-DR-cont-1} yields the corresponding lower bound
% \begin{equation}\label{eq:verifyH1-DR-cont-2}
%     \left(1+\|z + z'\|_2^2\right)^{\frac{1}{2}}
%     \geq \left(1+\|z \|_2^2\right)^{\frac{1}{2}} - \|z'\|_2,
%     \quad \text{for all } z,z'\in\R^{2d}.
% \end{equation}
% Now fix $c>0$, and assume that 
% $x,x',\omega,\omega'\in\R^d$ satisfy
% \begin{equation*}
%      \frac{\|x'\|_2}{\varphi(x, \omega)}
%     + \frac{\|\omega'\|_2}{\phi(x, \omega)}
%     \leq c,
% \end{equation*}
% so that, setting $z=(x,\omega)$ and $z'=(x',\omega')$, we have  
% \begin{equation}\label{eq:verifyH1-DR-cont-3}
%     \|z'\|_2 
%     = \|x'\|_2^2+\|\omega'\|_2^2
%     \leq \left(\|x'\|_2+\|\omega'\|_2\right)^{\frac{1}{2}}
%     \leq c \, \left(1+\|z \|_2^2\right)^{\frac{\rho}{2}}
%     \leq c \,  \left(1+\|z \|_2^2\right)^{\frac{1}{2}}.
% \end{equation}
% Using \eqref{eq:verifyH1-DR-cont-1} and \eqref{eq:verifyH1-DR-cont-2} in \eqref{eq:verifyH1-DR-cont-3},
% we get
% \begin{equation*}
%     1-c
%     \leq \frac{\left(1+\|z + z'\|_2^2\right)^{\frac{1}{2}}}{\left(1+\|z \|_2^2\right)^{\frac{1}{2}}}
%     \leq 1+c,
%     \quad \text{which implies} \quad 
%             C^{-1} 
%         \leq \frac{\phi^{-1}(z+z')}{\phi^{-1}(z)} \leq C
% \end{equation*}
% upon
% taking $c=1/2$ and $C=2^\rho$.
% A similar result holds for $\varphi^{-1}$, thereby establishing that $\phi^{-1}, \varphi^{-1}$ are $(\phi,\varphi)$-continuous.
% To show that $\phi^{-1}, \varphi^{-1}$ are $(1,1)$-temperate, we first start from Peetre's inequality (with $s=1/2$ and $\theta=z$, $\theta'=-z'$ in the formulation of \cite[(2.21)]{treves_introduction_1980},
% also see \cite[Lemma~4.3]{hinzmicrolocal})
% %$1+(t_1+t_2)^2 \leq 2(1+t_1^2)(1+t_2^2)$, valid for all $t_1,t_2\in \R$ (this can be verified easily by developing both sides).
% % Taking $t_1 = \|z'\|_2$ and $t_2 = \|z + z'\|_2$, we get for all $z,z'\in\R^{2d}$ 
% \begin{equation}\label{eq:Peetre}
%     \left(1+\|z\|_2^2\right)^{\frac{1}{2}}
%     % &= \left(1+\|z + z' - z'\|_2^2\right)^{\frac{1}{2}}\\
%     % &\leq \left(1+\left(\|z'\|_2 + \|z + z'\|_2\right)^2\right)^{\frac{1}{2}}
%     \leq \sqrt{2}\left(1+\|z'\|_2^2\right)^{\frac{1}{2}}\left(1+\|z + z'\|_2^2\right)^{\frac{1}{2}},
%     \quad \text{for all } z,z'\in\R^{2d},
% \end{equation}
% which we can rewrite according to 
% \begin{equation}\label{eq:DR-HA-hypo-tempstep}
%     \left(1+\|z + z'\|_2^2\right)^{-\frac{\rho}{2}}
%     \leq 2^{\frac{\rho}{2}}
%     \left(1+\|z\|_2^2\right)^{-\frac{\rho}{2}}\left(1+\|z'\|_2^2\right)^{\frac{\rho}{2}},
%     \quad \text{for all } z,z'\in\R^{2d}.
% \end{equation}
% Now, setting $z=(x,\omega)$ and $z'=(x',\omega')$, with $x,x',\omega,\omega'\in\R^d$, in  
% \eqref{eq:DR-HA-hypo-tempstep} and using the inequality $1+\|z'\|_2^2 \leq (1+\|x'\|_2+\|\omega'\|_2)^2 $ (which can be verified by developing both sides), 
% yields 
% \begin{equation*}
%     \phi^{-1}(x+x',\omega+\omega') 
%     \leq 2^{\frac{\rho}{2}} \,  \phi(x, \omega)^{-1} \left[1+ \|x'\|_2
%     + \|\omega'\|_2\right]^{\rho},
%     \quad \text{for all } x,x',\omega,\omega'\in\R^d.
% \end{equation*}
% A similar result holds for $\varphi$,
% which precisely means that the functions $\phi^{-1}, \varphi^{-1}$ are $(1,1)$-temperate. 
% Their boundedness follows immediately from $\phi^{-1}(z)
%     = \varphi^{-1}(z) \leq 1 $
% for all $z \in \R^{2d}$.
% \fi 

\textbf{Verifying (H2).}
To verify condition \emph{(H2)}, we start from the elementary bound $2t \le 1+t^2$, for $t\ge 0$.
Applying this estimate first with $t=\|x\|_2$ and then with $t=\|\omega\|_2$, where
$x,\omega\in\R^d$, and adding the resulting inequalities, yields
\[
1+\|x\|_2+\|\omega\|_2
\le
2\bigl(1+\|x\|_2^2+\|\omega\|_2^2\bigr),
\qquad x,\omega\in\R^d.
\]
Recalling that $\phi(x,\omega)=\varphi(x,\omega)
=(1+\|x\|_2^2+\|\omega\|_2^2)^{\rho/2}$, this bound can be rewritten as
\[
1+\|x\|_2+\|\omega\|_2
\le
2\bigl(\phi(x,\omega)\,\varphi(x,\omega)\bigr)^{1/\rho},
\qquad x,\omega\in\R^d.
\]
This is exactly condition~\eqref{eq:DR-H2} with $C=2^{-\rho}$ and $\zeta=\rho$.

% \iffalse 
% \textbf{Verifying (H2).}
% The condition (H2) is established by noting that $2 t\leq 1+t^2$, for all $t\geq 0$,
% applied successively to 
% $t=\|x\|_2$ and $t=\|\omega\|_2$ implies the bound
% \begin{equation*}
%     1+\|x\|_2 + \|\omega\|_2
%     \leq 2\left(1+\|x\|_2^2 + \|\omega\|_2^2\right)
%     = 2 \left(\phi(x, \omega)\, \varphi(x,\omega)\right)^{\frac{1}{\rho}},
%     \quad \text{for all } x,\omega \in \R^d,
% \end{equation*}
% which in turn yields \eqref{eq:DR-H2} with $C=1/2^\rho$ and $\zeta = \rho$.
% \fi 

\textbf{Verifying (W).}
By hypoellipticity of $\sigma$ (recall \eqref{eq:bound-def-hypoellip}), there exists
$R>0$ such that, for every multi-index $\alpha\in\N^{2d}$, the bound
\begin{equation}\label{eq:verify-W-hypo-DR-11}
    \bigl|\partial^\alpha \sigma(z)\bigr|
    \le
    C_\alpha\,\sigma(z)
    \bigl(1+\|z\|_2^2\bigr)^{-\frac{\rho|\alpha|}{2}}
\end{equation}
holds for all $z\in\R^{2d}$ with $\|z\|_2\ge R$, for some constant $C_\alpha>0$
depending on $\alpha$.
For fixed $\alpha$, the function
$z\mapsto \sigma(z)\bigl(1+\|z\|_2^2\bigr)^{-\frac{\rho|\alpha|}{2}}$ is continuous
and strictly positive, while $z\mapsto \partial^\alpha\sigma(z)$ is continuous
since $\sigma\in C^\infty(\R^{2d})$.
As the closed ball $\{z\in\R^{2d}:\|z\|_2\le R\}$ is compact, there exists a
constant $C'_\alpha>0$ such that
\begin{equation}\label{eq:verify-W-hypo-DR-12}
    \bigl|\partial^\alpha \sigma(z)\bigr|
    \le
    C'_\alpha\,\sigma(z)
    \bigl(1+\|z\|_2^2\bigr)^{-\frac{\rho|\alpha|}{2}},
    \qquad
    \|z\|_2 \le R.
\end{equation}
Combining \eqref{eq:verify-W-hypo-DR-11} and \eqref{eq:verify-W-hypo-DR-12}, and
recalling that $w=\sigma$ and $\phi=\varphi=(1+\|\cdot\|_2^2)^{\rho/2}$, we obtain,
for all $x,\omega\in\R^d$ and all multi-indices $\alpha,\beta\in\N^d$,
\begin{align}
    \bigl|\partial_x^\alpha \partial_\omega^\beta w(x,\omega)\bigr|
    &\le
    C_{\alpha,\beta}\,w(x,\omega)
    \bigl(1+\|(x,\omega)\|_2^2\bigr)^{-\frac{\rho(|\alpha|+|\beta|)}{2}}
    \label{eq:verify-W-hypo-1}\\
    &=
    C_{\alpha,\beta}
        \frac{w(x,\omega)}
        {\phi(x,\omega)^{|\alpha|}\,\varphi(x,\omega)^{|\beta|}}
    .
    \label{eq:verify-W-hypo-2}
\end{align}
This establishes condition~\eqref{eq:weight-class-DR} and hence shows that
$w\in\Gamma(\R^{2d};w,\phi,\varphi)$, as required for \emph{(W)}.




We show that the weight $w=\sigma$ is $(\phi,\varphi)$-temperate by a case
distinction on $z'$.
Fix $z,z'\in\R^{2d}$.

\medskip
\noindent\emph{Case 1:} 
    %$\|z'\|_2 \le \bigl(1+\|z\|_2^2\bigr)^{\rho/2}$.
    $\|z'\|_2 \leq \|z\|_2^\rho /2$.
Define the function
\[
    g \colon [0,1] \to \R,
    \qquad
    g(t) \coloneqq \ln\bigl(\sigma(z+t z')\bigr).
\]
Then
\begin{equation}\label{eq:proof-weight-temperate-DR1}
    g'(t)
    =
    \frac{\nabla \sigma(z+t z') \cdot z'}{\sigma(z+t z')},
    \qquad t\in[0,1].
\end{equation}
By hypoellipticity of $\sigma$, applying
\eqref{eq:verify-W-hypo-DR-11}--\eqref{eq:verify-W-hypo-DR-12} with $|\alpha|=1$
yields bounds on each first-order partial derivative of $\sigma$.
Combining the componentwise bounds corresponding to $|\alpha|=1$, there exists a
constant $K>0$ such that
\begin{equation}\label{eq:proof-weight-temperate-DR2}
    \|\nabla \sigma(z+t z')\|_2
    \le
    K\,\sigma(z+t z')
    \bigl(1+\|z+t z'\|_2^2\bigr)^{-\rho/2},
    \qquad t\in[0,1].
\end{equation}
Substituting this estimate into \eqref{eq:proof-weight-temperate-DR1} and applying
the Cauchy-Schwarz inequality yields
\begin{equation}\label{eq:proof-weight-temperate-DR3}
    |g'(t)|
    \le
    K\,\|z'\|_2\,\bigl(1+\|z+t z'\|_2^2\bigr)^{-\rho/2},
    \qquad t\in[0,1].
\end{equation}
Under the restriction of this case, $\|z'\|_2 \le \|z\|_2^\rho/2$, and assuming
$\|z\|_2\ge 1$, we estimate
\[
\|z+t z'\|_2
\ge
\|z\|_2 - t\|z'\|_2
\ge
\|z\|_2 - \tfrac12 \|z\|_2^\rho
\ge
\tfrac12 \|z\|_2,
\qquad t\in[0,1],
\]
where the last inequality uses $\rho\le 1$.
Inserting this bound into \eqref{eq:proof-weight-temperate-DR3} shows that the
right-hand side is uniformly bounded for $t\in[0,1]$. Hence, there exists a
constant $K_0>0$ such that
\begin{equation}\label{eq:g-uniform-bound}
    |g'(t)| \le K_0,
    \qquad t\in[0,1].
\end{equation}
% where the last inequality follows from the assumption of this case.
If $\|z\|_2 < 1$, then the assumption of this case implies
\(
\|z'\|_2 \le \tfrac12.
\)
Consequently, for all $t\in[0,1]$, the points $z+t z'$ remain in a fixed compact
subset of $\R^{2d}$.
Since $\sigma\in C^\infty(\R^{2d})$ is strictly positive, the function
$\nabla\sigma/\sigma$ is continuous and therefore bounded on this set.
% It follows that $|g'(t)|$ is uniformly bounded for all $t\in[0,1]$.
% Integrating \eqref{eq:proof-weight-temperate-DR3} 
It follows that $|g'(t)|$ is uniformly bounded for all $t\in[0,1]$, 
% \textcolor{red}{
which yields \eqref{eq:g-uniform-bound}.
% }
Integrating \eqref{eq:g-uniform-bound}
over the interval $[0,1]$ 
% \textcolor{red}{
gives
% }
\begin{equation}\label{eq:proof-weight-temperate-DR4}
    \left|\ln \left(\frac{\sigma(z+z')}{\sigma(z)}\right)\right|
    \le
    K_0,
\end{equation}
and therefore
\[
    \sigma(z+z') \le e^{K_0}\,\sigma(z).
\]
Writing \(z=(x,\omega)\) and \(z'=(x',\omega')\), and recalling that \(w=\sigma\), we hence get 
\[
w(x+x',\omega+\omega')
\le
e^{K_0}\,w(x,\omega)
\le
e^{K_0}\,w(x,\omega)\,
\bigl[1+\|x'\|_2\,\phi(x,\omega)+\|\omega'\|_2\,\varphi(x,\omega)\bigr],
\]
which is exactly \eqref{eq:def-phi-phi-temperate} with \(C=e^{K_0}\) and \(\zeta=1\).


% \iffalse 
% We first show that $w$ is $(\phi, \varphi)$-temperate. 
% To this end, we fix $z, z' \in \R^{2d}$ and split the analysis in two parts.
% First, we consider the case $\|z'\|_2 \leq \left(1+\|z\|_2^2\right)^{\rho/2}$.
% We introduce the function 
% \begin{equation*}
%     g \colon t \in [0,1]
%     \longmapsto \ln \left( \sigma(z+tz')\right),
% \end{equation*}
% so that 
% \begin{equation}\label{eq:proof-weight-temperate-DR1}
%     g'(t) 
%     = \frac{\nabla \sigma (z+tz')\cdot z'}{\sigma(z+tz')}, 
%     \quad \text{for all } t \in [0,1].
% \end{equation}
% By hypoellipticity of $\sigma$,
% we have (take $|\alpha|=1$ in \eqref{eq:verify-W-hypo-DR-11}--\eqref{eq:verify-W-hypo-DR-12})
% \begin{equation}\label{eq:proof-weight-temperate-DR2}
%     \left\| \nabla \sigma (z+tz') \right\|_2
%     \leq K \left\lvert  \sigma (z+tz') \right\rvert \left(1+\|z\|_2^2\right)^{-\frac{\rho}{2}} , 
%     \quad \text{for all } t \in [0,1],
% \end{equation}
% for some $K>0$.
% Combining \eqref{eq:proof-weight-temperate-DR1} with \eqref{eq:proof-weight-temperate-DR2} and applying the Cauchy-Schwarz inequality yields 
% \begin{equation}\label{eq:proof-weight-temperate-DR3}
%     \left\lvert g'(t) \right\rvert
%     \leq K \|z'\|_2 \left(1+\|z\|_2^2\right)^{-\frac{\rho}{2}}
%     \leq K, 
%     \quad \text{for all } t \in [0,1],
% \end{equation}
% where the second inequality follows from the initial assumption on $z$ and $z'$.
% Integrating both sides of \eqref{eq:proof-weight-temperate-DR3}
% between $0$ and $1$ gives 
% \begin{equation}\label{eq:proof-weight-temperate-DR4}
%     \left\lvert\ln \left( \frac{\sigma(z+z')}{\sigma(z)}\right)\right\rvert
%     \leq K
%     \quad \text{and thus} \quad 
%     \sigma(z+z') 
%     \leq e^{K}\sigma(z).
% \end{equation}
% Substituting $z=(x,\omega)$ and $z'=(x',\omega')$, with $x,x',\omega,\omega'\in\R^d$,
% and using that $w=\sigma$
% \eqref{eq:proof-weight-temperate-DR4} gives 
% \begin{equation*}
%     w(x+x',\omega+\omega') 
%     \leq e^{K}w(x,\omega) 
%     \leq e^{K}w(x,\omega) \left[1+ \|x'\|_2\, \phi(x, \omega)
%         + \|\omega'\|_2\, \varphi(x, \omega)\right],
% \end{equation*}
% which is \eqref{eq:def-phi-phi-temperate}
% with $C= e^{K}$ and $\zeta=1$.
% \fi 

\noindent\emph{Case 2:} $\|z'\|_2 \ge \|z\|_2^{\rho}/2$. By hypoellipticity, there exist constants $R>0$ and $C_1>0$ 
% \textcolor{red}{
with
% } 
\[
\sigma(z)\ge C_1\,\bigl(1+\|z\|_2^2\bigr)^{-m_+/2},
\qquad \|z\|_2\ge R.
\]
Since $\sigma$ is continuous and strictly positive on $\R^{2d}$, it attains a
strictly positive minimum on the closed ball
$B_R=\{z\in\R^{2d}:\|z\|_2\le R\}$.
% \textcolor{red}{
Hence, there exists $C_2>0$ satisfying
% }
\[
\sigma(z)\ge C_2\,\bigl(1+\|z\|_2^2\bigr)^{-m_+/2},
\qquad \|z\|_2\le R.
\]
These estimates imply that there exists a constant $K_1>0$ such that
\[
\sigma(z)\ge K_1\,\bigl(1+\|z\|_2^2\bigr)^{-m_+/2},
\qquad z\in\R^{2d}.
\]
Using this lower bound in conjunction with the corresponding upper bound from
\eqref{eq:bound-def-hypoellip}, we obtain
\[
    \frac{\sigma(z+z')}{\sigma(z)}
    \le
   K_2\,
    \frac{\bigl(1+\|z+z'\|_2^2\bigr)^{-m_-/2}}
         {\bigl(1+\|z\|_2^2\bigr)^{-m_+/2}},
\]
for some constant $K_2>0$.
Applying Peetre’s inequality \eqref{eq:Peetre} to the numerator gives
\[
    \frac{\sigma(z+z')}{\sigma(z)}
    \le
    K_3
    \bigl(1+\|z'\|_2^2\bigr)^{m_-/2}
    \bigl(1+\|z\|_2^2\bigr)^{(m_+-m_-)/2},
\]
for some constant $K_3>0$. Using the assumption of this case,
$\|z'\|_2 \ge \|z\|_2^{\rho}/2$, we can bound
$\bigl(1+\|z\|_2^2\bigr)^{(m_+-m_-)/2}
\le
C\,\bigl(1+\|z'\|_2^2\bigr)^{(m_+-m_-)/(2\rho)}$,
which leads to
\begin{equation}\label{eq:proof-weight-temperate-DR5}
    \sigma(z+z')
    \le
    K_4\,\sigma(z)
    \bigl(1+\|z'\|_2^2\bigr)^{\frac{m_-}{2}+\frac{m_+-m_-}{2\rho}},
\end{equation}
for some constant $K_4>0$. Writing $z=(x,\omega)$ and $z'=(x',\omega')$, we note that
\[
    1+\|z'\|_2^2
    =
    1+\|x'\|_2^2+\|\omega'\|_2^2
    \le
    \bigl(1+\|x'\|_2+\|\omega'\|_2\bigr)^2
    \le
    \bigl(1+\|x'\|_2\,\phi(x,\omega)
           +\|\omega'\|_2\,\varphi(x,\omega)\bigr)^2 .
\]
Recalling that $w=\sigma$, this implies
%%%t9[OK]: the exponent should be \frac{m_+-m_-}{\rho} instead of \frac{m_+-m_-}{2\rho}.
\[
    w(x+x',\omega+\omega')
    \le
    K_4\,w(x,\omega)
    \bigl[1+\|x'\|_2\,\phi(x,\omega)
          +\|\omega'\|_2\,\varphi(x,\omega)\bigr]^{\,m_-+\frac{m_+-m_-}{\rho}},
\]
which is precisely the $(\phi,\varphi)$-temperateness condition
\eqref{eq:def-phi-phi-temperate} with
$C=K_4$ and $\zeta = m_- + (m_+-m_-)/\rho$.

% \iffalse 
% \noindent\emph{Case 2:}
% % $\|z'\|_2 \geq \left(1+\|z\|_2^2\right)^{\rho/2}$.
%     $\|z'\|_2 \geq \|z\|_2^\rho /2$.
% Note that by similar continuity and strict positivity arguments as above, 
% we can extend the left part of \eqref{eq:bound-def-hypoellip} to the whole of $\R^d$, which gives 
% \begin{equation*}
%      \frac{\sigma(z+z')}{\sigma(z)}
%      \leq K_1 \frac{\left(1+\|z+z'\|_2^2\right)^{-\frac{m_{-}}{2}} }{\left(1+\|z\|_2^2\right)^{-\frac{m_{+}}{2}}}
%      \leq K_2 \left(1+\|z'\|_2^2\right)^{\frac{m_{-}}{2}} \left(1+\|z\|_2^2\right)^{\frac{m_{+}-m_{-}}{2}},
% \end{equation*}
% for some $K_1, K_2>0$,
% where the second step uses Peetre's inequality (recall \eqref{eq:Peetre}).
% Now, using the assumption on $z$ and $z'$, 
% we get 
% \begin{equation}\label{eq:proof-weight-temperate-DR5}
%      \sigma(z+z')
%      \leq \textcolor{blue}{K_3} \sigma(z) \left(1+\|z'\|_2^2\right)^{\frac{m_{-}}{2}+\frac{m_{+}-m_{-}}{\rho}},
% \end{equation}
% \textcolor{blue}{for some $K_3>0$.}
% We conclude by noting that, 
% $z=(x,\omega)$ and $z'=(x',\omega')$, with $x,x',\omega,\omega'\in\R^d$,
% \begin{equation*}
%    1+\|z'\|_2^2
%    = 1+\|x'\|_2^2 + \|\omega'\|_2^2
%    \leq \left(1+\|x'\|_2 + \|\omega'\|_2\right)^2
%    \leq \left(1+ \|x'\|_2\, \phi(x, \omega)
%         + \|\omega'\|_2\, \varphi(x, \omega)\right)^2
% \end{equation*}
% implies, together with \eqref{eq:proof-weight-temperate-DR5} and $w=\sigma$,
% the bound 
% \begin{equation*}
%     w(x+x',\omega+\omega') 
%     \leq K'' w(x,\omega) \left[1+ \|x'\|_2\, \phi(x, \omega)
%         + \|\omega'\|_2\, \varphi(x, \omega)\right]^{m_{-}+\frac{m_{+}-m_{-}}{2\rho}},
% \end{equation*}
% with $C= K''$ and $\zeta=m_{-}+(m_+-m_-)/(2\rho)$.
% \fi 




\textbf{Verifying (N).}
By hypoellipticity of $\sigma\in\hyposymbclassinv(\R^{2d})$
(recall \eqref{eq:bound-def-hypoellip}),
there exist constants $C_1,C_2,R>0$ such that
\begin{equation}\label{eq:verify-N-DR-hypo}
    C_1\bigl(1+\|z\|_2^2\bigr)^{-m_{+}/2}
    \le
    \sigma(z)
    \le
    C_2\bigl(1+\|z\|_2^2\bigr)^{-m_{-}/2},
    \qquad
    z\in\R^{2d}\ \text{with}\ \|z\|_2\ge R .
\end{equation}
%%%T9[OK]: This had m_ and m_+ swapped in the old version below.
With the choice $w=\sigma$, the bounds in \eqref{eq:verify-N-DR-hypo} imply
\begin{equation}\label{eq:verify-N-DR-hypo2}
    C_2^{-\frac{2\rho}{m_-}}\, w(z)^{\frac{2\rho}{m_-}}
    \le
    (1+\|z\|_2^2)^{-\rho}
    \le
    C_1^{-\frac{2\rho}{m_+}}\, w(z)^{\frac{2\rho}{m_+}},
    \qquad
    z\in\R^{2d}\ \text{with}\ \|z\|_2\ge R .
\end{equation}
% \iffalse
% \begin{equation}\label{eq:verify-N-DR-hypo2}
%     C_2^{-\frac{2\rho}{m_+}}\, w(z)^{\frac{2\rho}{m_+}}
%     \le
%     (1+\|z\|_2^2)^{-\rho}
%     \le
%     C_1^{-\frac{2\rho}{m_-}}\, w(z)^{\frac{2\rho}{m_-}},
%     \qquad
%     z\in\R^{2d}\ \text{with}\ \|z\|_2\ge R .
% \end{equation}
% \fi 
Since $w=\sigma$ is continuous and strictly positive, its restriction to the
closed ball $\{z\in\R^{2d}:\|z\|_2\le R\}$ attains a positive minimum and a finite
maximum. Define %%T9:There were two minima in the old version.
\[
    C_- \coloneqq \min_{\|z\|_2\le R} w(z),
    \qquad
    C_+ \coloneqq \max_{\|z\|_2\le R} w(z).
\]
Then, for all $z\in\R^{2d}$ with $\|z\|_2\le R$,
\begin{equation}\label{eq:verify-N-DR-hypo3}
    \left(1+R^2\right)^{-\rho}  C_+^{-\frac{2\rho}{m_-}}  w(z)^{\frac{2\rho}{m_-}}
    \leq \left(1+R^2\right)^{-\rho}
    \leq \left(1+\|z\|_2^2\right)^{-\rho} 
    \leq 1 
    \leq C_-^{-\frac{2\rho}{m_+}} w(z)^{\frac{2\rho}{m_+}} .
\end{equation}
Combining \eqref{eq:verify-N-DR-hypo2} and \eqref{eq:verify-N-DR-hypo3}, and setting
\[
    K' \coloneqq
    \min \left\{
        C_2^{-\frac{2\rho}{m_-}},
        (1+R^2)^{-\rho} C_+^{-\frac{2\rho}{m_-}}
    \right\},
    \qquad
    K \coloneqq
    \max \left\{
        C_1^{-\frac{2\rho}{m_+}},
        C_-^{-\frac{2\rho}{m_+}}
    \right\},
\]
$\gamma'=2\rho/m_-$, and $\gamma=2\rho/m_+$,
we obtain
\[
    K' w(z)^{\gamma'}
    \le
    \phi^{-1}(z)\,\varphi^{-1}(z)
    \le
    K w(z)^{\gamma},
    \qquad z\in\R^{2d},
\]
which is exactly condition~\eqref{eq:DR-N}.




% \iffalse
% \textbf{Verifying (N).}
% Invoking the hypoellipticity of $\sigma\in \hyposymbclassinv(\R^{2d})$ (recall \eqref{eq:bound-def-hypoellip}),
% we can find $C_1, C_2, R>0$ such that
% \begin{equation}\label{eq:verify-N-DR-hypo}
%     C_1\left(1+\|z\|_2^2\right)^{\frac{-m_{+}}{2}} 
%     \leq |\symbolnot(z)| 
%     \leq  C_2 \left(1+\|z\|_2^2\right)^{\frac{-m_{-}}{2}},
%     \quad \text{for all } z \in \R^{2d} \text{ with } \|z\|_2 \geq R.
% \end{equation}
% With our choice \eqref{eq:weight-choice-hypo},
% \eqref{eq:verify-N-DR-hypo} can equivalently be written as
% \begin{equation}\label{eq:verify-N-DR-hypo2}
%     C_2^{-\frac{2\rho}{m_+}} w(z)^{\frac{2\rho}{m_+}}  
%     \leq \left(1+\|z\|_2^2\right)^{-\rho} 
%     \leq C_1^{-\frac{2\rho}{m_-}} w(z)^{\frac{2\rho}{m_-}}, 
%     \quad \text{for all } z \in \R^{2d} \text{ with } \|z\|_2 \geq R.
% \end{equation}
% Moreover, the function $w$
% is continuous on the ball of radius $R$ in $\R^{2d}$, which is compact.
% Therefore, it attains its minimum and maximum (which are non-zero since $w$ is strictly positive), and, setting 
% \begin{equation*}
%     C_- \coloneqq \min\left\{w(z)\mid \|z\|_2 \leq R\right\}, 
%     \quad \text{and} \quad 
%     C_+ \coloneqq \min\left\{w(z)\mid \|z\|_2 \leq R\right\},
% \end{equation*}
% we obtain, for all  $z \in \R^{2d} $  with $\|z\|_2 \leq R$,
% \begin{equation}\label{eq:verify-N-DR-hypo3}
%     \left(1+R^2\right)^{-\rho}  C_+^{-\frac{2\rho}{m_+}}  w(z)^{\frac{2\rho}{m_+}}
%     \leq \left(1+R^2\right)^{-\rho}
%     \leq \left(1+\|z\|_2^2\right)^{-\rho} 
%     \leq 1 
%     \leq C_-^{-\frac{2\rho}{m_-}} w(z)^{\frac{2\rho}{m_-}} .
% \end{equation}
% Putting \eqref{eq:verify-N-DR-hypo2} and \eqref{eq:verify-N-DR-hypo3} together and setting 
% \begin{equation*}
%     K'=\min\left\{C_2^{-\frac{2\rho}{m_+}}, \left(1+R^2\right)^{-\rho}  C_+^{-\frac{2\rho}{m_+}}\right\}, 
%     \quad K=\max\left\{C_1^{-\frac{2\rho}{m_-}},  C_-^{-\frac{2\rho}{m_-}}\right\},
% \end{equation*}
% $\gamma' = 2\rho/m_+$, 
%  $\gamma=2\rho/m_-$
%  yields \eqref{eq:DR-N}.
% \fi 









% \iffalse
% As the symbol $\symbolnot$ is strictly positive, by assumption, its inverse $\symbolnot^{-1}$ is well-defined.
% The first step is now to apply Theorem~\ref{thm-shubin} to the symbol $\symbolnot^{-1}$.
% To this end, we note that, by \cite[Lemma~25.1 1) a)]{shubin_pseudodifferential_2001},  $\symbolnot \in \hyposymbclassinv(\R^{2d})$ implies $\symbolnot^{-1}\in \hyposymbclass(\R^{2d})$.
% Moreover, thanks to the assumption \eqref{eq:assumption-main-result}, 
% there exists $R>0$ and $C>0$ so that
% \begin{equation}\label{eq:verification-assump-inverse-tau}
%     \left|z\cdot\nabla\sigma^{-1}(z)\right|
%     =
%     \frac{\left|z\cdot\nabla\sigma(z)\right|}{\sigma^{2}(z)}
%     \ge
%     \frac{C\,\sigma(z)}{\sigma^{2}(z)}
%     =
%     C\,\sigma^{-1}(z),
% \end{equation}
% for all $z\in\R^{2d}$ with $\|z\|_2\ge R$.

% Consequently, all assumptions of Theorem~\ref{thm-shubin} are satisfied for the
% symbol $\sigma^{-1}$, and we obtain
% \begin{equation}\label{eq:application-main-result-proof}
%     \ecfup_{\sigma^{-1}}\!\left(\lambda^{-1}\right)
%     \sim
%     \volup_{\sigma^{-1}}\!\left(\lambda^{-1}\right),
%     \qquad \text{as } \lambda\to0.
% \end{equation}



% The second step builds on  \eqref{eq:application-main-result-proof} to characterize the entropy of a regularized inverse of the operator $T_{\symbolnot^{-1}}$.
% Next, we take $\lambda_* \geq 0$ outside the spectrum of $T_{\symbolnot^{-1}}$ and large enough for the operator $T_{\symbolnot^{-1}}+\lambda_* I$ to be strictly positive---this is possible by lower-boundedness of the spectrum of $T_{\symbolnot^{-1}}$, which follows from \eqref{eq:application-main-result-proof}. 
% (Recall from the discussion preceding Theorem~\ref{thm-shubin} that the spectrum of $T_{\symbolnot^{-1}}$ is discrete is equal to the set of its eigenvalues.)
% In particular,  $T_{\symbolnot^{-1}}+\lambda_* I$ has Weyl symbol $\symbolnot_*\coloneqq \symbolnot^{-1}+ \lambda_*$ and is invertible; we denote this inverse by $T_*=(T_{\symbolnot^{-1}}+\lambda_* I)^{-1}$.
% Specifically, we make two observations.
% First, we have
% \begin{equation*}
%     \int_{\mathbb{R}^{2d}} \mathbbm{1}_{\left\{\symbolnot^{-1}(x, \omega)\, < \, \lambda^{-1} \right\}} \diff  x\diff \omega 
%     =  \int_{\mathbb{R}^{2d}} \mathbbm{1}_{\{\symbolnot(x, \omega)\, > \, \lambda\}} \diff  x\diff \omega,
% \end{equation*}
% which immediately yields
% \begin{equation}\label{eq:transfer-vol}
%     \volup_{\symbolnot^{-1}}\left(\lambda^{-1}\right)
%     =V_{\symbolnot}(\lambda), 
%     \quad \text{for all } \lambda >0.
% \end{equation}
% It will be convenient for later purposes to also control the derivative $V_{\symbolnot}'(\cdot)$ of $V_{\symbolnot}(\cdot)$.
% This can be done as follows
% \begin{align}
%     V_{\symbolnot}'(\lambda)
%     = \left(\volup_{\symbolnot^{-1}}\left(\lambda^{-1}\right)\right)'
%     &= - \frac{\left(\volup_{\symbolnot^{-1}}\right)'\left(\lambda^{-1}\right)}{\lambda^2} \label{eq:shubin-derivative-bound1}\\
%     &= O_{\lambda\to 0} \left(\frac{\volup_{\symbolnot^{-1}}\left(\lambda^{-1}\right)}{\lambda}\right) 
%     = O_{\lambda\to 0} \left(\frac{V_{\symbolnot}(\lambda)}{\lambda}\right),\label{eq:shubin-derivative-bound2}
% \end{align}
% where the first and last equalities use \eqref{eq:transfer-vol}, and the third one is by application of \cite[Proposition~28.3]{shubin_pseudodifferential_2001}, whose application is justified by \eqref{eq:verification-assump-inverse-tau}.
% Second, we make use of the fact that the eigenvalues of the inverse of an operator are given by the inverse of the eigenvalues of this operator to get
% \begin{align*}
%     \# \left\{n\in\N^* \mid \lambda_n\left(T_{\symbolnot^{-1}}\right) < \lambda^{-1} - \lambda_* \right\}
%     &= \# \left\{n\in\N^* \mid \lambda_n\left(T_{\symbolnot^{-1}}+ \lambda_*I\right) < \lambda^{-1}\right\} \\
%     &=  \# \left\{n\in\N^* \mid \lambda_n\left(T_{\symbolnot^{-1}}+ \lambda_*I\right)^{-1} > \lambda\right\} \\
%     &=  \# \left\{n\in\N^* \mid \lambda_n\left(T_*\right) > \lambda \right\}.
% \end{align*}
% % where $\{\lambda_n(T_{\symbolnot^{-1}})\}_{n \in \N^*}$ and $\{\lambda_n(T_*)\}_{n \in \N^*}$ denote the eigenvalues of $T_{\symbolnot^{-1}}$ and $T_*$, respectively.
% In particular, we obtain 
% \begin{equation}\label{eq:transfer-ecf}
%     \ecfup_{\symbolnot^{-1}}\left(\lambda^{-1} - \lambda_*\right)
%     =M_{T_*}(\lambda), 
%     \quad \text{for all } \lambda \in (0, \lambda_*^{-1}).
% \end{equation}
% Combining \eqref{eq:application-main-result-proof}, \eqref{eq:transfer-vol}, and \eqref{eq:transfer-ecf}, yields the asymptotic equivalence
% \begin{equation}\label{eq:equivalence-ecf-volume-tstar}
%      M_{T_*}(\lambda)
%      \sim V_{\symbolnot}\left(\frac{\lambda}{1-\lambda\lambda_*}\right),
%       \quad \text{as } \lambda \to 0.
% \end{equation}
% Moreover, for all $\lambda\in (0, \lambda_*^{-1})$, there exists, by the mean value theorem, an element $\lambda_0\in ( \lambda, {\lambda}/(1-\lambda\lambda_*))$ for which 
% \begin{align*}
%     V_{\symbolnot}\left(\lambda\right)-V_{\symbolnot}\left(\frac{\lambda}{1-\lambda\lambda_*}\right)
%     &\leq C_1 \left\lvert V_{\symbolnot}'(\lambda_0) \right\rvert \left(\frac{\lambda}{1-\lambda\lambda_*} -\lambda\right)\\
%     &\leq C_2 \left\lvert V_{\symbolnot}(\lambda_0) \right\rvert \frac{\lambda}{\lambda_0} \left(\frac{1}{1-\lambda\lambda_*}-1 \right),
% \end{align*}
% where $C_1, C_2>0$ and we used \eqref{eq:shubin-derivative-bound1}--\eqref{eq:shubin-derivative-bound2}.
% In particular, we obtain 
% \begin{equation*}
%     V_{\symbolnot}\left(\lambda\right)-V_{\symbolnot}\left(\frac{\lambda}{1-\lambda\lambda_*}\right)
%     = o_{\lambda \to 0} \left(V_{\symbolnot}\left(\lambda\right)\right),
% \end{equation*}
% and, therefore,
% \begin{equation}\label{eq:volume-stability-final}
%     V_{\symbolnot}\left(\lambda\right)
%      \sim V_{\symbolnot}\left(\frac{\lambda}{1-\lambda\lambda_*}\right),
%       \quad \text{as } \lambda \to 0.
% \end{equation}
% Combining \eqref{eq:volume-stability-final} with \eqref{eq:equivalence-ecf-volume-tstar}, we can apply \cite[Theorem~2]{allard2025metricentropyminimaxrisk} (which is justified by regular variation of the volume which imply the condition (RC), cf the discussion in \cite[Appendix~C]{allard2025metricentropyminimaxrisk}) to asymptotically characterize the entropy of $T_*$ to the first order according to
% \begin{equation}\label{eq:asymp-entropy-inverse-operator}
%     H_{T_* }\left(\varepsilon\right)
%     \sim \int_{\varepsilon}^\infty 
%     \frac{V_{\symbolnot}(\lambda)}{\lambda} \diff \lambda, 
%     \quad \text{as } \varepsilon \to 0.    
% \end{equation}
% \fi



% We will assume, from now on, that $\text{Ker}(T_\symbolnot)=\{0\}$ so that $T_\symbolnot$ is invertible by \cite[Theorem~25.4]{shubin_pseudodifferential_2001}.
% This comes without loss of generality.
% Indeed, if it were not the case, one could consider an orthonormal eigenbasis $\{\psi_i\}_{i\in\N^*}$ of $L^2(\R^d)$ and build an operator $T_\psi$, which is diagonal in the basis $\{\psi_i\}_{i\in\N^*}$, with $i$-th diagonal element equal to $\exp{(-i)}$, if $\psi_i\in \text{Ker}(T_\symbolnot)$, and  $0$ otherwise.
% The operator $T_\symbolnot+T_\psi$ satisfies the same hypoellipticity properties as $T_\symbolnot$, has kernel $\text{Ker}(T_\symbolnot+T_\psi)=\{0\}$, and its entropy numbers (see below) are asymptotically equivalent to those of $T_\symbolnot$ due to the exponentially fast decay of the eigenvalues of $T_\psi$.

\color{black}

% \textbf{Proof that $H_{T^{-1}_{\symbolnot^{-1}}}(\varepsilon) 
% \sim H_\sigma(\varepsilon)$.}

% \subsection{Proof of Theorem~\ref{thm:main-result-2}}\label{sec:proof-main-result-2}


This completes the verification of conditions \emph{(H1)}–\emph{(N)} and hence
establishes \eqref{eq:main-res-volume-integral} for the Weyl quantization.
We next show that the same entropy asymptotics hold for the left and right
quantizations.
To this end, we first state a technical lemma, which is proved at the end of this
section.

\begin{lemma}\label{lem:technical-lemma}
Let $\rho_1>0$ and $m>0$.
Let $T_1$ and $T_2$ be compact linear operators on $L^2(\R^d)$ such that
\[
T_1 = T_2 (I + K) + K_{-\infty},
\]
for some $K\in \Psi^{-m}_{\rho_1}(\R^d)$ and $K_{-\infty}\in\Psi^{-\infty}(\R^d)$.
Assume further that there exists a strictly positive hypoelliptic symbol
$\sigma\in\hyponegclass(\R^{2d})$ whose associated volume function $V_\sigma$ is
regularly varying at zero, % with negative index
and that either $T_1$ or $T_2$ is the Weyl quantization
of $\sigma$.
Then %the metric entropies of $T_1$ and $T_2$ satisfy
\begin{equation}\label{eq:technical-lemma-result}
    H_{T_1}(\varepsilon)
    \sim
    H_{T_2}(\varepsilon),
    \qquad \varepsilon \to 0.
\end{equation}
\end{lemma}

% \iffalse 
% \begin{lemma}\label{lem:technical-lemma}
%     Let $\rho_1>0$ and $m>0$.
%     Let $T_1$ and $T_2$ be compact linear operators from $L^2(\R^d)$ to itself such that 
%     there exist $K\in \Psi^{-m}_{\rho_1}(\R^d)$ and  $K_{-\infty}\in\Psi^{-\infty}(\R^d)$ for which
%     $T_1 =  T_2 (I + K) + K_{-\infty}$. 
%     Let us further assume that there exists a strictly positive hypoelliptic symbol $\sigma\in\hyponegclass(\R^{2d})$, whose associated volume function $V_\sigma$ is regularly varying at zero,
%     and such that either $T_1$ or $T_2$ is the Weyl quantization of a $\sigma$.
%     Then, 
%     \begin{equation}\label{eq:technical-lemma-result}
%     H_{T_1}\left(\varepsilon \right)
%     \sim H_{T_2}\left(\varepsilon \right),
%     \quad \text{as } \varepsilon \to 0,
% \end{equation}
% \end{lemma}
% \fi 



% \noindent 
% With Lemma~\ref{lem:technical-lemma} at hand, we are ready to prove 
% the equivalence result for the difference quantizations as well as \eqref{eq:main-result-part2}.


\textbf{Invariance under change of quantization.}
We establish invariance of the entropy asymptotics under a change of
quantization by treating explicitly the Weyl–left case; the Weyl–right case
follows by the same arguments. Let $T_1$ and $T_2$ denote the left and Weyl quantizations of the symbol $\sigma$,
respectively, and let $\sigma_1$ be the Weyl symbol of $T_1$.
Since $T_1$ is the left quantization of a symbol
$\sigma\in\hyposymbclassinv(\R^{2d})$,
Proposition~25.1 in \cite{shubin_pseudodifferential_2001} implies that
the corresponding Weyl symbol $\sigma_1$ also belongs to
$\hyposymbclassinv(\R^{2d})$.
%By \cite[Proposition~25.1]{shubin_pseudodifferential_2001}, the symbol $\sigma_1$
%belongs to $\hyposymbclassinv(\R^{2d})$.

Since $T_2$ is the Weyl quantization of a hypoelliptic symbol in
$\hyposymbclassinv(\R^{2d})$, Theorem~25.1 in
\cite{shubin_pseudodifferential_2001} guarantees the existence of a right
parametrix $B_2$ with symbol $b_2\in\hyposymbclass(\R^{2d})$ such that
\begin{equation}\label{eq:parametrix-B-Tsigma}
    T_2 B_2 = I + S,
    \qquad S\in\Psi^{-\infty}(\R^d).
\end{equation}

% \textcolor{red}{
We may assume, without loss of generality, that the symbol $b_2$ is nowhere vanishing on $\R^{2d}$. 
Indeed, by hypoellipticity and the lower bound in \eqref{eq:bound-def-hypoellip}, there exists $R>0$ such that
$b_2(z)\neq 0$ for all $z \in \R^{2d}$ with $\|z\|_2 \ge R$. If $b_2$ vanishes at some point in the compact region
$\{z \in \R^{2d} \mid \|z\|_2 < R \}$, we can construct a function
$b_2' \in C^\infty(\R^{2d})$, supported in
$\{\, z \in \R^{2d} \mid \|z\|_2 \le 2R \,\}$, such that
\[
b_2'' \coloneqq b_2 + b_2' \in \hyposymbclass(\R^{2d})
\]
and $b_2''(z) \neq 0$ for all $z \in \R^{2d}$. Let $B_2'$ and $B_2''$ denote the Weyl quantizations of
$b_2'$ and $b_2''$, respectively. Replacing $B_2$ by $B_2''$ in \eqref{eq:parametrix-B-Tsigma}
yields
\[
T_2 B_2'' = I + S + T_2 B_2',
\]
where $T_2 B_2' \in \Psi^{-\infty}(\R^d)$ since $b_2'$ is compactly supported. Absorbing this term
into the remainder, we obtain a representation of the form
\[
T_2 B_2'' = I + S',
\qquad S' \in \Psi^{-\infty}(\R^d),
\]
with a parametrix whose symbol does not vanish on $\R^{2d}$.
Composing \eqref{eq:parametrix-B-Tsigma} on the right with $(T_1 - T_2)$ yields
\[
T_2 B_2 (T_1-T_2) = (T_1-T_2) + S(T_1-T_2).
\]
% }
\iffalse 
\textcolor{purple}{
    We can assume, without loss of generality, that $b_2(z)$ is non-zero for all $z\in \R^{2d}$.
    Indeed, by hypoellipticity and \eqref{eq:bound-def-hypoellip}, we are guaranteed that there exists $R>0$ such that $b_2(z) \neq 0$, whenever $\|z\|_2\geq R$. 
    If $b_2(z) = 0$ for some $z$ with $\|z\|_2< R$, it is possible to find $b'_2 \in C^\infty(\R^{2d})$, compactly supported in the ball of radius $2R$ of $\R^{2d}$, such that $b+b'\eqcolon b''_2 \in\hyposymbclass(\R^{2d})$ does not vanish in $\R^{2d}$.
    Therefore, upon replacing $B_2$ by $B''_2$ and $S$ by $T_2B'_2$ in \eqref{eq:parametrix-B-Tsigma}, with $B'_2\in\Psi^{-\infty}$ and $B''_2$ the Weyl quantizations of $b'_2$ and $b''_2$, respectively, 
    one can assume  that $b_2(z)$ is non-zero for all $z\in \R^{2d}$.
    Composing \eqref{eq:parametrix-B-Tsigma} %this identity 
    on the right with $(T_1-T_2)$ yields
}

\[
T_2 B_2 (T_1-T_2) = (T_1-T_2) + S(T_1-T_2).
\]
\fi 
Because $S\in\Psi^{-\infty}(\R^d)$ and $T_1-T_2$ is a pseudodifferential operator,
and since $\Psi^{-\infty}(\R^d)$ is a two-sided ideal in the
pseudodifferential calculus, the composition $S(T_1-T_2)$ belongs to
$\Psi^{-\infty}(\R^d)$; see, for example, \cite[Theorem~4.22]{hinzmicrolocal} or \cite[Theorem~23.6]{shubin_pseudodifferential_2001}.
%%%T9[OK]: please check the reference to Shubin, perhaps some other reference.
Letting
\[
S' \coloneqq -\,S(T_1-T_2)\in\Psi^{-\infty}(\R^d),
\]
we may therefore rewrite the above identity as
\[
T_2 B_2 (T_1-T_2) = (T_1-T_2) - S'.
\]
Rearranging terms yields
\begin{equation}\label{eq:t-perturb}
T_1
=
T_2\bigl(I + B_2(T_1-T_2)\bigr) + S'.
\end{equation}
%  where $S'\in\Psi^{-\infty}(\R^d)$ is a smoothing operator. 
With $K \coloneqq B_2(T_1-T_2)$ and $K_{-\infty}=S'$, we see that
\eqref{eq:t-perturb} is of the form required in
Lemma~\ref{lem:technical-lemma}.
Therefore, once it is shown that 
% \textcolor{red}{
$K\in\Psi^{-m}_{\rho_1}(\R^d)$ for some $\rho_1>0$ and $m>0$,
% }
the lemma applies and yields
\begin{equation}\label{eq:equiv-invariance-quant}
    H_{T_1}(\varepsilon)\sim H_{T_2}(\varepsilon),
    \qquad \varepsilon\to0.
\end{equation}




We next prove that in fact $K\in\Psi^{-2\rho}_{\rho}(\R^d)$.
This will be done by first showing that $\sigma^{-1}(\sigma_1-\sigma)\in \Gamma^{-2\rho}_{\rho}(\R^{2d})$. 
Recall that $\sigma_1$ denotes the Weyl symbol of the left quantization $T_1$.
By \cite[Theorem~23.3]{shubin_pseudodifferential_2001}, for every $N\in\N^*$ there
exists a remainder
$r_N\in\Gamma^{-m_- - 2(N+1)\rho}_{\rho}(\R^{2d})$ such that
\begin{equation}\label{eq:asymp-expansion-diff-symbols}
    \sigma_1(x,\omega)-\sigma(x,\omega)
    =
    \sum_{\substack{\alpha\in\N^d\\1\le|\alpha|\le N}}
    \frac{1}{\alpha!\,2^{|\alpha|}}
    \partial_\omega^\alpha\partial_x^\alpha\sigma(x,\omega)
    + r_N(x,\omega),
    \qquad x,\omega\in\R^d.
\end{equation}
Since $\sigma$ is strictly positive, division by $\sigma$ is well-defined, and
\eqref{eq:asymp-expansion-diff-symbols} yields
\begin{equation}\label{eq:finite-decom-Weyl-left}
    \frac{\sigma_1-\sigma}{\sigma}
    =
    \sum_{\substack{\alpha\in\N^d\\1\le|\alpha|\le N}}
    \frac{1}{\alpha!\,2^{|\alpha|}}
    \frac{\partial_\omega^\alpha\partial_x^\alpha\sigma}{\sigma}
    + \frac{r_N}{\sigma}.
\end{equation}
We treat the two terms on the right-hand side of \eqref{eq:finite-decom-Weyl-left} separately.
First, since $\sigma^{-1}\in\hyposymbclass(\R^{2d})$, application of
\cite[Lemma~25.1]{shubin_pseudodifferential_2001} allows us to conclude that
\[
\frac{r_N}{\sigma}
\in
\Gamma^{\,m_+ - m_- - 2(N+1)\rho}_{\rho}(\R^{2d}).
\]
Choosing $N$ sufficiently large ensures
\begin{equation}\label{eq:est-1}
\frac{r_N}{\sigma}\in \Gamma^{-2\rho}_{\rho}(\R^{2d}).
\end{equation}
Second, consider the derivative terms appearing in the finite sum in
\eqref{eq:finite-decom-Weyl-left}. For each multi-index $\alpha\neq 0$, the
Leibniz rule together with the derivative bounds in
\eqref{eq:bound-def-hypoellip} imply that
\[
\sigma^{-1}\,\partial_\omega^\alpha\partial_x^\alpha\sigma
\in
\Gamma^{-2\rho|\alpha|}_{\rho}(\R^{2d}).
\]
Since $|\alpha|\ge 1$, we have $-2\rho|\alpha|\le -2\rho$, and hence
\begin{equation}\label{eq:est-2}
\sigma^{-1}\,\partial_\omega^\alpha\partial_x^\alpha\sigma
\in
\Gamma^{-2\rho}_{\rho}(\R^{2d}).
\end{equation}
Combining \eqref{eq:est-1} and \eqref{eq:est-2} in \eqref{eq:finite-decom-Weyl-left} yields
\[
\sigma^{-1}(\sigma_1-\sigma)\in \Gamma^{-2\rho}_{\rho}(\R^{2d}),
\]
as claimed.

% \iffalse 
% In fact, we will show next that $K\in\Psi^{-2\rho}_{\rho}(\R^d)$. We first show that $\sigma^{-1}(\sigma_1-\sigma) \in \Gamma^{-2\rho}_{\rho}(\R^{2d})$.
% To this end, we apply \cite[Theorem~23.3]{shubin_pseudodifferential_2001} and use the hypoellipticity of $\sigma$
% to get that, for all $N\in\N^*$, there exists $r_N\in \Gamma^{-m_- - 2(N+1)\rho}_{\rho}(\R^{2d})$ such that
% \begin{equation}\label{eq:asymp-expansion-diff-symbols}
%     \sigma_1(x, \omega) - \sigma(x, \omega)
%     = \sum_{\substack{\alpha\in \N^{d}\\1\leq |\alpha|\leq N}} \frac{1}{\alpha!\,  2^{|\alpha|}} \partial^\alpha_\omega \partial_x^\alpha \sigma(x, \omega)
%     + r_N(x, \omega),
%     \quad \text{for all } x, \omega\in \R^d.
% \end{equation}
% Dividing both sides of \eqref{eq:asymp-expansion-diff-symbols} by $\sigma$ (recall that $\sigma$ is non-zero by assumption) yields 
% \begin{equation}\label{eq:finite-decom-Weyl-left}
%     \frac{\sigma_1(x, \omega) - \sigma(x, \omega)}{\sigma(x,\omega)}
%     = \sum_{\substack{\alpha\in \N^{d}\\1\leq |\alpha|\leq N}} \frac{1}{\alpha!\,  2^{|\alpha|}} 
%     \frac{\partial^\alpha_\omega \partial_x^\alpha \sigma(x, \omega)}{\sigma(x,\omega)}
%     + \frac{r_N(x, \omega)}{\sigma(x,\omega)},
%     \quad \text{for all } x, \omega\in \R^d.
% \end{equation}
% Similarly to \cite[Lemma~25.1]{shubin_pseudodifferential_2001},
% we verify that $r_N\in \Gamma^{-m_- - 2(N+1)\rho}_{\rho}(\R^{2d})$ together with $\symbolnot^{-1}\in \hyposymbclass(\R^{2d})$ 
% implies $\sigma^{-1}r_N\in \Gamma^{m_+-m_- - 2(N+1)\rho}_{\rho}(\R^{2d})$ and, therefore, we can choose $N$ large enough to have $\sigma^{-1}r_N\in \Gamma^{-2\rho}_{\rho}(\R^{2d})$.
% Next, one verifies, using the Leibniz rule and the second part of \eqref{eq:bound-def-hypoellip}, that $\sigma^{-1}\partial^\alpha_\omega \partial_x^\alpha \sigma \in \Gamma^{-2\rho|\alpha|}_{\rho}(\R^{2d})$, for all $\alpha \in \N^d$.
% We can therefore conclude from \eqref{eq:finite-decom-Weyl-left} that $\sigma^{-1}(\sigma_1-\sigma) \in \Gamma^{-2\rho}_{\rho}(\R^{2d})$.
% \fi 

Returning to the parametrix identity \(T_2 B_2 = I + S\), let \(s\) denote the Weyl
symbol of the % smoothing 
operator \(S\).
By the composition formula for Weyl quantization
\cite[Theorem~23.6]{shubin_pseudodifferential_2001} and the hypoellipticity of both
\(b_2\) and \(\sigma\), for every \(M\in\N\) there exist constants
\(C_{\alpha,\beta}>0\) and a remainder
\(r'_M\in\Gamma^{m_+ - m_- - 2(M+1)\rho}_{\rho}(\R^{2d})\) such that
\begin{align}
    1+s(z)
    &=
    \sum_{\substack{\alpha,\beta\in\N^d\\0\le|\alpha|+|\beta|\le M}}
    C_{\alpha,\beta}\,
    \partial_z^\alpha b_2(z)\,
    \partial_z^\beta \sigma(z)
    + r'_M(z)
    \label{eq:long-formula-1}\\
    &=
    b_2(z)\,\sigma(z)
    \left(
        1
        + \sum_{\substack{\alpha,\beta\in\N^d\\1\le|\alpha|+|\beta|\le M}}
        C_{\alpha,\beta}\,
        \frac{\partial_z^\alpha b_2(z)}{b_2(z)}\,
        \frac{\partial_z^\beta \sigma(z)}{\sigma(z)}
        + \frac{r'_M(z)}{b_2(z)\,\sigma(z)}
    \right),
    \label{eq:long-formula-2}
\end{align}
for all \(z\in\R^{2d}\). 
% \textcolor{red}{
Recall that $\sigma(z)$ is non-vanishing for all $z\in\R^{2d}$ by assumption, 
and that $b_2(z)$ has been chosen to be non-vanishing on $\R^{2d}$ without loss of generality.
% }

Choosing $M$ sufficiently large ensures that $r'_M\in\Gamma^{-2\rho}_{\rho}(\R^{2d})$.
Moreover, by the Leibniz rule and hypoellipticity,
\[
b_2^{-1}\,\partial^\alpha b_2 \in \Gamma^{-\rho|\alpha|}_{\rho}(\R^{2d}),
\qquad
\sigma^{-1}\,\partial^\beta \sigma \in \Gamma^{-\rho|\beta|}_{\rho}(\R^{2d}),
\]
for all multi-indices $\alpha,\beta$.
% \textcolor{red} {
Consequently, the expression in parentheses in \eqref{eq:long-formula-2}
is of the form $1+r'$ with $r'\in\Gamma^{-2\rho}_{\rho}(\R^{2d})$.
Since $r'(z)\to 0$ as $\|z\|_2\to\infty$, there exists $R>0$ such that
\[
|r'(z)|\le \tfrac12,
\qquad \text{for all } z\in\R^{2d} \text{ with } \|z\|_2 \ge R.
\]
We may now modify $r'$ on $\{z\in\R^{2d} \mid \|z\|_2<R \}$ as follows.
Let $\chi\in C_c^\infty(\R^{2d})$ satisfy
\[
\chi(z)=1 \quad \text{for } z\in\{z\in\R^{2d} \mid \|z\|_2\le R \},
\]
and
\[
\chi(z)=0 \quad \text{for } z\in\{z\in\R^{2d} \mid \|z\|_2\ge 2R \},
\]
and let $G\in C^\infty(\R)$ be a smooth truncation of the identity satisfying
\[
G(t)=t \quad \text{for } |t|\le \tfrac12,
\qquad
|G(t)|\le \tfrac12 \quad \text{for all } t\in\R.
\]
Define
\[
\widetilde r'(z)\coloneqq (1-\chi(z))\,r'(z)+\chi(z)\,G(r'(z)),
\qquad z\in\R^{2d}.
\]
Then
\[
\widetilde r' - r'
=
\chi\bigl(G(r')-r'\bigr)
\in C_c^\infty(\R^{2d}),
\]
and hence
\[
\widetilde r' \in \Gamma^{-2\rho}_{\rho}(\R^{2d}),
\]
since $r'\in \Gamma^{-2\rho}_{\rho}(\R^{2d})$ and
$C_c^\infty(\R^{2d})\subset \Gamma^{-\infty}(\R^{2d})
\subset \Gamma^{-2\rho}_{\rho}(\R^{2d})$.
Redefining
\[
1+\widetilde s \coloneqq b_2\,\sigma(1+\widetilde r'),
\]
we get 
\[
\widetilde s - s = b_2 \, \sigma(\widetilde r' - r')\in \Gamma^{-\infty}(\R^{2d}),
\]
since $\widetilde r'-r'\in C_c^\infty(\R^{2d})$.
Hence $\widetilde s\in\Gamma^{-\infty}(\R^{2d})$, while
$\widetilde r'\in\Gamma^{-2\rho}_{\rho}(\R^{2d})$ by construction, and the identity
\[
1+\widetilde s = b_2\,\sigma(1+\widetilde r')
\]
retains the form of \eqref{eq:long-formula-1}--\eqref{eq:long-formula-2}.
Replacing $(s,r')$ by $(\widetilde s,\widetilde r')$ in \eqref{eq:long-formula-1}--\eqref{eq:long-formula-2}, we thus preserve all structural
properties required in what follows—namely the identity, the symbol class of $r'$, and
the property of $s$—and we may therefore assume, without loss of generality, 
\[
|r'(z)|\le \tfrac12,
\qquad z\in\R^{2d}.
\]
It follows that
\[
b_2(z)=\sigma(z)^{-1}\,\frac{1+s(z)}{1+r'(z)},
\qquad z\in\R^{2d}.
\]
Moreover,
\[
(1+r')^{-1}=1+q,
\qquad
q=-\frac{r'}{1+r'}.
\]
Writing $q=r'\,H(r')$ with $H(t)\coloneqq -\frac{1}{1+t}$, and noting that $H$ is $C^\infty$ on an open neighborhood of
$[-\tfrac12,\tfrac12]$, we obtain $H(r')\in\Gamma^0_{\rho}(\R^{2d})$ and hence
\[
q\in\Gamma^{-2\rho}_{\rho}(\R^{2d}).
\]
Therefore,
\[
\frac{1+s}{1+r'} = (1+s)(1+q) = 1+\tilde r,
\qquad
\tilde r \in \Gamma^{-2\rho}_{\rho}(\R^{2d}),
\]
and consequently
\begin{equation}\label{eq:b2-symbol}
b_2=\sigma^{-1}(1+\tilde r),
\qquad \tilde r\in\Gamma^{-2\rho}_{\rho}(\R^{2d}).
\end{equation}
% }

\iffalse 
% \textcolor{red}{
Choosing $M$ sufficiently large ensures that $r'_M\in\Gamma^{-2\rho}_{\rho}(\R^{2d})$.
Moreover, by the Leibniz rule and hypoellipticity,
\[
b_2^{-1}\,\partial^\alpha b_2 \in \Gamma^{-\rho|\alpha|}_{\rho}(\R^{2d}),
\qquad
\sigma^{-1}\,\partial^\beta \sigma \in \Gamma^{-\rho|\beta|}_{\rho}(\R^{2d}),
\]
for all multi-indices $\alpha,\beta$.
\color{purple}
Consequently, the expression in parentheses in \eqref{eq:long-formula-2} is of the
form $1+r'$ with $r'\in\Gamma^{-2\rho}_{\rho}(\R^{2d})$.
Upon modifying $r'$ inside a ball of $\R^{2d}$,
one can assume without loss of generality that $|r'(z)|\leq 1/2$, for all $z\in\R^{2d}$.
We may thus write
\[
b_2(z)=\sigma(z)^{-1}\,\frac{1+s(z)}{1+r'(z)},
\qquad z\in\R^{2d}.
\]
% with $s\in\Gamma^{-\infty}(\R^{2d})$ and $r'\in\Gamma^{-2\rho}_{\rho}(\R^{2d})$.
We further have 
\[
(1+r')^{-1}
=1+ \sum_{j=1}^\infty (-r')^j
\eqcolon 1+q,
\qquad q\in\Gamma^{-2\rho}_{\rho}(\R^{2d}),
\]
where we are guaranteed that the series converges uniformly by $|r'(z)|\leq 1/2$, for all $z\in\R^{2d}$.
Therefore,
\[
\frac{1+s}{1+r'}=(1+s)(1+q)=1+\tilde r,
\qquad \tilde r \in \Gamma^{-2\rho}_{\rho}(\R^{2d}),
\]
and hence
\begin{equation}\label{eq:b2-symbol}
b_2=\sigma^{-1}(1+\tilde r),
\qquad \tilde r\in\Gamma^{-2\rho}_{\rho}(\R^{2d}).
\end{equation}
% }
%\color{black}
\fi 

\iffalse 
Choosing \(M\) sufficiently large ensures
\(r'_M\in\Gamma^{-2\rho}_{\rho}(\R^{2d})\).
Moreover, by the Leibniz rule and hypoellipticity,
\[
b_2^{-1}\,\partial^\alpha b_2 \in \Gamma^{-\rho|\alpha|}_{\rho}(\R^{2d}),
\qquad
\sigma^{-1}\,\partial^\beta \sigma \in \Gamma^{-\rho|\beta|}_{\rho}(\R^{2d}),
\]
for all multi-indices \(\alpha,\beta\). It follows that the expression in parentheses in \eqref{eq:long-formula-2} is of the
form \(1+r'\) with \(r'\in\Gamma^{-2\rho}_{\rho}(\R^{2d})\).
\textcolor{blue}{Therefore,}
\[
b_2(z)=\sigma(z)^{-1}\,\frac{1+s(z)}{1+r'(z)},
\textcolor{blue}{\quad z\in\R^{2d},}
\]
with \(s\in\Gamma^{-\infty}(\R^{2d})\) and \(r'\in\Gamma^{-2\rho}_{\rho}(\R^{2d})\),
\textcolor{blue}{and, since} the quotient \((1+s)/(1+r')\) can be written as \(1+\tilde r\) for some
\(\tilde r\in\Gamma^{-2\rho}_{\rho}(\R^{2d})\)\textcolor{blue}{,
it} follows that
\begin{equation}\label{eq:b2-symbol}
b_2=\sigma^{-1}(1+\tilde r),
\qquad \tilde r\in\Gamma^{-2\rho}_{\rho}(\R^{2d}).
\end{equation}
\fi 


Applying the same arguments to the operator
\(B_2(T_1-T_2)\) shows that its Weyl symbol is of the form
\begin{equation}\label{eq:b2-symbol-2}
    b_2(\sigma_1-\sigma)(1+r''),
\end{equation}
for some \(r''\in\Gamma^{-2\rho}_{\rho}(\R^{2d})\). Using \eqref{eq:b2-symbol} in \eqref{eq:b2-symbol-2} along with
\(\sigma^{-1}(\sigma_1-\sigma)\in\Gamma^{-2\rho}_{\rho}(\R^{2d})\), we obtain
% \iffalse 
% and the equality is a
% Substituting the above expression for \(b_2\) yields
% \[
% \sigma^{-1}(\sigma_1-\sigma)(1+\tilde{r})(1+r'').
% \]
% From the previous step we already know that
% \(\sigma^{-1}(\sigma_1-\sigma)\in\Gamma^{-2\rho}_{\rho}(\R^{2d})\).
% Since \(\Gamma^{-2\rho}_{\rho}\) is closed under multiplication by symbols of the
% form \(1+r\) with \(r\in\Gamma^{-2\rho}_{\rho}\), we conclude that
% \fi 
\[
b_2(\sigma_1-\sigma)(1+r'')=\sigma^{-1}(\sigma_1-\sigma)(1+\tilde{r})(1+r'') 
\in \Gamma^{-2\rho}_{\rho}(\R^{2d})
\]
and hence
\[
B_2(T_1-T_2)\in\Psi^{-2\rho}_{\rho}(\R^d).
\]
This establishes the required property for the operator
\(K = B_2(T_1 - T_2)\), 
and thus completes the proof of
the invariance under change of quantization stated in \eqref{eq:equiv-invariance-quant}.
% Lemma~\ref{lem:technical-lemma}.

% \iffalse 
% Coming back to the parametrix equation $T_2B_2=I+S$,
% we get, 
% from \cite[Theorem~23.6]{shubin_pseudodifferential_2001} and the hypoellipticity of both $b_2$ and $\sigma$,
% that 
% there exist a constant $C_{\alpha, \beta }>0$ and a symbol $r'_M \in \Gamma^{m_+-m_--2(M+1)\rho}_{\rho}(\R^{2d})$ for which
% \begin{align}
%     1+s(z)
%     &= \sum_{\substack{\alpha, \beta \in \N^d\\0\leq |\alpha|+ |\beta|\leq M}}
%     C_{\alpha, \beta }  \left({\partial^\alpha_z b_2(z)}\right)\left({\partial^\beta_z \sigma(z)} \right)
%     +r'_M(z) \label{eq:long-formula-1}\\
%     &= b_2 (z) \, \sigma (z)
%     \left(
%         1+ \sum_{\substack{\alpha, \beta \in \N^d\\1\leq |\alpha|+ |\beta|\leq M}}
%         C_{\alpha, \beta }  \frac{\partial^\alpha_z b_2(z)}{b_2(z)}\frac{\partial^\beta_z \sigma(z)}{\sigma(z)} +\frac{r'_M(z)}{b_2 (z) \, \sigma (z)}
%     \right),\label{eq:long-formula-2}
% \end{align}
% for all $z\in \R^{2d}$, where $s$ is the Weyl symbol of the operator $S$.
% Again, $r'_M \in \Gamma^{-2\rho}_{\rho}(\R^{2d})$ for $M$ large enough, and the Leibniz rule together with the hypoellipticity of $b_2$, resp. of $\sigma$, gives 
% $b_2^{-1}\partial^\alpha b_2 \in \Gamma^{-\rho|\alpha|}_{\rho}(\R^{2d})$, for all $\alpha \in \N^d$,
% resp. $\sigma^{-1} \partial^\beta \sigma \in \Gamma^{-\rho|\beta|}_{\rho}(\R^{2d})$, for all $\beta \in \N^d$.
% In particular, this shows that there exists 
% $r' \in \Gamma^{-2\rho}_{\rho}(\R^{2d})$ such that 
% $b_2 = \sigma^{-1} (1+r')$.
% A similar line of arguments as the one following \eqref{eq:long-formula-1}--\eqref{eq:long-formula-2} shows that the symbol of 
% $B_2(T_1-T_2)$ is of the form
% $b_2(\sigma_1-\sigma) (1+r'')$,
% and, therefore, of the form
% $\sigma^{-1}(\sigma-\sigma_2) (1+r') (1+r'')$.
% We have already seen in the previous paragraph that $\sigma^{-1}(\sigma-\sigma_2) \in \Gamma^{-2\rho}_{\rho}(\R^{2d})$.
% % Using \cite[Lemma~25.1]{shubin_pseudodifferential_2001},
% % it 
% It implies that $\sigma^{-1}(\sigma-\sigma_2) (1+r') (1+r'') \in \Gamma^{-2\rho}_{\rho}(\R^{2d})$
% and thus $B_2(T_1-T_2)\in \Psi^{-2\rho}_{\rho}(\R^d)$, which is the desired result.
% \fi 

% \textcolor{red}{
\textbf{Invertibility formula.}
We now prove the second statement of Theorem~\ref{thm:main-result}, namely
\[
H_{T_{\sigma^{-1}}^{-1}}(\varepsilon)\sim H_\sigma(\varepsilon),
\qquad \varepsilon\to0.
\]
To this end, we establish a parametrix relation between $T_\sigma$ and
$T_{\sigma^{-1}}$, showing that their composition differs from the identity by
a compact operator of strictly negative order.
Using the assumed invertibility of $T_{\sigma^{-1}}$, this relation allows us
to compare the entropy of $T_\sigma$ with that of $T_{\sigma^{-1}}^{-1}$.
% Recall that $\sigma\in\hyponegclass(\R^{2d})$ means
% $\sigma\in \mathrm{H}\Gamma^{-m_+,-m_-}_\rho(\R^{2d})$
% for some $\rho\in(0,1]$ and $m_+\ge m_->0$.
% By Lemma~25.1(1a) in \cite{shubin_pseudodifferential_2001},
% it follows that
% \[
% \sigma^{-1}\in \mathrm{H}\Gamma^{m_-,m_+}_\rho(\R^{2d}).
% \]
Applying the Weyl composition formula
\cite[Theorem~23.6]{shubin_pseudodifferential_2001}
to $\sigma^{-1}$ and $\sigma$ and following a similar line of arguments as the one leading to \eqref{eq:b2-symbol}, shows that
\iffalse 
\[
\sigma^{-1}\#\sigma
=
1+r,
\qquad
r\in \Gamma^{-2\rho}_\rho(\R^{2d}),
\]
so that
\[
T_{\sigma^{-1}}T_\sigma
=
I + T_r,
\]
where $T_r$ is a pseudodifferential operator of order $-2\rho$.
Since $2\rho>0$, $T_r$ is compact on $L^2(\R^d)$.
\textcolor{red}
\textbf{Invertibility formula.}
We now prove the second statement of Theorem~\ref{thm:main-result}, namely
\[
H_{T_{\sigma^{-1}}^{-1}}(\varepsilon)\sim H_\sigma(\varepsilon),
\qquad \varepsilon\to0.
\]
To this end, we establish a parametrix relation between \(T_\sigma\) and
\(T_{\sigma^{-1}}\), showing that their composition differs from the identity by
a compact operator of strictly lower order. Using the assumed invertibility of
\(T_{\sigma^{-1}}\), this parametrix relation allows us to relate the entropy of
\(T_\sigma\) to that of \(T_{\sigma^{-1}}^{-1}\).
\color{blue}
To make this precise, recall that $\sigma\in\hyponegclass(\R^{2d})$ means
$\sigma\in \mathrm{H}\Gamma^{-m_+,-m_-}_\rho(\R^{2d})$ for some
$\rho\in(0,1]$ and $m_+\ge m_->0$,
% Since $\sigma$ is strictly positive, Lemma~25.1(1a) in
% \cite{shubin_pseudodifferential_2001} implies that $\sigma^{-1}$ satisfies the
% corresponding hypoelliptic estimates outside a sufficiently large ball.
% On the complementary compact region, $\sigma^{-1}$ is smooth and bounded together
% with all its derivatives, and the hypoelliptic bounds extend to the whole of
% $\R^{2d}$ after adjusting constants.
% Consequently,
% \[
% \sigma^{-1}\in \mathrm{H}\Gamma^{m_-,m_+}_\rho(\R^{2d}).
% \]
so that  Lemma~25.1(1a) in
\cite{shubin_pseudodifferential_2001} implies 
$\sigma^{-1}\in \mathrm{H}\Gamma^{m_-,m_+}_\rho(\R^{2d})$.
Following a similar line of arguments as the one leading to \eqref{eq:b2-symbol},
we further get that the operator $T_{\symbolnot^{-1}}T_{\symbolnot }$ as Weyl symbol 
% We can now apply the Weyl composition formula
% \cite[Theorem~23.6]{shubin_pseudodifferential_2001} to the symbols
% $\sigma^{-1}$ and $\sigma$ to conclude that the product
% $T_{\sigma^{-1}}T_\sigma$ has Weyl symbol
\[
\sigma^{-1}\sigma\,(1+r)=1+r,
\]
with a remainder $r\in\Gamma^{-2\rho}_\rho(\R^{2d})$.
\fi 
%Consequently, there exists an operator
%\(R\in\Psi^{-2\rho}_\rho(\R^{d})\) such that
\begin{equation}\label{eq:prod-operators-identity-comp}
    T_{\sigma^{-1}}T_\sigma
    =
    I+R,
\end{equation}
with \(R\in\Psi^{-2\rho}_\rho(\R^{d})\).
% }

Since \(T_{\sigma^{-1}}\) is invertible by assumption, \eqref{eq:prod-operators-identity-comp}
can be rewritten as
\begin{equation}\label{eq:introduction-of-R2}
    T_\sigma
    =
    T_{\sigma^{-1}}^{-1}(I+R).
\end{equation}
By \cite[Theorem~24.4]{shubin_pseudodifferential_2001}, the operator
\(R\in\Psi^{-2\rho}_\rho(\R^{d})\) extends to a compact operator on \(L^2(\R^d)\).
Identity \eqref{eq:introduction-of-R2} is therefore of the form required in
Lemma~\ref{lem:technical-lemma}, with
\(T_1=T_\sigma\), \(T_2=T_{\sigma^{-1}}^{-1}\),
\(K=R\), and \(K_{-\infty}=0\).
Applying the lemma yields
\[
    H_{T_{\sigma^{-1}}^{-1}}(\varepsilon)
    \sim
    H_\sigma(\varepsilon),
    \qquad \varepsilon\to0,
\]
which is precisely \eqref{eq:main-result-part2}.

% \iffalse
% \textbf{Invertibility formula.}
% % We take $\lambda_* \geq 0$ outside the spectrum of $T_{\symbolnot^{-1}}$ and large enough for the operator $T_{\symbolnot^{-1}}+\lambda_* I$ to be strictly positive---this is possible by lower-boundedness of the spectrum of $T_{\symbolnot^{-1}}$.
% % %, which follows from \eqref{eq:application-main-result-proof}. 
% % (Recall from the discussion preceding Theorem~\ref{thm-shubin} that the spectrum of $T_{\symbolnot^{-1}}$ is discrete is equal to the set of its eigenvalues.)
% % In particular,  $T_{\symbolnot^{-1}}+\lambda_* I$ has Weyl symbol $\symbolnot_*\coloneqq \symbolnot^{-1}+ \lambda_*$ and is invertible; we denote this inverse by $T_*=(T_{\symbolnot^{-1}}+\lambda_* I)^{-1}$.
% % In the first step, we argue that 
% The operator $T_{\symbolnot^{-1}}$ is an almost-inverse of $T_{\symbolnot}$.
% Indeed, we get,
% \textcolor{blue}{from a similar line of arguments as the one following \eqref{eq:long-formula-1}--\eqref{eq:long-formula-2},}
% that the Weyl symbol $\tilde \symbolnot$ of the operator $T_{\symbolnot^{-1}}T_{\symbolnot }$ can be decomposed according to
% \begin{equation*}
%     \tilde \symbolnot 
%     =  \symbolnot^{-1} \symbolnot (1 + r) 
%     =  1+r,
% \end{equation*}
% for some $r\in \Gamma^{-2\rho}_\rho(\R^{2d})$.
% In particular, there exists an operator $R\in \Psi_{\rho}^{-2\rho}(\R^{d})$ satisfying
% \begin{equation}\label{eq:prod-operators-identity-comp}
%     T_{\symbolnot^{-1}}T_{\symbolnot}
%     = I + R,
% \end{equation}
% and, therefore, recalling that $T_{\symbolnot^{-1}}$ is invertible, we have
% \begin{equation}\label{eq:introduction-of-R2}
%     T_{\symbolnot}
%     =  T_{\symbolnot^{-1}}^{-1} (I + R).
% \end{equation}
% Upon application of \cite[Theorem~24.4]{shubin_pseudodifferential_2001}, we see that the operator $R$ is compact and from $L^2(\R^d)$ to itself.
% \textcolor{blue}{
%     We are left to apply Lemma~\ref{lem:technical-lemma} 
%     with $T_1=T_{\symbolnot}$, $T_2=T_{\symbolnot^{-1}}^{-1}$, $K=R$, and $K_{-\infty} = 0$,
%     whose assumption is satisfied by \eqref{eq:introduction-of-R2},
%     to obtain \eqref{eq:main-result-part2}.
% }
% \fi 

%%%T9: I wonder whether we couldn't use Lemma 5.1 to replace the arguments above on membership of sigma^-1 in HF, these arguments
%needed a split-up between outside a ball and inside a ball, Lemma 5.1 would simplify the arguments, could you please look into this.

\begin{proof}[Proof of Lemma~\ref{lem:technical-lemma}]
Let $k$ denote the Weyl symbol of $K$.
Since $K\in\Psi^{-m}_{\rho_1}(\R^d)$ for some $m>0$, we have
$k\in\Gamma^{-m}_{\rho_1}(\R^{2d})$, hence $k$ is of strictly negative order.
The constant symbol $1$ belongs to $\mathrm{H}\Gamma^{0,0}_{\rho_1}(\R^{2d})$.
Applying Shubin’s stability result for hypoelliptic symbols under lower-order
additive perturbations \cite[Lemma~25.1 (1c)]{shubin_pseudodifferential_2001} 
% with $\sigma\equiv 1$ and $r=k$ % <- here \sigma could be confused with the \sigma in the statement of the Lemma
yields 
\[
1+k \in \mathrm{H}\Gamma^{0,0}_{\rho_1}(\R^{2d}).
\]
Consequently, by \cite[Theorem~25.1]{shubin_pseudodifferential_2001},
the operator $I+K$ admits a right parametrix, that is, there exists a
pseudodifferential operator $P\in\Psi^0_{\rho_1}(\R^d)$ such that
\begin{equation}\label{eq:parametrix-formula}
    (I+K)P = I + R_{-\infty},
\end{equation}
for some residual operator $R_{-\infty}\in\Psi^{-\infty}(\R^d)$.
By the Weyl composition formula \cite[Theorem~23.6]{shubin_pseudodifferential_2001},
the Weyl symbol of $(I+K)P$ has the form
\[
(1+k)p + q,
\]
where $p$ is the Weyl symbol of $P$ and
$q\in\Gamma^{-2\rho_1}_{\rho_1}(\R^{2d})$.
On the other hand, by \eqref{eq:parametrix-formula}, the Weyl symbol of $(I+K)P$
is $1+r_{-\infty}$, where $r_{-\infty}$ is the symbol of $R_{-\infty}$.
Using that $k\in\Gamma^{-m}_{\rho_1}(\R^{2d})$ and comparing both expressions,
it follows that there exists a symbol
\[
r_1 \in \Gamma^{-\gamma}_{\rho_1}(\R^{2d}),
\qquad
\gamma \coloneqq \min\{2\rho_1,m\},
\]
such that
\[
p = 1 + r_1.
\]
Equivalently, in operator form,
\[
P = I + R_1,
\qquad
R_1 \in \Psi^{-\gamma}_{\rho_1}(\R^d).
\]
We now multiply the identity
\[
T_1 = T_2 (I+K) + K_{-\infty}
\]
from the right by $P$ and use \eqref{eq:parametrix-formula} to obtain
\begin{align*}
T_1 P
&= T_2 (I+K)P + K_{-\infty}P \\
&= T_2 (I+R_{-\infty}) + K_{-\infty}P \\
&= T_2 + T_2 R_{-\infty} + K_{-\infty}P.
\end{align*}
Rearranging terms yields
\begin{equation}\label{eq:introduction-of-R1}
    T_2
    =
    T_1 (I+R_1) + R'_{-\infty},
\end{equation}
where
\[
R'_{-\infty} \coloneqq -T_2 R_{-\infty} - K_{-\infty}P
\in \Psi^{-\infty}(\R^d),
\]
since $\Psi^{-\infty}(\R^d)$ is a two-sided ideal in the pseudodifferential
calculus. Equation \eqref{eq:introduction-of-R1} provides a representation of $T_2$
of the same structural form as the hypothesis of
Lemma~\ref{lem:technical-lemma}, namely
\[
T_2 = T_1 (I+R_1) + R'_{-\infty},
\]
with $R_1\in\Psi^{-\gamma}_{\rho_1}(\R^d)$ for some $\gamma>0$ and
$R'_{-\infty}\in\Psi^{-\infty}(\R^d)$.
Although the order $\gamma$ may differ from the original exponent $m$,
the lemma only requires the perturbation to be of strictly negative order.
This symmetry explains why Lemma~\ref{lem:technical-lemma} only requires that
either $T_1$ or $T_2$ be the Weyl quantization of~$\sigma$.
%Accordingly, the lemma may be proved assuming the representation
%$T_1 = T_2(I+K)+K_{-\infty}$, without loss of generality.




% \iffalse 
% \begin{proof}[Proof of Lemma~\ref{lem:technical-lemma}]
    

% \textcolor{blue}{By \cite[Lemma~5.1 f)]{shubin_pseudodifferential_2001},
% we know that $1+k \in \textrm{H}\Gamma^{0,0}_{\rho_1}(\R^{2d})$, where $k$ is the Weyl symbol of $K$.
% In particular, 
% \cite[Theorem~25.1]{shubin_pseudodifferential_2001} guarantees that
% the operator $I+K$ admits a right-parametrix, that is, a pseudodifferential operator $P\in\Psi^0_{\rho_1}(\R^d)$ satisfying
% \begin{equation}\label{eq:parametrix-formula}
%     (I+K)P = I + R_{-\infty},
% \end{equation}
% for some residual operator $R_{-\infty}\in\Psi^{-\infty}(\R^d)$.
% By the composition formula  \cite[Theorem 23.6]{shubin_pseudodifferential_2001}, 
% the Weyl symbol of $(I+K)P$ is of the form $(1+k)p + q$, 
% where $p$ is the Weyl symbol of $P$ and $q\in \Gamma^{-2\rho_1}_{\rho_1}(\R^{2d})$.
% But, by \eqref{eq:parametrix-formula}, the Weyl symbol of  $(I+K)P$ is also equal to $1+r_{-\infty}$, where $r_{-\infty}$ is the symbol of $R_{-\infty}$.
% Consequently, using $k\in \Gamma^{-m}_{\rho_1}(\R^{2d})$, $q\in \Gamma^{-2\rho_1}_{\rho_1}(\R^{2d})$, and $(1+k)p + q=1+r_{-\infty}$,
% we get that there exists a symbol $r_1\in \Gamma^{-\gamma}_{\rho_1}(\R^{2d})$, 
% with $\gamma \coloneqq \min\{2\rho_1, m\}$,
% such that $p = 1 + r_1$. 
% Equivalently in terms of operators, $P=I+R_1$, for some operator $R_1\in \Psi_{\rho_1}^{-\gamma}(\R^{d})$.
% % , with $\alpha\coloneqq \min\{2\rho, m_-\}$.
% Multiplying both sides of the identity $T_1 =  T_2 (I + K) + K_{-\infty}$ by $P$ from the right  and using \eqref{eq:parametrix-formula}, we get
% \begin{equation*}
%     T_1P
%     = T_2 (I + K) P + K_{-\infty} P
%     =  T_2 (I+R_{-\infty}) + K_{-\infty} P
%     = T_2 + T_2R_{-\infty} + K_{-\infty} P,
% \end{equation*}
% or, equivalently, 
% \begin{equation}\label{eq:introduction-of-R1}
%     T_2
%     =  T_1 (I + R_1) + R'_{-\infty},
% \end{equation}
% with $R'_{-\infty}\in\Psi^{-\infty}(\R^d)$. 
% % since it is the composition of elements in $\Psi^{-\infty}(\R^d)$ and $\Psi^{m_+}(\R^d)$.
% Therefore, without loss of generality, the roles of $T_1$ and $T_2$ may be interchanged in the statement of Lemma~\ref{lem:technical-lemma}.
% % , and we may assume $T_1=T_\sigma$ in what follows.
% }
% \fi 

We now turn to the comparison of the entropy numbers of $T_1$ and $T_2$.
Recall that the entropy numbers $\{e_n(T)\}_{n\in\N}$ of a compact operator
$T\colon L^2(\R^d)\to L^2(\R^d)$ are defined by
\begin{equation}\label{eq:def-entropy-nbs}
    e_n(T)
    \coloneqq
    \inf\bigl\{ \varepsilon>0 \,\big|\, H_T(\varepsilon)\le n \bigr\},
    \qquad n\in\N.
\end{equation}
We refer to
\cite{opideals1980pietsch,bookeigen1987pietsch,carlEntropyCompactnessApproximation1990,lorentzConstructiveApproximationAdvanced1996,edmundsFunctionSpacesEntropy1996}
for background on entropy numbers and their basic properties. Starting from the representation \eqref{eq:introduction-of-R1}, and using the
submultiplicativity and subadditivity of entropy numbers 
% together with their
% monotonicity in the index 
(see, for example,
\cite[Section~15.7]{lorentzConstructiveApproximationAdvanced1996} or
\cite[Section~1.3.1, Lemma~1% 
% (ii)--(iii)
]{edmundsFunctionSpacesEntropy1996}),
we obtain
\begin{equation}\label{eq:entropy-nb-arg-1}
    e_{\,n+5\lfloor n^{1/2}\rfloor}(T_2)
    \le
    e_n(T_1)\,e_{\lfloor n^{1/2}\rfloor}(I+R_1)
    +
    e_{4\lfloor n^{1/2}\rfloor}(R'_{-\infty}),
\end{equation}
together with
\begin{equation}\label{eq:entropy-nb-arg-2}
    e_{\lfloor n^{1/2}\rfloor}(I+R_1)
    \le
    e_{0}(I) + e_{\lfloor n^{1/2}\rfloor}(R_1),
    \qquad n\in\N.
\end{equation}
%%%T9[OK]: pls. check in the sentence above if we also refer to the monotonicity result in Lorentz or Edmunds, I included a reference to monotonicty.
Since $e_{0}(I)=\|I\|_{2\to2}=1$ 
% \cite[Section~1.3.1, Lemma~1(i)]{edmundsFunctionSpacesEntropy1996}
\cite[Section~15.7]{lorentzConstructiveApproximationAdvanced1996},
and $R_1$ by virtue of being compact satisfies
\[
\lim_{{j}\to\infty} e_{j}(R_1)=0,
\]
it follows, 
% by subadditivity of entropy numbers, 
by \eqref{eq:entropy-nb-arg-2},
that
\begin{equation}\label{eq:bound-entropy-identity}
    e_{\lfloor n^{1/2}\rfloor}(I+R_1)
    \le
    1+o(1),
    \qquad n\to\infty.
\end{equation}


% \iffalse

% We next establish a relationship between the entropy numbers of the operators $T_1$ and $T_2$.
% Recall that the entropy numbers $\{e_n(T)\}_{n\in\N}$ of a compact operator $T$ are defined as
% \begin{equation}\label{eq:def-entropy-nbs}
%     e_n(T)
%     \coloneqq \inf_{\varepsilon>0} \left\{ \varepsilon \mid H_T\left(\varepsilon \right) \leq n \right\},
%     \quad \text{for all } n\in \N.
% \end{equation}
% We refer to \cite{opideals1980pietsch,bookeigen1987pietsch,carlEntropyCompactnessApproximation1990,lorentzConstructiveApproximationAdvanced1996,edmundsFunctionSpacesEntropy1996} for overviews of the theory of entropy numbers.
% Using the sub-multiplicativity and sub-additivity of entropy numbers (see, e.g., \cite[Section~15.7]{lorentzConstructiveApproximationAdvanced1996} or \cite[Section~1.3.1, Lemma~1 (ii) and (iii)]{edmundsFunctionSpacesEntropy1996}), we get from \eqref{eq:introduction-of-R1} that
% \begin{equation}\label{eq:entropy-nb-arg-1}
%     e_{n+{3\lfloor {n}^{1/2} \rfloor}} \left(T_2\right)
%     \leq e_{n} \left(T_1\right) \,  e_{\lfloor {n}^{1/2} \rfloor} (I+R_1) + e_{2\lfloor {n}^{1/2} \rfloor}\left(R'_{-\infty}\right),
% \end{equation}
% and 
% \begin{equation}\label{eq:entropy-nb-arg-2}
%      e_{\lfloor {n}^{1/2} \rfloor} (I+R_1) 
%      \leq e_1 (I) + e_{\lfloor {n}^{1/2} \rfloor} (R_1),
% \end{equation}
% for all $n\in \N$.
% We observe that \cite[Section~1.3.1, Lemma~1 (i)]{edmundsFunctionSpacesEntropy1996} yields $e_1 (I)=\|I\|_{2\rightarrow 2}=1$
% and the compactness of $R_1$ implies $e_{\lfloor {n}^{1/2} \rfloor} (R_1) \to 0$, as $k\to\infty$ (indeed, a compact operator has finite entropy so that \eqref{eq:def-entropy-nbs} implies that its covering numbers go to zero when the index tends to infinity).
% \fi 


% Then, recalling that $T_*$ is invertible and applying once more the sub-multiplicativity of entropy numbers, we get 
% \begin{equation}\label{eq:entropy-nb-arg-3}
%     e_{(m+1) \lfloor {k}^{1/2} \rfloor} \left(R'_{-\infty}\right)
%     \leq e_{\lfloor {k}^{1/2} \rfloor} \left(T_*^m\right) \, e_{m\lfloor {k}^{1/2} \rfloor} \left(T_*^{-m} R'_{-\infty}\right).
% \end{equation}
% By \cite[Theorem~23.6]{shubin_pseudodifferential_2001},
% the operator $T_*^{-m} R'_{-\infty}$ belongs to $\Psi^{-\infty}(\R^d)$ since $T_*\in\Psi^{m_+}(\R^d)$ and $R'_{-\infty}\in\Psi^{-\infty}(\R^d)$.
% In particular, $T_*^{-m} R'_{-\infty}$ is compact operator hence we have  $e_{\lfloor {k}^{1/2} \rfloor} (T_*^{-m} R'_{-\infty}) \to 0$, as $k\to\infty$.
% Moreover, applying $m$ times the sub-multiplicativity of entropy numbers, yields \textcolor{blue}{EDIT ARGUMENT + CHANGE NOTATION $m$}
% \begin{equation}\label{eq:entropy-nb-arg-4}
%     e_{m\lfloor {k}^{1/2} \rfloor} \left(T_*^m\right)
%     \leq e_{\lfloor {k}^{1/2} \rfloor} \left(T_*\right)^m.
% \end{equation}
% Now, one can verify, from 
% % \eqref{eq:asymp-entropy-inverse-operator} and the  
% $\symbolnot \in \hyposymbclassinv(\R^{2d})$, 
% that the entropy of $T_*$ grows at most as $\varepsilon^{-d/m_-}$ and at least as $\varepsilon^{-d/m_+}$.
% (To see this, plug the bound \eqref{eq:bound-def-hypoellip} in the definition \eqref{eq:definition-volume-function} of $V_{\symbolnot}$.)
% This, in turn, implies 
% \begin{equation}\label{eq:entropy-nb-arg-5}
%     C_1 p^{-m_+/d}
%     \leq e_{p} \left(T_*\right)
%     \leq C_2 p^{-m_-/d},
%     \quad \text{for all } p\in \N,
% \end{equation}
% for some constants $C_1, C_2>0$.
% Consequently, choosing $m>2m_+/m_-$, we have  
% \begin{equation}\label{eq:entropy-nb-arg-6}
%     \left(e_{\lfloor {k}^{1/2} \rfloor} \left(T_*\right)\right)^m
%     = o_{k\to\infty}\left(e_{k+{(m+2)\lfloor {k}^{1/2} \rfloor}} \left(T_*\right)\right).
% \end{equation}

%%%T9[OK]: I think the domain here should be R^2d.






We now estimate the term
% \textcolor{red}{
$e_{4\lfloor n^{1/2}\rfloor}\!\left(R'_{-\infty}\right)$
% }
appearing in \eqref{eq:entropy-nb-arg-1}.
To this end, arbitrarily fix $s>0$ and consider the hypoelliptic symbol
$\bar\sigma\in \mathrm{H}\Gamma^{-s,-s}_{\rho_1}(\R^{2d})$
defined according to $\bar\sigma(z) = (1+\|z\|_2^2)^{-s/2}$, for all $z\in\R^{2d}$.
Let $T_{\bar\sigma}$ denote the Weyl quantization of $\bar\sigma$ and let $Q\in\Psi^{s}_{\rho_1}(\R^d)$ be a right parametrix, that is, 
\begin{equation*}
    T_{\bar\sigma} Q = I + R''_{-\infty}, 
    \quad  R''_{-\infty} \in \Psi^{-\infty}(\R^d).
\end{equation*}




One verifies that 
\begin{equation}\label{eq:proof-lemma-AAA}
    V_{\bar\sigma}(\lambda)
    =  \int_{\R^{2d}} \mathbbm{1}_{\{(1+\|z\|_2^2)^{-s/2}>\lambda\}}  \,dz
    \sim \omega_{2d}\,\lambda^{-\frac{2d}{s}}, 
    \quad  \lambda \to 0.
\end{equation}
Applying the entropy--volume relation \eqref{eq:main-res-volume-integral},
already established for Weyl quantization in \eqref{eq:mainr-result-for-Weyl-1}, to the operator
$T_{\bar\sigma}$ yields
\[
H_{T_{\bar\sigma}}(\varepsilon)
\sim
\int_\varepsilon^\infty \frac{V_{\bar\sigma}(\lambda)}{\lambda}\,d\lambda
\sim \frac{s \, \omega_{2d}}{2d}\varepsilon^{-\frac{2d}{s}},
\qquad \varepsilon\to0.
\]
By the definition of the entropy numbers \eqref{eq:def-entropy-nbs}, it follows that
\begin{equation}\label{eq:proof-lemma-AAB}
    e_j(T_{\bar\sigma})
    \sim
    \left(\frac{2d\,j}{s\,\omega_{2d}}\right)^{-\frac{s}{2d}},
    \qquad j\to\infty.
\end{equation}
Evaluating \eqref{eq:proof-lemma-AAB} at \(j=\lfloor n^{1/2}\rfloor\) yields
\begin{equation}\label{eq:poly-scaling-ent-nb}
    e_{\lfloor n^{1/2}\rfloor}(T_{\bar\sigma})
    \le
    C\,n^{-\frac{s}{4d}}\bigl(1+o(1)\bigr),
    \qquad n\to\infty,
\end{equation}
for a suitable constant \(C>0\).


\iffalse
By definition of the entropy numbers \eqref{eq:def-entropy-nbs}, this implies that
\begin{equation}\label{eq:proof-lemma-AAB}
    e_{\textcolor{blue}{j}}(T_{\bar\sigma})
\sim \left(\frac{2d{\textcolor{blue}{j}}}{s  \omega_{2d}}\right)^{-\frac{s}{2d}},
\qquad {\textcolor{blue}{j}}\to\infty.
\end{equation}
Evaluating this estimate at ${\textcolor{blue}{j}}=\lfloor n^{1/2}\rfloor$ yields bound
\begin{equation}\label{eq:poly-scaling-ent-nb}
    e_{\lfloor n^{1/2}\rfloor}(T_{\bar\sigma})
    \le
    C\,n^{-\frac{s}{4d}}\bigl(1+o(1)\bigr),
    \qquad n\to\infty,
\end{equation}
for a suitable constants $C>0$.
\fi 

We now return to the operator $R'_{-\infty}$.
Using the factorization
\[
R'_{-\infty} = \left(T_{\bar\sigma}Q - R''_{-\infty}\right)R'_{-\infty},
\]
and applying the subadditivity and submultiplicativity of entropy numbers, we obtain
\begin{align}
e_{4\lfloor n^{1/2}\rfloor}\!\left(R'_{-\infty}\right)
&\le
e_{\lfloor n^{1/2}\rfloor}(T_{\bar\sigma})\,
e_{\lfloor n^{1/2}\rfloor}\!\left(Q R'_{-\infty}\right)
+
e_{\lfloor n^{1/2}\rfloor}\!\left(R''_{-\infty}\right)
e_{\lfloor n^{1/2}\rfloor}\!\left(R'_{-\infty}\right)
\label{eq:bound-ent-split-1}\\
&\le
C\,n^{-\frac{s}{4d}}\bigl(1+o(1)\bigr)\,
e_{\lfloor n^{1/2}\rfloor}\!\left(Q R'_{-\infty}\right)
+
e_{\lfloor n^{1/2}\rfloor}\!\left(R''_{-\infty}\right)
e_{\lfloor n^{1/2}\rfloor}\!\left(R'_{-\infty}\right),
\label{eq:bound-ent-split-2}
\end{align}
where \eqref{eq:bound-ent-split-2} follows from \eqref{eq:poly-scaling-ent-nb}.
Since $Q\in\Psi^{s}_{\rho_1}(\R^d)$ and $R'_{-\infty}\in\Psi^{-\infty}(\R^d)$,
their composition $Q R'_{-\infty}$ belongs to $\Psi^{-\infty}(\R^d)$ and is therefore compact on $L^2(\R^d)$.
Moreover, $R''_{-\infty}\in\Psi^{-\infty}(\R^d)$ is also compact.
Consequently,
\[
\lim_{n\to\infty}
e_{\lfloor n^{1/2}\rfloor}\!\left(Q R'_{-\infty}\right)
=
0,
\qquad
\lim_{n\to\infty}
e_{\lfloor n^{1/2}\rfloor}\!\left(R''_{-\infty}\right)
=
0.
\]
Combining these limits with \eqref{eq:bound-ent-split-1}–\eqref{eq:bound-ent-split-2} yields
\begin{equation}\label{eq:R-fast-decay}
e_{4\lfloor n^{1/2}\rfloor}\!\left(R'_{-\infty}\right)
=
o \left(n^{-\frac{s}{4d}} + e_{\lfloor n^{1/2}\rfloor}\!\left(R'_{-\infty}\right)\right),
\qquad n\to\infty.
\end{equation}
\iffalse 
We now return to the operator $R'_{-\infty}$.
Using the factorization
\[
R'_{-\infty} = \left(T_{\bar\sigma}Q -R''_{-\infty} \right)R'_{-\infty},
\]
together with the subadditivity and submultiplicativity of entropy numbers, we obtain
\begin{align}
e_{4\lfloor n^{1/2}\rfloor}\!\left(R'_{-\infty}\right)
&\le
e_{\lfloor n^{1/2}\rfloor}(T_{\bar\sigma})\,
e_{\lfloor n^{1/2}\rfloor}\!\left(Q R'_{-\infty}\right)
+ e_{\lfloor n^{1/2}\rfloor}\!\left(R''_{-\infty}\right)e_{\lfloor n^{1/2}\rfloor}\!\left(R'_{-\infty}\right)
\label{eq:bound-ent-split-1}\\
&\le
C\,n^{-\frac{s}{4d}}\bigl(1+o(1)\bigr)\,
e_{\lfloor n^{1/2}\rfloor}\!\left(Q R'_{-\infty}\right)
+ e_{\lfloor n^{1/2}\rfloor}\!\left(R''_{-\infty}\right)e_{\lfloor n^{1/2}\rfloor}\!\left(R'_{-\infty}\right),
\label{eq:bound-ent-split-2}
\end{align}
where \eqref{eq:bound-ent-split-2} uses \eqref{eq:poly-scaling-ent-nb}. Since
$Q\in\Psi^{s}_{\rho_1}(\R^d)$ and $R'_{-\infty}\in\Psi^{-\infty}(\R^d)$,
their composition $Q R'_{-\infty}$ belongs to $\Psi^{-\infty}(\R^d)$
and is therefore compact on $L^2(\R^d)$.
Additionally, $R'_{-\infty}\in\Psi^{-\infty}(\R^d)$ is also compact.
Hence
\[
\lim_{n\to\infty}
e_{\lfloor n^{1/2}\rfloor}\!\left(Q R'_{-\infty}\right)
=
0
\quad \text{and} \quad 
\lim_{n\to\infty}
e_{\lfloor n^{1/2}\rfloor}\!\left(R''_{-\infty}\right)
=
0,
\]
and combining this with
\eqref{eq:bound-ent-split-1}–\eqref{eq:bound-ent-split-2} yields
\begin{equation}\label{eq:R-fast-decay}
e_{4\lfloor n^{1/2}\rfloor}\!\left(R'_{-\infty}\right)
=
o\left(n^{-\frac{s}{4d}}\right),
\qquad n\to\infty.
\end{equation}
\fi 
% \textcolor{red}{
Iterating \eqref{eq:R-fast-decay} and using that $s>0$ is arbitrary yields that, for every $q>0$,
\begin{equation}\label{eq:R-fast-decay-q}
e_{4\lfloor n^{1/2}\rfloor} \left(R'_{-\infty}\right)
=
o \left(n^{-q}\right),
\qquad n\to\infty.
\end{equation}
% }
\iffalse 
Since \(s>0\) was arbitrary, \eqref{eq:R-fast-decay} implies that the entropy
numbers of \(R'_{-\infty}\) decay faster than any polynomial in \(n\). More
precisely, for every \(q>0\),
\begin{equation}\label{eq:R-fast-decay-q}
    e_{4\lfloor n^{1/2}\rfloor}\!\left(R'_{-\infty}\right)
    =
    o \left(n^{-q}\right),
    \qquad n\to\infty.
\end{equation}
\fi
From now on, we assume without loss of generality that \(T_1=T_\sigma\), where
\(T_\sigma\) is the Weyl quantization of the hypoelliptic symbol \(\sigma\).
If instead \(T_2=T_\sigma\), the same argument applies after interchanging the
roles of \(T_1\) and \(T_2\) from the outset.




% \color{red}
An argument identical to that used to derive \eqref{eq:R-fast-decay-q} shows
that the entropy numbers of the % smoothing 
operator \(K_{-\infty}\) in the
parametrix relation
\begin{equation}\label{eq:parametrix-T1-T2}
T_1 = T_2(I+K)+K_{-\infty}
\end{equation}
also decay faster than any polynomial rate. More precisely, for every \(q>0\),
\begin{equation}\label{eq:K-fast-decay-q}
    e_{4\lfloor n^{1/2}\rfloor}\!\left(K_{-\infty}\right)
    =
    o \left(n^{-q}\right),
    \qquad n\to\infty.
\end{equation}
We next derive two-sided polynomial bounds for the entropy numbers of
\(T_1=T_\sigma\).
Since $\sigma \in \mathrm{H}\Gamma^{-}_{\rho}(\R^{2d})$, there exist exponents
$m_+ \ge m_- > 0$ such that the two-sided hypoellipticity bounds
\eqref{eq:bound-def-hypoellip} hold. Consequently, there exist constants
$c_1,c_2>0$ and $\lambda_0>0$ so that, for all $\lambda\in(0,\lambda_0)$,
\[
c_1 \lambda^{-2d/m_+}
\le
V_\sigma(\lambda)
\le
c_2 \lambda^{-2d/m_-}.
\]
Since \eqref{eq:main-res-volume-integral} has already been established for the
Weyl quantization, this yields
\[
H_{T_1}(\varepsilon)
\sim
\int_\varepsilon^\infty \frac{V_\sigma(\lambda)}{\lambda}\,d\lambda,
\qquad \varepsilon\to0,
\]
and therefore
\[
c_3 \, \varepsilon^{-2d/m_+}
\le
H_{T_1}(\varepsilon)
\le
c_4 \, \varepsilon^{-2d/m_-},
\qquad \varepsilon\in(0,\varepsilon_0),
\]
for suitable constants \(c_3,c_4,\varepsilon_0>0\).
By the definition of the entropy numbers, these bounds imply the existence of
constants \(c_5,c_6>0\) such that
\begin{equation}\label{eq:two-sided-polynomial-T1}
    c_5\, n^{-m_+/(2d)}
\lesssim
    e_n(T_1)
\lesssim
    c_6\, n^{-m_-/(2d)},
    \qquad n\to\infty.
\end{equation}
Choose \(q>m_+/(2d)\). Then \eqref{eq:R-fast-decay-q} and
\eqref{eq:two-sided-polynomial-T1} imply
\[
e_{4\lfloor n^{1/2}\rfloor}\!\left(R'_{-\infty}\right)
=
o \bigl(e_n(T_1)\bigr),
\qquad n\to\infty.
\]
Combining this with \eqref{eq:entropy-nb-arg-1} and
\eqref{eq:bound-entropy-identity}, we obtain
\begin{equation}\label{eq:inequality-entropy-nbs}
    e_{\,n+5\lfloor n^{1/2}\rfloor}(T_2)
    \le
    e_n(T_1)\bigl(1+o(1)\bigr),
    \qquad n\to\infty.
\end{equation}
Based on \eqref{eq:parametrix-T1-T2}, applying the same entropy-number calculus as before, we get
\begin{equation}\label{eq:entropy-nb-arg-1-bis}
    e_{\,n+5\lfloor n^{1/2}\rfloor}(T_1)
    \le
    e_n(T_2)\,e_{\lfloor n^{1/2}\rfloor}(I+K)
    +
    e_{4\lfloor n^{1/2}\rfloor}(K_{-\infty}).
\end{equation}
Since \(K\) is compact,
\[
e_{\lfloor n^{1/2}\rfloor}(I+K)\le 1+o(1),
\qquad n\to\infty.
\]
Moreover, by \eqref{eq:K-fast-decay-q},
\[
e_{4\lfloor n^{1/2}\rfloor}(K_{-\infty})
=
o \left(n^{-q}\right),
\]
for every $q>0$.
Using the lower bound in \eqref{eq:two-sided-polynomial-T1} together with
\eqref{eq:entropy-nb-arg-1-bis}, we infer that
\[
e_n(T_2)\gtrsim n^{-m_+/(2d)},
\qquad n\to\infty.
\]
In particular,
\[
e_{4\lfloor n^{1/2}\rfloor}(K_{-\infty})
=
o \bigl(e_n(T_2)\bigr),
\qquad n\to\infty.
\]
Reinserting this into \eqref{eq:entropy-nb-arg-1-bis}, we obtain
\begin{equation}\label{eq:inequality-entropy-nbs-bis}
    e_{\,n+5\lfloor n^{1/2}\rfloor}(T_1)
    \le
    e_n(T_2)\bigl(1+o(1)\bigr),
    \qquad n\to\infty.
\end{equation}
The preceding argument was carried out under the assumption \(T_1=T_\sigma\).
Since the parametrix relation assumed in Lemma~\ref{lem:technical-lemma} and the
derived reverse parametrix relation are invariant under exchanging \(T_1\) and
\(T_2\), the same reasoning applies when \(T_2=T_\sigma\).
Consequently, the estimates
\eqref{eq:inequality-entropy-nbs} and
\eqref{eq:inequality-entropy-nbs-bis}
hold irrespective of which of the two operators is the Weyl quantization of
\(\sigma\).


\iffalse
Since \(s>0\) was arbitrary, we conclude that $e_{4\lfloor n^{1/2}\rfloor}\!\left(R'_{-\infty}\right)$ decays faster than any power of \(n\).
For the hypoelliptic symbol \(\sigma\) appearing in the statement of
Lemma~\ref{lem:technical-lemma}, an argument analogous to that used in
\eqref{eq:proof-lemma-AAA}--\eqref{eq:proof-lemma-AAB}, combined with the upper bound in the left part of 
\eqref{eq:bound-def-hypoellip}, shows that the entropy numbers
\(\{e_j(T_\sigma)\}_{j\in\N}\) of its Weyl quantization \(T_\sigma\) decay
polynomially in \(j\). Consequently, irrespective of whether \(T_1=T_\sigma\)
or \(T_2=T_\sigma\), relation \eqref{eq:R-fast-decay} implies that the
contribution of the entropy numbers of \(R'_{-\infty}\) in
\eqref{eq:entropy-nb-arg-1} is asymptotically negligible. Combining
\eqref{eq:entropy-nb-arg-1} with \eqref{eq:bound-entropy-identity} therefore
yields
\begin{equation}\label{eq:inequality-entropy-nbs}
    e_{\,n+5\lfloor n^{1/2}\rfloor}(T_2)
    \le
    e_n(T_1)\bigl(1+o(1)\bigr),
    \qquad n\to\infty.
\end{equation}
By interchangeability of \(T_1\) and \(T_2\), we obtain the corresponding estimate
\begin{equation}\label{eq:inequality-entropy-nbs-bis}
    e_{\,n+5\lfloor n^{1/2}\rfloor}(T_1)
    \le
    e_n(T_2)\bigl(1+o(1)\bigr),
    \qquad n\to\infty.
\end{equation}
\fi 


\iffalse
Since $s>0$ was arbitrary, we can hence conclude that $e_{4\lfloor n^{1/2}\rfloor}\!\left(R'_{-\infty}\right)$
decays faster than any power of $n$. 
For the hypoelliptic symbol $\sigma$ in the statement of Lemma~\ref{lem:technical-lemma}, 
a similar argument as in \eqref{eq:proof-lemma-AAA}--\eqref{eq:proof-lemma-AAB},
combined with the left part of \eqref{eq:bound-def-hypoellip},
guarantees that the entropy numbers $\{e_j(T_\sigma)\}_{j\in\N}$ of its Weyl quantization $T_\sigma$ decay polynomially fast in $j$.
Therefore, irrespectively of whether $T_1=T_\sigma$ or $T_2 = T_\sigma$, 
we get from \eqref{eq:R-fast-decay} that 
the contribution of the entropy numbers of $R'_{-\infty}$ in \eqref{eq:entropy-nb-arg-1} is asymptotically negligible,
and combining \eqref{eq:entropy-nb-arg-1} with \eqref{eq:bound-entropy-identity} yields
\begin{equation}\label{eq:inequality-entropy-nbs}
    e_{\,n+\textcolor{blue}{5}\lfloor n^{1/2}\rfloor}(T_2)
    \le
    e_n(T_1) \left(1+o(1)\right),
    \qquad n\to\infty.
\end{equation}
Arguing by interchangeability of $T_1$ and $T_2$, we have the corresponding asymptotic bound
\begin{equation}\label{eq:inequality-entropy-nbs-bis}
    e_{\,n+\textcolor{blue}{5}\lfloor n^{1/2}\rfloor}(T_1)
    \le
    e_n(T_2) \left(1+o(1)\right),
    \qquad n\to\infty.
\end{equation}
\fi 



\color{black}

We now convert the entropy-number relations
\eqref{eq:inequality-entropy-nbs} and \eqref{eq:inequality-entropy-nbs-bis}
into corresponding asymptotic relations for the entropy functions.
Fix $\varepsilon>0$ and choose $n_\varepsilon\in\N$ such that
\begin{equation}\label{eq:definition-k-entropy}
    n_\varepsilon+5\lfloor n_\varepsilon^{1/2}\rfloor
    <
    H_{T_2}(\varepsilon)
    \le
  n_\varepsilon+1+5\lfloor (n_\varepsilon+1)^{1/2}\rfloor.
\end{equation}
By the definition of entropy numbers this implies
\[
e_{\,n_\varepsilon+5\lfloor n_\varepsilon^{1/2}\rfloor}(T_2)\ge \varepsilon .
\]
Combining this with \eqref{eq:inequality-entropy-nbs}, we obtain
\[
\varepsilon
\le
e_{\,n_\varepsilon+5\lfloor n_\varepsilon^{1/2}\rfloor}(T_2)
\le
e_{n_\varepsilon}(T_1)\bigl(1+o(1)\bigr),
\qquad \varepsilon\to0 .
\]
Hence there exists a function $\zeta:\R_+^*\to\R_+^*$ with
$\zeta(\varepsilon)\to1$ as $\varepsilon\to0$ such that
\begin{equation}\label{eq:relation-eps-entropy-nbs}
    \varepsilon\,\zeta(\varepsilon)
    <
    e_{n_\varepsilon}(T_1).
\end{equation}
Applying again the definition of entropy numbers yields
\begin{equation}\label{eq:useful-relation1}
    H_{T_1}\!\left(\varepsilon\,\zeta(\varepsilon)\right)
    \ge
    n_\varepsilon
    =
    H_{T_2}(\varepsilon)\bigl(1+o_{\varepsilon\to0}(1)\bigr).
\end{equation}
%A completely analogous argument based on
%\eqref{eq:inequality-entropy-nbs-bis} yields a function
%$\zeta':\R_+^*\to\R_+^*$ with $\zeta'(\varepsilon)\to1$ as $\varepsilon\to0$
%such that
%\begin{equation}\label{eq:useful-relation1-bis}
%    H_{T_1}\!\left(\varepsilon\,\zeta'(\varepsilon)\right)
%    \le
%    H_{T_2}(\varepsilon)\bigl(1+o_{\varepsilon\to0}(1)\bigr).
%\end{equation}
To obtain an upper bound for \(H_{T_1}\) in terms of \(H_{T_2}\), let $m_\varepsilon := \left\lceil H_{T_2}(\varepsilon) \right\rceil$.
Then $m_\varepsilon\to\infty$ as $\varepsilon\to0$, and by the definition of
entropy numbers we have
\[
e_{m_\varepsilon}(T_2)\le \varepsilon .
\]
Applying \eqref{eq:inequality-entropy-nbs-bis} with $n=m_\varepsilon$ yields
\[
e_{\,m_\varepsilon+5\lfloor m_\varepsilon^{1/2}\rfloor}(T_1)
\le
\varepsilon\bigl(1+o(1)\bigr),
\qquad \varepsilon\to0 .
\]
Consequently, there exists a function
$\zeta':\mathbb{R}_+^*\to\mathbb{R}_+^*$ with
$\zeta'(\varepsilon)\to1$ as $\varepsilon\to0$ such that
\begin{equation}\label{eq:useful-relation1-bis}
    H_{T_1}\!\left(\varepsilon\,\zeta'(\varepsilon)\right)
    \le
    m_\varepsilon+5\lfloor m_\varepsilon^{1/2}\rfloor.
\end{equation}
Since
\[
m_\varepsilon
=
H_{T_2}(\varepsilon)+O(1),
\qquad \varepsilon\to0,
\]
and $H_{T_2}(\varepsilon)\to\infty$, we obtain
\[
m_\varepsilon+5\lfloor m_\varepsilon^{1/2}\rfloor
=
H_{T_2}(\varepsilon)\bigl(1+o_{\varepsilon\to0}(1)\bigr).
\]
Substituting this relation into \eqref{eq:useful-relation1-bis} yields
\begin{equation}\label{eq:asymp-upper}
    H_{T_1} \left(\varepsilon\,\zeta'(\varepsilon)\right)
\le
H_{T_2}(\varepsilon)\bigl(1+o_{\varepsilon\to0}(1)\bigr).
\end{equation}


\iffalse 
Finally, we convert the asymptotic relation for entropy numbers \eqref{eq:inequality-entropy-nbs} into a corresponding asymptotic relation for the entropy of the corresponding operators.
To this end, fix $\varepsilon>0$ and let $n_\varepsilon\in\N$ be such that 
\begin{equation}\label{eq;definition-k-entropy}
    n_\varepsilon+{5\lfloor {n_\varepsilon}^{1/2} \rfloor}  
    < H_{T_2}\left(\varepsilon\right)
    \leq  n_\varepsilon+{5\lfloor {n_\varepsilon}^{1/2} \rfloor}+1.
\end{equation}
% In particular, by compactness of the operator $T_2$, we know that $n\to \infty$, as $\varepsilon\to 0$.
By definition of the entropy numbers, we get  $e_{n_\varepsilon+{5\lfloor {n_\varepsilon}^{1/2} \rfloor}} (T_{2}) \geq \varepsilon$,
which, using \eqref{eq:inequality-entropy-nbs}, implies the existence of a function $\zeta\colon \R_+^*\to \R_+^*$, 
satisfying $\zeta(\varepsilon)\to 1$ as $\varepsilon \to 0$, such that
\begin{equation}\label{eq:relation-eps-entropy-nbs}
    \varepsilon \, \zeta(\varepsilon) 
    <
    e_{n_\varepsilon} (T_1).
\end{equation}
Again using the definition of entropy numbers, we deduce from 
\eqref{eq;definition-k-entropy} and
\eqref{eq:relation-eps-entropy-nbs} that
\begin{equation}\label{eq:useful-relation1}
    H_{T_1}\left(\varepsilon \, \zeta(\varepsilon) \right) 
    \geq n_\varepsilon
    = H_{T_2}(\varepsilon)\left( 1+o_{\varepsilon\to 0}(1) \right).
\end{equation}
Similarly, starting from \eqref{eq:inequality-entropy-nbs-bis} instead of \eqref{eq:inequality-entropy-nbs},
there exists a function $\zeta'\colon \R_+^*\to \R_+^*$, 
satisfying $\zeta'(\varepsilon)\to 1$ as $\varepsilon \to 0$,
such that
\begin{equation}\label{eq:useful-relation1-bis}
    H_{T_1}\left(\varepsilon \, \zeta'(\varepsilon) \right) 
    \leq H_{T_2}(\varepsilon)\left( 1+o_{\varepsilon\to 0}(1) \right).
\end{equation}
\fi 

% \color{red}

From here on we assume, without loss of generality, that
$T_1=T_\sigma$. The case $T_2=T_\sigma$ is treated in the
same way after interchanging the roles of $T_1$ and $T_2$. Since \eqref{eq:main-res-volume-integral} has already been established
for the Weyl quantization, we obtain
\begin{align}
H_{T_1}\!\left(\varepsilon\,\zeta(\varepsilon)\right)
&\sim
\int_{\varepsilon\zeta(\varepsilon)}^\infty
\frac{V_\sigma(\lambda)}{\lambda}\,d\lambda
\label{eq:split-integral-zeta-proof-A1}\\
&=
\int_{\varepsilon}^{\infty}\frac{V_\sigma(\lambda)}{\lambda}\,d\lambda
+
\int_{\varepsilon\zeta(\varepsilon)}^{\varepsilon}
\frac{V_\sigma(\lambda)}{\lambda}\,d\lambda
\notag\\
&\sim
H_{T_1}(\varepsilon)
+
\int_{\varepsilon\zeta(\varepsilon)}^{\varepsilon}
\frac{V_\sigma(\lambda)}{\lambda}\,d\lambda,
\label{eq:split-integral-zeta-proof-B}
\end{align}
as $\varepsilon\to0$, where the second integral is understood with its
natural orientation. Using the monotonicity of $V_\sigma$, the remainder term in
\eqref{eq:split-integral-zeta-proof-A1}--\eqref{eq:split-integral-zeta-proof-B}
satisfies
\begin{equation}\label{eq:bound-remainder-volume-integral-A}
\left|
\int_{\varepsilon\zeta(\varepsilon)}^{\varepsilon}
\frac{V_\sigma(\lambda)}{\lambda}\,d\lambda
\right|
\le
\frac{V_\sigma \left(\varepsilon\,\zeta^*(\varepsilon)\right)}
     {\varepsilon\,\zeta^*(\varepsilon)}
\,
\bigl|\varepsilon-\varepsilon\,\zeta(\varepsilon)\bigr|
=
\frac{V_\sigma \left(\varepsilon\,\zeta^*(\varepsilon)\right)}
     {\zeta^*(\varepsilon)}
\,
|1-\zeta(\varepsilon)|,
\end{equation}
for all $\varepsilon>0$, where
$\zeta^*(\varepsilon)=\min\{1,\zeta(\varepsilon)\}$.
Since $V_\sigma$ is regularly varying at zero, Karamata’s theorem implies
\begin{equation}\label{eq:volume-integral-compp-A}
V_\sigma \left(\varepsilon\,\zeta^*(\varepsilon)\right)
=
O_{\varepsilon\to0} \left(
\int_{\varepsilon\zeta^*(\varepsilon)}^\infty
\frac{V_\sigma(\lambda)}{\lambda}\,d\lambda
\right).
\end{equation}
Because $\zeta(\varepsilon)\to1$ as $\varepsilon\to0$, it follows from
\eqref{eq:volume-integral-compp-A} that
\begin{equation}\label{eq:growth-control-volume-A}
V_\sigma \left(\varepsilon\,\zeta^*(\varepsilon)\right)
|1-\zeta(\varepsilon)|
=
o_{\varepsilon\to0}\!\left(
\int_{\varepsilon\zeta^*(\varepsilon)}^\infty
\frac{V_\sigma(\lambda)}{\lambda}\,d\lambda
\right).
\end{equation}
Combining
\textcolor{purple}{\eqref{eq:split-integral-zeta-proof-A1}--\eqref{eq:split-integral-zeta-proof-B}},
\eqref{eq:bound-remainder-volume-integral-A},
and \eqref{eq:growth-control-volume-A} therefore yields
\begin{equation}\label{eq:useful-relation3-A}
H_{T_1}\!\left(\varepsilon\,\zeta(\varepsilon)\right)
\sim
H_{T_1}(\varepsilon),
\qquad \varepsilon\to0 .
\end{equation}
The same argument, with $\zeta$ replaced by $\zeta'$, gives
\begin{equation}\label{eq:useful-relation3-bis-A}
H_{T_1}\!\left(\varepsilon\,\zeta'(\varepsilon)\right)
\sim
H_{T_1}(\varepsilon),
\qquad \varepsilon\to0 .
\end{equation}
Finally, combining
\eqref{eq:useful-relation1},
\eqref{eq:asymp-upper},
\eqref{eq:useful-relation3-A},
and \eqref{eq:useful-relation3-bis-A},
we conclude that
\[
H_{T_1}(\varepsilon)
\sim
H_{T_2}(\varepsilon),
\qquad \varepsilon\to0,
\]
which completes the proof of Lemma~\ref{lem:technical-lemma}.
\end{proof}


\iffalse 
From now on, we assume that $T_1$ is the Weyl quantization of $\sigma$ as per the assumption of Lemma~\ref{lem:technical-lemma}; again, there is no loss of generality here as $T_1$ and $T_2$ are interchangeable.
We observe that, owing to {\eqref{eq:main-res-volume-integral} (which has already been established for the Weyl quantization)}, 
we have 
\begin{align}
    H_{T_1}\left(\varepsilon\, \zeta(\varepsilon) \right)  
    \sim 
    \int_{\varepsilon\zeta(\varepsilon)}^\infty 
    \frac{V_{\symbolnot}(\lambda)}{\lambda} \diff \lambda
    &= \int_{\varepsilon}^\infty 
    \frac{V_{\symbolnot}(\lambda)}{\lambda} \diff \lambda
    + \int_{\varepsilon\zeta(\varepsilon)}^\varepsilon 
    \frac{V_{\symbolnot}(\lambda)}{\lambda} \diff \lambda \label{eq:split-integral-zeta-proof-A1}\\
    &\sim 
    H_{T_1}\left(\varepsilon \right) 
    + \int_{\varepsilon\zeta(\varepsilon)}^\varepsilon 
    \frac{V_{\symbolnot}(\lambda)}{\lambda} \diff \lambda,\label{eq:split-integral-zeta-proof-B}
\end{align}
as $\varepsilon\to 0$, 
where the second integral is understood with its natural orientation.
Using the monotonicity of $V_\sigma$, the remainder integral in
\eqref{eq:split-integral-zeta-proof-A1}--\eqref{eq:split-integral-zeta-proof-B} admits
the bound
% Using the monotonicity of $V_{\symbolnot}$,
% the second integral in the right-most side of \eqref{eq:split-integral-zeta-proof-A1}--\eqref{eq:split-integral-zeta-proof-B} can be bounded according to
\begin{equation}\label{eq:bound-remainder-volume-integral-A}
    \left\lvert\int_{\varepsilon\zeta(\varepsilon)}^\varepsilon 
    \frac{V_{\symbolnot}(\lambda)}{\lambda} \diff \lambda \right\rvert
    \leq  \frac{V_{\symbolnot}(\varepsilon \, \zeta^*(\varepsilon))}{\varepsilon \, \zeta^*(\varepsilon)}|\varepsilon - \varepsilon\, \zeta(\varepsilon)|
    = \frac{V_{\symbolnot}(\varepsilon \, \zeta^*(\varepsilon))}{\zeta^*(\varepsilon)}|1 - \zeta(\varepsilon)|,
\end{equation}
for all $ \varepsilon >0$,
where $\zeta^*(\varepsilon) \coloneqq \min \{1, \zeta(\varepsilon)\}$.
Moreover, by regular variation of $V_\sigma$ at zero and Karamata’s theorem for
integrals of regularly varying functions, we have
\begin{equation}\label{eq:volume-integral-compp-A}
    V_{\symbolnot}(\varepsilon \, \zeta^*(\varepsilon))
    % = - \int_{\varepsilon\zeta^*(\varepsilon)}^\infty 
    % V'_{\symbolnot}(\lambda) \diff \lambda
    = O_{\varepsilon \to 0}\left(\int_{\varepsilon\zeta^*(\varepsilon)}^\infty 
    \frac{V_{\symbolnot}(\lambda)}{\lambda} \diff \lambda\right).
\end{equation}
% (Recall that asymptotics are preserved through integration, see, e.g., \cite[Lemma~11]{allard2025metricentropyminimaxrisk}).
Since $\zeta(\varepsilon) \to 1$ as $\varepsilon\to 0$, it follows from \eqref{eq:volume-integral-compp-A} that 
\begin{equation}\label{eq:growth-control-volume-A}
    V_{\symbolnot}(\varepsilon \, \zeta^*(\varepsilon))|1 - \zeta(\varepsilon)|
    = o_{\varepsilon \to 0}\left(\int_{\varepsilon\zeta^*(\varepsilon)}^\infty 
    \frac{V_{\symbolnot}(\lambda)}{\lambda} \diff \lambda\right).
\end{equation}
Combining \eqref{eq:split-integral-zeta-proof-A1}--\eqref{eq:split-integral-zeta-proof-B}, \eqref{eq:bound-remainder-volume-integral-A}, and \eqref{eq:growth-control-volume-A}, 
we conclude that 
\begin{equation}\label{eq:useful-relation3-A}
    H_{T_1}\left(\varepsilon \, \zeta(\varepsilon)\right) 
    \sim H_{T_1}\left(\varepsilon\right),
    \quad \text{as }\varepsilon \to 0.
\end{equation}
% The same chain of arguments, starting from \eqref{eq:split-integral-zeta-proof-A1}--\eqref{eq:split-integral-zeta-proof-B} up to \eqref{eq:useful-relation3-A},
The same argument, with $\zeta$ replaced by $\zeta'$,
yields
\begin{equation}\label{eq:useful-relation3-bis-A}
    H_{T_1}\left(\varepsilon \, \zeta'(\varepsilon)\right) 
    \sim H_{T_1}\left(\varepsilon\right),
    \quad \text{as }\varepsilon \to 0.
\end{equation}
It therefore follows  from \eqref{eq:useful-relation1}, \eqref{eq:useful-relation1-bis}, \eqref{eq:useful-relation3-A}, and \eqref{eq:useful-relation3-bis-A} that
\begin{equation*}%\label{eq:entropy-bounds-part1}  
    H_{T_1}(\varepsilon)
    \sim H_{T_2}\left(\varepsilon \right),
     % \left( 1+o_{\varepsilon\to 0}(1) \right),
    \quad\text{as } \varepsilon\to 0,
\end{equation*}
% Conversely, a similar line of arguments starting from \eqref{eq:introduction-of-R2} instead of \eqref{eq:introduction-of-R1}, with the roles of $T_{\symbolnot}$ and $T_*$ reversed,  yields 
% \begin{equation}\label{eq:entropy-bounds-part2}
%      H_{\symbolnot}(\varepsilon)
%     \leq  H_{T_*}\left(\varepsilon\right)\left( 1+o_{\varepsilon\to 0}(1) \right).
% \end{equation}
% From \eqref{eq:entropy-bounds-part1} and \eqref{eq:entropy-bounds-part2}, we get 
% \begin{equation}\label{eq:equivalence-entropy-op-inverse-op}
%     H_{\symbolnot}(\varepsilon) 
%     \sim H_{T_*}\left(\varepsilon \right),
%     \quad \text{as } \varepsilon \to 0,
% \end{equation}
thereby completing the proof of Lemma~\ref{lem:technical-lemma}.
\fi 


\iffalse 
% We now estimate the term
% $e_{2\lfloor n^{1/2}\rfloor}\!\left(R'_{-\infty}\right)$
% appearing in \eqref{eq:entropy-nb-arg-1}.
% To this end, arbitrarily fix $m>0$ and consider a hypoelliptic symbol
% $\bar\sigma\in \mathrm{H}\Gamma^{m,m}_{\rho_1}(\R^{2d})$.
% Let $T_{\bar\sigma}$ denote its Weyl quantization and note that, by
% \cite[Corollary~23.3 and Theorem~26.3]{shubin_pseudodifferential_2001},
% $T_{\bar\sigma}$ has purely discrete spectrum.
% Choose $\mu\in\R$ outside the spectrum of $T_{\bar\sigma}$ and let 
% \[
% T_\mu \coloneqq T_{\bar\sigma}-\mu I .
% \]
% Since $T_{\bar\sigma}\in\Psi^{m}_{\rho_1}(\R^d)$ and $\mu I\in\Psi^{0}_{\rho_1}(\R^d)$,
% it follows that $T_\mu\in\Psi^{m}_{\rho_1}(\R^d)$.
% The choice of $\mu$ guarantees that $T_\mu$ is invertible; we denote
% its inverse (the resolvent) by
% \[
% R_\mu \coloneqq T_\mu^{-1}.
% \]
% \iffalse 
% By \cite[Corollary~23.3 and Theorem~26.3]{shubin_pseudodifferential_2001},
% $T_{\bar\sigma}$ has purely discrete spectrum.
% Choose $\mu\in\R$ outside the spectrum of $T_{\bar\sigma}$ and let
% \[
% T_\mu \coloneqq T_{\bar\sigma}-\mu I .
% \]
% Since $T_{\bar\sigma}\in\Psi^{m}_{\rho_1}(\R^d)$ and $\mu I\in\Psi^{0}_{\rho_1}(\R^d)$,
% the operator
% \[
% T_\mu \coloneqq T_{\bar\sigma}-\mu I
% \]
% belongs to $\Psi^{m}_{\rho_1}(\R^d)$. Choosing $\mu\in\R$ outside the spectrum of
% $T_{\bar\sigma}$ ensures that $T_\mu$ is invertible on $L^2(\R^d)$, and we denote
% its inverse (the resolvent) by
% \[
% R_\mu \coloneqq T_\mu^{-1}.
% \]

% Then $T_\mu$ is invertible, and we denote its resolvent by
% \[
% R_\mu \coloneqq T_\mu^{-1}.
% \]
% \fi 
%%%T9[OK]: Here I am not sure how the extension argument would follow from Theorem 25.4 in Shubin and whether we need it at all?

% Since $T_\mu\in \mathrm{H}\Gamma^{m,m}_{\rho_1}(\R^{2d})$ and
% $\ker T_\mu=\ker T_\mu^*=\{0\}$, \cite[Theorem~25.4]{shubin_pseudodifferential_2001} 
% shows that $R_\mu$ is a pseudodifferential operator with Weyl symbol
% $r_\mu\in \mathrm{H}\Gamma^{-m,-m}_{\rho_1}(\mathbb{R}^{2d})$.
% %By \cite[Theorem~25.4]{shubin_pseudodifferential_2001}, the resolvent 
% %$R_\mu$ is a pseudodifferential operator with symbol
% %$r_\mu\in \mathrm{H}\Gamma^{-m,-m}_{\rho_1}(\R^{2d})$. 
% Using the lower and upper bounds on the left side in \eqref{eq:bound-def-hypoellip},
% there exist constants $C_1,C_2>0$ such that, as $\lambda\to0$,
% \begin{align}
%     V_{r_\mu}(\lambda)
%     &\le
%     \left(
%         \int_{\R^{2d}}
%         \mathbbm{1}_{\{C_2(1+\|z\|_2^2)^{-m/2}>\lambda\}}
%         \,dz
%     \right)
%     \bigl(1+o(1)\bigr)
%     \label{eq:volume-rmu-up}\\
%     &=
%     \omega_{2d}
%     \left(\frac{C_2}{\lambda}\right)^{\frac{2d}{m}}
%     \bigl(1+o(1)\bigr),
%     \label{eq:volume-rmu-up-2}
% \end{align}
% and 
% \begin{align}
%     V_{r_\mu}(\lambda)
%     &\ge
%     \left(
%         \int_{\R^{2d}}
%         \mathbbm{1}_{\{C_1(1+\|z\|_2^2)^{-m/2}>\lambda\}}
%         \,dz
%     \right)
%     \bigl(1+o(1)\bigr)
%     \label{eq:volume-rmu-down}\\
%     &=
%     \omega_{2d}
%     \left(\frac{C_1}{\lambda}\right)^{\frac{2d}{m}}
%     \bigl(1+o(1)\bigr).
%     \label{eq:volume-rmu-down-2}
% \end{align}
% (The $o(1)$ terms account for the fact that the hypoelliptic bounds hold only
% outside a ball in phase space.) Taken together, these estimates show that $V_{r_\mu}$ is regularly varying at zero.

% Applying the entropy--volume relation \eqref{eq:main-res-volume-integral},
% which has already been established for Weyl quantization \textcolor{blue}{in \eqref{eq:mainr-result-for-Weyl-1}}, to the operator
% $R_\mu$ yields
% \[
% H_{R_\mu}(\varepsilon)
% \sim
% \int_\varepsilon^\infty \frac{V_{r_\mu}(\lambda)}{\lambda}\,d\lambda,
% \qquad \varepsilon\to0.
% \]
% Since \eqref{eq:volume-rmu-up}--\eqref{eq:volume-rmu-down-2} show that
% $V_{r_\mu}(\lambda)\asymp \lambda^{-2d/m}$ as $\lambda\to0$, the integral on the
% right-hand side behaves as
% \[
% \int_\varepsilon^\infty \lambda^{-2d/m-1}\,d\lambda
% \asymp
% \varepsilon^{-2d/m},
% \qquad \varepsilon\to0.
% \]
% Consequently,
% \[
% H_{R_\mu}(\varepsilon)\asymp \varepsilon^{-2d/m},
% \qquad \varepsilon\to0.
% \]
% By definition of the entropy numbers \eqref{eq:def-entropy-nbs}, this implies that
% \[
% e_{\textcolor{blue}{j}}(R_\mu)\asymp {\textcolor{blue}{j}}^{-m/(2d)},
% \qquad {\textcolor{blue}{j}}\to\infty.
% \]
% Evaluating this estimate at ${\textcolor{blue}{j}}=\lfloor n^{1/2}\rfloor$ yields the two-sided bound
% \begin{equation}\label{eq:poly-scaling-ent-nb}
%     C'\,n^{-\frac{m}{4d}}\bigl(1+o(1)\bigr)
%     \le
%     e_{\lfloor n^{1/2}\rfloor}(R_\mu)
%     \le
%     C''\,n^{-\frac{m}{4d}}\bigl(1+o(1)\bigr),
%     \qquad n\to\infty,
% \end{equation}
% for suitable constants $C',C''>0$.
% % \iffalse
% % Applying the entropy–volume relation \eqref{eq:main-res-volume-integral}
% % (which has already been established for Weyl quantization)
% % yields constants $C',C''>0$ such that
% % \begin{equation}\label{eq:poly-scaling-ent-nb}
% %     C'\,n^{-\frac{m}{4d}}\bigl(1+o(1)\bigr)
% %     \le
% %     e_{\lfloor n^{1/2}\rfloor}(R_\mu)
% %     \le
% %     C''\,n^{-\frac{m}{4d}}\bigl(1+o(1)\bigr),
% %     \qquad n\to\infty.
% % \end{equation}
% % \fi 
% We now return to the operator $R'_{-\infty}$.
% Using the factorization
% \[
% R'_{-\infty} = (R'_{-\infty}T_\mu)\,R_\mu,
% \]
% together with the submultiplicativity and monotonicity of entropy numbers, we obtain
% \begin{align}
% e_{2\lfloor n^{1/2}\rfloor}\!\left(R'_{-\infty}\right)
% &\le
% e_{\lfloor n^{1/2}\rfloor}(R_\mu)\,
% e_{\lfloor n^{1/2}\rfloor}\!\left(R'_{-\infty}T_\mu\right)
% \label{eq:bound-ent-split-1}\\
% &\le
% C''\,n^{-\frac{m}{4d}}\bigl(1+o(1)\bigr)\,
% e_{\lfloor n^{1/2}\rfloor}\!\left(R'_{-\infty}T_\mu\right),
% \label{eq:bound-ent-split-2}
% \end{align}
% where the last step employs \eqref{eq:poly-scaling-ent-nb}.
% Since $R'_{-\infty}\in\Psi^{-\infty}(\R^d)$ and
% $T_\mu\in\Psi^{m}_{\rho_1}(\R^d)$,
% their composition $R'_{-\infty}T_\mu$ belongs to $\Psi^{-\infty}(\R^d)$
% and is therefore compact on $L^2(\R^d)$.
% Hence
% \[
% \lim_{n\to\infty}
% e_{\lfloor n^{1/2}\rfloor}\!\left(R'_{-\infty}T_\mu\right)
% =
% 0,
% \]
% and combining this with
% \eqref{eq:bound-ent-split-1}–\eqref{eq:bound-ent-split-2} yields
% \begin{equation}\label{eq:R-fast-decay}
% e_{2\lfloor n^{1/2}\rfloor}\!\left(R'_{-\infty}\right)
% =
% o\left(n^{-\frac{m}{4d}}\right),
% \qquad n\to\infty.
% \end{equation}
% Since $m>0$ was arbitrary, we can hence conclude that $e_{2\lfloor n^{1/2}\rfloor}\!\left(R'_{-\infty}\right)$
% decays faster than any power of $n$. On the other hand, as either $T_1$ or $T_2$ is the Weyl quantization of a
% negative-order hypoelliptic symbol, the corresponding entropy number decays polynomially.
% Indeed, the same phase-space volume estimates used earlier imply the existence of
% constants \(c,C>0\) and exponents \(\alpha,\beta>0\) such that
% \[
% c\,n^{-\alpha}
% \;\le\;
% e_n(T_i)
% \;\le\;
% C\,n^{-\beta},
% \qquad n\to\infty,
% \]
% for the operator \(T_i\) that is the Weyl quantization of the hypoelliptic symbol.




\color{orange}
We now relate the entropy-number behavior of $T_1$ and $T_2$ by exploiting the
parametrix relations between the two operators.
The key observation is that polynomial entropy-number behavior of one of the
operators forces the same behavior for the other, which in turn allows
entropy-number comparisons to be converted into asymptotic relations for the
corresponding entropy functions.

Throughout this part of the proof, we assume that $T_1$ is the Weyl
quantization of the negative-order hypoelliptic symbol $\sigma$. The case in which $T_2$ is the Weyl quantization of $\sigma$ is treated in the same way,
with the roles of $T_1$ and $T_2$ interchanged. We start by noting that 
since $m>0$ was arbitrary in \eqref{eq:R-fast-decay}, the smoothing remainder
satisfies
\[
e_{2\lfloor n^{1/2}\rfloor}\!\left(R'_{-\infty}\right)
=o \left(n^{-p}\right)
\qquad\text{for every }p>0
\]
and hence decays faster than any polynomial in $n$.
Next, we explain why polynomial decay of $e_n(T_1)$ implies polynomial decay of
$e_n(T_2)$.
The parametrix construction yields the two equivalent operator representations
\begin{equation}\label{eq:parametrix-two}
T_2 = T_1(I+R_1)+R_{-\infty},
\qquad
T_1 = T_2(I+K)+K_{-\infty},
\end{equation}
where $R_1,K$ are of strictly negative order and $R_{-\infty},K_{-\infty}$ are smoothing.
Starting from the first identity and using the entropy-number estimates
\eqref{eq:entropy-nb-arg-1}--\eqref{eq:entropy-nb-arg-2}, the smoothing term contributes
super-polynomially small entropy numbers, while the polynomial upper bound on $e_n(T_1)$
precludes $e_n(T_2)$ from decaying slower than polynomially.
Conversely, starting from the second identity, the polynomial lower bound on $e_n(T_1)$
forces $e_n(T_2)$ to be at least polynomial, since otherwise the right-hand side would be
super-polynomially small.
Together, these two implications establish polynomial decay of $e_n(T_2)$.

Next, fix $\varepsilon>0$ and choose $n\in\N$ such that
\[
n+3\lfloor n^{1/2}\rfloor
< H_{T_2}(\varepsilon)
\le n+3\lfloor n^{1/2}\rfloor+1.
\]
Since $T_2$ is compact, $H_{T_2}(\varepsilon)\to\infty$ as $\varepsilon\to0$, and hence
$n\to\infty$ as $\varepsilon\to0$.
By definition of entropy numbers,
\[
e_{n+3\lfloor n^{1/2}\rfloor}(T_2)\ge \varepsilon
\quad\text{and}\quad
e_{n+3\lfloor n^{1/2}\rfloor+1}(T_2)\le \varepsilon.
\]

Since $V_\sigma$ is regularly varying at $0$ with negative index, there exist
$\rho>0$ and a slowly varying function $\ell$ such that
\[
V_\sigma(\lambda)\sim \lambda^{-\rho}\,\ell(1/\lambda),
\qquad \lambda\to0,
\]
see \cite{binghamRegularVariation1987}.
Moreover, the first part of Theorem~\ref{thm:main-result}, already established above,
yields the entropy representation
\[
H_{T_1}(\varepsilon)\sim \int_\varepsilon^\infty \frac{V_\sigma(\lambda)}{\lambda}\,d\lambda,
\qquad \varepsilon \to 0.
\]
Since the contribution of the integral over $(\lambda_0,\infty)$ is finite for any fixed
$\lambda_0>0$, the asymptotic behavior of $H_{T_1}$ as $\varepsilon \to 0$ is governed
by
\[
\int_\varepsilon^{\lambda_0} \lambda^{-\rho-1}\,\ell(1/\lambda)\,d\lambda.
\]
An application of Karamata’s theorem for integrals of regularly varying functions
yields
\begin{equation}\label{eq:HT1-varying}
H_{T_1}(\varepsilon)\sim \frac{1}{\rho}\,\varepsilon^{-\rho}\,\ell(1/\varepsilon),
\qquad \varepsilon \to 0,
\end{equation}
showing that $H_{T_1}$ is itself regularly varying at $0$ with index $\rho>0$.

The entropy numbers 
\[
e_n(T_1) =\inf\{\varepsilon>0:\ H_{T_1}(\varepsilon)\le n\},
\qquad n\in\mathbb{N},
\]
coincide with the generalized inverse of the nonincreasing function
$H_{T_1}$.
To invert the asymptotic relation \eqref{eq:HT1-varying} as $\varepsilon\to0$,
we recast it as a large-argument asymptotic by defining
\[
\Phi(t):=H_{T_1}(1/t),\qquad t\to\infty.
\]
Then $\Phi$ is nondecreasing and satisfies
\[
\Phi(t)=H_{T_1}(1/t)\sim \frac{1}{\rho}\,t^{\rho}\,\ell(t),
\qquad t\to\infty,
\]
so $\Phi$ is regularly varying at infinity with index $\rho$.
Let $\Psi$ denote its generalized inverse,
\[
\Psi(n):=\inf\{t>0:\ \Phi(t)\ge n\}.
\]
By standard inversion results for regularly varying monotone functions,
$\Psi$ is regularly varying at infinity with index $1/\rho$, and hence
\[
\Psi(n)\sim n^{1/\rho}\,L_1(n),
\qquad n\to\infty,
\]
for some slowly varying function $L_1$.
Since $e_n(T_1)=1/\Psi(n)$, this yields
\[
e_n(T_1)\sim n^{-1/\rho}\,\widetilde L(n),
\qquad n\to\infty,
\]
where $\widetilde L(n):=1/L_1(n)$ is slowly varying.
This representation allows one to control index shifts.
Let $k(n)=o(n)$.
Then
\[
\frac{e_{n+k(n)}(T_1)}{e_n(T_1)}
\sim
\Bigl(1+\frac{k(n)}{n}\Bigr)^{-1/\rho}
\frac{\widetilde L(n+k(n))}{\widetilde L(n)}.
\]
Since $\widetilde L$ is slowly varying, the Uniform Convergence Theorem
\cite[Theorem~1.2.1]{binghamRegularVariation1987} yields
$\widetilde L(n+k(n))/\widetilde L(n)\to 1$ for every $k(n)=o(n)$.
Consequently,
\[
\frac{e_{n+k(n)}(T_1)}{e_n(T_1)}\to 1,
\qquad n\to\infty.
\]
Although no symbol is assumed for $T_2$, the parametrix relations are sufficient to
transfer entropy–number asymptotics from $T_1$ to $T_2$.
Specifically, it follows from \eqref{eq:parametrix-two}, 
submultiplicativity and subadditivity of entropy numbers, and the fact that the smoothing remainders
have super-polynomially decaying entropy numbers that
\[
e_{n+k(n)}(T_2)\le e_n(T_1)\,(1+o(1))
\quad\text{and}\quad
e_{n+k(n)}(T_1)\le e_n(T_2)\,(1+o(1)),
\qquad n\to\infty.
\]
Now, using the second inequality and the shift stability of $e_n(T_1)$, we obtain
\[
e_n(T_2)\ge e_{n+k(n)}(T_1)\,(1+o(1))
= e_n(T_1)\,(1+o(1)).
\]
On the other hand, the first inequality together with the monotonicity of entropy
numbers yields
\[
e_n(T_2)\le e_{n-k(n)}(T_2)\le e_{n-2k(n)}(T_1)\,(1+o(1))
= e_n(T_1)\,(1+o(1)).
\]
Combining the two bounds shows that $e_n(T_2)\sim e_n(T_1)$.
Consequently, the index-shift stability of $e_n(T_1)$ transfers to $e_n(T_2)$, and
for every sequence $k(n)=o(n)$, we have 
\[
\frac{e_{n+k(n)}(T_2)}{e_n(T_2)}\to 1,
\qquad n\to\infty.
\]
Finally, setting $\zeta'_n = e_n(T_2)/e_{n+k(n)}(T_2)$ with $k(n)=3\lfloor n^{1/2}\rfloor+1$.
Since $e_{n+k(n)}(T_2)\le\varepsilon$ and $e_{n+k(n)}(T_2)/e_n(T_2)\to1$, we have
$\zeta'_n\to1$ and hence $e_n(T_2)\le \varepsilon\,\zeta'_n$.

% \iffalse 
% Consequently, since $e_{n+k(n)}(T_2)\le \varepsilon$ with
% $k(n)=3\lfloor n^{1/2}\rfloor+1$ and
% $e_{n+k(n)}(T_2)/e_n(T_2)\to 1$ as $n\to\infty$, there exists a sequence
% $\{\zeta'_n\}$ with $\zeta'_n\to1$ such that
% \[
% e_n(T_2)\le \varepsilon\,\zeta'_n.
% \]
% \fi 

We now use the representation $T_1=T_2(I+K)+K_{-\infty}$ to compare entropy numbers in the
opposite direction.
Applying the same entropy-number calculus as in
\eqref{eq:entropy-nb-arg-1}--\eqref{eq:entropy-nb-arg-2}
to this representation yields
\[
e_{n+3\lfloor n^{1/2}\rfloor}(T_1)
\le e_n(T_2)\,e_{\lfloor n^{1/2}\rfloor}(I+K)
+ e_{2\lfloor n^{1/2}\rfloor}(K_{-\infty}).
\]
Since $K$ is compact, $e_{\lfloor n^{1/2}\rfloor}(I+K)\le 1+o(1)$ as $n\to\infty$.
Moreover, $K_{-\infty}$ is smoothing, hence
$e_{2\lfloor n^{1/2}\rfloor}(K_{-\infty})=o(n^{-p})$ for every $p>0$.
Because $e_n(T_2)$ decays polynomially as established above, the smoothing term is negligible relative to
$e_n(T_2)$, and we obtain
\[
e_{n+3\lfloor n^{1/2}\rfloor}(T_1)\le e_n(T_2)\,\bigl(1+o(1)\bigr),
\qquad n\to\infty.
\]
Inserting the bound $e_n(T_2)\le \varepsilon\,\zeta'_n$ and absorbing the factor
$(1+o(1))$ into $\zeta'_n$ (redefining $\zeta'_n$ if necessary) yields
\[
e_{n+3\lfloor n^{1/2}\rfloor}(T_1)\le \varepsilon\,\zeta'_n.
\]
By definition of entropy numbers and the monotonicity of $H_{T_1}$, this implies
\[
H_{T_1}(\varepsilon\zeta'_n)\le n+3\lfloor n^{1/2}\rfloor.
\]
Using again that $H_{T_2}(\varepsilon)\to\infty$ as $\varepsilon\to0$, the choice of $n$ made above yields
\[
n+3\lfloor n^{1/2}\rfloor
= H_{T_2}(\varepsilon)\,\bigl(1+o_{\varepsilon\to0}(1)\bigr),
\]
and therefore
\begin{equation}\label{eq:upper-bound-H}
H_{T_1}(\varepsilon\zeta'_n)
\;\le\;
H_{T_2}(\varepsilon)\,\bigl(1+o_{\varepsilon\to0}(1)\bigr).
\end{equation}
We next derive the complementary inequality.
Starting from 
\[
e_{n+3\lfloor n^{1/2}\rfloor}(T_2)
\le e_n(T_1)\,\bigl(1+o(1)\bigr),
\qquad n\to\infty
\]
and combining with $e_{n+3\lfloor n^{1/2}\rfloor}(T_2)\ge\varepsilon$, yields a
sequence $\{\zeta_n\}$ with $\zeta_n\to1$ such that
\[
\varepsilon\,\zeta_n \le e_n(T_1).
\]
Invoking again the definition of entropy numbers and the monotonicity of $H_{T_1}$, we
obtain
\[
H_{T_1}(\varepsilon\zeta_n)\ge H_{T_1}\bigl(e_n(T_1)\bigr)\ge n
= H_{T_2}(\varepsilon)\,\bigl(1+o_{\varepsilon\to0}(1)\bigr),
\]
that is,
\begin{equation}\label{eq:lower-bound-H}
H_{T_1}(\varepsilon\zeta_n)
\;\ge\;
H_{T_2}(\varepsilon)\,\bigl(1+o_{\varepsilon\to0}(1)\bigr).
\end{equation}


% \iffalse 
% We start from the swapped entropy-number estimate
% \begin{equation}\label{eq:inequality-entropy-nbs-case2}
% e_{n+3\lfloor n^{1/2}\rfloor}(T_1)
% \le
% e_n(T_2)\,\bigl(1+o(1)\bigr),
% \qquad n\to\infty.
% \end{equation}
% Fix $\varepsilon>0$ and set
% \[
% n \coloneqq \left\lceil H_{T_2}(\varepsilon)\right\rceil .
% \]
% Then $H_{T_2}(\varepsilon)\le n$, and by definition of entropy numbers this implies
% \[
% e_n(T_2)\le \varepsilon .
% \]
% Inserting this bound into \eqref{eq:inequality-entropy-nbs-case2} yields
% \[
% e_{n+3\lfloor n^{1/2}\rfloor}(T_1)
% \le
% \varepsilon\,\bigl(1+o(1)\bigr).
% \]
% Hence there exists a sequence $\{\zeta_n'\}$ with $\zeta_n'\to1$ such that
% \[
% e_{n+3\lfloor n^{1/2}\rfloor}(T_1)\le \varepsilon\,\zeta_n'.
% \]
% By definition of entropy numbers,
% \[
% H_{T_1}(\varepsilon\,\zeta_n')
% \le
% n+3\lfloor n^{1/2}\rfloor .
% \]
% Since $T_2$ is compact, $H_{T_2}(\varepsilon)\to\infty$ as $\varepsilon\to0$, and
% therefore $n=\lceil H_{T_2}(\varepsilon)\rceil\to\infty$ as $\varepsilon\to0$.
% Moreover, $n=H_{T_2}(\varepsilon)+\theta_\varepsilon$ with
% $\theta_\varepsilon\in[0,1)$, so that
% \[
% n+3\lfloor n^{1/2}\rfloor
% =
% H_{T_2}(\varepsilon)
% +
% O(1)
% +
% O\!\left(\sqrt{H_{T_2}(\varepsilon)}\right),
% \qquad \varepsilon\to0.
% \]
% Dividing by $H_{T_2}(\varepsilon)$ and using $H_{T_2}(\varepsilon)\to\infty$ yields
% \[
% n+3\lfloor n^{1/2}\rfloor
% =
% H_{T_2}(\varepsilon)\,\bigl(1+o_{\varepsilon\to0}(1)\bigr).
% \]
% Consequently,
% \[
% H_{T_1}(\varepsilon\,\zeta_n')
% \le
% H_{T_2}(\varepsilon)\,\bigl(1+o_{\varepsilon\to0}(1)\bigr).
% \]
% \fi 

% \iffalse 
% Applying the corresponding entropy-number comparison
% \eqref{eq:inequality-entropy-nbs-case2} to this swapped representation, and
% repeating the same entropy-number–to–entropy conversion as in
% \eqref{eq:def-n-entropy}--\eqref{eq:relation-eps-entropy-nbs}, yields a sequence
% $\{\zeta_n'\}_{n\in\N}$ with $\zeta_n'\to1$ as $n\to\infty$ such that
% \begin{equation}\label{eq:useful-relation1-bis}
%     H_{T_1}\!\left(\varepsilon\,\zeta_n'\right)
%     \;\le\;
%     H_{T_2}(\varepsilon)\,\bigl(1+o_{\varepsilon\to0}(1)\bigr).
% \end{equation}



% As shown above, the structural hypothesis is invariant under exchanging
% $T_1$ and $T_2$, and the corresponding entropy-number inequality obtained
% from \eqref{eq:inequality-entropy-nbs} after swapping the operators reads
% (cf.~\eqref{eq:68})
% \begin{equation}\label{eq:inequality-entropy-nbs-swapped}
%     e_{n+3\lfloor n^{1/2}\rfloor}(T_1)
%     \le
%     e_n(T_2)\,\bigl(1+o(1)\bigr),
%     \qquad n\to\infty.
% \end{equation}

% Fix $\varepsilon>0$ and choose $n\in\N$ such that
% \begin{equation}\label{eq:def-n-entropy-swapped}
%     n+3\lfloor n^{1/2}\rfloor
%     < H_{T_1}(\varepsilon)
%     \le n+3\lfloor n^{1/2}\rfloor+1.
% \end{equation}
% Since $T_1$ is compact, $H_{T_1}(\varepsilon)\to\infty$ as $\varepsilon\to0$,
% and hence $n\to\infty$ as $\varepsilon\to0$.
% By definition of entropy numbers, this implies
% \[
% e_{n+3\lfloor n^{1/2}\rfloor}(T_1)\ge \varepsilon.
% \]
% Combining this with \eqref{eq:inequality-entropy-nbs-swapped} yields a sequence
% $\{\zeta_n'\}_{n\in\N}$ with $\zeta_n'\to1$ as $n\to\infty$ such that
% \[
% \varepsilon\,\zeta_n' \le e_n(T_2).
% \]

% Using again the definition of entropy numbers, we conclude that
% \begin{equation}\label{eq:useful-relation1-bis}
%     H_{T_2}\!\left(\varepsilon\,\zeta_n'\right)
%     \ge
%     n
%     =
%     H_{T_1}(\varepsilon)\,\bigl(1+o_{\varepsilon\to0}(1)\bigr).
% \end{equation}
% \fi 

% \iffalse 
% Again, using the definition of entropy numbers, we deduce from \eqref{eq;definition-k-entropy} and \eqref{eq:relation-eps-entropy-nbs} that
% \begin{equation}\label{eq:useful-relation1}
%     H_{T_1}\left(\varepsilon\zeta_n \right) 
%     \geq H_{T_1}\left(e_{n} (T_1) \right) 
%     \geq n
%     = H_{T_2}(\varepsilon)\left( 1+o_{\varepsilon\to 0}(1) \right).
% \end{equation}
% By interchangeability of $T_1$ and $T_2$, \eqref{eq:useful-relation1} can be reversed, that is,
% there exists a sequence $\{\zeta_j\}_{j\in\N}$, satisfying $\zeta_j\to 1$ as $j\to\infty$,
% such that
% \begin{equation}\label{eq:useful-relation1-bis}
%     H_{T_1}\left(\varepsilon\zeta_n' \right) 
%     \leq H_{T_2}(\varepsilon)\left( 1+o_{\varepsilon\to 0}(1) \right).
% \end{equation}
% % Now, using the fact that the entropy of the operator $T_{\symbolnot}$ goes to infinity as $\varepsilon \to 0$, we get
% % \begin{equation}\label{eq:useful-relation2}
% %     H_{\symbolnot}(\varepsilon) + (m+2)\sqrt{H_{\symbolnot}(\varepsilon)}
% %     \sim H_{\symbolnot}(\varepsilon) ,
% %     \quad \text{as }\varepsilon \to 0.
% % \end{equation}
% \fi 
%%%T9: Please check carefully the back and forth between T_1 and T_2 throughout and make sure to explain whenever we use this argument that we only need %statements that are invariant to the roles of T_1 and T_2.

Next, we observe that, owing to \eqref{eq:main-res-volume-integral} (which has already
been established for the Weyl quantization), we have
\begin{align}
H_{T_1}(\varepsilon\zeta_n)
&\sim
\int_{\varepsilon\zeta_n}^{\infty}\frac{V_\sigma(\lambda)}{\lambda}\,d\lambda \notag\\
&=
H_{T_1}(\varepsilon)
+
\int_{\varepsilon\zeta_n}^{\varepsilon}\frac{V_\sigma(\lambda)}{\lambda}\,d\lambda,
\label{eq:split-integral-zeta-proof-A}
\end{align}
where the second integral is understood with its natural orientation.
%Equivalently,
%\begin{align}
%H_{T_1}(\varepsilon\zeta_n)
%&\sim
%H_{T_1}(\varepsilon)
%+
%\int_{\min\{\varepsilon,\,\varepsilon\zeta_n\}}^{\max\{\varepsilon,\,\varepsilon\zeta_n\}}
%\frac{V_\sigma(\lambda)}{\lambda}\,d\lambda,
%\label{eq:split-integral-zeta-proof-B}
%\end{align}
%as \(\varepsilon\to0\). 
Using the monotonicity of $V_\sigma$, the remainder integral in
\eqref{eq:split-integral-zeta-proof-A} admits
the bound
\begin{equation}\label{eq:bound-remainder-volume-integral}
\left|
\int_{\varepsilon\zeta_n}^{\varepsilon}\frac{V_\sigma(\lambda)}{\lambda}\,d\lambda
\right|
\le
\frac{V_\sigma(\varepsilon\zeta_n^*)}{\varepsilon\zeta_n^*}
\,|\varepsilon-\varepsilon\zeta_n|
=
\frac{V_\sigma(\varepsilon\zeta_n^*)}{\zeta_n^*}\,|1-\zeta_n|,
\end{equation}
for all $\varepsilon>0$, where $\zeta_n^*:=\min\{1,\zeta_n\}$.
Moreover, by regular variation of $V_\sigma$ at zero and Karamata’s theorem for
integrals of regularly varying functions, we have
\begin{equation}\label{eq:volume-integral-compp}
V_\sigma(\varepsilon\zeta_n^*)
=
O_{\varepsilon\to0} \left(
\int_{\varepsilon\zeta_n^*}^{\infty}
\frac{V_\sigma(\lambda)}{\lambda}\,d\lambda
\right).
\end{equation}
Since $\zeta_n\to1$ as $n\to\infty$, it follows from
\eqref{eq:volume-integral-compp} that
\begin{equation}\label{eq:growth-control-volume}
V_\sigma(\varepsilon\zeta_n^*)\,|1-\zeta_n|
=
o_{\varepsilon\to0} \left(
\int_{\varepsilon\zeta_n^*}^{\infty}
\frac{V_\sigma(\lambda)}{\lambda}\,d\lambda
\right).
\end{equation}
Combining
\eqref{eq:split-integral-zeta-proof-A}, 
\eqref{eq:bound-remainder-volume-integral}, and
\eqref{eq:growth-control-volume}, we conclude that
\begin{equation}\label{eq:useful-relation3}
H_{T_1}(\varepsilon\zeta_n)
\sim
H_{T_1}(\varepsilon),
\qquad \varepsilon\to0.
\end{equation}
The same argument, with $\zeta_n$ replaced by $\zeta_n'$, yields
\begin{equation}\label{eq:useful-relation3-bis}
H_{T_1}(\varepsilon\zeta_n')
\sim
H_{T_1}(\varepsilon),
\qquad \varepsilon\to0.
\end{equation}
It therefore follows from
\eqref{eq:upper-bound-H}, \eqref{eq:lower-bound-H}, \eqref{eq:useful-relation3}, and \eqref{eq:useful-relation3-bis} that
\begin{equation}\label{eq:entropy-bounds-part1}
H_{T_1}(\varepsilon)
\sim
H_{T_2}(\varepsilon),
\qquad \varepsilon\to0,
\end{equation}
thereby completing the proof of Lemma~\ref{lem:technical-lemma}.
\end{proof}
\fi 

\color{black}

% \iffalse 
% Next, we observe that, owing to \textcolor{blue}{\eqref{eq:main-res-volume-integral} (which has already been proven for the Weyl quantization)}, 
% we have 
% \begin{align}
%     H_{T_1}\left(\varepsilon\zeta_n \right)  
%     \sim 
%     \int_{\varepsilon\zeta_n}^\infty 
%     \frac{V_{\symbolnot}(\lambda)}{\lambda} \diff \lambda
%     &= \int_{\varepsilon}^\infty 
%     \frac{V_{\symbolnot}(\lambda)}{\lambda} \diff \lambda
%     + \int_{\varepsilon\zeta_n}^\varepsilon 
%     \frac{V_{\symbolnot}(\lambda)}{\lambda} \diff \lambda \label{eq:split-integral-zeta-proof-A}\\
%     &\sim 
%     H_{T_1}\left(\varepsilon \right) 
%     + \int_{\varepsilon\zeta_n}^\varepsilon 
%     \frac{V_{\symbolnot}(\lambda)}{\lambda} \diff \lambda,\label{eq:split-integral-zeta-proof-B}
% \end{align}
% as $\varepsilon\to 0$.
% Using the monotonicity of $V_{\symbolnot}$,
% the second integral in \eqref{eq:split-integral-zeta-proof-A}--\eqref{eq:split-integral-zeta-proof-B} can be bounded according to
% \begin{equation}\label{eq:bound-remainder-volume-integral}
%     \left\lvert\int_{\varepsilon\zeta_n}^\varepsilon 
%     \frac{V_{\symbolnot}(\lambda)}{\lambda} \diff \lambda \right\rvert
%     \leq  \frac{V_{\symbolnot}(\varepsilon \zeta^*_n)}{\varepsilon \zeta^*_n}|\varepsilon - \varepsilon\zeta_n|
%     = \frac{V_{\symbolnot}(\varepsilon \zeta^*_n)}{\zeta^*_n}|1 - \zeta_n|,
% \end{equation}
% for all $ \varepsilon >0$,
% where $\zeta^*_n \coloneqq \min \{1, \zeta_n\}$.
% Moreover, 
% % integrating the asymptotic bound \eqref{eq:shubin-derivative-bound1}--\eqref{eq:shubin-derivative-bound2} from $\varepsilon\zeta_k'$ to $\infty$ yields 
% we get, by regular variation of $V_\sigma$, that 
% \begin{equation}\label{eq:volume-integral-compp}
%     V_{\symbolnot}(\varepsilon \zeta^*_n)
%     = \int_{\varepsilon\zeta^*_n}^\infty 
%     V'_{\symbolnot}(\lambda) \diff \lambda
%     = O_{\varepsilon \to 0}\left(\int_{\varepsilon\zeta^*_n}^\infty 
%     \frac{V_{\symbolnot}(\lambda)}{\lambda} \diff \lambda\right).
% \end{equation}
% (Recall that asymptotics are preserved through integration, see, e.g., \cite[Lemma~11]{allard2025metricentropyminimaxrisk}).
% Using $\zeta_n \to 1$ as $n\to\infty$, \eqref{eq:volume-integral-compp} implies 
% \begin{equation}\label{eq:growth-control-volume}
%     V_{\symbolnot}(\varepsilon \zeta^*_n)|1 - \zeta_n|
%     = o_{\varepsilon \to 0}\left(\int_{\varepsilon\zeta^*_n}^\infty 
%     \frac{V_{\symbolnot}(\lambda)}{\lambda} \diff \lambda\right).
% \end{equation}
% Putting \eqref{eq:split-integral-zeta-proof-A}--\eqref{eq:split-integral-zeta-proof-B}, \eqref{eq:bound-remainder-volume-integral}, and \eqref{eq:growth-control-volume} together yields 
% \begin{equation}\label{eq:useful-relation3}
%     H_{T_1}\left(\varepsilon\zeta_n\right) 
%     \sim H_{T_1}\left(\varepsilon\right),
%     \quad \text{as }\varepsilon \to 0.
% \end{equation}
% The same chain of arguments, starting from \eqref{eq:split-integral-zeta-proof-A}--\eqref{eq:split-integral-zeta-proof-B} up to \eqref{eq:useful-relation3},
% yields
% \begin{equation}\label{eq:useful-relation3-bis}
%     H_{T_1}\left(\varepsilon\zeta'_n\right) 
%     \sim H_{T_1}\left(\varepsilon\right),
%     \quad \text{as }\varepsilon \to 0.
% \end{equation}
% It hence follows  from \eqref{eq:useful-relation1}--\eqref{eq:useful-relation1-bis} and \eqref{eq:useful-relation3}--\eqref{eq:useful-relation3-bis} that
% \begin{equation}\label{eq:entropy-bounds-part1}  
%     H_{T_1}(\varepsilon)
%     \sim H_{T_2}\left(\varepsilon \right),
%      % \left( 1+o_{\varepsilon\to 0}(1) \right),
%     \quad\text{as } \varepsilon\to 0,
% \end{equation}
% % Conversely, a similar line of arguments starting from \eqref{eq:introduction-of-R2} instead of \eqref{eq:introduction-of-R1}, with the roles of $T_{\symbolnot}$ and $T_*$ reversed,  yields 
% % \begin{equation}\label{eq:entropy-bounds-part2}
% %      H_{\symbolnot}(\varepsilon)
% %     \leq  H_{T_*}\left(\varepsilon\right)\left( 1+o_{\varepsilon\to 0}(1) \right).
% % \end{equation}
% % From \eqref{eq:entropy-bounds-part1} and \eqref{eq:entropy-bounds-part2}, we get 
% % \begin{equation}\label{eq:equivalence-entropy-op-inverse-op}
% %     H_{\symbolnot}(\varepsilon) 
% %     \sim H_{T_*}\left(\varepsilon \right),
% %     \quad \text{as } \varepsilon \to 0,
% % \end{equation}
% thereby finishing the proof of Lemma~\ref{lem:technical-lemma}.
% \end{proof}
% \fi 









\subsection{Proof of Corollary~\ref{cor:ME-Main}}\label{sec:proof-ME-Main}

From the definition \eqref{eq:definition-volume-function} of $V_{\symbolnot}$, we have
\begin{align*}
\int_{\varepsilon}^{\infty}\frac{V_{\symbolnot}(\lambda)}{\lambda}\,d\lambda
&=
\int_{\varepsilon}^{\infty}
\int_{\mathbb{R}^{2d}}
\frac{\mathbbm{1}_{\{\symbolnot(x,\omega)>\lambda\}}}{\lambda}
\,dx\,d\omega\,d\lambda \\
&=
\int_{0}^{\infty}
\int_{\mathbb{R}^{2d}}
\frac{\mathbbm{1}_{\{\symbolnot(x,\omega)>\lambda>\varepsilon\}}}{\lambda}
\,dx\,d\omega\,d\lambda ,
\qquad \varepsilon>0.
\end{align*}
Since the integrand is nonnegative, Tonelli’s % Fubini’s 
theorem yields
\[
\int_{0}^{\infty}
\int_{\mathbb{R}^{2d}}
\frac{\mathbbm{1}_{\{\symbolnot(x,\omega)>\lambda>\varepsilon\}}}{\lambda}
\,dx\,d\omega\,d\lambda
=
\int_{\mathbb{R}^{2d}}
\int_{0}^{\infty}
\frac{\mathbbm{1}_{\{\symbolnot(x,\omega)>\lambda>\varepsilon\}}}{\lambda}
\,d\lambda\,dx\,d\omega, \qquad \varepsilon>0.
\]
For fixed $(x,\omega)\in\mathbb{R}^{2d}$, the inner integral can be computed explicitly:
\[
\int_{0}^{\infty}
\frac{\mathbbm{1}_{\{\symbolnot(x,\omega)>\lambda>\varepsilon\}}}{\lambda}
\,d\lambda
=
\ln_{+}\!\left(\frac{\symbolnot(x,\omega)}{\varepsilon}\right),
\]
where $\ln_{+}(t)=\max\{\ln (t),0\}$.
Consequently,
\[
\int_{\varepsilon}^{\infty}\frac{V_{\symbolnot}(\lambda)}{\lambda}\,d\lambda
=
\int_{\mathbb{R}^{2d}}
\ln_{+}\!\left(\frac{\symbolnot(x,\omega)}{\varepsilon}\right)
\,dx\,d\omega ,
\qquad \varepsilon>0.
\]
The claim now follows directly from Theorem~\ref{thm:main-result}.


% \iffalse
% \subsection{Proof of Corollary~\ref{cor:ME-Main}}\label{sec:proof-ME-Main}

% From the definition \eqref{eq:definition-volume-function} of $V_{\symbolnot}$, we get 
% \begin{align*}
%         \int_{\varepsilon}^\infty \frac{V_{\symbolnot}(\lambda)}{\lambda} \diff \lambda
%         &= \int_{\varepsilon}^\infty \int_{\mathbb{R}^{2d}} \frac{\mathbbm{1}_{\{\symbolnot(x, \omega)\, > \, \lambda\}}}{\lambda} \diff  x\diff \omega  \diff \lambda \\
%         &= \int_{0}^\infty \int_{\mathbb{R}^{2d}} \frac{\mathbbm{1}_{\{\symbolnot(x, \omega)\, > \, \lambda>\varepsilon\}}}{\lambda} \diff  x\diff \omega  \diff \lambda, 
%         \quad \text{for all } \varepsilon > 0.
% \end{align*}
% Applying Fubini's theorem yields 
% \begin{equation*}
%         \int_{0}^\infty \int_{\mathbb{R}^{2d}} \frac{\mathbbm{1}_{\{\symbolnot(x, \omega)\, > \, \lambda>\varepsilon\}}}{\lambda} \diff  x\diff \omega  \diff \lambda
%         =  \int_{\mathbb{R}^{2d}} \int_{0}^\infty \frac{\mathbbm{1}_{\{\symbolnot(x, \omega)\, > \, \lambda>\varepsilon\}}}{\lambda}  \diff \lambda \diff  x\diff \omega , 
%         \quad \text{for all } \varepsilon > 0.
% \end{equation*}
% We can compute the inner integral to obtain
% \begin{equation*}
%      \int_{0}^\infty \frac{\mathbbm{1}_{\{\symbolnot(x, \omega)\, > \, \lambda>\varepsilon\}}}{\lambda}  \diff \lambda
%      =  \int_{\varepsilon}^{\symbolnot(x, \omega)} \frac{\mathbbm{1}_{\{\symbolnot(x, \omega)\, > \varepsilon\}}}{\lambda}  \diff \lambda
%      = \ln_{+}\left(\frac{\symbolnot(x, \omega)}{\varepsilon}\right),
% \end{equation*}
% for all $\varepsilon > 0$ and $(x,\omega)\in\mathbb{R}^{2d}$.
% Therefore, 
% \begin{equation*}
%     \int_{\varepsilon}^\infty \frac{V_{\symbolnot}(\lambda)}{\lambda} \diff \lambda
%     = \int_{\R^{2d}} \ln_{+}\left(\frac{\symbolnot(x, \omega)}{\varepsilon}\right) \diff  x \diff \omega,
%     \quad \text{for all } \varepsilon > 0,
% \end{equation*}
% and the desired result follows by application of Theorem~\ref{thm:main-result}.
% \fi 

\subsection{Proof of Corollary~\ref{cor:MR-Main}}\label{sec:proof-MR-Main}

\iffalse 
We combine the theory of regular variation
(see, e.g., \cite{binghamRegularVariation1987})
with Theorem~\ref{thm:main-result} to derive the asymptotics of the
type-$\tau$ integrals \eqref{eq:type-tau-integrals}.
Specifically, for $\tau=2$ and $\tau=3$, definition \eqref{eq:type-tau-integrals} gives
\[
I_2(\varepsilon)
=
\int_{\varepsilon}^{\infty}
M_{\symbolnot}(\lambda)\,\lambda^{-2}\,d\lambda,
\qquad
I_3(\varepsilon)
=
\int_{\varepsilon}^{\infty}
M_{\symbolnot}(\lambda)\,\lambda^{-3}\,d\lambda.
\]
By Theorem~\ref{thm:main-result}, we have
\[
H_{\symbolnot}(\varepsilon)
\sim
\int_{\varepsilon}^{\infty}
\frac{V_{\symbolnot}(\lambda)}{\lambda}\,d\lambda,
\qquad \varepsilon\to 0,
\]
for all quantizations. Since the regular variation of $V_{\symbolnot}$
implies condition (RC), we may invoke \cite[Theorem~2]{allard2025metricentropyminimaxrisk}
to obtain
\[
H_{\symbolnot}(\varepsilon)\sim I_1(\varepsilon).
\]
It hence follows that
\[
I_1(\varepsilon)
\sim
\int_{\varepsilon}^{\infty}
\frac{V_{\symbolnot}(\lambda)}{\lambda}\,d\lambda,
\qquad \varepsilon\to 0.
\]
\fi 


We combine the theory of regular variation
(see, e.g., \cite{binghamRegularVariation1987})
with Theorem~\ref{thm:main-result} to derive the asymptotics of the
type-$\tau$ integrals \eqref{eq:type-tau-integrals}.
Specifically, for $\tau=2$ and $\tau=3$, definition \eqref{eq:type-tau-integrals} gives
\[
I_2(\varepsilon)
=
\int_{\varepsilon}^{\infty}
M_{\symbolnot}(\lambda)\,\lambda^{-2}\,d\lambda,
\qquad
I_3(\varepsilon)
=
\int_{\varepsilon}^{\infty}
M_{\symbolnot}(\lambda)\,\lambda^{-3}\,d\lambda.
\]
By Theorem~\ref{thm:main-result}, we have
\[
H_{\symbolnot}(\varepsilon)
\sim
\int_{\varepsilon}^{\infty}
\frac{V_{\symbolnot}(\lambda)}{\lambda}\,d\lambda,
\qquad \varepsilon\to 0,
\]
for all quantizations. Since $V_{\symbolnot}$ is regularly varying at zero,
condition~(RC) holds and \cite[Theorem~2]{allard2025metricentropyminimaxrisk} yields
\[
H_{\symbolnot}(\varepsilon)\sim I_1(\varepsilon),
\]
which, in turn, implies
\[
I_1(\varepsilon)
\sim
\int_{\varepsilon}^{\infty}
\frac{V_{\symbolnot}(\lambda)}{\lambda}\,d\lambda,
\qquad \varepsilon\to 0.
\]
% \textcolor{red}{
Using integration by parts, we obtain
\[
I_2(\varepsilon)
=
\varepsilon^{-1}I_1(\varepsilon)
-
\int_{\varepsilon}^{\infty}
\frac{I_1(\lambda)}{\lambda^2}\,d\lambda,
\]
and
\[
I_3(\varepsilon)
=
\varepsilon^{-2}I_1(\varepsilon)
-
2\int_{\varepsilon}^{\infty}
\frac{I_1(\lambda)}{\lambda^3}\,d\lambda.
\]
Since $I_1$ is regularly varying at zero, another application of
Karamata’s theorem implies
\[
I_2(\varepsilon)
\sim
\int_{\varepsilon}^{\infty}
\frac{V_{\symbolnot}(\lambda)}{\lambda^{2}}\,d\lambda,
\qquad
I_3(\varepsilon)
\sim
\int_{\varepsilon}^{\infty}
\frac{V_{\symbolnot}(\lambda)}{\lambda^{3}}\,d\lambda,
\qquad \varepsilon\to 0.
\]
% }
%Since $V_{\symbolnot}$ is regularly varying at zero,
%Karamata’s theorem then implies
%\[
%I_2(\varepsilon)
%\sim
%\int_{\varepsilon}^{\infty}
%\frac{V_{\symbolnot}(\lambda)}{\lambda^{2}}\,d\lambda,
%\qquad
%I_3(\varepsilon)
%\sim
%\int_{\varepsilon}^{\infty}
%\frac{V_{\symbolnot}(\lambda)}{\lambda^{3}}\,d\lambda,
%\qquad
%\varepsilon\to 0.
%\]
\iffalse 
We combine the theory of regular variation
(see, e.g., \cite{binghamRegularVariation1987})
with Theorem~\ref{thm:main-result} to derive the asymptotics of the
type-$\tau$ integral \eqref{eq:type-tau-integrals}.
Specifically, for $\tau=2$ and $\tau=3$, the definition \eqref{eq:type-tau-integrals} gives
$I_2(\varepsilon)=\int_{\varepsilon}^{\infty}M_{\symbolnot}(\lambda)\lambda^{-2}d\lambda$ and $I_3(\varepsilon)=\int_{\varepsilon}^{\infty}M_{\symbolnot}(\lambda)\lambda^{-3}d\lambda$.
\fi 
% As shown in the proof of Theorem~\ref{thm:main-result}, we have
% $M_{\textcolor{blue}{\symbolnot}}(\lambda)\sim V_{\symbolnot}(\lambda)$ as $\lambda\to0$, and hence
%%%T12: I am not sure I understand this, we need asymptotic equivalence of $M_{\textcolor{blue}{\symbolnot}}(\lambda)\sim V_{\symbolnot}(\lambda)$,
%further Theorem 1 specifies entropy so unclear why it is invoked here. Also unclear why we need Karamata here. I feel that the flow is broken here.
% \color{red}???
\iffalse 
\color{black}
\textcolor{blue}{By Theorem~\ref{thm:main-result}, we have
$I_1(\varepsilon)\sim\int_{\varepsilon}^{\infty}V_{\textcolor{blue}{\symbolnot}}(\lambda)\lambda^{-1}d\lambda$ as $\lambda\to0$, and hence, by Karamata's theorem for functions regularly varying at zero,}
\[
I_2(\varepsilon)\sim
\int_{\varepsilon}^{\infty}\frac{V_{\symbolnot}(\lambda)}{\lambda^2}\,d\lambda
\quad \text{and} \quad 
\textcolor{blue}{I_3(\varepsilon)\sim
\int_{\varepsilon}^{\infty}\frac{V_{\symbolnot}(\lambda)}{\lambda^3}\,d\lambda,}
\quad \varepsilon \to 0.
\]
\fi 
In particular, by \cite[Theorem~4]{allard2025metricentropyminimaxrisk},
\begin{equation}\label{eq:proof-risk-1}
R_{\symbolnot}(\kappa)
\sim
\kappa^2\,\varepsilon_\kappa
\int_{\varepsilon_\kappa}^{\infty}\frac{V_{\symbolnot}(\lambda)}{\lambda^2}\,d\lambda,
\qquad \kappa\to0,
\end{equation}
where the critical radius $\varepsilon_\kappa$ is determined by
\begin{equation}\label{eq:proof-risk-4}
\kappa^2
\int_{\varepsilon_\kappa}^{\infty}
V_{\symbolnot}(\lambda)
\left(\frac{2}{\lambda^3}-\frac{1}{\lambda^2\varepsilon_\kappa}\right)
\,d\lambda
=1,
\qquad \kappa>0.
\end{equation}
We next compute the integrals appearing in \eqref{eq:proof-risk-1} and
\eqref{eq:proof-risk-4} explicitly.
From the definition \eqref{eq:definition-volume-function} of $V_{\symbolnot}$,
\begin{align*}
\int_{\varepsilon}^{\infty}\frac{V_{\symbolnot}(\lambda)}{\lambda^2}\,d\lambda
&=
\int_{\varepsilon}^{\infty}
\int_{\mathbb{R}^{2d}}
\frac{\mathbbm{1}_{\{\symbolnot(x,\omega)>\lambda\}}}{\lambda^2}
\,dx\,d\omega\,d\lambda \\
&=
\int_{0}^{\infty}
\int_{\mathbb{R}^{2d}}
\frac{\mathbbm{1}_{\{\symbolnot(x,\omega)>\lambda>\varepsilon\}}}{\lambda^2}
\,dx\,d\omega\,d\lambda .
\end{align*}
Since the integrand is nonnegative, Tonelli’s theorem allows us to exchange the
order of integration, giving
\[
\int_{\mathbb{R}^{2d}}
\int_{0}^{\infty}
\frac{\mathbbm{1}_{\{\symbolnot(x,\omega)>\lambda>\varepsilon\}}}{\lambda^2}
\,d\lambda\,dx\,d\omega .
\]
For fixed $(x,\omega)\in\mathbb{R}^{2d}$, the inner integral can be evaluated explicitly:
\[
\int_{0}^{\infty}
\frac{\mathbbm{1}_{\{\symbolnot(x,\omega)>\lambda>\varepsilon\}}}{\lambda^2}
\,d\lambda
=
\left(\frac{1}{\varepsilon}-\frac{1}{\symbolnot(x,\omega)}\right)_+ .
\]
Consequently,
\begin{equation}\label{eq:proof-risk-2}
\int_{\varepsilon}^{\infty}\frac{V_{\symbolnot}(\lambda)}{\lambda^2}\,d\lambda
=
\int_{\mathbb{R}^{2d}}
\left(\frac{1}{\varepsilon}-\frac{1}{\symbolnot(x,\omega)}\right)_+
\,dx\,d\omega .
\end{equation}
Substituting \eqref{eq:proof-risk-2} into \eqref{eq:proof-risk-1} yields
\eqref{eq:main-res-risk-1-corr}.

A completely analogous computation gives
\begin{equation}\label{eq:proof-risk-3}
\int_{\varepsilon}^{\infty}\frac{V_{\symbolnot}(\lambda)}{\lambda^3}\,d\lambda
=
\int_{\mathbb{R}^{2d}}
\left(\frac{1}{2\varepsilon^2}-\frac{1}{2\symbolnot(x,\omega)^2}\right)_+
\,dx\,d\omega .
\end{equation}
Combining \eqref{eq:proof-risk-2} and \eqref{eq:proof-risk-3}, we obtain
\begin{align*}
\int_{\varepsilon}^{\infty}
V_{\symbolnot}(\lambda)
\left(\frac{2}{\lambda^3}-\frac{1}{\lambda^2\varepsilon}\right)
\,d\lambda
&=
\int_{\mathbb{R}^{2d}}
\left[
\left(\frac{1}{\varepsilon^2}-\frac{1}{\symbolnot(x,\omega)^2}\right)_+
-
\left(\frac{1}{\varepsilon^2}-\frac{1}{\symbolnot(x,\omega)\varepsilon}\right)_+
\right]
dx\,d\omega \\
&=
\int_{\mathbb{R}^{2d}}
\frac{1}{\symbolnot(x,\omega)}
\left(\frac{1}{\varepsilon}-\frac{1}{\symbolnot(x,\omega)}\right)_+
\,dx\,d\omega .
\end{align*}
Inserting this expression into \eqref{eq:proof-risk-4} yields
\eqref{eq:main-res-risk-2-corr}, thereby completing the proof.


% \iffalse
% \subsection{Proof of Corollary~\ref{cor:MR-Main}}\label{sec:proof-MR-Main}



% We can use the theory of regular variations (see, e.g., \cite{binghamRegularVariation1987} and \cite[Appendix~C]{allard2025metricentropyminimaxrisk}) together with Theorem~\ref{thm:main-result} to get the asymptotics of the type-$2$ 
% %and type-$3$ integrals 
% integral \eqref{eq:type-tau-integrals} according to 
% \begin{equation*}
%     I_2(\varepsilon) 
%     \sim \int_{\varepsilon}^\infty \frac{V_{\symbolnot}(\lambda)}{\lambda^2} \diff \lambda,
%     \quad \text{as } \varepsilon \to 0.
% \end{equation*}
% In particular, applying \cite[Theorem~4]{allard2025metricentropyminimaxrisk}
% we get 
% \begin{equation}\label{eq:proof-risk-1}
%     R_\symbolnot (\kappa)
%     \sim \kappa^2 \varepsilon_\kappa  \int_{\varepsilon_\kappa}^\infty \frac{V_{\symbolnot}(\lambda)}{\lambda^2} \diff \lambda,
%     \quad \text{as } \kappa \to 0,
% \end{equation}
% where the critical radius $\varepsilon_\kappa$ satisfies
% \begin{equation}\label{eq:proof-risk-4}
%     \kappa^2 \int_{\varepsilon_\kappa}^\infty V_{\symbolnot}(\lambda) \left(\frac{2}{\lambda^3}- \frac{1}{\lambda^2 \varepsilon_\kappa} \right) \diff \lambda = 1, 
%     \quad \text{for all } \kappa>0.
% \end{equation}


% From the definition \eqref{eq:definition-volume-function} of $V_{\symbolnot}$, we get 
% \begin{align*}
%         \int_{\varepsilon}^\infty \frac{V_{\symbolnot}(\lambda)}{\lambda^2} \diff \lambda
%         &= \int_{\varepsilon}^\infty \int_{\mathbb{R}^{2d}} \frac{\mathbbm{1}_{\{\symbolnot(x, \omega)\, > \, \lambda\}}}{\lambda^2} \diff  x\diff \omega  \diff \lambda \\
%         &= \int_{0}^\infty \int_{\mathbb{R}^{2d}} \frac{\mathbbm{1}_{\{\symbolnot(x, \omega)\, > \, \lambda>\varepsilon\}}}{\lambda^2} \diff  x\diff \omega  \diff \lambda, 
%         \quad \text{for all } \varepsilon > 0.
% \end{align*}
% Applying Fubini's theorem yields 
% \begin{equation*}
%         \int_{0}^\infty \int_{\mathbb{R}^{2d}} \frac{\mathbbm{1}_{\{\symbolnot(x, \omega)\, > \, \lambda>\varepsilon\}}}{\lambda^2} \diff  x\diff \omega  \diff \lambda
%         =  \int_{\mathbb{R}^{2d}} \int_{0}^\infty \frac{\mathbbm{1}_{\{\symbolnot(x, \omega)\, > \, \lambda>\varepsilon\}}}{\lambda^2}  \diff \lambda \diff  x\diff \omega , 
%         \quad \text{for all } \varepsilon > 0.
% \end{equation*}
% We can compute the inner integral to obtain
% \begin{equation*}
%      \int_{0}^\infty \frac{\mathbbm{1}_{\{\symbolnot(x, \omega)\, > \, \lambda>\varepsilon\}}}{\lambda^2}  \diff \lambda
%      =  \int_{\varepsilon}^{\symbolnot(x, \omega)} \frac{\mathbbm{1}_{\{\symbolnot(x, \omega)\, > \varepsilon\}}}{\lambda^2}  \diff \lambda
%      = \left(\frac{1}{\varepsilon}- \frac{1}{\symbolnot(x, \omega)}\right)_+,
% \end{equation*}
% for all $\varepsilon > 0$ and $(x,\omega)\in\mathbb{R}^{2d}$.
% Therefore, 
% \begin{equation}\label{eq:proof-risk-2}
%     \int_{\varepsilon}^\infty \frac{V_{\symbolnot}(\lambda)}{\lambda^2} \diff \lambda
%     = \int_{\R^{2d}} \left(\frac{1}{\varepsilon}- \frac{1}{\symbolnot(x, \omega)}\right)_+ \diff  x \diff \omega,
%     \quad \text{for all } \varepsilon > 0.
% \end{equation}
% Plugging \eqref{eq:proof-risk-2} in \eqref{eq:proof-risk-1} yields \eqref{eq:main-res-risk-1-corr}.



% Similar arguments as the ones employed to establish \eqref{eq:proof-risk-2} yield
% \begin{equation}\label{eq:proof-risk-3}
%     \int_{\varepsilon}^\infty \frac{V_{\symbolnot}(\lambda)}{\lambda^3} \diff \lambda
%     = \int_{\R^{2d}} \left(\frac{1}{2\varepsilon^2}- \frac{1}{2\symbolnot(x, \omega)^2}\right)_+ \diff  x \diff \omega,
%     \quad \text{for all } \varepsilon > 0.
% \end{equation}
% Combining \eqref{eq:proof-risk-2} and \eqref{eq:proof-risk-3} yields
% \begin{align*}
%     \int_{\varepsilon}^\infty V_{\symbolnot}(\lambda) \left(\frac{2}{\lambda^3}- \frac{1}{\lambda^2 \varepsilon} \right) \diff \lambda
%     &= \int_{\R^{2d}} \left(\frac{1}{\varepsilon^2}- \frac{1}{\symbolnot(x, \omega)^2}\right)_+ - \left(\frac{1}{\varepsilon^2}- \frac{1}{\symbolnot(x, \omega)\varepsilon }\right)_+  \diff  x \diff \omega \\
%     &= \int_{\R^{2d}} \frac{1}{\symbolnot(x, \omega)} \left(\frac{1}{\varepsilon}- \frac{1}{\symbolnot(x, \omega)}\right)_+  \diff  x \diff \omega,
% \end{align*}
% for all $\varepsilon>0$, which, when inserted in \eqref{eq:proof-risk-4}, yields \eqref{eq:main-res-risk-2-corr}.

% \fi 



\subsection{Proof of Theorem~\ref{thm:sob-sch}}\label{sec:proof-thm-sob-sch}

By \cite[Proposition~25.4]{shubin_pseudodifferential_2001}, $\ballsch$ is compact,
and therefore bounded, in $L^2(\mathbb R^d)$.
We next show that $\ker \tsch=\{0\}$.
If $\tsch\psi=0$ for some nonzero $\psi\in L^2(\R^d)$, then
$\mu\psi\in\ballsch$ for all $\mu\in\R$, which contradicts boundedness of
$\ballsch$ in $L^2(\R^d)$.
Since $\tsch$ is self-adjoint, $\ker \tsch^*=\ker \tsch=\{0\}$, and
\cite[Theorem~25.4]{shubin_pseudodifferential_2001} therefore guarantees the
existence of an inverse $\tsch^{-1}$,
which can be extended via \cite[Theorem~24.4]{shubin_pseudodifferential_2001} to a compact operator 
% \color{red}
on 
% \color{black} 
$L^2(\mathbb R^d)$.
% on $L^2(\mathbb R^d)$, which is again
% a pseudodifferential operator. 
Setting $g=\tsch f$ in \eqref{eq:definition-ballsch}, we obtain
\[
\ballsch
=
\{\tsch^{-1}g:\ \|g\|_{L^2(\mathbb R^d)}\le1\}
=
\tsch^{-1}\mathcal B_2,
\]
where $\mathcal B_2$ denotes the unit ball in $L^2(\mathbb R^d)$.
Consequently,
\begin{equation}\label{eq:equality-metric-entropy-op-entropy-sch}
H\!\left(\varepsilon;\ballsch,\|\cdot\|_{L^2(\mathbb R^d)}\right)
=
H_{\tsch^{-1}}(\varepsilon),
\qquad \varepsilon>0.
\end{equation}

To apply Theorem~\ref{thm:main-result}, we study the volume function
\begin{equation}\label{eq:integral-volume-sch-sob1}
V_{\symbolnotsch^{-1}}(\lambda)
=
\int_{\mathbb R^{2d}}
\mathbbm{1}_{\{(1+c\|x\|_2^2+(2\pi\|\omega\|_2)^2)^{-s/2}>\lambda\}}
\,dx\,d\omega,
\qquad \lambda>0.
\end{equation}
Performing the change of variables $x\mapsto x'/\sqrt c$ and
$\omega\mapsto\omega'/(2\pi)$ yields
\begin{align}
V_{\symbolnotsch^{-1}}(\lambda)
&=
\frac{1}{(2\pi\sqrt c)^d}
\int_{\mathbb R^{2d}}
\mathbbm{1}_{\{(1+\|x'\|_2^2+\|\omega'\|_2^2)^{-s/2}>\lambda\}}
\,dx'\,d\omega' \label{eq:volume-integral-sch-formula1}\\
&=
\frac{1}{(2\pi\sqrt c)^d}
\int_{\mathbb R^{2d}}
\mathbbm{1}_{\{\|x'\|_2^2+\|\omega'\|_2^2\, <\, \lambda^{-2/s}-1\}}
\,dx'\,d\omega'. \label{eq:volume-integral-sch-formula2}
\end{align}
Hence $V_{\symbolnotsch^{-1}}(\lambda)$ is proportional to the volume of a ball of
radius $(\lambda^{-2/s}-1)^{1/2}$ in $\mathbb R^{2d}$, and therefore
\[
V_{\symbolnotsch^{-1}}(\lambda)
=
\frac{\omega_{2d}}{(2\pi\sqrt c)^d}(\lambda^{-2/s}-1)^d
\sim
\frac{\omega_{2d}}{(2\pi\sqrt c)^d}\lambda^{-2d/s},
\qquad \lambda\to0.
\]
A direct calculation further shows that
\begin{equation}\label{eq:integral-sob-vol-1}
\int_{\varepsilon}^{\infty}
\frac{V_{\symbolnotsch^{-1}}(\lambda)}{\lambda}\,d\lambda
\textcolor{blue}{\sim}
\frac{s\,\omega_{2d}}{2d(2\pi\sqrt c)^d}\,\varepsilon^{-2d/s},
\qquad \varepsilon \textcolor{blue}{\to} 0.
\end{equation}
The symbol $\symbolnotsch^{-1}$ satisfies the assumptions of
Theorem~\ref{thm:main-result}.
Applying successively
\eqref{eq:equality-metric-entropy-op-entropy-sch},
\eqref{eq:main-result-part2},
\eqref{eq:main-res-volume-integral},
and \eqref{eq:integral-sob-vol-1}, we obtain
\begin{align*}
H\!\left(\varepsilon;\ballsch,\|\cdot\|_{L^2(\mathbb R^d)}\right)
&=
H_{\tsch^{-1}}(\varepsilon)
\sim
H_{\symbolnotsch^{-1}}(\varepsilon)\\
&\sim
\int_{\varepsilon}^{\infty}
\frac{V_{\symbolnotsch^{-1}}(\lambda)}{\lambda}\,d\lambda
\sim
\frac{s\,\omega_{2d}}{2d(2\pi\sqrt c)^d}\,\varepsilon^{-2d/s},
\qquad \varepsilon\to0,
\end{align*}
which establishes \eqref{eq:pinsker-entropy}.


Finally, by \cite[Appendix~A, Table~(iv)]{allard2025metricentropyminimaxrisk}, metric entropy scaling of the form 
\[
H(\varepsilon)\sim \frac{\mathfrak c\,\varepsilon^{-\alpha}}{\alpha},
\qquad \varepsilon\to0,
\]
with $\mathfrak c,\alpha>0$, is equivalent to minimax risk scaling
\[
R_\kappa
\sim
\frac{\alpha+2}{\alpha}
\left(
\frac{\mathfrak c\,\alpha\,\kappa^2}{(\alpha+1)(\alpha+2)}
\right)^{\!\frac{2}{\alpha+2}},
\qquad \kappa\to0.
\]
In view of \eqref{eq:pinsker-entropy}, for $\ballsch$ we may choose
\[
\mathfrak c=\frac{\omega_{2d}}{(2\pi\sqrt c)^d},
\qquad
\alpha=\frac{2d}{s},
\]
which yields
\[
R_\kappa(\ballsch)
\sim
\frac{d+s}{d}
\left(
\frac{d\,s\,\omega_{2d}\,\kappa^2}
{(2\pi\sqrt c)^d(d+s)(2d+s)}
\right)^{\!\frac{s}{d+s}},
\qquad \kappa\to0.
\]
This establishes \eqref{eq:pinsker-unbounded} and completes the proof of
Theorem~\ref{thm:sob-sch}.





% \iffalse 
% \subsection{Proof of Theorem~\ref{thm:sob-sch}}\label{sec:proof-thm-sob-sch}




% It follows from \cite[Proposition~25.4]{shubin_pseudodifferential_2001}, 
% that $\ballsch $ is compact in $W^{0,2}_{\mathcal{S}} = L^2(\R^d)$.
% In particular, the operator $\tsch$ has the trivial kernel $\{0\}$ by \cite[Theorem~25.4]{shubin_pseudodifferential_2001} and is thus invertible.
% Indeed, assuming, by way of contradiction, that there exists a nonzero $\psi\in L^2(\R^d)$ such that $\tsch\psi=0$,
% we have $\mu \psi \in \ballsch $, for all $\mu\in\R$, so that $\ballsch $ cannot be bounded, let alone compact, in $L^2(\R^d)$.
% Setting $g=\tsch f$ in \eqref{eq:definition-ballsch}, we obtain
% \begin{align*}
%     \ballsch
%     = \left\{\tsch^{-1} g \mid \left\| g \right\|_{L^2(\R^d)} \leq 1 \right\}
%     = \tsch^{-1}\mathcal{B}_2,
% \end{align*}
% where $\mathcal{B}_2$ denotes the unit ball in $L^2(\R^d)$.
% Consequently, %%%T5[OK]: holds for all eps >0?
% \begin{equation}\label{eq:equality-metric-entropy-op-entropy-sch}
%     H\left(\varepsilon;\ballsch, \|\cdot\|_{L^2(\R^d)}\right)
%     = H_{\tsch^{-1}}\left(\varepsilon\right),
%     \quad \varepsilon > 0.
% \end{equation} 
% With a view towards the application of Theorem~\ref{thm:main-result}, 
% we study the asymptotic properties of
% \begin{equation}\label{eq:integral-volume-sch-sob1}
%     V_{\symbolnotsch^{-1}}(\lambda)
%     =  \int_{\R^{2d}} \mathbbm{1}_{\left\{{\left(1+c \|x\|_2^2+(2\pi \|\omega\|_2)^2\right)^{-s/2}}\, >\, \lambda\right\}} \diff x \diff \omega,
% \end{equation}
% for all $\lambda >0$.
% To this end, we perform the change of variables $x \to x'/\sqrt{c}$ and $\omega \to \omega'/(2\pi)$ resulting in
% \begin{align}
%     V_{\symbolnotsch^{-1}}(\lambda)
%     &= \frac{1}{(2 \pi\sqrt{c})^{d}} \int_{\R^{2d}} \mathbbm{1}_{\left\{{\left(1+ \|x'\|_2^2+\|\omega'\|_2^2\right)^{-s/2}} \, > \, \lambda\right\}} \diff x' \diff \omega' \label{eq:volume-integral-sch-formula1} \\
%     &= \frac{1}{(2 \pi\sqrt{c})^{d}} \int_{\R^{2d}} \mathbbm{1}_{\left\{\|x'\|_2^2+\|\omega'\|_2^2 \, < \,  \lambda^{-2/s}-1\right\}} \diff x' \diff \omega'. \label{eq:volume-integral-sch-formula2}
% \end{align}
% It is apparent from \eqref{eq:volume-integral-sch-formula1}--\eqref{eq:volume-integral-sch-formula2} that $V_{\symbolnotsch^{-1}}(\lambda)$ is proportional to the volume of a ball of radius $(\lambda^{-2/s}-1)^{1/2}$ in $\R^{2d}$.
% Specifically, we have 
% \begin{equation*}
%     V_{\symbolnotsch^{-1}}(\lambda)
%     = \frac{\omega_{2d}}{(2 \pi\sqrt{c})^{d}} \left(\lambda^{-2/s}-1\right)^{d}
%     \sim \frac{\omega_{2d}}{(2 \pi\sqrt{c})^{d}} \lambda^{-2d/s},
%     \quad \text{as } \lambda \to 0.
% \end{equation*}
% \color{blue}
% A direct calculation further shows that 
% \begin{equation}\label{eq:integral-sob-vol-1}
%     \int_{\varepsilon}^\infty \frac{V_{\symbolnotsch^{-1}}(\lambda)}{\lambda} \diff \lambda 
%     = \frac{s \, \omega_{2d}}{2d(2 \pi\sqrt{c})^{d}} \varepsilon^{-2d/s}, 
%     \quad \text{for all } \varepsilon>0.
% \end{equation}
% One can easily verify that $\symbolnotsch^{-1}$ satisfies the assumptions of Theorem~\ref{thm:main-result}. Applying successively \eqref{eq:equality-metric-entropy-op-entropy-sch},
% \eqref{eq:main-result-part2},
% \eqref{eq:main-res-volume-integral},
% and \eqref{eq:integral-sob-vol-1} , yields
% %%%T5[OK]: removed ln(2)
% \begin{align*}
%     H\left(\varepsilon;\ballsch, \|\cdot\|_{L^2(\R^d)}\right)
%     = H_{\tsch^{-1}}\left(\varepsilon\right)
%     &\sim H_{{\symbolnotsch}^{-1}}\left(\varepsilon\right)\\
%     &\sim     \int_{\varepsilon}^\infty \frac{V_{\symbolnotsch^{-1}}(\lambda)}{\lambda} \diff \lambda 
%     \sim \frac{s \, \omega_{2d}}{2d(2 \pi\sqrt{c})^{d}} \varepsilon^{-2d/s},
%     \quad \text{as } \varepsilon \to 0,
% \end{align*}
% thereby establishing \eqref{eq:pinsker-entropy}.
% \color{black}

% Note that one gets from \cite[Theorem~7]{allard2025metricentropyminimaxrisk} (see also \cite[Section~3 and Appendix~A]{allard2025metricentropyminimaxrisk}) that 
% the metric entropy of an ellipsoid scaling as 
% \begin{equation*}
%     H(\varepsilon) \sim \frac{\mathfrak{c}\, \varepsilon^{-\alpha}}{\alpha},
%     \quad \text{as } \varepsilon \to 0,
% \end{equation*}
% for $c,\alpha>0$, is equivalent to the minimax risk scaling according to 
% \begin{equation*}
%     R_\kappa \sim \frac{\alpha+2}{\alpha} \left(\frac{\mathfrak{c}\,\alpha\,\kappa^2}{(\alpha+1)(\alpha+2)}\right)^{\frac{2}{2+\alpha}},
%     \quad \text{as } \kappa\to 0.
% \end{equation*}
% By \eqref{eq:pinsker-entropy}, for the ellipsoid $\ballsch$, we can choose 
% \begin{equation*}
%     \mathfrak{c}= \frac{\omega_{2d}}{(2 \pi\sqrt{c})^{d}} \quad \text{and} \quad \alpha = \frac{2d}{s},
% \end{equation*}
% so that 
% \begin{equation*}
%     R_\kappa \left(\ballsch\right) 
%     \sim \frac{d+s }{d} \left(\frac{d\, s\, \omega_{2d}\, \kappa^2}{(2 \pi\sqrt{c})^{d}(d+s)(2d+s)}\right)^{\frac{s}{d+s }},
%     \quad \text{as } \kappa \to 0,
% \end{equation*}
% which establishes \eqref{eq:pinsker-unbounded} and finishes the proof of Theorem~\ref{thm:sob-sch}.
% \fi 



\subsection{Proof of Theorem~\ref{thm:sob-weighted}}\label{sec:proof-sob-weighted}

%%%T10[0K]: The reference here is to bounded domains, need a reference to unbounded domains with decaying weights, also the ref. was Theorem 3.17, should be %2.17. Perhaps Adams Fournier. and shouldn't it be Kondrachev Rellich?
%%%T12: why do we need Exercise 2.12 here?
It follows from Rellich's compactness theorem (see \cite[Theorem~2.17 and Exercise~2.12]{hinzmicrolocal}) that $\ballweight$ is compact in $L^2(\R^d)$.
Proceeding as in the proof of Theorem~\ref{thm:sob-sch}, with
$\symbolnotsobw$ in place of $\symbolnotsch$,
we reduce the problem to determining the asymptotic behavior of
\begin{equation}\label{eq:integral-volume-weighted-sob1}
V_{\symbolnotsobw^{-1}}(\lambda)
=
\frac{1}{(2\pi\sqrt{c})^{d}}
\int_{\mathbb{R}^{2d}}
\mathbbm{1}_{\{(1+\|\omega'\|_2^2)^{-s/2}(1+\|x'\|_2^2)^{-r/2}>\lambda\}}
\,dx'\,d\omega',
\qquad \lambda>0.
\end{equation}
We rewrite the integral in polar coordinates to get
\[
\int_0^\infty\!\!\int_0^\infty
\mathbbm{1}_{\{(1+u_0^2)^{-s/2}(1+v_0^2)^{-r/2}>\lambda\}}
\,s_{d-1}^2 (u_0v_0)^{d-1}\,du_0\,dv_0,
\]
where $s_{d-1}$ denotes the surface measure of the unit sphere in $\mathbb{R}^d$.
With the change of variables $u=u_0^d$ and $v=v_0^d$, we get
\begin{equation}\label{eq:double-integral-sob-var-ch}
    V_{\symbolnotsobw^{-1}}(\lambda)
    =
    \frac{\omega_d^2}{(2\pi\sqrt{c})^{d}}
    \int_0^\infty\!\!\int_0^\infty
    \mathbbm{1}_{\{(1+u^{2/d})^{-s/2}(1+v^{2/d})^{-r/2}>\lambda\}}
    \,du\,dv,
\end{equation}
using $\omega_d = s_{d-1}/d$.


% \color{red}

Let us first assume that $r<s$. For $v\ge0$, the inequality
$(1+u^{2/d})^{-s/2}(1+v^{2/d})^{-r/2}>\lambda$
has solutions in $u$ only if
$\lambda(1+v^{2/d})^{r/2} < 1 $. Solving this inequality for $v$
gives $v < (\lambda^{-2/r}-1)^{d/2}$, and we therefore define
\begin{equation}\label{eq:def-v-lambda}
    v_\lambda \coloneqq (\lambda^{-2/r}-1)^{d/2}.
\end{equation}
Fix $v\in[0,v_\lambda)$. Then the condition
$(1+u^{2/d})^{-s/2}(1+v^{2/d})^{-r/2}>\lambda$
is equivalent to
$1+u^{2/d}<\lambda^{-2/s}(1+v^{2/d})^{-r/s}$,
and hence to
\[
u < \bigl(\lambda^{-2/s}(1+v^{2/d})^{-r/s}-1\bigr)^{d/2}.
\]
Consequently, integrating first with respect to $u$ gives
\[
V_{\symbolnotsobw^{-1}}(\lambda)
=
\frac{\omega_d^2}{(2\pi\sqrt{c})^{d}}
\int_0^{v_\lambda}
\left(
\lambda^{-2/s}(1+v^{2/d})^{-r/s}-1
\right)^{d/2}
\,dv,
\qquad \lambda>0.
\]
After the change of variables $t = v/v_\lambda$, the integral becomes
\begin{equation}\label{eq:vsigma-integral}
V_{\symbolnotsobw^{-1}}(\lambda)
=
\frac{\omega_d^2\, v_\lambda}{(2\pi\sqrt{c})^{d}}
\int_0^{1} u_\lambda(t)\,dt,
\qquad \lambda>0,
\end{equation}
where
\begin{equation}\label{eq:def-u-lambda}
u_\lambda(t)
\coloneqq
\left(
\lambda^{-2/s}\left(1+(t v_\lambda)^{2/d}\right)^{-r/s}-1
\right)^{d/2},
% \textcolor{red}{
    \qquad t\in[0,1).
% }
\end{equation}
% \textcolor{red}{
Note that the integral \eqref{eq:vsigma-integral} includes the endpoint $t=1$, although the admissible region corresponds to $t\in[0,1)$. This is justified since the singleton $\{1\}$ has Lebesgue measure zero and $u_\lambda$ extends continuously to $[0,1]$ with $u_\lambda(1)=0$.
% }
From the definition of $v_\lambda$ in \eqref{eq:def-v-lambda} we have $v_\lambda \sim \lambda^{-d/r}$, as 
$\lambda \to 0$. Letting $\gamma \coloneqq 2r/(ds)$ and using \eqref{eq:def-u-lambda}, we obtain the
pointwise limit
\[
u_\lambda(t) \longrightarrow
\left(t^{-\gamma}-1\right)^{d/2},
\qquad \lambda\to0,
\quad t\in(0,1).
\]
Since $r<s$, we have $\gamma<2/d$, and therefore the function
$t\mapsto (t^{-\gamma}-1)^{d/2}$ is integrable on $(0,1)$.
Moreover,
\[
\left(1+(t v_\lambda)^{2/d}\right)^{-r/s}
\le
t^{-\gamma}\left(1+v_\lambda^{2/d}\right)^{-r/s}
=
t^{-\gamma}\lambda^{2/s},
\qquad \lambda>0,\; t\in(0,1),
\]
which implies
\[
u_\lambda(t)
\le
\left(t^{-\gamma}-1\right)^{d/2}.
\]
Since $u_\lambda(t)\to (t^{-\gamma}-1)^{d/2}$ pointwise on $(0,1)$ as
$\lambda\to0$, and $|u_\lambda(t)|\le (t^{-\gamma}-1)^{d/2}$ with the latter
integrable on $(0,1)$, the hypotheses of the dominated convergence theorem are
satisfied. Consequently,
\[
\int_0^{1} u_\lambda(t)\,dt
\;\longrightarrow\;
\int_0^{1} (t^{-\gamma}-1)^{d/2}\,dt,
\qquad \lambda\to0 .
\]
Using the representation of $V_{\symbolnotsobw^{-1}}(\lambda)$ derived above
and applying the dominated convergence theorem, we obtain
\begin{equation}\label{eq:first-asymp-vol-r-leq-s}
V_{\symbolnotsobw^{-1}}(\lambda)
=
\frac{\omega_d^2\, v_\lambda}{(2\pi\sqrt{c})^{d}}
\int_0^{1} u_\lambda(t)\,dt
\;\sim\;
\frac{\omega_d^2\, v_\lambda}{(2\pi\sqrt{c})^{d}}
\int_0^{1} (t^{-\gamma}-1)^{d/2} dt,
\qquad \lambda \to 0 .
\end{equation}
Changing variables according to $t_0=t^\gamma$ in the integral on the
right-hand side of \eqref{eq:first-asymp-vol-r-leq-s}, we obtain
\begin{equation}\label{eq:binet-1}
    \int_0^{1} (t^{-\gamma}-1)^{d/2} dt
    =
    \frac{1}{\gamma}
    \int_0^{1}
    t_0^{\frac{1}{\gamma}-1-\frac{d}{2}}
    (1-t_0)^{d/2}
    \,dt_0
    =
    \frac{1}{\gamma}
    B \left(\frac{1}{\gamma}-\frac{d}{2}, \frac{d}{2}+1\right),
\end{equation}
where $B$ denotes Euler’s Beta function. Using the relation between the Beta
and Gamma functions (see \cite[Chapter~2]{artin}), 
we further get 
\begin{equation}\label{eq:binet-2}
    \frac{1}{\gamma}
    B \left(\frac{1}{\gamma}-\frac{d}{2}, \frac{d}{2}+1\right)
    =
    \frac{
        \Gamma \left(\frac{1}{\gamma}-\frac{d}{2}\right)
        \Gamma \left(\frac{d}{2}+1\right)
    }{
        \gamma\,\Gamma \left(\frac{1}{\gamma}+1\right)
    }
    =
    \frac{
        d\,\Gamma \left(\frac{d(s-r)}{2r}\right)\Gamma \left(\frac{d}{2}\right)
    }{
        2\,\Gamma \left(\frac{ds}{2r}\right)
    }.
\end{equation}
Combining \eqref{eq:def-v-lambda}, \eqref{eq:first-asymp-vol-r-leq-s},
\eqref{eq:binet-1}, and \eqref{eq:binet-2}, we conclude that
\begin{equation}\label{eq:final-asymp-vol-r-leq-s}
    V_{\symbolnotsobw^{-1}}(\lambda)
    \sim
    \frac{
        d\,\omega_d^2\,
        \Gamma \left(\frac{d(s-r)}{2r}\right)
        \Gamma \left(\frac{d}{2}\right)
    }{
        2(2\pi\sqrt{c})^{d}\,
        \Gamma \left(\frac{ds}{2r}\right)
    }
    \lambda^{-d/r},
    \qquad \lambda \to 0,
\end{equation}
in the case $r<s$.

By symmetry, the case $r>s$ is treated in the same way after interchanging the
roles of $r$ and $s$ throughout the preceding argument. In particular, the
auxiliary quantities $v_\lambda$, $u_\lambda$, and $\gamma$ are redefined with
$r$ and $s$ swapped, so that $\gamma=2s/(dr)$. This yields
\begin{equation}\label{eq:final-asymp-vol-r-geq-s}
    V_{\symbolnotsobw^{-1}}(\lambda)
    \sim \frac{d \,\omega_d^2\,  \Gamma\left(\frac{d(r-s)}{2s}\right) \Gamma \left( \frac{d}{2} \right)}{2(2\pi\sqrt{c})^{d}\, \Gamma \left(\frac{dr}{2s}\right)}\lambda^{-d/s},
    \quad \lambda \to 0.
\end{equation}

\iffalse 
For $r=s$, it is convenient to introduce the comparison integral
\begin{equation}\label{eq:integral-after-normalization-exponents}
\bar V_{\symbolnotsobw^{-1}}(\lambda)
=
\frac{\omega_d^2}{(2\pi\sqrt{c})^{d}}
\int_0^\infty\!\!\int_0^\infty
\mathbbm{1}_{\{(1+u)^{-s/d}(1+v)^{-s/d}>\lambda\}}
\,du\,dv .
\end{equation}
Fix $\lambda>0$ and $v\ge0$. For $u\ge0$, the inequality $(1+u)^{-s/d}(1+v)^{-s/d}>\lambda$ 
is equivalent to $u<\frac{\lambda^{-d/s}}{1+v}-1$. Hence admissible values of $u$ exist if and only if $0\le v<\lambda^{-d/s}-1$.
Integrating first with respect to $u$ and then with respect to $v$ yields
\begin{align*}
\int_0^\infty\!\!\int_0^\infty
\mathbbm{1}_{\{(1+u)^{-s/d}(1+v)^{-s/d}>\lambda\}}
\,du\,dv
&=
\int_0^{\lambda^{-d/s}-1}
\left(\frac{\lambda^{-d/s}}{1+v}-1\right)\,dv \\
&=
\frac{d}{s}\,\lambda^{-d/s}\ln(\lambda^{-1})
-
\bigl(\lambda^{-d/s}-1\bigr).
\end{align*}
Substituting this expression into \eqref{eq:integral-after-normalization-exponents}
gives
\begin{equation}\label{eq:vbar-asymp1}
\bar V_{\symbolnotsobw^{-1}}(\lambda)
\sim
\frac{d\,\omega_d^2}{s(2\pi\sqrt{c})^{d}}
\,\lambda^{-d/s}\ln(\lambda^{-1}),
\qquad \lambda\to0 .
\end{equation}
We now show that $V_{\symbolnotsobw^{-1}}(\lambda)$ has the same
asymptotic behavior. Observe that
\[
(1+u^{2/d})^{-s/2}\sim(1+u)^{-s/d},
\qquad u\to\infty .
\]
Consequently, for every fixed $\eta \in (0,1)$ there exists $R_\eta>0$ such that
\begin{equation}\label{eq:r-eta-inequalities}
(1-\eta)(1+u^{2/d})^{-s/2}
\le
(1+u)^{-s/d}
\le
(1+\eta)(1+u^{2/d})^{-s/2},
\qquad u\ge R_\eta ,
\end{equation}
and the same estimates hold with $u$ replaced by $v$. Using \eqref{eq:r-eta-inequalities}, we can compare the indicator constraints
in \eqref{eq:integral-after-normalization-exponents} and
\eqref{eq:double-integral-sob-var-ch} in the regime where both variables are large.
Indeed, whenever $u,v\ge R_\eta$, \eqref{eq:r-eta-inequalities} and its version with $u$ replaced by $v$ imply that
\[
(1-\eta)^2 (1+u^{2/d})^{-s/2}(1+v^{2/d})^{-s/2}
\le
(1+u)^{-s/d}(1+v)^{-s/d}
\le
(1+\eta)^2 (1+u^{2/d})^{-s/2}(1+v^{2/d})^{-s/2}.
\]
Consequently, for $u,v\ge R_\eta$, we obtain the set inclusions
\begin{equation}\label{eq:set-inclusion-1}
\{(1+u^{2/d})^{-s/2}(1+v^{2/d})^{-s/2}>\lambda\}
\subseteq
\{(1+u)^{-s/d}(1+v)^{-s/d}>(1-\eta)^2\lambda\}
\end{equation}
and
\begin{equation}\label{eq:set-inclusion-2}
\{(1+u)^{-s/d}(1+v)^{-s/d}>(1+\eta)^2\lambda\}
\subseteq
\{(1+u^{2/d})^{-s/2}(1+v^{2/d})^{-s/2}>\lambda\}.
\end{equation}
We next use these inclusions to compare the defining integrals of
$V_{\symbolnotsobw^{-1}}(\lambda)$ and $\bar V_{\symbolnotsobw^{-1}}(\lambda)$.
To this end, we split the double integral defining
$V_{\symbolnotsobw^{-1}}(\lambda)$ according to 
\[
V_{\symbolnotsobw^{-1}}(\lambda)
=
\frac{\omega_d^2}{(2\pi\sqrt{c})^{d}}
\left(
\int_{0}^{R_\eta}\!\!\int_{0}^{R_\eta}
+
2\int_{R_\eta}^{\infty}\!\!\int_{0}^{R_\eta}
+
\int_{R_\eta}^{\infty}\!\!\int_{R_\eta}^{\infty}
\right)
\mathbbm{1}_{\{(1+u^{2/d})^{-s/2}(1+v^{2/d})^{-s/2}>\lambda\}}
\,du\,dv .
\]
The first integral satisfies
\[
\int_{0}^{R_\eta}\!\!\int_{0}^{R_\eta}
\mathbbm{1}_{\{(1+u^{2/d})^{-s/2}(1+v^{2/d})^{-s/2}>\lambda\}}
\,du\,dv
\longrightarrow
R_\eta^{2},
\qquad \lambda\to0 .
\]
Next, for the mixed region we have
\begin{align*}
\int_{R_\eta}^{\infty}\!\!\int_{0}^{R_\eta}
\mathbbm{1}_{\{(1+u^{2/d})^{-s/2}(1+v^{2/d})^{-s/2}>\lambda\}}
\,du\,dv
&\le
R_\eta
\int_{R_\eta}^{\infty}
\mathbbm{1}_{\{(1+v^{2/d})^{-s/2}>\lambda\}}\,dv \\
&=
O_{\lambda\to0} \left(\lambda^{-d/s}\right).
\end{align*}
For the tail region $u,v\ge R_\eta$, we use
the set inclusions \eqref{eq:set-inclusion-1} and
\eqref{eq:set-inclusion-2} to get
\[
\int_{R_\eta}^{\infty}\!\!\int_{R_\eta}^{\infty}
\mathbbm{1}_{\{(1+u)^{-s/d}(1+v)^{-s/d}>(1+\eta)^2\lambda\}}
\,du\,dv
\;\le\;
\int_{R_\eta}^{\infty}\!\!\int_{R_\eta}^{\infty}
\mathbbm{1}_{\{(1+u^{2/d})^{-s/2}(1+v^{2/d})^{-s/2}>\lambda\}}
\,du\,dv
\]
and
\[
\int_{R_\eta}^{\infty}\!\!\int_{R_\eta}^{\infty}
\mathbbm{1}_{\{(1+u^{2/d})^{-s/2}(1+v^{2/d})^{-s/2}>\lambda\}}
\,du\,dv
\;\le\;
\int_{R_\eta}^{\infty}\!\!\int_{R_\eta}^{\infty}
\mathbbm{1}_{\{(1+u)^{-s/d}(1+v)^{-s/d}>(1-\eta)^2\lambda\}}
\,du\,dv .
\]
Combining these bounds with \eqref{eq:limit-double-int-simple-form} gives
\[
(1+\eta)^{-2d/s}
\le
\liminf_{\lambda\to0}
\frac{
\displaystyle
\int_{R_\eta}^{\infty}\!\!\int_{R_\eta}^{\infty}
\mathbbm{1}_{\{(1+u^{2/d})^{-s/2}(1+v^{2/d})^{-s/2}>\lambda\}}
\,du\,dv
}{
\frac{d}{s}\lambda^{-d/s}\ln(\lambda^{-1})
}
\]
and
\[
\limsup_{\lambda\to0}
\frac{
\displaystyle
\int_{R_\eta}^{\infty}\!\!\int_{R_\eta}^{\infty}
\mathbbm{1}_{\{(1+u^{2/d})^{-s/2}(1+v^{2/d})^{-s/2}>\lambda\}}
\,du\,dv
}{
\frac{d}{s}\lambda^{-d/s}\ln(\lambda^{-1})
}
\le
(1-\eta)^{-2d/s}.
\]
Since $\eta\in(0,1)$ is arbitrary, letting $\eta\downarrow0$ yields
\begin{equation}\label{eq:limit-double-int-simple-form-2}
\int_{R_\eta}^{\infty}\!\!\int_{R_\eta}^{\infty}
\mathbbm{1}_{\{(1+u^{2/d})^{-s/2}(1+v^{2/d})^{-s/2}>\lambda\}}
\,du\,dv
\sim
\frac{d}{s}\lambda^{-d/s}\ln(\lambda^{-1}),
\qquad \lambda\to0.
\end{equation}
Reinserting the bounded and mixed regions, whose total contribution is
$O(\lambda^{-d/s})$, we conclude that
\begin{equation}\label{eq:asymp-vol-case-r-equals-s}
V_{\symbolnotsobw^{-1}}(\lambda)
\sim
\frac{d\,\omega_d^2}{s(2\pi\sqrt{c})^{d}}
\lambda^{-d/s}\ln\!\left(\lambda^{-1}\right),
\qquad \lambda\to0.
\end{equation}



Using the set inclusions \eqref{eq:set-inclusion-1} and
\eqref{eq:set-inclusion-2} on the region $u,v\ge R_\eta$, and recalling that
the complementary region $\{u<R_\eta\}\cup\{v<R_\eta\}$ contributes only
$O(\lambda^{-d/s})$, we obtain
\[
\bar V_{\symbolnotsobw^{-1}}\!\bigl((1+\eta)^2\lambda\bigr)
+ O(\lambda^{-d/s})
\;\le\;
V_{\symbolnotsobw^{-1}}(\lambda)
\;\le\;
\bar V_{\symbolnotsobw^{-1}}\!\bigl((1-\eta)^2\lambda\bigr)
+ O(\lambda^{-d/s}),
\qquad \lambda\to0 .
\]
Since
\[
\lambda^{-d/s}
=
o\!\left(\lambda^{-d/s}\ln(\lambda^{-1})\right),
\qquad \lambda\to0,
\]
and $\bar V_{\symbolnotsobw^{-1}}(\lambda)\sim
\frac{d\,\omega_d^2}{s(2\pi\sqrt{c})^d}\lambda^{-d/s}\ln(\lambda^{-1})$,
it follows that
\[
V_{\symbolnotsobw^{-1}}(\lambda)
\sim
\bar V_{\symbolnotsobw^{-1}}(\lambda),
\qquad \lambda\to0.
\]


Since $\bar V_{\symbolnotsobw^{-1}}(\lambda)$ is regularly varying at
zero with index $-d/s$, rescaling the argument by the constant
$(1\pm\eta)^2$ does not change its leading-order asymptotics.
Letting $\eta\downarrow0$ therefore yields the asymptotic equivalence
\[
V_{\symbolnotsobw^{-1}}(\lambda)
\sim
\bar V_{\symbolnotsobw^{-1}}(\lambda),
\qquad \lambda\to0 .
\]

Combining this relation with \eqref{eq:vbar-asymp1} finally gives
\begin{equation}\label{eq:asymp-vol-case-r-equals-s}
V_{\symbolnotsobw^{-1}}(\lambda)
\sim
\frac{d\,\omega_d^2}{s(2\pi\sqrt{c})^{d}}
\lambda^{-d/s}\ln(\lambda^{-1}),
\qquad \lambda\to0 .
\end{equation}




Therefore, in the region $u,v\ge R_\eta$ the domains of integration
defining $V_{\symbolnotsobw^{-1}}(\lambda)$ and
$\bar V_{\symbolnotsobw^{-1}}(\lambda)$ coincide up to the replacement
of $\lambda$ by $(1\pm\eta)^2\lambda$.

It remains to control the contribution of the complementary region
$\{u<R_\eta\}\cup\{v<R_\eta\}$.
As in the analysis of $\bar V_{\symbolnotsobw^{-1}}(\lambda)$ above,
one obtains
\[
\int_{R_\eta}^{\infty}\!\!\int_0^{R_\eta}
\mathbbm{1}_{\{(1+u^{2/d})^{-s/2}(1+v^{2/d})^{-s/2}>\lambda\}}
\,du\,dv
=
O_{\lambda\to0}\!\left(\lambda^{-d/s}\right),
\]
and the same bound holds for the symmetric region
$u\in[0,R_\eta],\,v\in[R_\eta,\infty)$.
The contribution of the bounded region
$u,v\in[0,R_\eta]$ converges to $R_\eta^2$ as $\lambda\to0$.

Hence the entire contribution of
$\{u<R_\eta\}\cup\{v<R_\eta\}$ is of order
$O(\lambda^{-d/s})$, whereas the leading term obtained in
\eqref{eq:vbar-asymp1} behaves as
$\lambda^{-d/s}\ln(\lambda^{-1})$.
Since $\ln(\lambda^{-1})\to\infty$ as $\lambda\to0$, the latter
dominates the former, and the asymptotic behavior is therefore
governed entirely by the region $u,v\ge R_\eta$.


Combining these observations with \eqref{eq:vbar-asymp1} yields
\[
V_{\symbolnotsobw^{-1}}(\lambda)
\sim
\bar V_{\symbolnotsobw^{-1}}(\lambda),
\qquad \lambda\to0 .
\]
Consequently,
\begin{equation}\label{eq:asymp-vol-case-r-equals-s}
V_{\symbolnotsobw^{-1}}(\lambda)
\sim
\frac{d\,\omega_d^2}{s(2\pi\sqrt{c})^{d}}
\lambda^{-d/s}\ln\!\left(\lambda^{-1}\right),
\qquad \lambda\to0 .
\end{equation}

\iffalse
For $r=s$, it is convenient to introduce the comparison integral
\begin{equation}\label{eq:integral-after-normalization-exponents}
\bar V_{\symbolnotsobw^{-1}}(\lambda)
=
\frac{\omega_d^2}{(2\pi\sqrt{c})^{d}}
\int_0^\infty\!\!\int_0^\infty
\mathbbm{1}_{\{(1+u)^{-s/d}(1+v)^{-s/d}>\lambda\}}
\,du\,dv .
\end{equation}
Fix $\lambda>0$ and $v\ge0$. For $u\ge0$, the inequality
$(1+u)^{-s/d}(1+v)^{-s/d}>\lambda$ is equivalent to
$u<\frac{\lambda^{-d/s}}{1+v}-1$.
Hence admissible values of $u$ exist if and only if
$0\le v<\lambda^{-d/s}-1$. 
Integrating first with respect to $u$ and then with respect to $v$, we obtain
\begin{align*}
\int_0^\infty\!\!\int_0^\infty
\mathbbm{1}_{\{(1+u)^{-s/d}(1+v)^{-s/d}>\lambda\}}
\,du\,dv
&=
\int_0^{\lambda^{-d/s}-1}
\left(\frac{\lambda^{-d/s}}{1+v}-1\right)\,dv \\
&=
\frac{d}{s}\,\lambda^{-d/s}\ln(\lambda^{-1})
-
\bigl(\lambda^{-d/s}-1\bigr).
\end{align*}
Substituting this identity into \eqref{eq:integral-after-normalization-exponents}
yields
\begin{equation}\label{eq:vbar-asymp1}
\bar V_{\symbolnotsobw^{-1}}(\lambda)
\sim
\frac{d\,\omega_d^2}{s(2\pi\sqrt{c})^{d}}
\,\lambda^{-d/s}\ln(\lambda^{-1}),
\qquad \lambda\to0.
\end{equation}
\iffalse 
If $r=s$, it is convenient
to introduce the comparison integral
\begin{equation}\label{eq:integral-after-normalization-exponents}
\bar V_{\symbolnotsobw^{-1}}(\lambda)
=
\frac{\omega_d^2}{(2\pi\sqrt{c})^{d}}
\int_0^\infty\!\!\int_0^\infty
\mathbbm{1}_{\{(1+u)^{-s/d}(1+v)^{-s/d}>\lambda\}}
\,du\,dv.
\end{equation}
Fix $\lambda>0$ and $v\ge0$. For $u\ge0$, the inequality
\[
(1+u)^{-\frac{s}{d}}(1+v)^{-\frac{s}{d}}>\lambda
\quad \text{is equivalent to} \quad 
u<\frac{\lambda^{-\frac{d}{s}}}{1+v}-1.
\]
Therefore, for admissible values of $u$ to exist, we must have
$0 \le v < \lambda^{-d/s}-1$.
Integrating first in $u$ and then in $v$ therefore gives
\begin{align*}
    \int_0^\infty\!\!\int_0^\infty
\mathbbm{1}_{\{(1+u)^{-s/d}(1+v)^{-s/d}>\lambda\}}
\,du\,dv
&=
\int_0^{\lambda^{-\frac{d}{s}}-1}
\left(\frac{\lambda^{-\frac{d}{s}}}{1+v}-1\right)\,dv \\
&= \frac{d\,\lambda^{-\frac{d}{s}} }{s} \ln\left(\lambda^{-1}\right)
- \left(\lambda^{-\frac{d}{s}}-1\right).
\end{align*}
When injected in \eqref{eq:integral-after-normalization-exponents}, 
this yields 
\begin{equation}\label{eq:vbar-asymp1}
\bar V_{\symbolnotsobw^{-1}}(\lambda)
\sim
\frac{d\, \omega_d^2}{s\, (2\pi\sqrt{c})^{d}}
\, \lambda^{-\frac{d}{s}}\ln\left(\lambda^{-1}\right),
\qquad \lambda\to0.
\end{equation}
\fi 
We next show that $V_{\symbolnotsobw^{-1}}(\lambda)$ and
$\bar V_{\symbolnotsobw^{-1}}(\lambda)$ are asymptotically equivalent.
To this end, we first observe that
\[
(1+u^{2/d})^{-s/2} \sim (1+u)^{-s/d},
\qquad u\to\infty.
\]
Consequently, for every fixed $\eta>0$ there exists $R_\eta>0$ such that
\begin{equation}\label{eq:r-eta-inequalities}
(1-\eta)(1+u^{2/d})^{-s/2}
\le (1+u)^{-s/d}
\le (1+\eta)(1+u^{2/d})^{-s/2},
\qquad u\ge R_\eta .
\end{equation}
The same estimates hold with $u$ replaced by $v$. We next decompose the double integral defining $\bar V_{\symbolnotsobw^{-1}}(\lambda)$
according to whether $u$ and $v$ lie below or above the threshold $R_\eta$.
Writing the integral as a sum over these regions yields
\[
\bar V_{\symbolnotsobw^{-1}}(\lambda)
=
\frac{\omega_d^2}{(2\pi\sqrt{c})^{d}}
\left(
\int_{0}^{R_\eta}\!\!\int_{0}^{R_\eta}
+
2\int_{R_\eta}^{\infty}\!\!\int_{0}^{R_\eta}
+
\int_{R_\eta}^{\infty}\!\!\int_{R_\eta}^{\infty}
\right)
\mathbbm{1}_{\{(1+u)^{-s/d}(1+v)^{-s/d}>\lambda\}}
\,du\,dv .
\]
First observe that
\[
\int_{0}^{R_\eta}\!\!\int_{0}^{R_\eta}
\mathbbm{1}_{\{(1+u)^{-s/d}(1+v)^{-s/d}>\lambda\}}
\,du\,dv
\longrightarrow R_\eta^{2},
\qquad \lambda\to0 .
\]
Indeed, for fixed $u,v\in[0,R_\eta]$ the quantity
$(1+u)^{-s/d}(1+v)^{-s/d}$ is strictly positive, so the indicator
eventually equals $1$ as $\lambda\to0$. Next, for the mixed region we have
\begin{align*}
\int_{R_\eta}^{\infty}\!\!\int_{0}^{R_\eta}
\mathbbm{1}_{\{(1+u)^{-s/d}(1+v)^{-s/d}>\lambda\}}
\,du\,dv
&\le
R_\eta
\int_{R_\eta}^{\infty}
\mathbbm{1}_{\{(1+v)^{-s/d}>\lambda\}}\,dv \\
&=
O_{\lambda\to0} \left(\lambda^{-d/s}\right).
\end{align*}
Consequently,
\[
\bar V_{\symbolnotsobw^{-1}}(\lambda)
=
\frac{\omega_d^{2}}{(2\pi\sqrt{c})^{d}}
\left(
\int_{R_\eta}^{\infty}\!\!\int_{R_\eta}^{\infty}
\mathbbm{1}_{\{(1+u)^{-s/d}(1+v)^{-s/d}>\lambda\}}
\,du\,dv
\right)
+
O_{\lambda\to0} \left(\lambda^{-d/s}\right).
\]
Combining this decomposition with the asymptotic formula
\eqref{eq:vbar-asymp1} yields
\begin{equation}\label{eq:limit-double-int-simple-form}
\int_{R_\eta}^{\infty}\!\int_{R_\eta}^{\infty}
\mathbbm{1}_{\{(1+u)^{-s/d}(1+v)^{-s/d}>\lambda\}}
\,du\,dv
\sim
\frac{d}{s}\lambda^{-d/s}\ln \left(\lambda^{-1}\right),
\qquad \lambda\to 0.
\end{equation}
\iffalse 
First, we observe that
\begin{equation*}
    \int_0^{R_\eta}\!\!\int_0^{R_\eta}
    \mathbbm{1}_{\{(1+u)^{-s/d}(1+v)^{-s/d}>\lambda\}}
    \,du\,dv
    \xrightarrow{} R_\eta^2,
    \quad \text{as }\lambda \to 0,
\end{equation*}
and that 
\begin{align*}
    \int_{R_\eta}^\infty\!\!\int_0^{R_\eta}
    \mathbbm{1}_{\{(1+u)^{-s/d}(1+v)^{-s/d}>\lambda\}}
    \,du\,dv
    \leq R_\eta \int_{R_\eta}^\infty \mathbbm{1}_{\{(1+v)^{-s/d}>\lambda\}}
    \,dv = O_{\lambda \to 0} \left(\lambda^{-\frac{d}{s}}\right).
\end{align*}
Therefore, we get 
\begin{equation*}
    \bar V_{\symbolnotsobw^{-1}}(\lambda)
    = \frac{\omega_d^2}{(2\pi\sqrt{c})^{d}}
    \left(
    \int_{R_\eta}^\infty\!\!\int_{R_\eta}^\infty 
    \mathbbm{1}_{\{(1+u)^{-s/d}(1+v)^{-s/d}>\lambda\}}
    \,du\,dv\right) 
    + O_{\lambda \to 0} \left(\lambda^{-\frac{d}{s}}\right).
\end{equation*}
Together with \eqref{eq:vbar-asymp1}, this implies 
\begin{equation}\label{eq:limit-double-int-simple-form}
    \int_{R_\eta}^\infty\!\!\int_{R_\eta}^\infty 
    \mathbbm{1}_{\{(1+u)^{-s/d}(1+v)^{-s/d}>\lambda\}}
    \,du\,dv
    \sim 
    \frac{d}{s}\lambda^{-\frac{d}{s}}\ln\left(\lambda^{-1}\right),
    \quad \text{as } \lambda\to0.
\end{equation}
\fi 
Using the definition of $R_\eta$ in \eqref{eq:r-eta-inequalities}, we obtain
\[
\int_{R_\eta}^{\infty}\!\int_{R_\eta}^{\infty}
\mathbbm{1}_{\{(1+u^{2/d})^{-s/2}(1+v^{2/d})^{-s/2}\,>\,\lambda\}}
\,du\,dv
\le
\int_{R_\eta}^{\infty}\!\int_{R_\eta}^{\infty}
\mathbbm{1}_{\{(1+u)^{-s/d}(1+v)^{-s/d}\,>\,(1-\eta)^2\lambda\}}
\,du\,dv
\]
and
\[
\int_{R_\eta}^{\infty}\!\int_{R_\eta}^{\infty}
\mathbbm{1}_{\{(1+u^{2/d})^{-s/2}(1+v^{2/d})^{-s/2}\,>\,\lambda\}}
\,du\,dv
\ge
\int_{R_\eta}^{\infty}\!\int_{R_\eta}^{\infty}
\mathbbm{1}_{\{(1+u)^{-s/d}(1+v)^{-s/d}\,>\,(1+\eta)^2\lambda\}}
\,du\,dv .
\]
Combining these bounds with \eqref{eq:limit-double-int-simple-form} yields
\[
(1+\eta)^{-2d/s}
\le
\liminf_{\lambda\to0}
\frac{
\displaystyle
\int_{R_\eta}^{\infty}\!\!\int_{R_\eta}^{\infty}
\mathbbm{1}_{\{(1+u^{2/d})^{-s/2}(1+v^{2/d})^{-s/2}>\lambda\}}
\,du\,dv
}{
\frac{d}{s}\lambda^{-d/s}\ln(\lambda^{-1})
}
\]
and
\[
\limsup_{\lambda\to0}
\frac{
\displaystyle
\int_{R_\eta}^{\infty}\!\!\int_{R_\eta}^{\infty}
\mathbbm{1}_{\{(1+u^{2/d})^{-s/2}(1+v^{2/d})^{-s/2}>\lambda\}}
\,du\,dv
}{
\frac{d}{s}\lambda^{-d/s}\ln(\lambda^{-1})
}
\le
(1-\eta)^{-2d/s}.
\]
Since $\eta>0$ is arbitrary, letting $\eta\downarrow0$ gives
\begin{equation}\label{eq:limit-double-int-simple-form-2}
\int_{R_\eta}^{\infty}\!\!\int_{R_\eta}^{\infty}
\mathbbm{1}_{\{(1+u^{2/d})^{-s/2}(1+v^{2/d})^{-s/2}>\lambda\}}
\,du\,dv
\sim
\frac{d}{s}\lambda^{-d/s}\ln(\lambda^{-1}),
\qquad \lambda\to0.
\end{equation}
\iffalse
Now, using the definition of $R_\eta$ in \eqref{eq:r-eta-inequalities}, we have 
\begin{equation*}
    \int_{R_\eta}^\infty\!\!\int_{R_\eta}^\infty 
    \mathbbm{1}_{\{(1+u^{2/d})^{-s/2}(1+v^{2/d})^{-s/2}>\lambda\}}
    \,du\,dv
    \leq \int_{R_\eta}^\infty\!\!\int_{R_\eta}^\infty 
    \mathbbm{1}_{\{(1+u)^{-s/d}(1+v)^{-s/d}>(1-\eta)^2\lambda\}}
    \,du\,dv
\end{equation*}
and
\begin{equation*}
    \int_{R_\eta}^\infty\!\!\int_{R_\eta}^\infty 
    \mathbbm{1}_{\{(1+u^{2/d})^{-s/2}(1+v^{2/d})^{-s/2}>\lambda\}}
    \,du\,dv
    \geq \int_{R_\eta}^\infty\!\!\int_{R_\eta}^\infty 
    \mathbbm{1}_{\{(1+u)^{-s/d}(1+v)^{-s/d}>(1+\eta)^2\lambda\}}
    \,du\,dv.
\end{equation*}
From the asymptotic behavior \eqref{eq:limit-double-int-simple-form}, we thus get 
\begin{equation*}
    (1+\eta)^{-\frac{2d}{s}}
    \leq  \lim_{\lambda\to 0} \int_{R_\eta}^\infty\!\!\int_{R_\eta}^\infty 
     \frac{\mathbbm{1}_{\{(1+u^{2/d})^{-s/2}(1+v^{2/d})^{-s/2}>\lambda\}}}{\frac{d}{s}\lambda^{-\frac{d}{s}}\ln\left(\lambda^{-1}\right)}
    \,du\,dv
    \leq (1-\eta)^{-\frac{2d}{s}},
\end{equation*}
and, choosing $\eta$ arbitrarily small yields
\begin{equation}\label{eq:limit-double-int-simple-form-2}
    \int_{R_\eta}^\infty\!\!\int_{R_\eta}^\infty 
    \mathbbm{1}_{\{(1+u^{2/d})^{-s/2}(1+v^{2/d})^{-s/2}>\lambda\}}
    \,du\,dv
    \sim 
    \frac{d}{s}\lambda^{-\frac{d}{s}}\ln\left(\lambda^{-1}\right),
    \quad \text{as } \lambda\to0.
\end{equation}
\fi 
A decomposition of the double integral defining $V_{\symbolnotsobw^{-1}}(\lambda)$
analogous to that used for $\bar V_{\symbolnotsobw^{-1}}(\lambda)$ shows that 
\begin{equation}\label{eq:asymp-vol-case-r-equals-s}
V_{\symbolnotsobw^{-1}}(\lambda)
\sim
\frac{d\,\omega_d^2}{s(2\pi\sqrt{c})^{d}}
\lambda^{-d/s}\ln \left(\lambda^{-1}\right),
\qquad \lambda\to0 .
\end{equation}
\iffalse
A decomposition of the double integral defining $V_{\symbolnotsobw^{-1}}(\lambda)$
analogous to that used for $\bar V_{\symbolnotsobw^{-1}}(\lambda)$ yields
\[
\int_0^{R_\eta}\!\!\int_0^{R_\eta}
\mathbbm{1}_{\{(1+u^{2/d})^{-s/2}(1+v^{2/d})^{-s/2}>\lambda\}}
\,du\,dv
\longrightarrow
R_\eta^2,
\qquad \lambda\to0,
\]
and
\begin{align*}
\int_{R_\eta}^{\infty}\!\!\int_0^{R_\eta}
\mathbbm{1}_{\{(1+u^{2/d})^{-s/2}(1+v^{2/d})^{-s/2}>\lambda\}}
\,du\,dv
&\le
R_\eta
\int_{R_\eta}^{\infty}
\mathbbm{1}_{\{(1+v^{2/d})^{-s/2}>\lambda\}}
\,dv \\
&=
O_{\lambda\to0}\!\left(\lambda^{-d/s}\right).
\end{align*}
The same estimate holds for the symmetric mixed region. Consequently,
\[
V_{\symbolnotsobw^{-1}}(\lambda)
=
\frac{\omega_d^2}{(2\pi\sqrt{c})^{d}}
\left(
\int_{R_\eta}^{\infty}\!\!\int_{R_\eta}^{\infty}
\mathbbm{1}_{\{(1+u^{2/d})^{-s/2}(1+v^{2/d})^{-s/2}>\lambda\}}
\,du\,dv
\right)
+
O_{\lambda\to0}\!\left(\lambda^{-d/s}\right).
\]
Combining this decomposition with \eqref{eq:limit-double-int-simple-form-2}, we conclude that
\begin{equation}\label{eq:asymp-vol-case-r-equals-s}
V_{\symbolnotsobw^{-1}}(\lambda)
\sim
\frac{d\,\omega_d^2}{s(2\pi\sqrt{c})^{d}}
\lambda^{-d/s}\ln\!\left(\lambda^{-1}\right),
\qquad \lambda\to0.
\end{equation}
\fi 
\iffalse 
A similar splitting of the double integral as was done for $\bar V_{\symbolnotsobw^{-1}}(\lambda)$
yields 
\begin{equation*}
    \int_0^{R_\eta}\!\!\int_0^{R_\eta}\mathbbm{1}_{\{(1+u^{2/d})^{-s/2}(1+v^{2/d})^{-s/2}>\lambda\}}
    \,du\,dv
    \xrightarrow{} R_\eta^2,
    \quad \text{as }\lambda \to 0,
\end{equation*}
and  
\begin{align*}
    \int_{R_\eta}^\infty\!\!\int_0^{R_\eta} \mathbbm{1}_{\{(1+u^{2/d})^{-s/2}(1+v^{2/d})^{-s/2}>\lambda\}}
    \,du\,dv
    \leq R_\eta \int_{R_\eta}^\infty \mathbbm{1}_{\{(1+v^{2/d})^{-s/2}>\lambda\}}
    \,dv 
    = O_{\lambda \to 0} \left(\lambda^{-\frac{d}{s}}\right).
\end{align*}
Therefore, 
\begin{equation*}
    V_{\symbolnotsobw^{-1}}(\lambda)
    = \frac{\omega_d^2}{(2\pi\sqrt{c})^{d}}
    \left(
    \int_{R_\eta}^\infty\!\!\int_{R_\eta}^\infty 
    \mathbbm{1}_{\{(1+u^{2/d})^{-s/2}(1+v^{2/d})^{-s/2}>\lambda\}}
    \,du\,dv\right) 
    + O_{\lambda \to 0} \left(\lambda^{-\frac{d}{s}}\right),
\end{equation*}
which implies, using \eqref{eq:limit-double-int-simple-form-2}, that 
\begin{equation}\label{eq:asymp-vol-case-r-equals-s}
    V_{\symbolnotsobw^{-1}}(\lambda)
    \sim \frac{d\,\omega_d^2}{s\,(2\pi\sqrt{c})^{d}}
    \lambda^{-\frac{d}{s}}\ln\left(\lambda^{-1}\right),
    \quad \text{as } \lambda\to0.
\end{equation}
\fi 
% \iffalse 
% In order to simplify the analysis of the asymptotic behavior of
% $V_{\symbolnotsobw^{-1}}(\lambda)$ as $\lambda \to 0$, it is convenient
% to introduce the comparison integral
% \begin{equation}\label{eq:integral-after-normalization-exponents}
% \bar V_{\symbolnotsobw^{-1}}(\lambda)
% =
% \frac{\omega_d^2}{(2\pi\sqrt{c})^{d}}
% \int_0^\infty\!\!\int_0^\infty
% \mathbbm{1}_{\{(1+u)^{-s/d}(1+v)^{-r/d}>\lambda\}}
% \,du\,dv.
% \end{equation}
% We next show that $V_{\symbolnotsobw^{-1}}(\lambda)$ and
% $\bar V_{\symbolnotsobw^{-1}}(\lambda)$ are equivalent up to rescaling of the argument.
% Consequently, the leading-order asymptotic
% behavior of $V_{\symbolnotsobw^{-1}}(\lambda)$ can be determined from the
% simpler integral $\bar V_{\symbolnotsobw^{-1}}(\lambda)$.
% To establish this equivalence, we compare the defining inequalities in the two
% volume functions. In particular, we relate the expressions
% $(1+u^{2/d})^{-s/2}$ and $(1+u)^{-s/d}$ (and similarly in $v$)
% appearing in their respective definitions.
% For $d \ge 2$, the elementary inequalities
% \[
% (1+y)^{2/d} \le 1+y^{2/d}
% \le 2^{1-2/d}(1+y)^{2/d},
% \qquad y>0,
% \]
% imply the inclusions
% \begin{align*}
% \{(1+u)^{-s/d}(1+v)^{-r/d}
%   > 2^{\frac{d-2}{2d}(s+r)}\lambda\}
% &\subseteq
% \{(1+u^{2/d})^{-s/2}(1+v^{2/d})^{-r/2}
%   > \lambda\}, \\[1mm]
% \{(1+u^{2/d})^{-s/2}(1+v^{2/d})^{-r/2}
%   > \lambda\}
% &\subseteq
% \{(1+u)^{-s/d}(1+v)^{-r/d}
%   > \lambda\}.
% \end{align*}
% These inclusions in turn yield the volume bounds. For $d\ge2$ and all $\lambda>0$,
% \begin{equation}\label{eq:volume-notions-comparisons1}
% \bar V_{\symbolnotsobw^{-1}}\!\left(2^{\frac{d-2}{2d}(s+r)}\lambda\right)
% \le V_{\symbolnotsobw^{-1}}(\lambda)
% \le \bar V_{\symbolnotsobw^{-1}}(\lambda).
% \end{equation}
% For $d=1$ and all $\lambda>0$, an analogous argument yields
% \begin{equation}\label{eq:volume-notions-comparisons2}
% \bar V_{\symbolnotsobw^{-1}}(\lambda)
% \le V_{\symbolnotsobw^{-1}}(\lambda)
% \le \bar V_{\symbolnotsobw^{-1}}\!\left(2^{-(s+r)/2}\lambda\right).
% \end{equation}
% We now analyze the double integral defining $\bar V_{\symbolnotsobw^{-1}}(\lambda)$ in
% \eqref{eq:integral-after-normalization-exponents} as $\lambda\to0$.
% Fix $\lambda>0$ and $v\ge0$. For $u\ge0$, the inequality
% \[
% (1+u)^{-s/d}(1+v)^{-r/d}>\lambda
% \]
% is equivalent to
% \[
% 1+u<\lambda^{-d/s}(1+v)^{-r/s}.
% \]
% Therefore, the admissible values of $u$ are given by
% \[
% 0\le u<\frac{\lambda^{-d/s}}{(1+v)^{r/s}}-1.
% \]
% Moreover, for admissible values of $u$ to exist, we must have
% \[
% 0 \le v < \lambda^{-d/r}-1.
% \]
% Integrating first in $u$ and then in $v$ therefore gives
% \[
% \int_0^\infty\!\!\int_0^\infty
% \mathbbm{1}_{\{(1+u)^{-s/d}(1+v)^{-r/d}>\lambda\}}
% \,du\,dv
% =
% \int_0^{\lambda^{-d/r}-1}
% \Bigl(\frac{\lambda^{-d/s}}{(1+v)^{r/s}}-1\Bigr)\,dv.
% \]
% If $s=r$, this yields 
% \begin{equation}\label{eq:vbar-asymp1}
% \bar V_{\symbolnotsobw^{-1}}(\lambda)
% \sim
% \frac{\omega_d^2}{(2\pi\sqrt{c})^{d}}
% \frac{d}{s}\lambda^{-d/s}\ln(\lambda^{-1}),
% \qquad \lambda\to0.
% \end{equation}
% For $r\neq s$, a direct computation delivers 
% \begin{equation}\label{eq:vbar-asymp2}
% \bar V_{\symbolnotsobw^{-1}}(\lambda)
% \sim
% \frac{\omega_d^2}{(2\pi\sqrt{c})^{d}}
% \begin{cases}
% \frac{s}{r-s}\lambda^{-d/s}, & r>s,\\[0.2cm]
% \frac{r}{s-r}\lambda^{-d/r}, & r<s,
% \end{cases}
% \qquad \lambda\to0.
% \end{equation}
% \fi
\fi 
\fi 
For $r=s$, we first observe that
\[
(1+u^{2/d})^{-s/2}\sim(1+u)^{-s/d},
\qquad u\to\infty .
\]
Hence, for every fixed $\eta\in(0,1)$ there exists $R_\eta>0$ such that
\begin{equation}\label{eq:r-eta-inequalities}
(1-\eta)(1+u^{2/d})^{-s/2}
\le
(1+u)^{-s/d}
\le
(1+\eta)(1+u^{2/d})^{-s/2},
\qquad u\ge R_\eta ,
\end{equation}
and the same estimates hold with $u$ replaced by $v$. Now, whenever
$u,v\ge R_\eta$, \eqref{eq:r-eta-inequalities} and its
version with $u$ replaced by $v$ imply that
\begin{align*}
    (1-\eta)^2 (1+u^{2/d})^{-s/2}(1+v^{2/d})^{-s/2}
&\le
(1+u)^{-s/d}(1+v)^{-s/d} \\
&\le
(1+\eta)^2 (1+u^{2/d})^{-s/2}(1+v^{2/d})^{-s/2}.
\end{align*}
Consequently, for $u,v\ge R_\eta$, we obtain the set inclusions
\begin{equation}\label{eq:set-inclusion-1}
\{(1+u^{2/d})^{-s/2}(1+v^{2/d})^{-s/2}>\lambda\}
\subseteq
\{(1+u)^{-s/d}(1+v)^{-s/d}>(1-\eta)^2\lambda\}
\end{equation}
and
\begin{equation}\label{eq:set-inclusion-2}
\{(1+u)^{-s/d}(1+v)^{-s/d}>(1+\eta)^2\lambda\}
\subseteq
\{(1+u^{2/d})^{-s/2}(1+v^{2/d})^{-s/2}>\lambda\}.
\end{equation}
We next decompose the double integral \eqref{eq:double-integral-sob-var-ch} defining
$V_{\symbolnotsobw^{-1}}(\lambda)$ according to
\[
V_{\symbolnotsobw^{-1}}(\lambda)
=
\frac{\omega_d^2}{(2\pi\sqrt{c})^{d}}
\left(
\int_{0}^{R_\eta}\!\!\int_{0}^{R_\eta}
+
2\int_{R_\eta}^{\infty}\!\!\int_{0}^{R_\eta}
+
\int_{R_\eta}^{\infty}\!\!\int_{R_\eta}^{\infty}
\right)
\mathbbm{1}_{\{(1+u^{2/d})^{-s/2}(1+v^{2/d})^{-s/2}>\lambda\}}
\,du\,dv .
\]
The first integral satisfies
\begin{equation}\label{eq:first-integral}
\int_{0}^{R_\eta}\!\!\int_{0}^{R_\eta}
\mathbbm{1}_{\{(1+u^{2/d})^{-s/2}(1+v^{2/d})^{-s/2}>\lambda\}}
\,du\,dv
\longrightarrow
R_\eta^2,
\qquad \lambda\to0 .
\end{equation}
Next, for the mixed region we have
\begin{align}\label{eq:second-integral}
\int_{R_\eta}^{\infty}\!\int_{0}^{R_\eta}
\mathbbm{1}_{\{(1+u^{2/d})^{-s/2}(1+v^{2/d})^{-s/2}>\lambda\}}
\,du\,dv \nonumber
&\le
R_\eta
\int_{R_\eta}^{\infty}
\mathbbm{1}_{\{(1+v^{2/d})^{-s/2}>\lambda\}}\,dv \\
&=
O_{\lambda\to0} \left(\lambda^{-d/s}\right).
\end{align}
It remains to analyze the tail region
\[
\int_{R_\eta}^{\infty}\!\!\int_{R_\eta}^{\infty}
\mathbbm{1}_{\{(1+u^{2/d})^{-s/2}(1+v^{2/d})^{-s/2}>\lambda\}}
\,du\,dv .
\]
To this end, we introduce the tail part of the corresponding 
comparison integral
\[
\bar I_\eta(\lambda)
:=
\int_{R_\eta}^{\infty}\!\!\int_{R_\eta}^{\infty}
\mathbbm{1}_{\{(1+u)^{-s/d}(1+v)^{-s/d}>\lambda\}}
\,du\,dv .
\]
For $\lambda$ sufficiently small, the admissible region in this tail integral is given by
\[
R_\eta \le v < \frac{\lambda^{-d/s}}{1+R_\eta}-1,
\qquad
R_\eta \le u < \frac{\lambda^{-d/s}}{1+v}-1,
\]
and therefore
\[
\bar I_\eta(\lambda)
=
\int_{R_\eta}^{\frac{\lambda^{-d/s}}{1+R_\eta}-1}
\left(\frac{\lambda^{-d/s}}{1+v}-1-R_\eta\right)\,dv
=
\frac{d}{s}\lambda^{-d/s}\ln(\lambda^{-1})
+ O_{\lambda\to0} \left(\lambda^{-d/s}\right),
\]
so that
\begin{equation}\label{eq:tail-comparison-integral}
\bar I_\eta(\lambda)
\sim
\frac{d}{s}\lambda^{-d/s}\ln(\lambda^{-1}),
\qquad \lambda\to0 .
\end{equation}
We now use the set inclusions \eqref{eq:set-inclusion-1} and
\eqref{eq:set-inclusion-2} to get
\[
\bar I_\eta \bigl((1+\eta)^2\lambda\bigr)
\;\le\;
\int_{R_\eta}^{\infty}\!\int_{R_\eta}^{\infty}
\mathbbm{1}_{\{(1+u^{2/d})^{-s/2}(1+v^{2/d})^{-s/2}>\lambda\}}
\,du\,dv
\]
and
\[
\int_{R_\eta}^{\infty}\!\int_{R_\eta}^{\infty}
\mathbbm{1}_{\{(1+u^{2/d})^{-s/2}(1+v^{2/d})^{-s/2}>\lambda\}}
\,du\,dv
\;\le\;
\bar I_\eta \bigl((1-\eta)^2\lambda\bigr).
\]
By \eqref{eq:tail-comparison-integral}, this implies
\[
(1+\eta)^{-2d/s}
\le
\liminf_{\lambda\to0}
\frac{
\displaystyle
\int_{R_\eta}^{\infty}\!\!\int_{R_\eta}^{\infty}
\mathbbm{1}_{\{(1+u^{2/d})^{-s/2}(1+v^{2/d})^{-s/2}>\lambda\}}
\,du\,dv
}{
\frac{d}{s}\lambda^{-d/s}\ln(\lambda^{-1})
}
\]
and
\[
\limsup_{\lambda\to0}
\frac{
\displaystyle
\int_{R_\eta}^{\infty}\!\!\int_{R_\eta}^{\infty}
\mathbbm{1}_{\{(1+u^{2/d})^{-s/2}(1+v^{2/d})^{-s/2}>\lambda\}}
\,du\,dv
}{
\frac{d}{s}\lambda^{-d/s}\ln(\lambda^{-1})
}
\le
(1-\eta)^{-2d/s}.
\]
Since $\eta\in(0,1)$ is arbitrary, letting $\eta\downarrow0$ yields
\begin{equation}\label{eq:limit-double-int-simple-form-2}
\int_{R_\eta}^{\infty}\!\!\int_{R_\eta}^{\infty}
\mathbbm{1}_{\{(1+u^{2/d})^{-s/2}(1+v^{2/d})^{-s/2}>\lambda\}}
\,du\,dv
\sim
\frac{d}{s}\lambda^{-d/s}\ln(\lambda^{-1}),
\qquad \lambda\to0.
\end{equation}
Finally, combining \eqref{eq:first-integral}, \eqref{eq:second-integral}, and \eqref{eq:limit-double-int-simple-form-2}, we can conclude that
\begin{equation}\label{eq:asymp-vol-case-r-equals-s}
V_{\symbolnotsobw^{-1}}(\lambda)
\sim
\frac{d\,\omega_d^2}{s(2\pi\sqrt{c})^{d}}
\lambda^{-d/s}\ln \left(\lambda^{-1}\right),
\qquad \lambda\to0.
\end{equation}
We now apply Theorem~\ref{thm:main-result} with
$\symbolnot=\symbolnotsobw^{-1}$ to obtain
\begin{equation}\label{eq:appli-thm1-sob2}
    H_{\symbolnotsobw^{-1}}(\varepsilon)
    \sim
    \int_\varepsilon^\infty
    \frac{V_{\symbolnotsobw^{-1}}(\lambda)}{\lambda}\,d\lambda,
    \qquad \varepsilon\to0.
\end{equation}
Combining \eqref{eq:final-asymp-vol-r-leq-s},
\eqref{eq:final-asymp-vol-r-geq-s},
\eqref{eq:asymp-vol-case-r-equals-s},
and \eqref{eq:appli-thm1-sob2}, we arrive at
\[
    H_{\symbolnotsobw^{-1}}(\varepsilon)
    \sim
    \frac{\omega_d^2}{(2\pi\sqrt{c})^{d}}
    \begin{dcases}
        \Xi_{r,s,d}\,\varepsilon^{-d/\min\{r,s\}},
        & \text{if } r\neq s,\\[.25cm]
        \varepsilon^{-d/s}\ln(\varepsilon^{-1}),
        & \text{if } r=s,
    \end{dcases}
    \qquad \varepsilon\to0,
\]
where
\[
    \Xi_{r,s,d}
    \coloneqq
    \frac{
        \min\{r,s\}\,
        \Gamma \left(\frac{d|s-r|}{2\min\{r,s\}}\right)
        \Gamma \left(\frac{d}{2}\right)
    }{
        2\,
        \Gamma \left(\frac{d\max\{r,s\}}{2\min\{r,s\}}\right)
    }.
\]
Finally,
\[
H \left(\varepsilon;\mathcal{B}_{\mathcal{W}}^{(s,r)},\|\cdot\|_{L^2(\mathbb R^d)}\right)
\sim
H_{\symbolnotsobw^{-1}}(\varepsilon),
\qquad \varepsilon \to 0,
\]
which completes the proof.
\iffalse
We now apply Theorem~\ref{thm:main-result} with
$\symbolnot=\symbolnotsobw^{-1}$ to obtain
% By
% \eqref{eq:volume-notions-comparisons1}--\eqref{eq:volume-notions-comparisons2},
% there exist constants $c_1,c_2>0$ such that, for all $\lambda>0$,
% \[
% \bar V_{\symbolnotsobw^{-1}}(c_1\lambda)
% \le
% V_{\symbolnotsobw^{-1}}(\lambda)
% \le
% \bar V_{\symbolnotsobw^{-1}}(c_2\lambda).
% \]
% Integrating these inequalities over $(\varepsilon,\infty)$ yields
% \[
% \int_\varepsilon^\infty
% \frac{\bar V_{\symbolnotsobw^{-1}}(c_1\lambda)}{\lambda}\,d\lambda
% \le
% \int_\varepsilon^\infty
% \frac{V_{\symbolnotsobw^{-1}}(\lambda)}{\lambda}\,d\lambda
% \le
% \int_\varepsilon^\infty
% \frac{\bar V_{\symbolnotsobw^{-1}}(c_2\lambda)}{\lambda}\,d\lambda.
% \]
% A change of variables gives
% \begin{equation}\label{eq:v-double-ineq}
% \int_{c_1\varepsilon}^\infty
% \frac{\bar V_{\symbolnotsobw^{-1}}(\lambda)}{\lambda}\,d\lambda
% \le
% \int_\varepsilon^\infty
% \frac{V_{\symbolnotsobw^{-1}}(\lambda)}{\lambda}\,d\lambda
% \le
% \int_{c_2\varepsilon}^\infty
% \frac{\bar V_{\symbolnotsobw^{-1}}(\lambda)}{\lambda}\,d\lambda.
% \end{equation}
% By Theorem~\ref{thm:main-result},
\begin{equation}\label{eq:appli-thm1-sob2}
    H_{\symbolnotsobw^{-1}}(\varepsilon)
    \sim
    \int_\varepsilon^\infty
    \frac{V_{\symbolnotsobw^{-1}}(\lambda)}{\lambda}\,d\lambda,
    \qquad \varepsilon\to0.
\end{equation}
Thus, combining the asymptotics \eqref{eq:final-asymp-vol-r-leq-s}, \eqref{eq:final-asymp-vol-r-geq-s}, and \eqref{eq:asymp-vol-case-r-equals-s} with \eqref{eq:appli-thm1-sob2} gives 
\begin{equation*}
    H_{\symbolnotsobw^{-1}}(\varepsilon)
    \sim \frac{\omega_d^2}{(2\pi\sqrt{c})^{d}}
    \begin{dcases}
        \Xi_{r,s,d}\, \varepsilon^{-\frac{d}{\min\{r,s\}}}, 
        \quad &\text{if } r\neq s,\\[.25cm]           
        \varepsilon^{-\frac{d}{s}}\ln\left(\varepsilon^{-1}\right), \quad &\text{if } r=s,
    \end{dcases}
    \quad \text{as } \varepsilon\to 0,
\end{equation*}
where 
\begin{equation*}
     \Xi_{r,s,d}
     \coloneqq \frac{\min\{r,s\} \, \Gamma\left(\frac{d \, |s-r| }{2 \min\{r,s\}}\right) \Gamma \left( \frac{d}{2} \right)}{2\Gamma \left(\frac{d\max\{r,s\} }{2\min\{r,s\} }\right)}.
\end{equation*}
% With
% \[
% f(\varepsilon)
% \coloneqq
% \int_\varepsilon^\infty
% \frac{\bar V_{\symbolnotsobw^{-1}}(\lambda)}{\lambda}\,d\lambda,
% \]
% the inequalities \eqref{eq:v-double-ineq} imply that
% \[
% f(c_1\varepsilon)
% \le
% H_{\symbolnotsobw^{-1}}(\varepsilon)
% \le
% f(c_2\varepsilon),
% \qquad \varepsilon\to0.
% \]
% It follows that $f(\varepsilon)$ is regularly varying at $0$ with index $-d/\min\{r,s\}$; in particular, 
% \[
% \frac{f(c\, \varepsilon)}{f(\varepsilon)}
% \longrightarrow
% c^{-d/\min\{r,s\}},
% \qquad \varepsilon\to0.
% \]
% Hence, as $\varepsilon \to 0$,
% \[
% f(\varepsilon)
% \;\lesssim\;
% H_{\symbolnotsobw^{-1}}(\varepsilon)
% \;\lesssim\;
% 2^{\frac{s+r}{2\min\{r,s\}}}\,
% f(\varepsilon),
% \quad \text{if } d=1,
% \]
% and
% \[
% 2^{-\frac{(d-2)(s+r)}{2\min\{r,s\}}}\,
% f(\varepsilon)
% \;\lesssim\;
% H_{\symbolnotsobw^{-1}}(\varepsilon)
% \;\lesssim\;
% f(\varepsilon),
% \quad \text{if } d \ge 2.
% \]
% Finally, using \eqref{eq:vbar-asymp1} and \eqref{eq:vbar-asymp2}, we obtain
% \[
% f(\varepsilon)
% \sim
% \frac{\omega_{d}^2}{d(2 \pi\sqrt{c})^{d}}
% \begin{cases}
% \dfrac{s^2}{r-s}\,\varepsilon^{-d/s}, & r>s, \\[.25cm]
% d\,\varepsilon^{-d/s}\ln(\varepsilon^{-1}), & r=s, \\[.25cm]
% \dfrac{r^2}{s-r}\,\varepsilon^{-d/r}, & r<s,
% \end{cases}
% \qquad \varepsilon\to0.
% \]
The proof is finalized by noting that
\[
H \left(\varepsilon;\mathcal{B}_{\mathcal{W}}^{(s,r)},\|\cdot\|_{L^2(\mathbb R^d)}\right)
\sim
H_{\symbolnotsobw^{-1}}(\varepsilon),
\qquad \varepsilon \to 0.
\]
\fi
\color{black}

%Combining the preceding relations gives the stated entropy asymptotics
%for $\mathcal{B}_{\mathcal{W}}^{(s,r)}$ and completes the proof of
%Theorem~\ref{thm:sob-weighted}.


% \iffalse
% A direct application of Theorem~\ref{thm:main-result} with
% $\symbolnot=\symbolnotsobw^{-1}$, combined with 
% \eqref{eq:volume-notions-comparisons1}--\eqref{eq:volume-notions-comparisons2}, yields
% two-sided bounds for the entropy $H_{T^{-1}_{\symbolnotsobw}}(\varepsilon)$ in terms of
% \[
% f(\varepsilon)\;\coloneqq\;\int_{\varepsilon}^{\infty}
% \frac{\bar V_{\symbolnotsobw^{-1}}(\lambda)}{\lambda}\,d\lambda.
% \]
% More precisely, for $d=1$ we have
% \[
% f(\varepsilon)
% \;\le\;
% H_{T^{-1}_{\symbolnotsobw}}(\varepsilon)
% \;\le\;
% 2^{\frac{s+r}{2\min\{r,s\}}}\,f(\varepsilon),
% \]
% whereas for $d\ge 2$ we obtain
% \[
% 2^{-\frac{(d-2)(s+r)}{2\min\{r,s\}}}\,f(\varepsilon)
% \;\le\;
% H_{T^{-1}_{\symbolnotsobw}}(\varepsilon)
% \;\le\;
% f(\varepsilon).
% \]

% Furthermore, the asymptotic evaluation of $\bar V_{\symbolnotsobw^{-1}}(\lambda)$
% as $\lambda\to0$ derived above yields
% \[
% f(\varepsilon)
% \sim
% \int_{\varepsilon}^{\infty}\frac{\bar V_{\symbolnotsobw^{-1}}(\lambda)}{\lambda}\,d\lambda
% \sim
% \frac{\omega_{d}^2}{d(2 \pi\sqrt{c})^{d}}
% \begin{cases}
% \frac{s^2}{r-s}\,\varepsilon^{-d/s}, & r>s, \\[.25cm]
% d\,\varepsilon^{-d/s}\ln(\varepsilon^{-1}), & r=s, \\[.25cm]
% \frac{r^2}{s-r}\,\varepsilon^{-d/r}, & r<s,
% \end{cases}
% \qquad \varepsilon\to0.
% \]
% Combining the preceding bounds with this asymptotic expression for $f(\varepsilon)$
% yields the desired entropy asymptotics and completes the proof of
% Theorem~\ref{thm:sob-weighted}.
% \fi 

% \iffalse 
% By Theorem~\ref{thm:main-result}, applied to
% $\symbolnotsobw^{-1}$ together with
% \eqref{eq:volume-notions-comparisons1}--\eqref{eq:volume-notions-comparisons2},
% we obtain entropy asymptotics of the form
% \[
% H_{T^{-1}_{\symbolnotsobw}}(\varepsilon)
% \sim
% \int_\varepsilon^\infty
% \frac{\bar V_{\symbolnotsobw^{-1}}(\lambda)}{\lambda}\,d\lambda,
% \qquad \varepsilon\to0.
% \]
% Evaluating this integral yields
% \[
% H(\varepsilon)
% \sim
% \frac{\omega_d^2}{d(2\pi\sqrt{c})^{d}}
% \begin{cases}
% \frac{s^2}{r-s}\varepsilon^{-d/s}, & r>s,\\[0.2cm]
% d\,\varepsilon^{-d/s}\ln(\varepsilon^{-1}), & r=s,\\[0.2cm]
% \frac{r^2}{s-r}\varepsilon^{-d/r}, & r<s,
% \end{cases}
% \qquad \varepsilon\to0,
% \]
% which is the desired result and completes the proof of
% Theorem~\ref{thm:sob-weighted}.
% \fi 


% \iffalse 
% Therefore, following the beginning of the proof of Theorem~\ref{thm:sob-sch} until \eqref{eq:volume-integral-sch-formula1}--\eqref{eq:volume-integral-sch-formula2} with $\symbolnotsobw$ in place of $\symbolnotsch$, we reduce the problem to characterizing the asymptotic behavior of
% \begin{equation}\label{eq:integral-volume-weighted-sob1}
%     V_{\symbolnotsobw^{-1}}(\lambda)
%     =  \frac{1}{(2 \pi\sqrt{c})^{d}} \int_{\R^{2d}} \mathbbm{1}_{\left\{{\left(1+\|\omega'\|_2^2\right)^{-s/2}}{\left(1+\|x'\|_2^2\right)^{-r/2}} \, > \, \lambda\right\}} \diff x' \diff \omega'.
% \end{equation}
% To this end,
% we first rewrite the integral in \eqref{eq:integral-volume-weighted-sob1} in polar coordinates according to
% \begin{equation*}
%     \int_{0}^\infty \int_{0}^\infty  \mathbbm{1}_{\left\{{\left(1+u_0^{2}\right)^{-s/2}}{\left(1+v_0^{2}\right)^{-r/2}}\, > \, \lambda\right\}} s_{d-1}^2 \left(u_0v_0\right)^{d-1} \diff u_0 \diff v_0,
% \end{equation*}
% where $s_{d-1}^2$ denotes the $(d-1)$-dimensional Hausdorff measure of the surface of the unit ball in $\R^d$.
% Further, employing the change of variables $u = u_0^{d}$ and $v= v_0^{d}$,
% we obtain
% \begin{equation*}
%     V_{\symbolnotsobw^{-1}}(\lambda)
%     = \frac{s_{d-1}^2}{(2 \pi\sqrt{c})^{d} d^2}
%     \int_{0}^\infty \int_{0}^\infty  \mathbbm{1}_{\left\{{\left(1+u^{2/d}\right)^{-s/2}}{\left(1+v^{2/d}\right)^{-r/2}}\, > \, \lambda\right\}} \diff u \diff v,
% \end{equation*}
% which, upon recalling the classical formula $ \omega_d = s_{d-1}/d$, can be rewritten as
% \begin{equation*}
%     V_{\symbolnotsobw^{-1}}(\lambda)
%     = \frac{\omega_{d}^2}{(2 \pi\sqrt{c})^{d}}
%     \int_{0}^\infty \int_{0}^\infty  \mathbbm{1}_{\left\{{\left(1+u^{2/d}\right)^{-s/2}}{\left(1+v^{2/d}\right)^{-r/2}}\, >\, \lambda\right\}} \diff u \diff v.
% \end{equation*}


% It turns out that it will be more convenient to study the slightly different integral
% \begin{equation}\label{eq:integral-after-normalization-exponents}
%   \bar  V_{\symbolnotsobw^{-1}}(\lambda)
%    \coloneqq  \frac{\omega_{d}^2}{(2 \pi\sqrt{c})^{d}} \int_{0}^\infty \int_{0}^\infty  \mathbbm{1}_{\left\{{\left(1+u\right)^{-s/d}}{\left(1+v\right)^{-r/d}}\, > \, \lambda\right\}} \diff u \diff v,
% \end{equation}
% defined for all $\lambda >0$. 
% We thus proceed to relate $ V_{\symbolnotsobw^{-1}}$ and $\bar  V_{\symbolnotsobw^{-1}}$.
% To this end, we first observe, in the case $d\geq 2$, that
% \begin{equation}\label{eq:useful-inequalities}
%     (1+y)^{2/d}
%     \leq 1+y^{2/d}
%     \leq 2^{1-2/d} (1+y)^{2/d},
%    \quad \text{for all } y>0, 
% \end{equation}
% where the lower bound is standard
% and the upper bound is deduced from the concavity of the function $t\mapsto t^{2/d}$, for $t>0$.
% Applying the bounds in \eqref{eq:useful-inequalities} with $y=u$ and $y=v$,
% we get--still in the case $d\geq 2$--the inclusions
% \begin{align*}
%     &\left\{{\left(1+u\right)^{-s/d}}{\left(1+v\right)^{-r/d}}>2^{\frac{d-2}{d}(s+r)} \lambda\right\}\\
%     &\subseteq \left\{{\left(1+u^{2/d}\right)^{-s/2}}{\left(1+v^{2/d}\right)^{-r/2}}>\lambda\right\} 
%     \subseteq \left\{{\left(1+u\right)^{-s/d}}{\left(1+v\right)^{-r/d}}>\lambda\right\},
% \end{align*}
% which readily translate into the inequalities
% \begin{equation}\label{eq:volume-notions-comparisons1}
%     \bar V_{\symbolnotsobw^{-1}} \left(2^{\frac{d-2}{d}(s+r)}\lambda\right)
%     \leq V_{\symbolnotsobw^{-1}} \left(\lambda\right)
%     \leq \bar V_{\symbolnotsobw^{-1}} \left(\lambda\right),
%     \quad \text{for all } \lambda >0.
% \end{equation}
% In the case $d=1$, a similar analysis yields the inequalities
% \begin{equation}\label{eq:volume-notions-comparisons2}
%     \bar V_{\symbolnotsobw^{-1}} \left(\lambda\right)
%     \leq V_{\symbolnotsobw^{-1}} \left(\lambda\right)
%     \leq \bar V_{\symbolnotsobw^{-1}} \left(2^{-({s+r})}\lambda\right),
%     \quad \text{for all } \lambda >0.
% \end{equation}



% We now fix $\lambda >0$ and  study the integral appearing in the definition \eqref{eq:integral-after-normalization-exponents} of $\bar V_{\symbolnotsobw^{-1}}(\lambda) $.
% Specifically, we first perform successive integrations in the domain of integration according to 
% \begin{equation*}
%     \int_{0}^\infty \int_{0}^\infty  \mathbbm{1}_{\left\{{\left(1+u\right)^{-s/d}}{\left(1+v\right)^{-r/d}}>\lambda\right\}} \diff u \diff v
%     =\int_{0}^{\lambda^{-d/r}-1} \left( \int_{0}^{\frac{\lambda^{-d/s}}{(1+v)^{r/s}}-1}  \diff u \right) \diff v. 
% \end{equation*}
% The inner integral can be calculated explicitly, yielding
% \begin{equation*}
%     \int_{0}^{\lambda^{-d/r}-1} \left( \int_{0}^{\frac{\lambda^{-d/s}}{(1+v)^{r/s}}-1}  \diff u \right) \diff v = \int_{0}^{\lambda^{-d/r}-1} \frac{\lambda^{-d/s}}{(1+v)^{r/s}} \diff v - \left(\lambda^{-d/r}-1\right).
% \end{equation*}
% We now distinguish between the cases $s=r$ and $s\neq r$.
% First consider the case $s=r$, for which the integral
% \begin{equation*}
%      \int_{0}^{\lambda^{-d/r}-1} \frac{\lambda^{-d/s}}{1+v} \diff v
%      = \frac{d}{r} \lambda^{-d/s} \ln\left(\lambda^{-1}\right)
% \end{equation*}
% asymptotically dominates the term $\lambda^{-d/r}-1$.
% Next assume that $r\neq s$ and observe that
% \begin{equation}\label{eq:calculation-integral-asymp-lambda}
%     \int_{0}^{\lambda^{-d/r}-1} \frac{\lambda^{-d/s}}{(1+v)^{r/s}} \diff v - \lambda^{-d/r}+1
%     =  \frac{\lambda^{-d/r}-\lambda^{-d/s}}{1-r/s}- \lambda^{-d/r}+1.
% \end{equation}
% The right-hand side in \eqref{eq:calculation-integral-asymp-lambda} satisfies the following asymptotic behavior
% \begin{equation*}
%      \frac{\lambda^{-d/r}-\lambda^{-d/s}}{1-r/s}- \lambda^{-d/r} 
%      \sim
%     \begin{cases}
%         \frac{s}{r-s}\lambda^{-d/s}, \quad &\text{if } r>s, \\[.25cm]
%         \frac{r}{s-r}\lambda^{-d/r}, \quad &\text{if } r<s,
%     \end{cases}
%     \quad \quad  \text{as } \lambda \to 0.
% \end{equation*}
% Putting everything together, we have proven that 
% \begin{equation*}
%     \bar V_{\symbolnotsobw^{-1}} \left(\lambda\right) 
%     \sim \frac{\omega_{d}^2}{(2 \pi\sqrt{c})^{d}}
%     \begin{cases}
%         \frac{s}{r-s}\lambda^{-d/s}, \quad &\text{if } r>s, \\[.25cm]
%         \frac{d}{s}\lambda^{-d/s} \ln\left(\lambda^{-1}\right), \quad &\text{if } r=s, \\[.25cm]
%         \frac{r}{s-r}\lambda^{-d/r}, \quad &\text{if } r<s,
%     \end{cases}
%     \quad \quad  \text{as } \lambda \to 0.
% \end{equation*}


% \textcolor{blue}{A direct application of Theorem~\ref{thm:main-result}, with $\symbolnot=\symbolnotsobw^{-1}$,
% combined with the relationship \eqref{eq:volume-notions-comparisons1}--\eqref{eq:volume-notions-comparisons2}
% between $ V_{\symbolnotsobw^{-1}}$ and $\bar V_{\symbolnotsobw^{-1}}$,  
% yields}
% \begin{equation*}
%     f(\varepsilon)
%     \leq H\left(\varepsilon;T^{-1}_{\symbolnotsobw}\right)
%     \leq 2^{\frac{s+r}{\min\{r,s\}}} f(\varepsilon)
% \end{equation*}
% and 
% \begin{equation*}
%     2^{-\frac{(d-2)(s+r)}{\min\{r,s\}}} f(\varepsilon)
%     \leq H\left(\varepsilon;T^{-1}_{\symbolnotsobw}\right)
%     \leq f(\varepsilon),
% \end{equation*}
% if $d=1$ and $d\geq 2$ respectively, 
% where 
% \color{blue}
% \begin{equation*}
%     f(\varepsilon) 
%     \sim  \int_{\varepsilon}^\infty \frac{\bar V_{\symbolnotsobw^{-1}} \left(\lambda\right) }{\lambda} \diff \lambda
%     \sim \frac{\omega_{d}^2}{d  (2 \pi\sqrt{c})^{d}}
%     \begin{cases}
%         \frac{s^2}{r-s}\varepsilon^{-d/s}, \quad &\text{if } r>s, \\[.25cm]
%         d \,  \varepsilon^{-d/s} \ln\left(\varepsilon^{-1}\right), \quad &\text{if } r=s, \\[.25cm]
%         \frac{r^2}{s-r}\varepsilon^{-d/r}, \quad &\text{if } r<s,
%     \end{cases}
%     \quad\quad  \text{as } \varepsilon \to 0.
% \end{equation*}
% \color{black}
% This is the desired result, which concludes the proof of Theorem~\ref{thm:sob-weighted}.
% %%%T5[OK]: I removed ln(2), please double-check.
% \fi 




%%%T10[OK]: please update ref. 6, accepted, 

\section*{Acknowledgments}

The authors thank A.~Künzi for useful suggestions.


\bibliography{main}

\appendix


\section{Review of Spectral Theory for Pseudodifferential Operators}
\label{sec:review-lit}

This appendix reviews the spectral-theoretic material required for the
development of our results.
We begin by recalling the Kohn--Nirenberg representation of a
pseudodifferential operator.
With the Fourier transform defined by
\begin{equation}\label{eq:definition-fourier-transform}
    \hat f(\omega)
    =
    \int_{\R^d} f(y)\, e^{-2\pi i\, y\cdot \omega}\,\diff y,
    \qquad \omega\in\R^d,
\end{equation}
the operator $T_\sigma$ associated with a symbol $\sigma$ admits the integral
representation
\begin{equation}\label{eq:kohn-nirenberg-def2}
    (T_\sigma f)(x)
    =
    \int_{\R^{2d}} \sigma(x,\omega)\,
    e^{2\pi i (x-y)\cdot\omega}\,
    f(y)\,\diff y\,\diff\omega,
    \qquad x\in\R^d.
\end{equation}

Following standard practice (see, e.g.,
\cite[Definition~23.1]{shubin_pseudodifferential_2001}),
we assume throughout that the symbols under consideration satisfy appropriate
regularity and growth conditions.
Specifically, for $N\in\N^*$ (typically $N=2d$ or $N=3d$ below), $m\in\R$, and
$\rho\in(0,1]$, we define the symbol class
$\psidosymbclass(\R^N)$ as the set of real-valued functions
$\sigma\in C^\infty(\R^N)$ such that, for every multi-index
$\alpha\in\N^N$, there exists a constant $C_\alpha>0$ with
\begin{equation}\label{eq:decay-cond-symbol-classs}
    \bigl|\partial^\alpha \sigma(z)\bigr|
    \le
    C_\alpha\,\bigl(1+\|z\|_2^2\bigr)^{\frac{m-\rho|\alpha|}{2}},
    \qquad z\in\R^N.
\end{equation}
%We define the order of a symbol $\sigma$ as the infimum over all $m\in\R$ such that $\sigma\in\psidosymbclass(\R^N)$.
%%%T8: I removed the definition of order as we only use "positive order" and "negative order" and those two quantities are
%defined below. If we keep the general definition of order here, the reader may wonder why below we do not just say that symbols of positive order have %order that is positive with order as was defined here.

%\section{Review of Spectral Theory for Pseudodifferential Operators}\label{sec:review-lit}

% \iffalse 
% This section contains a review of the material needed in the development of our results.
% We first note that the 
% Kohn-Nirenberg representation \eqref{eq:kohn-nirenberg-def} of a pseudodifferential operator can be reformulated, by developing the Fourier transform 
% \begin{equation}\label{eq:definition-fourier-transform}
%     \hat f (\omega) 
%     = \int_{\R^d} f(y) e^{-2i\pi y \cdot \omega } \diff y,
%     \quad \text{for all } \omega \in \R^d,
% \end{equation}
% according to
% \begin{equation}\label{eq:kohn-nirenberg-def2}
%     \left(T_{\symbolnot} f\right) (x) 
%     = \int_{\R^{2d}} \symbolnot(x, \omega)   \, e^{2i\pi (x-y) \cdot \omega} f (y) \diff y \diff \omega,
%     \quad \text{for all } x \in \R^d.
% \end{equation}
% Following standard practice (see, e.g., \cite[Definition~23.1]{shubin_pseudodifferential_2001}), we assume in what follows
% that the symbols $\symbolnot$  considered satisfy regularity conditions.
% Specifically, given $N\in\N^*$ (typically taken to be $2d$ or $3d$ below), $m\in\R$, and $\rho\in(0,1]$, we define the {symbol class} $\psidosymbclass(\R^N)$ 
% as the class of real-valued functions $\symbolnot \in C^\infty(\R^N)$ such that, for every multi-index $\alpha\in\N^N$, there exists a $C_\alpha>0$ so that
% \begin{equation}\label{eq:decay-cond-symbol-classs}
%     \left\lvert \partial^\alpha \symbolnot(z) \right\rvert
%     \leq C_{\alpha} \left(1+\|z\|^2_2\right)^{\frac{m-\rho|\alpha|}{2}},
%     \quad \text{for all } z\in \R^N.
% \end{equation}
% \fi 


For later use, it is convenient to work with a slightly more general integral
representation than \eqref{eq:kohn-nirenberg-def2}, allowing the symbol to depend
on the integration variable $y$ as well.
Specifically, we consider operators of the form
\begin{equation}\label{eq:amplitude-integral-representation}
    (T_{\bar\sigma} f)(x)
    =
    \int_{\R^{2d}} \bar\sigma(x,y,\omega)\,
    e^{2\pi i (x-y)\cdot\omega}\,
    f(y)\,\diff y\,\diff\omega,
\end{equation}
where $\bar\sigma\in\psidosymbclass(\R^{3d})$ is referred to as the
amplitude of the operator. The precise definition of an amplitude varies somewhat across the literature;
see, for example, \cite[Definition~23.3]{shubin_pseudodifferential_2001} and
\cite[Definition~2.1]{treves2013psido}.
The choice adopted here—namely, requiring $\bar\sigma$ to belong to
$\psidosymbclass(\R^{3d})$—is slightly more restrictive than
\cite[Definition~23.3]{shubin_pseudodifferential_2001}, but this distinction plays
no role in the arguments below. Since $\bar\sigma\in\psidosymbclass(\R^{3d})$, it satisfies 
\eqref{eq:decay-cond-symbol-classs} with $N=3d$.
Operators of the form \eqref{eq:amplitude-integral-representation} with amplitudes $\bar\sigma\in\psidosymbclass(\R^{3d})$ define the
class $\psidoopclass(\R^d)$.
Finally, we set
\[
\Psi^{-\infty}
\;\coloneqq\;
\bigcap_{m\in\R}\psidoopclass(\R^d),
\]
which can be shown to be independent of the parameter $\rho$; see the discussion following
\cite[Definition~23.4]{shubin_pseudodifferential_2001}.

The amplitude formulation \eqref{eq:amplitude-integral-representation} provides a
unified framework encompassing several standard quantization schemes.
The Kohn--Nirenberg representation \eqref{eq:kohn-nirenberg-def2} is recovered when
the amplitude is independent of the variable \(y\), that is, when there exists a
symbol \(\sigma_L\in\Gamma_\rho^m(\R^{2d})\) such that
\[
\bar\sigma(x,y,\omega)=\sigma_L(x,\omega).
\]
In this case, \(\sigma_L\) is referred to as the Kohn--Nirenberg symbol (or
left symbol) of the operator \(T_{\bar\sigma}\).
If \(\sigma_W\in\Gamma_\rho^m(\R^{2d})\) and
\[
\bar\sigma(x,y,\omega)=\sigma_W \left(\tfrac{x+y}{2},\omega\right),
\]
then \(T_{\bar\sigma}\) coincides with the Weyl quantization of the symbol
\(\sigma_W\), which is therefore called the Weyl symbol of the operator.
Similarly, when the amplitude depends only on the second spatial variable,
\[
\bar\sigma(x,y,\omega)=\sigma_R(y,\omega),
\qquad \sigma_R\in\Gamma_\rho^m(\R^{2d}),
\]
the function \(\sigma_R\) is referred to as the right symbol of
\(T_{\bar\sigma}\).
Conversely, every operator \(T_{\bar\sigma}\in\Psi_\rho^m(\R^d)\) admits equivalent
representations in terms of a left symbol, a Weyl symbol, and a right symbol; see
\cite[Theorem~23.1]{shubin_pseudodifferential_2001}.


A central theme in the spectral theory of pseudodifferential operators is that the
asymptotic behavior of the eigenvalue-counting function can be characterized in
terms of the associated symbol, independently of the chosen quantization
(Kohn--Nirenberg, Weyl, or right).
We highlight two results of this type.
The first, stated as Theorem~\ref{thm-shubin}, is a classical Weyl-type result for
hypoelliptic pseudodifferential operators and is included here for context.
The second, Theorem~\ref{thm:Dauge-Robert}, concerns compact pseudodifferential
operators and provides the spectral characterization used in the proof of our
main result.

To state Theorem~\ref{thm-shubin}, we first introduce the notion of a
hypoelliptic symbol; see, for example,
\cite[Definition~25.1]{shubin_pseudodifferential_2001}.
A real-valued function $\sigma\in C^\infty(\R^{2d})$ is said to be hypoelliptic if
there are constants $C_1,C_2,\rho,R>0$ and real numbers $m_+\ge m_-$ such that,
for every multi-index $\alpha\in\N^{2d}$, one can find a constant $C_\alpha>0$
for which
\begin{equation}\label{eq:bound-def-hypoellip}
    C_1\bigl(1+\|z\|_2^2\bigr)^{\frac{m_-}{2}}
    \le
    |\sigma(z)|
    \le
    C_2\bigl(1+\|z\|_2^2\bigr)^{\frac{m_+}{2}},
    \quad\text{and}\quad
    \frac{|\partial^\alpha\sigma(z)|}{|\sigma(z)|}
    \le
    C_\alpha\bigl(1+\|z\|_2^2\bigr)^{-\frac{\rho|\alpha|}{2}},
\end{equation}
for all $z\in\R^{2d}$ with $\|z\|_2\ge R$.
We denote this class of symbols by $\hyposymbclass(\R^{2d})$.
By direct comparison of \eqref{eq:bound-def-hypoellip} with the estimate
\eqref{eq:decay-cond-symbol-classs}, every hypoelliptic symbol is a symbol in the
sense of $\Gamma_\rho^{m_+}(\R^{2d})$, that is,
$\hyposymbclass(\R^{2d})\subset\Gamma_\rho^{m_+}(\R^{2d})$.
A hypoelliptic operator is a pseudodifferential operator whose symbol is hypoelliptic.


Depending on the asymptotic behavior of the symbol at infinity,
hypoelliptic symbols exhibit qualitatively different spectral properties.
We therefore distinguish between symbols that grow at infinity and symbols
that decay at infinity.


\medskip

\noindent
\textbf{Positive-order hypoelliptic symbols.}
A hypoelliptic symbol $\sigma$ is said to be of positive order if it belongs
to the class
\[
\hypoposclass(\R^{2d})
\;\coloneqq\;
\bigcup_{\rho\in(0,1]}
\ \bigcup_{\substack{m_+,m_->0\\ m_+\,\ge\, m_-}}
\mathrm{H}\Gamma^{m_-,m_+}_\rho(\R^{2d}).
\]
That is, $\sigma$ is of positive order if there exist parameters $\rho\in(0,1]$
and real numbers $m_+\ge m_->0$ such that
\[
\sigma \in \mathrm{H}\Gamma^{m_-,m_+}_\rho(\R^{2d}).
\]
For symbols in this class, we have
\[
|\sigma(z)| \to \infty,
\qquad \text{as } \|z\|_2 \to \infty.
\]

% \iffalse 
% \textbf{Positive-order hypoelliptic symbols.}
% \textcolor{blue}{A hypoelliptic symbol $\sigma$ is said to be of positive order whenever there exist parameters
% $\rho\in(0,1]$ and real numbers $m_+\ge m_->0$ such that $\sigma \in \mathrm{H}\Gamma^{m_-,m_+}_\rho(\R^{2d})$.}
% We denote the class of positive-order hypoelliptic symbols by
% \[
% \hypoposclass(\R^{2d})
% \;\coloneqq\;
% \bigcup_{\rho\in(0,1]}\ \bigcup_{\substack{m_+,m_->0\\ m_+\ge m_-}}
% \mathrm{H}\Gamma^{m_-,m_+}_\rho(\R^{2d}).
% \]
% \textcolor{blue}{If $\sigma$ is of positive order, then $|\sigma(z)| \to \infty$, as $\|z\|_2 \to \infty$.}
% \fi 

\noindent 
\textbf{Negative-order hypoelliptic symbols.}
A hypoelliptic symbol $\sigma$ is said to be of negative order if it belongs
to the class
\begin{equation}\label{eq:negative-order-def-class}
\hyponegclass(\R^{2d})
\;\coloneqq\;
\bigcup_{\rho\in(0,1]}
\ \bigcup_{\substack{m_+,m_->0\\ m_+\,\ge\, m_-}}
\mathrm{H}\Gamma^{-m_+,-m_-}_\rho(\R^{2d}).
\end{equation}
That is, $\sigma$ is of negative order if there exist parameters $\rho\in(0,1]$
and real numbers $m_+\ge m_->0$ such that
\[
\sigma \in \mathrm{H}\Gamma^{-m_+,-m_-}_\rho(\R^{2d}).
\]
For symbols in this class, we have
\[
|\sigma(z)| \to 0,
\qquad \text{as } \|z\|_2 \to \infty.
\]

% \iffalse 
% \noindent
% \textbf{Negative-order hypoelliptic symbols.}
% \textcolor{blue}{A hypoelliptic symbol $\sigma$ is said to be of negative order whenever there exist parameters
% $\rho\in(0,1]$ and real numbers $m_+\ge m_->0$ such that $\sigma \in \mathrm{H}\Gamma^{-m_+,-m_-}_\rho(\R^{2d})$.}
% We denote the class of negative-order hypoelliptic symbols by
% \begin{equation}\label{eq:negative-order-def-class}
%     \hyponegclass(\R^{2d})
% \;\coloneqq\;
% \bigcup_{\rho\in(0,1]}\ \bigcup_{\substack{m_+,m_->0\\ m_+\ge m_-}}
% \mathrm{H}\Gamma^{-m_+,-m_-}_\rho(\R^{2d}).
% \end{equation}
% \textcolor{blue}{If $\sigma$ is of negative order, then $|\sigma(z)| \to 0$, as $\|z\|_2 \to \infty$.}
% \fi 


% \iffalse
% Depending on the sign of the exponents $m_-$ and $m_+$, 
% a hypoelliptic symbol can have qualitatively different behaviors.
% We introduce the classes 
% of positive- and negative-order hypoelliptic symbols 
% according to
% \begin{align}
%     \hypoposclass (\R^{2d}) &\coloneqq \left\{\sigma \mid \sigma \in \hyposymbclass(\R^{2d}) \text{ for some } \rho\in(0,1], \text{ and } m_+\ge m_->0  \right\} \quad \text{and} \nonumber \\
%     \hyponegclass (\R^{2d}) &\coloneqq  \left\{\sigma \mid \sigma \in \hyposymbclassinv(\R^{2d}) \text{ for some } \rho\in(0,1], \text{ and }  m_+\ge m_->0  \right\}, \label{eq:negative-order-def-class}
% \end{align}
% respectively.
% \fi 


The classical spectral theory of hypoelliptic pseudodifferential operators
predominantly addresses the 
% case $m_->0$
positive-order case. %, corresponding to symbols whose
%magnitude $|\sigma(z)|$ tends to infinity as $\|z\|_2\to\infty$.
This stands in contrast to the negative-order setting of the present paper, where the symbols of
interest decay at infinity, so that the volume integral in
\eqref{eq:definition-volume-function} is finite.
Nevertheless, in order to place our results in context and to recall the
classical Weyl-type asymptotics available in the 
positive-order case, it is
convenient to introduce counterparts of the eigenvalue-counting function
\eqref{eq:definition-ecf} and of the volume function
\eqref{eq:definition-volume-function}, which we denote by a superscript
$\uparrow$. 
Specifically, for a Weyl symbol $\sigma_W\in\hypoposclass(\R^{2d})$, we define the
upward eigenvalue-counting function of the pseudodifferential operator
$T_{\sigma_W}$ by
\[
    \ecfup_{\sigma_W}(\lambda)
    \;\coloneqq\;
    \#\bigl\{n\in\N^* : \lambda_n(T_{\sigma_W}) < \lambda\bigr\},
    \qquad \lambda>0,
\]
where $\{\lambda_n(T_{\sigma_W})\}_{n\in\N^*}$
% \textcolor{red}{
denotes the eigenvalues
% } % denotes the spectrum 
of $T_{\sigma_W}$.
It follows from \cite[Corollary~23.3 and Theorem~26.3]{shubin_pseudodifferential_2001}
that $T_{\sigma_W}$ has purely discrete spectrum.
The upward volume associated with the same symbol $\sigma_W$ is defined by
\begin{equation}\label{eq:definition-upward-volume}
    \volup_{\sigma_W}(\lambda)
    \;\coloneqq\;
    \int_{\R^{2d}} \mathbbm{1}_{\{\sigma_W(x,\omega) \, < \, \lambda\}}
    \,\diff x\,\diff\omega,
    \qquad \lambda>0.
\end{equation}

With these notions in place, a classical Weyl-type result in the spectral theory
of hypoelliptic pseudodifferential operators can be stated as follows.

% \begin{theorem}[{\cite[Theorem~30.1]{shubin_pseudodifferential_2001}}]
% \label{thm-shubin}
% Let $\sigma$ be a Weyl symbol in $\hypoposclass(\R^{2d})$.
% That is, assume there exist parameters $\rho\in(0,1]$ and real numbers
% $m_+\ge m_->0$ for which $\sigma\in \mathrm{H}\Gamma^{m_-,m_+}_\rho(\R^{2d})$.
% Assume moreover that there exist constants $R>0$ and $C>0$ so that 
% \begin{equation}\label{eq:assumption-shubin}
%     |z\cdot\nabla\sigma(z)| \ge C\,|\sigma(z)|,
%     \qquad \text{for all } z\in\R^{2d}\text{ with }\|z\|_2\ge R.
% \end{equation}
% Then the upward eigenvalue-counting function $\ecfup_\sigma$ of the corresponding
% Weyl-quantized operator satisfies
% \[
%     \ecfup_\sigma(\lambda)\sim \volup_\sigma(\lambda),
%     \qquad \text{as } \lambda\to\infty.
% \]
% \end{theorem}

\begin{theorem}[{\cite[Theorem~30.1]{shubin_pseudodifferential_2001}}]
\label{thm-shubin}
Let $\sigma\in\hypoposclass(\R^{2d})$.
Assume that there exist constants $R>0$ and $C>0$ such that
\begin{equation}\label{eq:assumption-shubin}
    |z\cdot\nabla\sigma(z)| \ge C\,|\sigma(z)|,
    \qquad \text{for all } z\in\R^{2d}\text{ with }\|z\|_2\ge R.
\end{equation}
Then the upward eigenvalue-counting function $\ecfup_\sigma$ of the
corresponding Weyl-quantized operator satisfies
\[
    \ecfup_\sigma(\lambda)\sim \volup_\sigma(\lambda),
    \qquad \text{as } \lambda\to\infty.
\]
\end{theorem}

% \iffalse
% \begin{theorem}[{\cite[Theorem~30.1]{shubin_pseudodifferential_2001}}]
% \label{thm-shubin}
% \textcolor{blue}{Let $\sigma$ be a 
% % Weyl 
% symbol in $\hypoposclass(\R^{2d})$.
% % That is, assume there exist parameters $\rho\in(0,1]$ and real numbers
% % $m_+\ge m_->0$ for which $\sigma\in \mathrm{H}\Gamma^{m_-,m_+}_\rho(\R^{2d})$.
% Assume % moreover 
% that there exist constants $R>0$ and $C>0$ so that }
% \begin{equation}\label{eq:assumption-shubin}
%     |z\cdot\nabla\sigma(z)| \ge C\,|\sigma(z)|,
%     \qquad \text{for all } z\in\R^{2d}\text{ with }\|z\|_2\ge R.
% \end{equation}
% Then the upward eigenvalue-counting function $\ecfup_\sigma$ of the corresponding
% Weyl-quantized operator satisfies
% \[
%     \ecfup_\sigma(\lambda)\sim \volup_\sigma(\lambda),
%     \qquad \text{as } \lambda\to\infty.
% \]
% \end{theorem}
% \fi 






% \iffalse 
% \begin{theorem}[{\cite[Theorem~30.1]{shubin_pseudodifferential_2001}}]
% \label{thm-shubin}
% Let $\rho\in(0,1]$ and let $m_+,m_->0$ satisfy $m_+\ge m_-$.
% Let $\sigma\in\mathrm{H}\Gamma^{m_-,m_+}_\rho(\R^{2d})$ be a Weyl symbol.
% Assume that there exist constants $R>0$ and $C>0$ such that
% \begin{equation}\label{eq:assumption-shubin}
%     |z\cdot\nabla\sigma(z)| \ge C\,|\sigma(z)|,
%     \qquad \text{for all } z\in\R^{2d}\text{ with }\|z\|_2\ge R.
% \end{equation}
% Then the upward eigenvalue-counting function $\ecfup_\sigma$ of the
% corresponding Weyl-quantized operator satisfies
% \[
% \ecfup_\sigma(\lambda)\sim \volup_\sigma(\lambda),
% \qquad \text{as } \lambda\to\infty.
% \]
% \end{theorem}
% \fi 
% %%%T5: I added here that sigma has to be a Weyl symbol as per Shubin.
% \iffalse 
% \begin{theorem}[{\cite[Theorem~30.1]{shubin_pseudodifferential_2001}}] \label{thm-shubin}
%     Let $\rho\in (0,1]$, let $m_+,m_->0$ be such that $m_+ \geq m_-$,
%     and let $\symbolnot \in \hyposymbclass(\R^{2d})$ be such that 
%     there exist $R>0$ and $C>0$ for which
%         \begin{equation}\label{eq:assumption-shubin}
%         |z \cdot \nabla \symbolnot(z)| \geq C |\symbolnot(z)|,
%         \quad \text{for all } z \in \R^{2d} \text{ with } \|z\|_2 \geq R.
%         \end{equation}
%     Then, we have $\ecfup_{\symbolnot}(\lambda) \sim \volup_\symbolnot(\lambda)$, as $\lambda \to \infty$.
%  \end{theorem}
% \fi

\noindent
Note that our definition of the upward volume in
\eqref{eq:definition-upward-volume} differs from that used in
\cite[Theorem~30.1]{shubin_pseudodifferential_2001} by a multiplicative factor of
$(2\pi)^d$.
This discrepancy reflects the different normalization conventions for the
Fourier transform; see \eqref{eq:definition-fourier-transform} and
\cite[eq.~(1.2)]{shubin_pseudodifferential_2001}.



The second result we invoke is due to Dauge and Robert \cite{dauge1987weyl}.
Their work concerns compact pseudodifferential operators of negative order and
establishes a Weyl-type asymptotic formula for the eigenvalue-counting function,
including a remainder estimate. Earlier results in this direction are due to Birman and Solomyak
\cite{birman_solomyak_1980_anisotropic}, who analyzed compact pseudodifferential
operators with symbols that are homogeneous or quasi-homogeneous with respect to
the phase variable.
A typical example of such behavior is given by symbols $\sigma$ for which there
exists $\zeta<0$ such that
\[
\sigma(x,t\omega)=t^{\zeta}\sigma(x,\omega),
\]
for all $t\ge1$, $x\in\R^d$, and $\omega\in\R^d$ with $\|\omega\|_2$ sufficiently
large. Dauge and Robert relax the requirement of exact homogeneity and work under more
general decay assumptions at infinity, thereby extending Weyl-type spectral
asymptotics to a broader class of compact pseudodifferential operators. This
flexibility makes their result applicable to the negative-order symbol regime
considered in the present paper.

Dauge--Robert formulate their results in terms of general weight functions
\(\phi,\varphi\in C^\infty(\R^{2d})\) taking values in \(\R_+^*\); see
\cite[Definition~1.1]{dauge1987weyl}.
In this framework, a function \(w\in C^\infty(\R^{2d})\) with values in \(\R_+^*\)
is called \((\phi,\varphi)\)-continuous if there exist constants \(c,C>0\) such that
\[
    C^{-1}
    \le
    \frac{w(x+x',\omega+\omega')}{w(x,\omega)}
    \le
    C
\]
for all \(x,x',\omega,\omega'\in\R^d\) satisfying
\[
    \frac{\|x'\|_2}{\varphi(x,\omega)}
    +
    \frac{\|\omega'\|_2}{\phi(x,\omega)}
    \le c .
\]
The weight \(w\) is said to be \((\phi,\varphi)\)-temperate if there exist constants
\(C,\zeta>0\) so that
\begin{equation}\label{eq:def-phi-phi-temperate}
    w(x+x',\omega+\omega')
    \le
    C\, w(x,\omega)
    \bigl[1+\|x'\|_2\,\phi(x,\omega)
          +\|\omega'\|_2\,\varphi(x,\omega)\bigr]^{\zeta},
\end{equation}
for all \(x,x',\omega,\omega'\in\R^d\).

Further, given a \((\phi,\varphi)\)-temperate weight \(w\), Dauge--Robert define the symbol
class \(\Gamma(\R^{2d};w,\phi,\varphi)\) as the set of real-valued functions
\(\sigma\in C^\infty(\R^{2d})\) such that, for every pair of multi-indices
\(\alpha,\beta\in\N^{d}\), one has
\begin{equation}\label{eq:weight-class-DR}
    \bigl|\partial_x^\alpha \, \partial_\omega^\beta \, \sigma(x,\omega)\bigr|
    \;\leq C_{\alpha,\beta}
    \frac{w(x,\omega)}
         {\phi(x,\omega)^{|\alpha|}\,\varphi(x,\omega)^{|\beta|}},
    \qquad x,\omega\in\R^d,
\end{equation}
with a constant $C_{\alpha,\beta} > 0$.

% \iffalse 
% Dauge-Robert state their results in terms of general weight functions $\phi, \varphi \in C^\infty(\R^{2d})$ with values in $\R_+^*$.
% To do so, they introduce the following two notions:
% \begin{itemize}
%     \item a function $w \in C^\infty(\R^{2d})$ with values in $\R_+^*$ is called $(\phi, \varphi)$-continuous if there exists $c, C>0$ such that 
%     \begin{equation*}
%         C^{-1} 
%         \leq \frac{w(x+x',\omega+\omega')}{w(x, \omega)} \leq C
%         \quad \text{for all }x,x',\omega,\omega'\in\R^d \text{ with }\frac{\|x'\|_2}{\varphi(x, \omega)}
%         + \frac{\|\omega'\|_2}{\phi(x, \omega)}
%         \leq  c ;
%     \end{equation*}

%     \item a function $w\in C^\infty(\R^{2d})$ with values in $\R_+^*$ is called $(\phi, \varphi)$-temperate if there exists $C, \zeta>0$ such that 
%     \begin{equation}\label{eq:def-phi-phi-temperate}
%         w(x+x',\omega+\omega') 
%         \leq C w(x, \omega) \left[1+ \|x'\|_2\, \phi(x, \omega)
%         + \|\omega'\|_2\, \varphi(x, \omega)\right]^\zeta,
%     \end{equation}
%     for all $x,x',\omega,\omega'\in\R^d$.
% \end{itemize}
% Given a $(\phi, \varphi)$-temperate weight $w$, 
% they define, on the template of \eqref{eq:decay-cond-symbol-classs}, 
% a more general class of symbols parameterized by $w$, $\phi$, and $\varphi$.
% Specifically, the class $\Gamma(\R^{2d}; w, \phi, \varphi)$ 
% consists of real-valued functions $\symbolnot \in C^\infty(\R^{2d})$ such that, 
% for every multi-indices $\alpha, \beta \in\N^{2d}$, there exists a $C_{\alpha,\beta}>0$ so that
% \begin{equation}\label{eq:weight-class-DR}
%     \left \lvert\partial^\alpha_x \, \partial^\beta_\omega \, \symbolnot(x, \omega) \right \rvert 
%     \leq C_{\alpha,\beta} \left \lvert \frac{w(x, \omega)}{\phi(x, \omega) ^{|\alpha|}\, \varphi(x,\omega)^{|\beta|}} \right \rvert, 
%     \quad \text{for all } x,\omega \in \R^d.
% \end{equation}
% \fi 


In order to control the spectral asymptotics of compact pseudodifferential
operators with symbol decay, Dauge--Robert impose a set of structural conditions
on the weight functions \(w,\phi,\varphi\) governing localization and scaling in
phase space. For ease of reference, we reproduce these assumptions below, using
the same labels as in \cite{dauge1987weyl}.

\begin{citemize}
    \item[(H1)]
    The functions \(\phi^{-1}\) and \(\varphi^{-1}\) are
    \((\phi,\varphi)\)-continuous, \((1,1)\)-temperate, and bounded on
    \(\R^{2d}\).

    \item[(H2)]
    There exist constants \(C,\zeta>0\) such that
    \begin{equation}\label{eq:DR-H2}
        C\,(1+\|x\|_2+\|\omega\|_2)^{\zeta}
        \le \phi(x,\omega)\,\varphi(x,\omega),
        \qquad x,\omega\in\R^{d}.
    \end{equation}

    \item[(W)]
    The weight \(w\) belongs to the symbol class
    \(\Gamma(\R^{2d};w,\phi,\varphi)\) and is
    \((\phi,\varphi)\)-temperate.

    \item[(N)]
    There exist constants \(K,K',\gamma,\gamma'>0\) such that
    \begin{equation}\label{eq:DR-N}
        K'\,w(x,\omega)^{\gamma'}
        \le \phi^{-1}(x,\omega)\,\varphi^{-1}(x,\omega)
        \le K\,w(x,\omega)^{\gamma},
        \qquad x,\omega\in\R^{d}.
    \end{equation}
\end{citemize}

Under assumptions \emph{(H1)}–\emph{(N)}, Dauge–Robert establish the following Weyl-type
asymptotic formula for the eigenvalue-counting function of compact
pseudodifferential operators whose symbol
$\sigma$ belongs to $\Gamma(\R^{2d}; w, \phi, \varphi)$.

% In order to state result, 
% we need to introduce the following two quantities, namely
% \begin{equation*}
%     M_\sigma^+(\lambda) 
%     \coloneqq  \# \left\{n\in\N^* \mid \lambda_n \geq \lambda \right\}
%     \quad \text{and} \quad 
%     M_\sigma^-(\lambda) 
%     \coloneqq  \# \left\{n\in\N^* \mid -\lambda_n \geq \lambda \right\},
%     \quad \text{for all } \lambda>0,
% \end{equation*}
% and
% \begin{equation*}
%     V^+_{\symbolnot}(\lambda) 
%     \coloneqq \int_{\mathbb{R}^{2d}} \mathbbm{1}_{\{\symbolnot(x, \omega)\, > \, \lambda\}} \diff  x\diff \omega
%     \quad \text{and} \quad 
%     V^-_{\symbolnot}(\lambda) 
%     \coloneqq \int_{\mathbb{R}^{2d}} \mathbbm{1}_{\{-\symbolnot(x, \omega)\, > \, \lambda\}} \diff  x\diff \omega,
% \end{equation*}
% for all $\lambda>0$,
% which distinguish the sign of the eigenvalues and the symbol in  \eqref{eq:definition-ecf} and \eqref{eq:definition-volume-function}, respectively.


\begin{theorem}[{\cite[Theorem~1.3]{dauge1987weyl}}]\label{thm:Dauge-Robert}
    Let $w, \phi, \varphi \in C^\infty(\R^{2d})$
    be weight functions satisfying (H1), (H2), (W), and (N).
    Let $\sigma\in\Gamma(\R^{2d}; w, \phi, \varphi)$ 
    be non-negative
    and assume that the maps $\lambda \mapsto V_\sigma(\lambda)$ % ,
    % $\lambda \mapsto V^-_\sigma(\lambda)$, 
    and $\lambda \mapsto V_w(\lambda)$ are regularly varying at zero with $V_w (\lambda) = O_{\lambda\to 0} \left(V_\sigma (\lambda)\right)$.
    %%%T12[0K]: Is the first part of the statement, i.e., the one on M_{|T_\sigma|} really contained in Dauge-Robert Theorem 1.3?
    Then, $M_{|T_\sigma|} (\lambda) \sim M_{T_\sigma} (\lambda) \sim V_{\sigma} (\lambda)$, as $\lambda\to 0$, where $T_\sigma$ is the Weyl quantization of $\sigma$.
    % \begin{equation*}
    %     % M^-_\sigma (\lambda) = V^-_\sigma (\lambda) + o_{\lambda\to\infty} \left(V_w (\lambda)\right)
    %     % \quad \text{and} \quad
    %     M_\sigma (\lambda) = V_\sigma (\lambda) + o_{\lambda\to\infty} \left(V_w (\lambda)\right).
    % \end{equation*}
\end{theorem}

%%%T6: Do Dauge-Robert assume non-negative sigma as well?

\noindent
We remark that both Theorems~\ref{thm-shubin} and~\ref{thm:Dauge-Robert} are stated
here in slightly simplified forms relative to their original counterparts in
\cite[Theorem~30.1]{shubin_pseudodifferential_2001} and
\cite[Theorem~1.3]{dauge1987weyl}, respectively.
More precisely, \cite[Theorem~30.1]{shubin_pseudodifferential_2001} replaces the
condition~\eqref{eq:assumption-shubin} by the weaker requirement
\begin{equation}\label{eq:assumption-shubin2}
    |z \cdot \nabla \sigma(z)|
    \ge
    C\,|\sigma(z)|^{1-\delta},
    \qquad
    \text{for all } z \in \R^{2d} \text{ with } \|z\|_2 \ge R,
\end{equation}
for some $\delta \in [0,\rho')$, where $\rho'>0$ is a parameter depending on the
symbol $\sigma$ (with the universal lower bound $\rho' \ge \rho/m_+$).
Under this assumption, Shubin establishes the refined asymptotic expansion
\[
    \ecfup_\sigma(\lambda)
    =
    \volup_\sigma(\lambda)
    \left(
        1
        +
        O_{\lambda\to\infty}
        \bigl(\lambda^{\delta-\rho'+\eta}\bigr)
    \right),
    \qquad
    \text{for every } \eta>0.
\]
Similarly, \cite[Theorem~1.3]{dauge1987weyl} provides a remainder estimate of the
form
\[
    M_\sigma(\lambda)
    =
    V_\sigma(\lambda)
    +
    O_{\lambda\to0}
    \left(
        \lambda^{\zeta}\,V_w(\lambda)
    \right),
    \qquad
    \text{for some } \zeta>0,
\]
under their original set of assumptions. %%%T6[OK]: I added "under their ..." as otherwise it would be unclear under which assumptions their result holds.
In the present paper, only first-order asymptotics are required, and the finer
remainder estimates available in the original results are not needed.
For this reason, the simplified formulations given in
Theorems~\ref{thm-shubin} and~\ref{thm:Dauge-Robert} are sufficient for our
purposes. 
% Finally, 
Additionally, we note that the regular-variation assumption imposed in
Theorem~\ref{thm:Dauge-Robert} is stronger than the original condition~(T) in
\cite{dauge1987weyl}. This strengthening is introduced here in order to enable the
subsequent connection between eigenvalue-counting functions and metric entropy,
and therefore does not entail any loss of generality in the context of our
results. Finally, the non-negativity assumption on the symbol $\sigma$ can in principle be
removed. In \cite{dauge1987weyl}, this is achieved by separating the asymptotic
analysis of the positive and negative parts of the spectrum of $T_\sigma$. 
%%%T12[OK]: mention that $M_{|T_\sigma|} (\lambda) \sim M_{T_\sigma} (\lambda) + O(1), would also resolve my question on the theorem 2.13 of Dauge-Robert above.
When
$\sigma$ is strictly positive and of negative order, the negative spectrum 
% is finite and 
does not contribute to the leading asymptotics,
and $M_{|T_\sigma|} (\lambda) \sim M_{T_\sigma} (\lambda)$. Since our analysis
focuses exclusively on the leading-order behavior, the formulation of
Theorem~\ref{thm:Dauge-Robert} for non-negative symbols is compatible with the
assumptions of Theorem~\ref{thm:main-result}.

% \iffalse 
% \textcolor{blue}{Finally, 
% one can dispense with the non-negativity assumption on the symbol $\sigma$;
% this is done in \cite{dauge1987weyl} by separating the analysis of the positive and negative parts of the spectrum of the corresponding operator $T_\sigma$.
% In view of the strict-positivity assumption in % the statement of
% Theorem~\ref{thm:main-result} (cf. the subsequent discussion), 
% there is no loss of generality in our formulation of Theorem~\ref{thm:Dauge-Robert}.}
% \fi 




% \iffalse 
% We remark that both Theorems~\ref{thm-shubin} and \ref{thm:Dauge-Robert} are actually weaker versions of \cite[Theorem~30.1]{shubin_pseudodifferential_2001} and \cite[Theorem~1.3]{dauge1987weyl}, respectively.
% Concretely, \cite[Theorem~30.1]{shubin_pseudodifferential_2001} applies under \eqref{eq:assumption-shubin} replaced by
% \begin{equation}\label{eq:assumption-shubin2}
%     |z \cdot \nabla \symbolnot(z)| \geq C \symbolnot(z)^{1-\delta},
%     \quad \text{for all } z \in \R^{2d} \text{ with } \|z\|_2 \geq R,
% \end{equation}
% for some $\delta\in[0, \rho')$, where $\rho'\geq \rho / m_+ $ is a parameter that depends on the symbol $\symbolnot$,
% and establishes the stronger asymptotic equivalence
%  \begin{equation*} 
%     \ecfup_{\symbolnot}(\lambda)
%     =\volup_\symbolnot(\lambda)\left(1+O_{\lambda\to\infty}\left(\lambda^{\delta-\rho^{\prime}+\eta}\right)\right),
%     \quad \text{for all } \eta>0.
% \end{equation*}
% Likewise, \cite[Theorem~1.3]{dauge1987weyl} actually characterizes the remainder term according to 
% \begin{equation*}
%     M_\sigma (\lambda) = V_\sigma (\lambda) + O_{\lambda\to 0} \left(\lambda^\zeta \,  V_w (\lambda)\right),
%     \quad \text{for some } \zeta>0.
% \end{equation*}
% However, in the present paper, we do not need to go beyond first-order terms, 
% which is why our formulations of these results are sufficient.
% Similarly, the regular-variation assumption in Theorem~\ref{thm:Dauge-Robert} is stronger than the original assumption (T) in \cite{dauge1987weyl}, 
% but is needed to relate the metric entropy and the eigenvalue-counting function, so that there is no issue on losing generality here.
% \fi 


% % Finally, we conclude this appendix by introducing a subclass of hypoelliptic symbols that is useful in the main body.
% % We say that $\sigma\in C^\infty(\R^{2d})$ is factorially hypoelliptic if 
% % there exists $\alpha_0, \beta_0, R, C >0$ such that 
% %     \begin{equation}\label{eq:fact-hypo1}
% %         \sigma(x,\omega)\geq C\left(1+\|x\|^2_2\right)^{\frac{-\alpha_0}{2}}\left(1+\|\omega\|^2_2\right)^{\frac{-\beta_0}{2}},
% %         \quad \text{for all } x,\omega\in \R^d  \text{ with } \|(x,\omega)\|_2 \ge R,
% %     \end{equation}
% %     and such that for every multi-indices $\alpha, \beta \in\N^{2d}$, there exists a $C_{\alpha,\beta}>0$ for which
% %     \begin{equation}\label{eq:fact-hypo2}
% %         \left\lvert \partial^\alpha_x \, \partial^\beta_\omega \, \symbolnot(x, \omega) \right\rvert
% %         \leq C_{\alpha,\beta} \left(1+\|x\|^2_2\right)^{\frac{-\alpha_0-\rho|\alpha|}{2}}\left(1+\|\omega\|^2_2\right)^{\frac{-\beta_0-\rho|\beta|}{2}},
% %         \quad \text{for all } x,\omega\in \R^d.
% %     \end{equation}



\end{document}
